\documentclass[11pt]{article}

\usepackage[english]{babel}
\usepackage[top=1in,bottom=1.2in,left=1in,right=1in]{geometry}
\usepackage{amssymb,amsmath,amsthm,mathtools}
\usepackage{microtype}
\usepackage{newtxtext,newtxmath}
\usepackage{graphicx}
\usepackage{placeins}
\usepackage[colorlinks=true,linkcolor=blue,citecolor=red,urlcolor=blue]{hyperref}
\hypersetup{pdfauthor={}}

\newtheorem{theorem}{Theorem}[section]
\newtheorem{proposition}[theorem]{Proposition}
\newtheorem{lemma}[theorem]{Lemma}
\newtheorem{corollary}[theorem]{Corollary}
\theoremstyle{remark}
\newtheorem{remark}[theorem]{Remark}

\numberwithin{equation}{section}

\newcommand{\E}{\mathbb E}
\newcommand{\Pbb}{\mathbb P}
\newcommand{\R}{\mathbb R}
\newcommand{\cN}{\mathcal N}
\newcommand{\eps}{\varepsilon}
\newcommand{\wh}{\widehat}
\newcommand{\op}{\mathrm{op}}
\newcommand{\Id}{\mathrm{Id}}
\DeclareMathOperator{\Var}{Var}
\DeclareMathOperator{\supp}{supp}

\title{Radial Phase Transitions for Gaussian Mixture NPMLEs\\
in Fixed and Growing Dimensions}
\author{}
\date{}

\begin{document}

\maketitle

\begin{abstract}
Can a fitted Gaussian mixture distinguish a single Gaussian population from
latent heterogeneity? Without regularization, the support of the mixture
nonparametric maximum likelihood estimator (NPMLE) diverges even when the data
come from a single Gaussian component. We ask whether fitting wider
Gaussian components regularizes this overfitting. The observations are
i.i.d.\ spherical Gaussians with covariance $(1-\eps_n)I_d$, while the fitted
mixture uses kernels with covariance $I_d$.

In dimension two, the transition occurs when $\eps_n\log n$ is of constant
order. The limiting profile is explicit and combines the largest point of a
radial Poisson edge process with the aggregate contribution of observations
just inside the edge. For every sequence $d_n\ge3$, fixed or growing, a single
exact Gamma-quantile normalization describes the transition. When the
normalized coordinate tends to \(s\), the NPMLE is unique with probability
tending to one, and its support size minus one converges to
\(\operatorname{Pois}(e^{-s})\). Thus the one-atom probability tends to
\(\exp\{-e^{-s}\}\). For fixed $d\ge3$, the Gamma-quantile normalization
reduces to
\(\eps_n=(d/2-1)(\log\log n)/\log n+O(1/\log n)\). Under a common coupling
of the data, the support counts across each critical window converge jointly
to a Poisson tail-count process. We also identify the fitted edge masses and,
in fixed dimension, their directions.
\end{abstract}

\medskip
\noindent\textbf{Note.} This draft was produced through an iterative process in which Mark Sellke repeatedly asked GPT-5.6 in OpenAI Codex to prove, refine, and write up further extensions.
\medskip

\setcounter{tocdepth}{2}
\tableofcontents

\section{Introduction}
\label{sec:introduction}

A natural way to look for latent heterogeneity is to fit a mixture and inspect
the number of fitted components. For the Gaussian mixture nonparametric
maximum likelihood estimator, this suggests a basic question: when the data
come from one spherical Gaussian, does maximum likelihood select a single
component mean? Our companion paper shows that it does not when the fitted
component covariance matches the data-generating covariance
\cite{supportDivergenceCompanion}; in fact, the fitted support diverges in
probability. Here we ask whether variance inflation can regularize this
overfitting. We determine precisely how much inflation is
needed in every dimension $d\ge2$, including dimensions that grow with the
sample size.

Let $d=d_n\ge2$, let $Z_1,\ldots,Z_n$ be independent
$\cN(0,I_d)$ vectors, and observe
\begin{equation}
\label{eq:model}
    X_i=\sqrt{1-\eps_n}\,Z_i,
    \qquad 0\le\eps_n<1.
\end{equation}
We fit location mixtures of $\cN(\mu,I_d)$ densities, whose component
covariance is therefore larger than the data covariance. For a probability
measure $G$ on $\R^d$, write
\[
    q_G(x)=\int_{\R^d}\phi_d(x-u)\,dG(u),
    \qquad
    \ell_n(G)=\sum_{i=1}^n\log q_G(X_i),
\]
where $\phi_d$ is the $\cN(0,I_d)$ density. A mixture NPMLE is any maximizer
$\wh G_n$ of $\ell_n$. Because multivariate NPMLEs need not be unique
\cite{soloff2024multivariate}, we
measure their support complexity by
\begin{equation}
\label{eq:minimum-support-size}
    \wh K_n
    =\min\bigl\{|\supp(G)|:G\text{ maximizes }\ell_n\bigr\}.
\end{equation}
Lindsay's support-reduction theorem gives a maximizer with at most \(n\) atoms,
so the minimum in \eqref{eq:minimum-support-size} is attained
\cite{lindsay1983geometryA}. Nonuniqueness does not affect the one-atom event:
Proposition~\ref{prop:one-atom-unique} shows that, almost surely, any one-atom
maximizer is the unique NPMLE.  The minimum-support convention is likewise
immaterial in the critical regimes below: our proof shows that the NPMLE is
unique with probability tending to one.

For each sample, the one-atom event is monotone in the variance defect.
Increasing $\eps_n$ damps every likelihood perturbation away from the sample
mean, so once a one-atom fit is optimal, it remains optimal under any larger
defect. We can therefore study a samplewise critical value of $\eps_n$;
Lemmas~\ref{lem:variance-monotonicity}
and~\ref{lem:critical-defect-interval} show that this critical value is well
defined. The asymptotic law of the critical defect changes sharply between
dimensions two and three: in
dimension two, observations just inside the radial edge have a nonvanishing
collective effect, whereas in dimensions at least three the radial extremes
alone govern the transition.

\subsection{Main results}
\label{sec:intro-results}

We first state the exceptional two-dimensional limit. Put
\begin{equation}
\label{eq:d2-intro-parameters}
    \kappa_n=\frac{\eps_n}{2(1-\eps_n)},
    \qquad
    \rho_n=\sqrt{2\log n},
    \qquad
    \lambda_n=\kappa_n\rho_n^2.
\end{equation}
For $\lambda>0$, let $x_\lambda$ be the unique real number satisfying
\begin{equation}
\label{eq:x-lambda-definition}
    \sup_{r\in\R}
       \left\{\Phi(-r)+\exp\left(x_\lambda-\frac{r^2}{2}\right)\right\}
    =e^\lambda,
\end{equation}
where $\Phi$ and $\phi$ are the standard normal distribution function and
density, respectively, and set
\begin{equation}
\label{eq:Q2-definition}
    Q_2(\lambda)=\exp\{-e^{-x_\lambda}\}.
\end{equation}
The term $\Phi(-r)$ in \eqref{eq:x-lambda-definition} is the limiting
contribution of observations just inside the radial edge, while the
exponential term comes from the largest radial observation. The function
\(Q_2\) extends continuously and strictly increasingly to \([0,\infty]\), with
\(Q_2(0)=0\) and \(Q_2(\infty)=1\); see Lemma~\ref{lem:q-two-curve}.

For dimensions at least three, an exact radial normalization covers all
growth regimes. Let $Y_n$ have the $\operatorname{Gamma}(d_n/2,1)$
distribution, with density $f_n$ and upper tail $\overline F_n$. Define
\begin{equation}
\label{eq:ba-intro}
    n\overline F_n(b_n)=1,
    \qquad
    a_n=\frac{\overline F_n(b_n)}{f_n(b_n)},
\end{equation}
and put
\begin{equation}
\label{eq:exact-intro-coordinate}
    \gamma_n=(1-\eps_n)^{-1},
    \qquad
    \sigma_n^2=1-\frac1n,
    \qquad
    R_n^\star=\frac{(n-1)^2}{n(n-2)}\log(n-1),
    \qquad
    s_n^\star
    =\frac{\gamma_nR_n^\star/\sigma_n^2-b_n}{a_n}.
\end{equation}
Thus $b_n$ is the \(1-1/n\) quantile of the radial energy \(|Z|^2/2\), and
$a_n$ is its extreme-value scale. The correction $R_n^\star$ is asymptotically
negligible in fixed and moderately growing dimension, but changes the limiting
coordinate when $d_n\asymp n^2\log n$.

We state the transition directly as a limit law for the support size; the
one-atom probability is then the probability that \(\widehat K_n-1\) equals
zero. Here \(\operatorname{Pois}(\mu)\) denotes the Poisson law with mean
\(\mu\).
\begin{theorem}
\label{thm:main-support-law}
Under either critical-window hypothesis below, the NPMLE is unique with
probability tending to one.  For the model \eqref{eq:model}, the support-size
variable in \eqref{eq:minimum-support-size} has the following limits.
\begin{enumerate}
\item If $d=2$ and $\lambda_n\to\lambda\in(0,\infty)$, then
\begin{equation}
\label{eq:main-d2-support-law}
    \wh K_n-1\Rightarrow\operatorname{Pois}(e^{-x_\lambda}),
    \qquad
    \Pbb\{\wh K_n=1\}\to Q_2(\lambda).
\end{equation}
Moreover, $\wh K_n=O_{\Pbb}(1)$ whenever
$\liminf_n\lambda_n>0$.
\item If $d_n\ge3$ and $s_n^\star\to s\in\R$, then
\begin{equation}
\label{eq:main-highd-support-law}
    \wh K_n-1\Rightarrow\operatorname{Pois}(e^{-s}),
    \qquad
    \Pbb\{\wh K_n=1\}\to\exp\{-e^{-s}\}.
\end{equation}
Moreover, $\wh K_n=O_{\Pbb}(1)$ whenever
$\liminf_n s_n^\star> -\infty$, and $\wh K_n=1$ with probability tending
to one if $s_n^\star\to\infty$.
\end{enumerate}
\end{theorem}

For fixed $d\ge3$, the Gamma tail expansion gives
\[
    b_n=\log n+\left(\frac d2-1\right)\log\log n
        -\log\Gamma(d/2)+o(1),
    \qquad a_n\to1.
\]
The critical scaling forces \(\eps_n\to0\) and
\(\eps_n^2\log n=o(1)\). Thus $s_n^\star\to s$ is equivalent to
\begin{equation}
\label{eq:fixed-d-intro-window}
    \eps_n\log n-\left(\frac d2-1\right)\log\log n
        +\log\Gamma(d/2)\longrightarrow s.
\end{equation}
The transition therefore occurs at scale $(\log\log n)/\log n$ in fixed
dimensions $d\ge3$, but at scale $1/\log n$ in dimension two.

The exact normalization also gives explicit scales when $d_n$ grows. For
$d_n=o(\log n)$ with $d_n\to\infty$, the critical defect satisfies
\begin{equation}
\label{eq:sublog-intro-critical}
    \eps_{c,n}
    \sim\frac{d_n}{2\log n}
       \log\frac{2\log n}{d_n},
\end{equation}
and the transition window has width asymptotic to $1/\log n$. In larger
dimensions, the implicit Gamma quantile \eqref{eq:ba-intro} is simpler and more
accurate than separate expansions for different growth rates. At scale
$d_n\asymp n^2\log n$, the $O(n^{-1})$ correction in $R_n^\star$ shifts the
Gumbel coordinate by order one, and the exact formula incorporates this shift.

The two fixed-parameter limits in Theorem~\ref{thm:main-support-law} describe
one variance defect at a time. In fact, a single radial Poisson cloud governs
the entire regularization path. The path theorem below
makes precise the sense in which an extreme observation creates a fitted atom
and then loses it as the fitted variance is increased. Let \(\mathcal X\) be a
Poisson point process on \(\R\) with intensity \(e^{-x}\,dx\). Couple each path
through one standardized sample \(Z_1,\ldots,Z_n\). Write
\(\wh K_n(z_{1:n})\) for \eqref{eq:minimum-support-size} computed from the data
\(z_{1:n}\), and define
\begin{align}
\label{eq:intro-d2-path-parametrization}
 \eps_{n,2}(\lambda)
 &=\frac{2\lambda}{\rho_n^2+2\lambda},
 &K_{n,2}(\lambda)
 &=\wh K_n\bigl(\sqrt{1-\eps_{n,2}(\lambda)}\,Z_{1:n}\bigr),\\
\label{eq:intro-highd-path-parametrization}
 \eps_{n,\ge3}(s)
 &=1-\frac{R_n^\star}{\sigma_n^2(b_n+a_ns)},
 &K_{n,\ge3}(s)
 &=\wh K_n\bigl(\sqrt{1-\eps_{n,\ge3}(s)}\,Z_{1:n}\bigr).
\end{align}
The parametrization in \eqref{eq:intro-highd-path-parametrization} applies to
any sequence \(d_n\ge3\); both defects lie in
\([0,1)\) uniformly
on the compact intervals considered below, once \(n\) is large.

\begin{theorem}
\label{thm:regularization-path}
For every compact \(J\Subset(0,\infty)\), in dimension two,
\begin{equation}
\label{eq:intro-d2-path-limit}
 \{K_{n,2}(\lambda)-1:\lambda\in J\}
 \Longrightarrow
 \{\mathcal X((x_\lambda,\infty)):\lambda\in J\}
 \qquad\text{in }D(J).
\end{equation}
For every compact \(I\subset\R\) and every sequence \(d_n\ge3\),
\begin{equation}
\label{eq:intro-highd-path-limit}
 \{K_{n,\ge3}(s)-1:s\in I\}
 \Longrightarrow
 \{\mathcal X((s,\infty)):s\in I\}
 \qquad\text{in }D(I).
\end{equation}
Here both Skorokhod spaces carry the \(J_1\) topology. In either case, with
probability tending to one, the NPMLE is unique simultaneously throughout the
specified interval.
\end{theorem}

The map \(\lambda\mapsto x_\lambda\) is a smooth increasing bijection; write
\(\Lambda\) for its inverse. The limiting path is a decreasing step function:
a Poisson point \(x\) disappears at \(\lambda=\Lambda(x)\) in dimension two
and at \(s=x\) in dimensions at least three. Proposition
\ref{prop:path-masses} gives the corresponding mass profiles and spatial marks.

Samplewise monotonicity turns these probability limits into limit laws for the
critical variance defect itself. For a standardized sample $z_{1:n}$, define
\[
    \eps_*(z_{1:n})
    =\inf\{\eps\in[0,1):\sqrt{1-\eps}\,z_{1:n}
       \text{ has a one-atom NPMLE}\}.
\]
In dimension two, the corresponding variable
\(\rho_n^2\eps_*(Z_{1:n})/\{2(1-\eps_*(Z_{1:n}))\}\) has limiting
distribution function \(Q_2\). For \(d_n\ge3\), define
\[
 S_{*,n}
 =\frac{R_n^\star/
     [\sigma_n^2\{1-\eps_*(Z_{1:n})\}]-b_n}{a_n}.
\]
Then \(S_{*,n}\) converges in distribution to a standard Gumbel variable.

\subsection{Ideas of the proof}
\label{sec:intro-proof-ideas}

The first-order optimality criterion for a one-atom fit can be written as a
random field. After centering the observations and rescaling displacements
from the sample mean, this field becomes a sum of unit-width Gaussian bumps.
The direction and radius of $Z_i-\bar Z_n$ determine the center and height of
the $i$th bump. A sufficiently large observation can therefore push the field
above its critical level and force an additional fitted atom.

In dimension two, observations just below the radial maximum have an order-one
aggregate contribution. In a local window centered at the maximizing
direction, this contribution converges to $\Phi(-r)$ at radial offset $r$.
Angular separation leaves at most one extreme observation in the same window,
which yields \eqref{eq:x-lambda-definition}. In dimensions at least three, the
aggregate inward contribution vanishes at the transition. The radial
exceedances converge to a Poisson process with intensity $e^{-x}\,dx$, and
each observation above the critical radial-energy threshold contributes one
additional fitted atom.

Although the limiting law is universal for $d_n\ge3$, proving that no other
field crossing occurs requires two forms of control. In the lower-dimensional
regimes, radial edge points are well separated, and spatial empirical-process
estimates control the contribution of the remaining observations. When the
dimension is large enough that the centered observations form an
approximately regular simplex, an entropy dual replaces Euclidean covering
and reduces the local optimization to an ordered mean-field Potts problem.
Section~\ref{sec:support} then converts the surviving field crossings into
fitted atoms. For the path limit, we differentiate the fixed-parameter local
likelihood expansions in the variance parameter. After the dimension-dependent
normalization of the edge masses, each mass crosses zero with a derivative
bounded away from zero. The support identity therefore continues uniformly
across the critical window.

\subsection{Related literature}
\label{sec:related-literature}

The mixture NPMLE goes back to Robbins and Kiefer--Wolfowitz
\cite{robbins1950generalization,kiefer1956consistency}. Lindsay developed its
likelihood geometry and finite-support reduction
\cite{lindsay1983geometryA}. In empirical Bayes, the fitted mixing
distribution estimates the prior; representative treatments include
Jiang--Zhang~\cite{jiang2009general}, Koenker--Mizera
\cite{koenker2014convex}, Dicker--Zhao~\cite{dicker2016high}, and the
multivariate theory of Soloff--Guntuboyina--Sen
\cite{soloff2024multivariate}.

Recent work distinguishes statistical accuracy from the combinatorial
complexity of an NPMLE\@. Polyanskiy and Wu proved logarithmic upper bounds
on the support size in one-dimensional subgaussian models and asked what
happens under a single Gaussian component~\cite{polyanskiy2020self}. Our
companion paper proves that the support diverges under every fixed finite
mixing distribution \cite{supportDivergenceCompanion}. Polyanskiy and Sellke
study certified computation and generic support behavior for multivariate
Gaussian mixtures \cite{polyanskiySellke2025}. A separate companion paper
treats one-dimensional variance misspecification, where the critical scale is
polynomial and the limit is governed by a two-sided compensated Poisson field
\cite{varianceD1Companion}.

A separate line of Gaussian-mixture NPMLE theory concerns density estimation
or empirical-Bayes risk; see Zhang~\cite{zhang2009generalized},
Saha--Guntuboyina~\cite{saha2020nonparametric}, Doss et al.
\cite{doss2020optimal}, and more recent stability and adaptivity results
\cite{zhangVolgushev2026adaptivity,GhoshGuntuboyinaMukherjeeTran2026,
ChenWu2026RegretHellinger}. Smooth mixture NPMLEs provide a closely related
form of variance inflation for estimation and inference
\cite{KimSen2026SmoothNPMLE}. In the present problem the variance inflation
vanishes with $n$, and we study the support of the likelihood maximizer rather
than its estimation risk.

Mixture homogeneity tests compare a point-mass null with an unrestricted
mixture alternative; see Gu--Koenker--Volgushev
\cite{guKoenkerVolgushev2018} and the references therein. Our criterion asks
for more than rejection or acceptance at a likelihood-ratio threshold: it
requires the unrestricted optimizer itself to collapse to one atom. Because
this event is monotone in the variance inflation, each sample has a critical
inflation level, analogous in spirit to the critical bandwidth in modality
testing
\cite{silverman1981multimodality,hallYork2001calibration}.

\section{Likelihood Geometry and Radial Normalization}
\label{sec:geometry}

Both Theorems~\ref{thm:main-support-law} and
\ref{thm:regularization-path} rest on two dimension-free reductions. The KKT
criterion identifies one-atomicity with a supremum of a Gaussian-bump field
and makes its monotonicity in the variance defect explicit. The Gamma
normalization then puts the radial extremes on a common Poisson scale for
every sequence \(d_n\ge2\). This section establishes the KKT and radial
reductions and the exact centering identities used in all later regimes.

\subsection{Notation and one-atom likelihood geometry}\label{sec:mv-one-atom-geometry}

The deterministic part of both main theorems is the reduction from mixture
likelihood to a Gaussian-bump field. This subsection derives that KKT
criterion, proves monotonicity in the variance defect, and records the exact
single-contact calculation used later in the simplex regime. We first fix the
notation used throughout the paper.

All asymptotic notation is as $n\to\infty$. For a deterministic positive
sequence $r_n$, the relation $U_n=O_{\Pbb}(r_n)$ means that
$|U_n|/r_n$ is tight, while $U_n=o_{\Pbb}(r_n)$ means that
$U_n/r_n\to0$ in probability. We write $r_n\lesssim s_n$ when
$r_n\le C s_n$ for a constant independent of $n$, and use
$\gtrsim,\asymp,\ll$, and $\gg$ analogously. We write
$U_n\xrightarrow{\Pbb}U$ for convergence in
probability and $U_n\Rightarrow U$ for convergence in distribution. A
positive random sequence satisfies $U_n=n^{o_{\Pbb}(1)}$ when, for every
$\eta>0$, $\Pbb\{n^{-\eta}\le U_n\le n^\eta\}\to1$. A
subscript on asymptotic notation records allowed dependence: for example,
$O_{\Pbb,R}(1)$ is tight for each fixed $R$. In iterated limits,
$o_A(1)$ denotes a deterministic quantity that tends to zero as
$A\to\infty$ after the $n\to\infty$ limit. Constants implicit in
$O_m(\cdot)$ may depend on the fixed integer $m$, and other constants may
change from line to line. Point processes carry the vague topology, and a
standard Gumbel variable has distribution function
$s\mapsto\exp\{-e^{-s}\}$.

Let \(d=d_n\ge2\), let \(Z_1,\ldots,Z_n\) be independent
\(\cN(0,I_d)\) vectors, and observe
\[
  X_i=\sqrt{1-\eps_n}\,Z_i,
  \qquad 0\le\eps_n<1.
\]
We fit location mixtures of \(\cN(\mu,I_d)\) densities. The only one-atom
maximum-likelihood candidate is the point mass at the sample mean. Writing
\(\bar Z_n=n^{-1}\sum_iZ_i\) and
\(\widetilde Z_i=Z_i-\bar Z_n\), the directional-derivative criterion says
that this candidate is optimal exactly when
\begin{equation}
\label{eq:one-atom-event}
  E_n^{(1)}
  :=
  \left\{\sup_{t\in\R^d}D_n(t)\le1\right\},
  \qquad
  D_n(t)
  =
  \frac1n\sum_{i=1}^n
  \exp\left\{t\cdot\widetilde Z_i-\frac{\gamma_n|t|^2}{2}\right\},
\end{equation}
where
\begin{equation}
\label{eq:gamma-definition}
  \kappa_n=\frac{\eps_n}{2(1-\eps_n)},
  \qquad
  \gamma_n=\frac1{1-\eps_n}=1+2\kappa_n.
\end{equation}
Indeed, if \(a\) denotes displacement of a fitted component mean from the
sample mean, the usual directional derivative is
\[
  \frac1n\sum_i
  \exp\left\{a\cdot(X_i-\bar X_n)-\frac{|a|^2}{2}\right\}.
\]
The change of variables \(t=\sqrt{1-\eps_n}\,a\) gives
\eqref{eq:one-atom-event}. Notice that \(D_n(0)=1\), so the issue is whether
the field rises above its value at the origin.

Although a multivariate mixture NPMLE need not be unique, this causes no
ambiguity for the event studied here.  The following dimension-free fact
shows that \(E_n^{(1)}\) is almost surely the event that the NPMLE is unique
and one-atomic.

\begin{proposition}
\label{prop:one-atom-unique}
Let \(n\ge2\), and let \(X_1,\ldots,X_n\) be i.i.d.\ from a nondegenerate
spherical Gaussian distribution in \(\R^d\). The probability that
\(\delta_{\bar X_n}\) and a different mixing measure both maximize the
Gaussian-mixture likelihood is zero.
\end{proposition}

\begin{proof}
If \(\delta_{\bar X_n}\) and \(G\ne\delta_{\bar X_n}\) both maximize the
likelihood, strict concavity in the vector of fitted values implies
\(q_G(X_i)=q_{\delta_{\bar X_n}}(X_i)\) for every \(i\).  Since the one-atom
fit is optimal, its physical-coordinate directional field
\[
 \mathcal D_{X,n}(a)
 =\frac1n\sum_{i=1}^n
 \exp\left\{a\cdot(X_i-\bar X_n)-\frac{|a|^2}{2}\right\}
\]
satisfies \(\mathcal D_{X,n}(a)\le1\).
Moreover,
\[
 \int \mathcal D_{X,n}(u-\bar X_n)\,dG(u)
 =\frac1n\sum_i\frac{q_G(X_i)}
 {q_{\delta_{\bar X_n}}(X_i)}=1.
\]
Thus \(\mathcal D_{X,n}(a)=1\) for some \(a\ne0\).

Write \(Y=(X_1-\bar X_n,\ldots,X_n-\bar X_n)=RU\) in polar form in the subspace
\(\{(y_1,\ldots,y_n):\sum_i y_i=0\}\). Conditional on its direction \(U\),
the radius \(R\) has a density. After setting \(t=Ra\), the
directional-derivative field is
\[
 \mathcal D_{R,U}(t)
 =\frac1n\sum_i
 \exp\left\{t\cdot U_i-\frac{|t|^2}{2R^2}\right\}.
\]
For fixed \(U\) and \(t\ne0\), this is strictly increasing in \(R\).
Consequently, for each \(U\), at most one radius can admit both a one-atom
maximizer and a nonzero contact: at any larger radius that same contact
would make the directional derivative exceed one. Thus, conditional on
\(U\), the radius belongs to an at-most-singleton exceptional set, which has
probability zero because \(R\mid U\) has a density. Integration over \(U\)
proves the claim.
\end{proof}

The one-atom event is monotone in the variance defect. This elementary fact
allows all sharp-window limits below to be converted into laws for a
samplewise critical shrinkage.
\begin{lemma}
\label{lem:variance-monotonicity}
Fix $z_1,\ldots,z_n\in\R^d$, and let
$x_i^{(\eps)}=\sqrt{1-\eps}\,z_i$ for $0\le\eps<1$. If the sample based on
$x_1^{(\eps_1)},\ldots,x_n^{(\eps_1)}$ admits a one-atom NPMLE and
$\eps_2\ge\eps_1$, then the sample based on
$x_1^{(\eps_2)},\ldots,x_n^{(\eps_2)}$ also admits a one-atom NPMLE.
\end{lemma}

\begin{proof}
Let $\bar z=n^{-1}\sum_i z_i$ and
\[
 K_n(t)=\frac1n\sum_{i=1}^n
 \exp\left\{t\cdot(z_i-\bar z)-\frac{|t|^2}{2}\right\}.
\]
After the change of variables $t=a\sqrt{1-\eps}$, the one-atom directional
field is
\[
 e^{-\kappa(\eps)|t|^2}K_n(t),
 \qquad
 \kappa(\eps)=\frac{\eps}{2(1-\eps)}.
\]
Because $\kappa$ is increasing, the supremum of this field is nonincreasing
in $\eps$. The directional-derivative criterion proves the claim.
\end{proof}

The next lemma shows that the monotone one-atom set is nonempty and closed,
hence has the form \([\eps_*,1)\). Its endpoint \(\eps_*\) is the samplewise
critical defect used in the introduction.
\begin{lemma}
\label{lem:critical-defect-interval}
For every fixed sample $z_1,\ldots,z_n\in\R^d$, the set of
$\eps\in[0,1)$ for which
$\sqrt{1-\eps}\,z_1,\ldots,\sqrt{1-\eps}\,z_n$ admits a one-atom NPMLE is
a nonempty interval of the form $[\eps_*,1)$.
\end{lemma}

\begin{proof}
Let $y_i=z_i-\bar z$ and $M=\max_i|y_i|$. The empirical random variable
taking the values $t\cdot y_i$ has mean zero and range of length at most
$2M|t|$. Hoeffding's lemma gives
\[
 \log\left(\frac1n\sum_i e^{t\cdot y_i-|t|^2/2}\right)
 \le \frac{M^2-1}{2}|t|^2.
\]
Hence the one-atom field
$e^{-\kappa|t|^2}n^{-1}\sum_i e^{t\cdot y_i-|t|^2/2}$ is at most one
for every $t$ once $\kappa\ge(M^2-1)_+/2$, proving nonemptiness. Lemma
\ref{lem:variance-monotonicity} makes the set an upper interval. Finally, if
$\eps_m$ belongs to the set and $\eps_m\to\eps<1$, then the one-atom
criterion holds at $\eps$ by pointwise convergence of its field for every
$t$. The interval is therefore closed in $[0,1)$ and contains its infimum.
\end{proof}

For the spatial analysis, it is useful to rewrite \(D_n\) as a Gaussian
radial-basis field whose centers and heights are determined by the sample.
This representation will allow the radial and angular contributions to be
studied separately.

\begin{lemma}
\label{lem:bump-representation}
Set
\[
  w=\sqrt{\gamma_n}\,t,
  \qquad
  q_i=\frac{\widetilde Z_i}{\sqrt{\gamma_n}},
  \qquad
  h_i=\frac1n\exp\left\{\frac{|q_i|^2}{2}\right\}.
\]
Define the bump-coordinate field
\[
  B_n(w):=D_n(w/\sqrt{\gamma_n}).
\]
Then, exactly,
\begin{equation}
\label{eq:bump-representation}
  B_n(w)
  =
  \sum_{i=1}^n h_i
  \exp\left\{-\frac{|w-q_i|^2}{2}\right\}.
\end{equation}
In particular, the \(i\)th bump has maximum \(h_i\), attained at \(w=q_i\),
and
\begin{equation}
\label{eq:one-point-threshold}
  h_i>1
  \quad\Longleftrightarrow\quad
  \frac{|\widetilde Z_i|^2}{2}>\gamma_n\log n.
\end{equation}
\end{lemma}

\begin{proof}
Completing the square gives
\[
  t\cdot\widetilde Z_i-\frac{\gamma_n|t|^2}{2}
  =
  w\cdot q_i-\frac{|w|^2}{2}
  =
  \frac{|q_i|^2}{2}-\frac{|w-q_i|^2}{2}.
\]
Summing proves \eqref{eq:bump-representation}; maximizing one summand proves
\eqref{eq:one-point-threshold}.
\end{proof}

The bump representation isolates the basic mechanism. A radial extreme creates a
localized bump whose height is close to one. Other observations contribute a
diffuse background, and one-atomicity fails precisely when the full bump field
exceeds one. At the transition the centers capable of producing a crossing have
radius approximately \(\sqrt{2\log n}\) in \(w\)-coordinates, even when the
original dimension is much larger than \(\log n\).

There is also an exact dual formulation which will be decisive in
Section~\ref{sec:explicit-regime}. Let
\(\Delta_n=\{p\in[0,1]^n:\sum_i p_i=1\}\), let
\(\mathbf u_n=(1/n,\ldots,1/n)\), and write
\[
  K(p\Vert \mathbf u_n)=\sum_{i=1}^n p_i\log(np_i)
\]
for relative entropy, with \(0\log0=0\).

\begin{lemma}
\label{lem:entropy-dual}
For every centered collection \(q_1,\ldots,q_n\in\R^d\),
\begin{equation}
\label{eq:entropy-dual}
  \log\sup_{w\in\R^d}B_n(w)
  =
  \sup_{p\in\Delta_n}
  \left\{
    \frac12\left|\sum_{i=1}^np_iq_i\right|^2
    -K(p\Vert \mathbf u_n)
  \right\}.
\end{equation}
The expression on the right is always nonnegative, since it vanishes at
\(p=\mathbf u_n\).
\end{lemma}

\begin{proof}
The Gibbs variational formula gives, for every \(w\),
\[
  \log\left\{\frac1n\sum_i e^{w\cdot q_i}\right\}
  =
  \sup_{p\in\Delta_n}
  \left\{w\cdot\sum_i p_iq_i-K(p\Vert \mathbf u_n)\right\}.
\]
Subtract \(|w|^2/2\), take the supremum over \(w\), and use
\(\sup_w\{w\cdot v-|w|^2/2\}=|v|^2/2\). Finally,
\(\sum_iq_i=0\), so \(p=\mathbf u_n\) gives zero.
\end{proof}

The point-mass choice \(p=e_i\) in \eqref{eq:entropy-dual} gives the
one-bump obstruction \eqref{eq:one-point-threshold}.  A slightly diffuse
competitor gives an exact and strictly stronger finite-sample threshold.

\begin{proposition}
\label{prop:potts-radial-obstruction}
For \(n\ge3\), define
\begin{equation}
\label{eq:potts-radial-threshold}
 R_n^\star
 =
 \frac{(n-1)^2}{n(n-2)}\log(n-1)
 =
 \log n-\frac1n+O\left(\frac{\log n}{n^2}\right).
\end{equation}
If
\[
 \max_i\frac{|q_i|^2}{2}>R_n^\star,
\]
then \(\sup_wB_n(w)>1\). Consequently, \(E_n^{(1)}\) implies
\(\max_i|q_i|^2/2\le R_n^\star\).
\end{proposition}

\begin{proof}
Fix \(i\), and define \(p\in\Delta_n\) by
\[
 p_i=\frac{n-1}{n},
 \qquad
 p_j=\frac1{n(n-1)}\quad(j\ne i).
\]
Since \(\sum_jq_j=0\),
\[
 \sum_jp_jq_j=\frac{n-2}{n-1}q_i,
 \qquad
 K(p\Vert \mathbf u_n)=\frac{n-2}{n}\log(n-1).
\]
The expression in braces in \eqref{eq:entropy-dual} is therefore positive
exactly when
\[
 \frac{|q_i|^2}{2}>
 \frac{(n-1)^2}{n(n-2)}\log(n-1).
\]
The expansion in \eqref{eq:potts-radial-threshold} follows from Taylor
expansion of \(\log(n-1)\).
\end{proof}

\subsection{Gamma extremes and empirical centering}\label{sec:mv-gamma-normalization}

The normalization in Theorem~\ref{thm:main-support-law} is chosen so that the
largest centered radial energies have an order-one Poisson limit. This
subsection proves that limit uniformly over arbitrary dimension sequences and
quantifies the effect of replacing \(Z_i\) by \(Z_i-\bar Z_n\).

Put
\[
  L_n=\log n,
  \qquad
  k_n=\frac{d_n}{2},
\]
and let \(Y_n\) have the \(\Gamma(k_n,1)\) distribution. Write \(f_n\) and
\(\overline F_n\) for its density and upper tail. Define \(b_n\) and \(a_n\)
by
\begin{equation}
\label{eq:b-a-definition}
  n\overline F_n(b_n)=1,
  \qquad
  a_n=\frac{\overline F_n(b_n)}{f_n(b_n)}.
\end{equation}
Thus \(b_n\) is the exact radial \(1-1/n\) quantile and \(a_n\) is the
reciprocal hazard there. Unlike the fixed-dimensional expansion of \(b_n\),
these definitions remain valid without specifying a relation between \(d_n\)
and \(n\).

We first verify that \(a_n\) is the correct extreme-value scale for every
triangular array of Gamma shapes. The proof is included because uniformity in
\(k_n\) is essential here.

\begin{lemma}
\label{lem:gamma-hazard}
For any sequence \(k_n\ge1\), the quantile \(b_n\) in
\eqref{eq:b-a-definition} satisfies
\[
  \frac{b_n-k_n}{\sqrt{k_n}}\longrightarrow\infty.
\]
Moreover,
\begin{equation}
\label{eq:hazard-asymptotic}
  a_n
  =
  \frac{b_n}{b_n-k_n+1}\{1+o(1)\},
\end{equation}
and, uniformly for \(x\) in compact subsets of \(\R\),
\begin{equation}
\label{eq:gamma-tail-ratio}
  \frac{\overline F_n(b_n+a_nx)}{\overline F_n(b_n)}
  \longrightarrow e^{-x}.
\end{equation}
\end{lemma}

\begin{proof}
We may argue along subsequences. If \(k_n\) remains bounded, then
\(b_n\to\infty\), and the first assertion is immediate. If \(k_n\to\infty\),
the central limit theorem gives
\[
  \Pbb\{Y_n>k_n+M\sqrt{k_n}\}\longrightarrow1-\Phi(M)
\]
for every fixed \(M\), where \(\Phi\) is the standard normal distribution
function. Since \(\overline F_n(b_n)=1/n\to0\), eventually
\(b_n>k_n+M\sqrt{k_n}\). Letting \(M\to\infty\) proves the first claim.

For \(y>k_n-1\), changing variables \(z=y+u\) in the Gamma tail gives the
exact identity
\begin{equation}
\label{eq:tail-density-ratio}
  \frac{\overline F_n(y)}{f_n(y)}
  =
  \int_0^\infty
  \exp\left\{-u+(k_n-1)\log\left(1+\frac uy\right)\right\}\,du.
\end{equation}
Set
\(\delta_y=1-(k_n-1)/y=(y-k_n+1)/y\). The inequality
\(\log(1+v)\le v\) shows that the last integral is at most
\(\delta_y^{-1}\). After the substitution \(v=\delta_yu\), its exponent is
\[
  -v-\frac{k_n-1}{2y^2\delta_y^2}v^2
  +O\left(\frac{k_n}{y^3\delta_y^3}v^3\right)
\]
on every bounded \(v\)-interval. At \(y=b_n\),
\[
  \frac{k_n}{b_n^2\delta_{b_n}^2}
  =
  \frac{k_n}{(b_n-k_n+1)^2}
  \longrightarrow0.
\]
Dominated convergence in \eqref{eq:tail-density-ratio}, using the upper
envelope \(e^{-v}\), therefore yields
\(a_n\delta_{b_n}\to1\), which is \eqref{eq:hazard-asymptotic}.

If \(y=b_n+O(a_n)\), then
\[
  \frac{\delta_y}{\delta_{b_n}}=1+o(1),
  \qquad
  \frac{\overline F_n(y)}{f_n(y)}=a_n\{1+o(1)\},
\]
uniformly over the indicated range. The first relation follows by
differentiating \(\delta_y\) and using
\(k_n/(b_n-k_n+1)^2\to0\); the second follows because the dominated-convergence
argument applied to \eqref{eq:tail-density-ratio} is uniform for
\(y=b_n+O(a_n)\). Since
\((\log\overline F_n)'(y)=-f_n(y)/\overline F_n(y)\), integration from
\(b_n\) to \(b_n+a_nx\) gives \eqref{eq:gamma-tail-ratio}.
\end{proof}

The triangular-array maximum therefore has the usual Gumbel law. No
dimension-dependent centering expansion is needed.

\begin{proposition}
\label{prop:gamma-max}
Let \(Y_{n,1},\ldots,Y_{n,n}\) be independent
\(\Gamma(k_n,1)\) variables. Then, for every \(x\in\R\),
\begin{equation}
\label{eq:gumbel-unscaled}
  \Pbb\left\{
    \max_{1\le i\le n}\frac{Y_{n,i}-b_n}{a_n}\le x
  \right\}
  \longrightarrow
  \exp\{-e^{-x}\}.
\end{equation}
More generally, the number of indices satisfying
\(Y_{n,i}>b_n+a_nx\) converges to a Poisson variable with mean \(e^{-x}\).
\end{proposition}

\begin{proof}
By Lemma~\ref{lem:gamma-hazard},
\(n\overline F_n(b_n+a_nx)\to e^{-x}\). The exceedance count is binomial
with success probability asymptotic to \(e^{-x}/n\), proving the Poisson
limit. Its probability of being zero gives \eqref{eq:gumbel-unscaled}.
\end{proof}

Throughout the radial analysis, put
\begin{equation}
\label{eq:residual-variance}
 \sigma_n^2=1-\frac1n.
\end{equation}
The exact transition coordinate and the associated centered radial excesses
are
\begin{equation}
\label{eq:exact-transition-coordinate}
 s_n^\star
 =\frac{\gamma_nR_n^\star/\sigma_n^2-b_n}{a_n},
 \qquad
 \xi_{n,i}^\star
 =\frac{|Z_i-\bar Z_n|^2/(2\sigma_n^2)-b_n}{a_n}-s_n^\star,
 \qquad
 \delta_n^\star=\frac{\sigma_n^2a_n}{\gamma_n}.
\end{equation}
Thus, for \(q_i=(Z_i-\bar Z_n)/\sqrt{\gamma_n}\),
\begin{equation}
\label{eq:exact-support-radius}
 \frac{|q_i|^2}{2}=R_n^\star+\delta_n^\star\xi_{n,i}^\star.
\end{equation}
These exact quantities will be used in every theorem statement. Some proofs
temporarily replace \(R_n^\star\) by \(\log n\); the comparison is recorded
once in Lemma~\ref{lem:exact-old-edge-coordinate}.

The exact Gamma normalization agrees with the moving-shell normalization in
every fixed dimension. The following bridge is used when the bounded- and
growing-dimensional arguments are combined.
\begin{lemma}
\label{lem:fixed-d-normalization-bridge}
Fix \(d\ge3\), put \(k=d/2\), and define
\[
 b_{n,d}=\log n+(k-1)\log\log n-\log\Gamma(k).
\]
Then \(a_n\to1\), \(b_n-b_{n,d}\to0\), and
\(R_n^\star/\sigma_n^2=\log n+o(1)\). If \(s_n^\star=O(1)\), then
\begin{equation}
\label{eq:fixed-d-coordinate-bridge}
 s_n^\star
 =\frac{\gamma_n\log n-b_n}{a_n}+o(1)
 =\frac{\eps_n}{1-\eps_n}\log n
   -(k-1)\log\log n+\log\Gamma(k)+o(1).
\end{equation}
Moreover, uniformly over \(1\le i\le n\),
\begin{equation}
\label{eq:fixed-d-excess-bridge}
 \frac{|Z_i-\bar Z_n|^2/(2\sigma_n^2)-b_n}{a_n}
 =\frac{|Z_i|^2}{2}-b_{n,d}+o_{\Pbb}(1).
\end{equation}
\end{lemma}

\begin{proof}
The fixed-shape Gamma tail expansion gives
\[
 \overline F_n(y)=\frac{y^{k-1}e^{-y}}{\Gamma(k)}\{1+O(y^{-1})\},
 \qquad y\to\infty.
\]
Substitution in \(n\overline F_n(b_n)=1\) yields
\(b_n=b_{n,d}+o(1)\), and Lemma~\ref{lem:gamma-hazard} gives
\(a_n\to1\). The expansion of \(R_n^\star\) in
\eqref{eq:potts-radial-threshold} gives
\(R_n^\star/\sigma_n^2=\log n+o(1)\). Since \(s_n^\star=O(1)\) forces
\(\gamma_n=O(1)\), the first equality in
\eqref{eq:fixed-d-coordinate-bridge} follows. The second uses
\(\gamma_n-1=\eps_n/(1-\eps_n)\).

Finally, \(\max_i|Z_i|=O_{\Pbb}(\sqrt{\log n})\) and
\(|\bar Z_n|=O_{\Pbb}(n^{-1/2})\). Hence
\[
 \max_i\bigl||Z_i-\bar Z_n|^2-|Z_i|^2\bigr|
 =O_{\Pbb}\left(\sqrt{\frac{\log n}{n}}\right)=o_{\Pbb}(1),
\]
while multiplication by \(\sigma_n^{-2}=1+O(n^{-1})\) changes the maximum
energy by \(O_{\Pbb}(\log n/n)=o_{\Pbb}(1)\). Together with
\(a_n\to1\) and \(b_n-b_{n,d}\to0\), this proves
\eqref{eq:fixed-d-excess-bridge}.
\end{proof}

The Karush--Kuhn--Tucker (KKT) field uses centered observations. The next
condition is sufficient
for empirical centering to be negligible on the Gamma extreme scale:
\begin{equation}
\label{eq:centering-condition}
  \frac{\sqrt{b_nd_n/n}}{a_n}\longrightarrow0.
\end{equation}
It is automatic throughout \(d_n=o(\log n)\).  In high dimension the direct
norm bound behind this condition is wasteful; the leave-one-out estimate
below reaches \(d_n\log n=o(n^2)\).

\begin{lemma}
\label{lem:centered-radial-max}
Assume \eqref{eq:centering-condition}. Then
\begin{equation}
\label{eq:centered-uniform-energy-equivalence}
  \max_i
  \left|
    \frac{|Z_i-\bar Z_n|^2-|Z_i|^2}{2a_n}
  \right|
  \longrightarrow_{\Pbb}0.
\end{equation}
In particular,
\begin{equation}
\label{eq:centered-max-equivalence}
  \max_i
  \frac{|Z_i-\bar Z_n|^2/2-b_n}{a_n}
  -
  \max_i
  \frac{|Z_i|^2/2-b_n}{a_n}
  \longrightarrow_{\Pbb}0.
\end{equation}
Consequently, the centered maximum has the Gumbel limit
\eqref{eq:gumbel-unscaled}.
\end{lemma}

\begin{proof}
Proposition~\ref{prop:gamma-max} and tightness imply
\(\max_i|Z_i|^2/2=b_n+O_{\Pbb}(a_n)\). Lemma~\ref{lem:gamma-hazard} also
implies \(a_n=O(b_n)\), so \(\max_i|Z_i|=O_{\Pbb}(\sqrt{b_n})\). Since
\(|\bar Z_n|=O_{\Pbb}(\sqrt{d_n/n})\),
\[
  \max_i
  \left|
    \frac{|Z_i-\bar Z_n|^2-|Z_i|^2}{2a_n}
  \right|
  \le
  \frac{\max_i|Z_i|\,|\bar Z_n|+|\bar Z_n|^2/2}{a_n}
  =o_{\Pbb}(1)
\]
by \eqref{eq:centering-condition}. This proves
\eqref{eq:centered-uniform-energy-equivalence}; taking maxima gives
\eqref{eq:centered-max-equivalence}.
\end{proof}

The norm bound in Lemma~\ref{lem:centered-radial-max} controls
\(Z_i\cdot\bar Z_n\) by the product of two norms and is unnecessarily
wasteful when \(d_n\) is large. A
leave-one-out decomposition retains the cancellation in this inner product.

\begin{lemma}
\label{lem:centered-radial-max-high-d}
Uniformly over \(d_n\ge1\),
\begin{equation}
\label{eq:centered-energy-leave-one-out}
 \max_{1\le i\le n}
 \left|
  \frac{|Z_i-\bar Z_n|^2-|Z_i|^2}{2}
 \right|
 =
 O_{\Pbb}\left\{
  \frac{d_n}{n}
  +\sqrt{\frac{d_nL_n}{n}}
  +\frac{L_n}{\sqrt n}
 \right\}.
\end{equation}
Consequently, if \(L_n=o(d_n)\) and \(d_nL_n=o(n^2)\), then
\eqref{eq:centered-max-equivalence} holds.
\end{lemma}

\begin{proof}
Write
\[
 Z_i\cdot\bar Z_n
 =
 \frac{|Z_i|^2}{n}
 +Z_i\cdot\bar Z_{-i},
 \qquad
 \bar Z_{-i}=\frac1n\sum_{j\ne i}Z_j.
\]
The two vectors in the second inner product are independent, and
\(\bar Z_{-i}\sim\mathcal N(0,(n-1)n^{-2}I_{d_n})\).  Its moment generating
function therefore gives the subexponential bound
\[
 \max_i|Z_i\cdot\bar Z_{-i}|
 =
 O_{\Pbb}\left\{
  \sqrt{\frac{d_nL_n}{n}}+\frac{L_n}{\sqrt n}
 \right\}.
\]
The standard chi-square tail bound and a union bound also give
\[
 \max_i|Z_i|^2
 =
 O_{\Pbb}\{d_n+\sqrt{d_nL_n}+L_n\},
 \qquad
 |\bar Z_n|^2=O_{\Pbb}\{(d_n+1)/n\}.
\]
Substituting these estimates into
\[
 \frac{|Z_i-\bar Z_n|^2-|Z_i|^2}{2}
 =-Z_i\cdot\bar Z_n+\frac{|\bar Z_n|^2}{2}
\]
proves \eqref{eq:centered-energy-leave-one-out}.

In the stated range,
\(a_n\sim\sqrt{d_n/(4L_n)}\).  Dividing the three terms in
\eqref{eq:centered-energy-leave-one-out} by \(a_n\) gives, respectively,
\[
 O\left(\frac{\sqrt{d_nL_n}}n\right),\qquad
 O\left(\frac{L_n}{\sqrt n}\right),\qquad
 O\left(\frac{L_n^{3/2}}{\sqrt{d_nn}}\right),
\]
all of which tend to zero.  This proves
\eqref{eq:centered-max-equivalence}.
\end{proof}

The two norm-based centering lemmas do not use the exact marginal
distribution of a centered residual. In fact,
\[
 \frac{Z_i-\bar Z_n}{\sigma_n}\sim\mathcal N(0,I_{d_n})
\]
for every \(i\).  The residuals are dependent, but their correlations are
weak enough that the extreme exceedance counts still have Poisson limits.
The next proposition uses this exact marginal law to reach dimensions beyond
the range of Lemma~\ref{lem:centered-radial-max-high-d}.

\begin{proposition}
\label{prop:direct-centered-gamma-max}
Suppose
\[
 L_n=o(d_n).
\]
For every fixed \(x\in\mathbb R\),
\begin{equation}
\label{eq:direct-centered-gumbel}
 \Pbb\left\{
  \max_i
  \frac{|Z_i-\bar Z_n|^2/(2\sigma_n^2)-b_n}{a_n}
  \le x
 \right\}
 \longrightarrow e^{-e^{-x}}.
\end{equation}
More generally, the number of indices for which
\[
 \frac{|Z_i-\bar Z_n|^2}{2\sigma_n^2}>b_n+a_nx
\]
converges to \(\operatorname{Pois}(e^{-x})\).
\end{proposition}

\begin{proof}
Fix \(m\ge1\). We prove, uniformly for \(x\) in compact subsets of
\(\mathbb R\), that
\begin{equation}
\label{eq:centered-joint-tail-factorization}
 \Pbb\{Y_1>y_n,\ldots,Y_m>y_n\}
 =
 \{1+o(1)\}\Pbb\{Y_1>y_n\}^m,
 \qquad y_n=b_n+a_nx,
\end{equation}
where
\[
 V_i=\frac{Z_i-\bar Z_n}{\sigma_n},
 \qquad
 Y_i=\frac{|V_i|^2}{2}.
\]
The proposition will then follow by the method of factorial moments.

We first record the Gamma estimates needed below. Put \(k=k_n=d_n/2\),
\(L=L_n=\log n\), and
\[
 h(\delta)=\delta-\log(1+\delta),\qquad \delta>-1.
\]
If \(Y\sim\Gamma(k,1)\) and \(y=k(1+\delta)\), where
\(\delta\to0\) and \(k\delta^2\to\infty\), then
\begin{equation}
\label{eq:centered-gamma-moderate-tail}
 \Pbb\{Y>y\}
 =
 \frac{1+o(1)}{\delta\sqrt{2\pi k}}
 \exp\{-kh(\delta)\}.
\end{equation}
Writing \(f_k\) for the density of \(Y\), Stirling's formula gives
\[
 f_k\{k(1+\delta)\}
 =
 \frac{1+o(1)}{(1+\delta)\sqrt{2\pi k}}
 \exp\{-kh(\delta)\}.
\]
Moreover, the exact identity
\[
 \frac{\Pbb\{Y>y\}}{f_k(y)}
 =
 \int_0^\infty
 \exp\left\{-v+(k-1)\log\left(1+\frac{v}{y}\right)\right\}\,dv
\]
shows that
\begin{equation}
\label{eq:centered-gamma-mills}
 \frac{\Pbb\{Y>y\}}{f_k(y)}
 =
 \frac{y}{y-k+1}\{1+o(1)\}
 =
 \frac{1+\delta}{\delta}\{1+o(1)\}.
\end{equation}
For completeness, set \(\eta=(y-k+1)/y\) and substitute \(s=\eta v\)
in the integral. Its exponent becomes
\[
 -s-\frac{k-1}{2y^2\eta^2}s^2
 +O\left(\frac{k}{y^3\eta^3}s^3\right)
\]
on bounded \(s\)-intervals, while
\(\log(1+v/y)\le v/y\) supplies the dominating function \(e^{-s}\).
Since \(k/(y^2\eta^2)\sim(k\delta^2)^{-1}\to0\), dominated convergence
proves \eqref{eq:centered-gamma-mills}, and hence
\eqref{eq:centered-gamma-moderate-tail}.

Let
\[
 \delta_n=\frac{b_n-k}{k}.
\]
The central limit theorem and \(\Pbb\{Y>b_n\}=n^{-1}\) imply
\(\delta_n\sqrt{k}\to\infty\). The Chernoff bound
\[
 \Pbb\{Y>k(1+\delta)\}\le e^{-kh(\delta)}
\]
gives \(kh(\delta_n)\le L=o(k)\), and therefore \(\delta_n\to0\).
Applying \eqref{eq:centered-gamma-moderate-tail} at \(b_n\) yields
\[
 L
 =
 kh(\delta_n)+\log(\delta_n\sqrt{2\pi k})+o(1).
\]
Because \(h(\delta)\sim\delta^2/2\), this implies
\begin{equation}
\label{eq:centered-gamma-local-scales}
 \delta_n\sim\sqrt{\frac{2L}{k}},
 \qquad
 b_n-k\sim\sqrt{2kL},
 \qquad
 a_n\sim\frac1{\delta_n}\sim\sqrt{\frac{k}{2L}}.
\end{equation}
If \(y_n=b_n+a_nx=k(1+\delta_{n,x})\), then, uniformly for bounded \(x\),
\[
 \frac{\delta_{n,x}}{\delta_n}=1+o(1).
\]
Using \(h'(\delta)=\delta/(1+\delta)\),
\eqref{eq:centered-gamma-moderate-tail}, and
\(a_n\delta_n/(1+\delta_n)\to1\), we obtain
\[
 k\{h(\delta_{n,x})-h(\delta_n)\}
 =
 x+O\left(\frac{a_n^2}{k}\right)+o(1)
 =
 x+o(1).
\]
Consequently,
\begin{equation}
\label{eq:centered-gamma-local-tail-ratio}
 p_n(x):=\Pbb\{Y_i>y_n\}
 =
 \frac{e^{-x}+o(1)}{n}.
\end{equation}

We next compare the joint law of \(Y_1,\ldots,Y_m\) with \(m\)
independent Gamma variables. For each Euclidean coordinate, the Gaussian
vector \((V_1,\ldots,V_m)\) has covariance
\[
 A=A_{n,m}
 =
 \frac n{n-1}I_m-\frac1{n-1}\mathbf1\mathbf1^\top
 =
 I_m+C,
\]
where
\[
 C_{ii}=0,
 \qquad
 C_{ij}=-\frac1{n-1}\quad(i\ne j),
 \qquad
 \|C\|_{\mathrm{op}}+\|C\|_{\mathrm F}=O_m(n^{-1}).
\]
Set
\[
 t=t_n=1-\frac{k}{y_n},
 \qquad
 \beta=\beta_n=1-t=\frac{k}{y_n}.
\]
By \eqref{eq:centered-gamma-local-scales},
\begin{equation}
\label{eq:centered-tilt-scale}
 t\to0,
 \qquad
 t\sqrt{d_n}=(2+o(1))\sqrt{L_n},
 \qquad
 d_nt^2=(4+o(1))L_n.
\end{equation}

Let \(M_{n,m}^{\mathrm{dep}}(t)\) be the joint moment-generating function
of \((Y_1,\ldots,Y_m)\) at the common argument \(t\). The Gaussian
quadratic-form formula gives
\[
 M_{n,m}^{\mathrm{dep}}(t)
 =
 \det(I_m-tA)^{-d_n/2}.
\]
For \(m\) independent copies,
\[
 M_{n,m}^{\mathrm{ind}}(t)=\beta^{-md_n/2}.
\]
Since \(\operatorname{tr}C=0\),
\begin{align}
\label{eq:centered-tilt-normalizer}
 \log\frac{M_{n,m}^{\mathrm{dep}}(t)}
               {M_{n,m}^{\mathrm{ind}}(t)}
 &=
 -\frac{d_n}{2}
 \log\det\left(I_m-\frac{t}{\beta}C\right) \notag\\
 &=
 O_m\left(\frac{d_nt^2}{n^2}\right)
 =
 O_m\left(\frac{L_n}{n^2}\right)
 =
 o(1).
\end{align}

Under the dependent exponential tilt, each coordinate vector is Gaussian
with covariance
\[
 \Sigma^{\mathrm{dep}}
 =
 (A^{-1}-tI_m)^{-1}
 =
 A(I_m-tA)^{-1}.
\]
Under the independent tilt the corresponding covariance is
\[
 \Sigma^{\mathrm{ind}}=\beta^{-1}I_m.
\]
We write \(\E_t^{\mathrm{dep}}\) and \(\E_t^{\mathrm{ind}}\) for expectation
under these two tilted laws, and use \(\E_t^\star\) when the same formula
applies to either one.
Writing \(f_t(a)=a/(1-ta)\) and expanding \(f_t(I_m+C)\) at \(I_m\)
gives
\[
 \Sigma^{\mathrm{dep}}
 =
 \beta^{-1}I_m+\beta^{-2}C+R_{n,m},
 \qquad
 \|R_{n,m}\|_{\mathrm{op}}
 \le C_m t\|C\|_{\mathrm{op}}^2.
\]
In particular,
\begin{equation}
\label{eq:centered-tilted-covariance}
 (\Sigma^{\mathrm{dep}})_{ii}
 =
 \beta^{-1}+O_m(tn^{-2}),
 \qquad
 (\Sigma^{\mathrm{dep}})_{ij}
 =
 -\frac{\beta^{-2}}{n-1}+O_m(tn^{-2})
 \quad(i\ne j).
\end{equation}
The eigenvalues of both tilted covariance matrices therefore lie in a
fixed compact subset of \((0,\infty)\).

We now prove the local limit estimate needed for the tilted overshoots.
Let \(f_{n,m}^{\mathrm{dep}}\) and \(f_{n,m}^{\mathrm{ind}}\) denote the
respective densities of
\[
 W_n
 =
 \frac1{\sqrt{d_n}}
 (Y_1-y_n,\ldots,Y_m-y_n)
\]
under the dependent and independent tilted laws. We claim that
\begin{equation}
\label{eq:centered-tilted-local-limit}
 \sup_{w\in\mathbb R^m}
 \left|f_{n,m}^{\star}(w)-\varphi_m(w)\right|
 \longrightarrow0,
 \qquad
 \varphi_m(w)=\pi^{-m/2}e^{-|w|^2},
\end{equation}
for \(\star\in\{\mathrm{dep},\mathrm{ind}\}\).

To prove the claim, let \(\Sigma^\star\) be the appropriate tilted
covariance, let \(D_z=\operatorname{diag}(z_1,\ldots,z_m)\), and define
\[
 B_z^\star=(\Sigma^\star)^{1/2}D_z(\Sigma^\star)^{1/2}.
\]
The characteristic function of \(W_n\) is
\begin{equation}
\label{eq:centered-tilted-characteristic-function}
 \chi_{n,m}^{\star}(z)
 =
 \exp\left\{-\frac{iy_n}{\sqrt{d_n}}\sum_{j=1}^mz_j\right\}
 \det\left(
 I_m-\frac{i}{\sqrt{d_n}}\Sigma^\star D_z
 \right)^{-d_n/2}.
\end{equation}
For fixed \(z\), expansion of the logarithm gives
\begin{align*}
 \log\chi_{n,m}^{\star}(z)
 &=
 \frac{i\sqrt{d_n}}2
 \sum_{j=1}^m
 \left\{(\Sigma^\star)_{jj}-\beta^{-1}\right\}z_j\\
 &\quad
 -\frac14
 \operatorname{tr}(\Sigma^\star D_z\Sigma^\star D_z)
 +O_m\left(\frac{|z|^3}{\sqrt{d_n}}\right).
\end{align*}
The linear term vanishes identically in the independent case and is
\(O_m(t\sqrt{d_n}|z|/n^2)=o(1)\) in the dependent case by
\eqref{eq:centered-tilt-scale} and
\eqref{eq:centered-tilted-covariance}. Since
\(\Sigma^\star\to I_m\), it follows that
\[
 \chi_{n,m}^{\star}(z)\longrightarrow e^{-|z|^2/4}
\]
for every fixed \(z\).

We next give an integrable domination uniform in the two tilted laws.
Since \(B_z^\star\) is real symmetric,
\begin{equation}
\label{eq:centered-characteristic-modulus}
 |\chi_{n,m}^{\star}(z)|
 =
 \det\left(
 I_m+\frac{(B_z^\star)^2}{d_n}
 \right)^{-d_n/4}.
\end{equation}
If \(\mu_1,\ldots,\mu_m\) are the eigenvalues of \(B_z^\star\), the
uniform lower eigenvalue bound for \(\Sigma^\star\) gives
\[
 \sum_{j=1}^m\mu_j^2
 =
 \|B_z^\star\|_{\mathrm F}^2
 \ge c_m|z|^2.
\]
Hence
\begin{equation}
\label{eq:centered-characteristic-global-bound}
 |\chi_{n,m}^{\star}(z)|
 \le
 \left(1+\frac{c_m|z|^2}{d_n}\right)^{-d_n/4}.
\end{equation}
Choose \(c_0>0\) sufficiently small. If
\(|z|\le c_0\sqrt{d_n}\), then all
\(\mu_j^2/d_n\le1\); using \(\log(1+v)\ge v/2\) in
\eqref{eq:centered-characteristic-modulus} yields
\begin{equation}
\label{eq:centered-characteristic-gaussian-bound}
 |\chi_{n,m}^{\star}(z)|\le e^{-c_m'|z|^2}.
\end{equation}
On the complementary region, polar coordinates and
\eqref{eq:centered-characteristic-global-bound} give
\begin{align*}
 &\int_{|z|>c_0\sqrt{d_n}}
 \left(1+\frac{c_m|z|^2}{d_n}\right)^{-d_n/4}\,dz\\
 &\qquad\le
 C_m d_n^{m/2}(1+c_m'')^{-d_n/8}
 \mathrm B\left(\frac m2,\frac{d_n}{8}-\frac m2\right)
 =o(1).
\end{align*}
Here \(\mathrm B\) denotes the Euler beta function.
Thus the characteristic functions are uniformly integrable. Pointwise
convergence, \eqref{eq:centered-characteristic-gaussian-bound}, and the
last tail estimate imply
\[
 \int_{\mathbb R^m}
 \left|\chi_{n,m}^{\star}(z)-e^{-|z|^2/4}\right|\,dz
 \longrightarrow0.
\]
Fourier inversion proves \eqref{eq:centered-tilted-local-limit}. Notice
that this argument treats the dependent and independent tilted densities
separately; both converge uniformly to the same density \(\varphi_m\), and
in particular
\[
 f_{n,m}^{\mathrm{dep}}(0),
 \ f_{n,m}^{\mathrm{ind}}(0)
 \longrightarrow
 \varphi_m(0)=\pi^{-m/2}.
\]

It remains to compare the tilted overshoot integrals. For
\(\star\in\{\mathrm{dep},\mathrm{ind}\}\), put
\[
 J_{n,m}^{\star}
 =
 \E_t^\star\left[
  \exp\left\{-t\sum_{j=1}^m(Y_j-y_n)\right\}
  \mathbf1_{\{Y_j>y_n,\ 1\le j\le m\}}
 \right].
\]
Writing \(q_n=t\sqrt{d_n}\), the density of \(W_n\) gives
\[
 q_n^mJ_{n,m}^{\star}
 =
 \int_{\mathbb R_+^m}
 e^{-\sum_jz_j}
 f_{n,m}^{\star}(z/q_n)\,dz.
\]
By \eqref{eq:centered-tilt-scale}, \(q_n\to\infty\). The uniform local
limit theorem supplies a common bound on the densities, so dominated
convergence yields
\[
 q_n^mJ_{n,m}^{\star}\longrightarrow\varphi_m(0).
\]
Therefore
\[
 \frac{J_{n,m}^{\mathrm{dep}}}{J_{n,m}^{\mathrm{ind}}}\longrightarrow1.
\]
The exponential-tilting identities are
\[
 \Pbb\{Y_1>y_n,\ldots,Y_m>y_n\}
 =
 M_{n,m}^{\mathrm{dep}}(t)e^{-tmy_n}J_{n,m}^{\mathrm{dep}}
\]
and
\[
 p_n(x)^m
 =
 M_{n,m}^{\mathrm{ind}}(t)e^{-tmy_n}J_{n,m}^{\mathrm{ind}}.
\]
Together with \eqref{eq:centered-tilt-normalizer}, they prove
\eqref{eq:centered-joint-tail-factorization}.

Finally, let
\[
 N_n(x)=\sum_{i=1}^n\mathbf1_{\{Y_i>b_n+a_nx\}}.
\]
For an integer-valued variable \(N\), write
\(\{N\}_m=N(N-1)\cdots(N-m+1)\), and similarly
\((n)_m=n(n-1)\cdots(n-m+1)\).
For every fixed \(m\), exchangeability,
\eqref{eq:centered-joint-tail-factorization}, and
\eqref{eq:centered-gamma-local-tail-ratio} give
\[
 \E\{N_n(x)\}_m
 =
 (n)_m\Pbb\{Y_1>y_n,\ldots,Y_m>y_n\}
 \longrightarrow e^{-mx}
\]
for every fixed \(m\). These are the factorial moments of
\(\operatorname{Pois}(e^{-x})\), so
\(N_n(x)\Rightarrow\operatorname{Pois}(e^{-x})\). Taking the probability
that \(N_n(x)=0\) proves \eqref{eq:direct-centered-gumbel}.
\end{proof}


\section{The Two-Dimensional Sharp Window}
\label{sec:d2}
\label{sec:d-two-sharp-window}

This section establishes the field estimates for the two-dimensional clauses
of Theorems~\ref{thm:main-support-law} and
\ref{thm:regularization-path}. At the critical scale, the KKT field near its
maximizing direction contains both a single extreme radial spike and a
nonvanishing contribution from observations just inside the edge. We identify
their limiting profile, confine every possible crossing to a compact radial
window, and prove a uniform decomposition used for support recovery in
Section~\ref{sec:support}.

Throughout this section let \(Z_1,\dots,Z_n\) be i.i.d.\ \(\cN(0,I_2)\), let
\(\rho_n=\sqrt{2\log n}\), set \(\bar Z_n=n^{-1}\sum_iZ_i\), and put
\[
    H_n(t)=\frac1n\sum_{i=1}^n\exp\{t\cdot Z_i-|t|^2/2\},
    \qquad
    D_n(t)=e^{-t\cdot\bar Z_n}e^{-\kappa_n|t|^2}H_n(t),
    \qquad
    \lambda_n=\kappa_n\rho_n^2 .
\]
The one-atom event is \(E_n^{(1)}=\{\sup_tD_n(t)\le1\}\).
We write \(\Phi\) and \(\phi\) for the standard normal distribution
function and density.

\subsection{The limiting profile}

The KKT supremum chooses its direction after seeing the sample. A fixed
angular interval of width \(1/\rho_n\) contains a point with bounded radial
excess with probability \(O(1/\rho_n)\), but the maximizing interval can be
centered on such a point after the sample is observed. This selected extreme
creates one random spike. The many observations whose radial energies lie
just below \(\log n\) produce the deterministic background.

Write
\[
    Z_i=R_iU_i,\qquad U_i=e^{i\Theta_i}\in S^1,\qquad
    \xi_i=\frac{R_i^2}{2}-\log n ,
\]
and let \(\operatorname{ang}(\phi-\psi)\in(-\pi,\pi]\) denote the principal
signed angular displacement from direction \(\psi\) to direction \(\phi\). We
also write \(\operatorname{ang}(u,v)\) for the same displacement from \(v\) to
\(u\), and set \(Y_i(v)=\rho_n\operatorname{ang}(U_i,v)\). The global radial edge process
\[
    \Xi_n=\sum_{i=1}^n\delta_{(\xi_i,\Theta_i)}
\]
converges to a Poisson point process \(\Xi=\sum_j\delta_{(\xi_j,\Theta_j)}\)
on \(\R\times\mathbb T\), where \(\mathbb T=\R/2\pi\mathbb Z\), with intensity
\[
    e^{-x}\,dx\,\frac{d\phi}{2\pi}.
\]
For a fixed deterministic direction \(\psi\), however, the locally rescaled
process
\[
    \sum_i
    \delta_{\left(\xi_i,\ \rho_n\operatorname{ang}(\Theta_i-\psi)\right)}
\]
has intensity
\[
    \frac{1+o(1)}{2\pi\rho_n}\,e^{-x}\,dx\,dy
\]
on compact \(x,y\)-sets, and hence converges to the empty process. The
nontrivial local limit is obtained only after conditioning on, and recentering
around, a radial edge point.

For \(t=(\rho_n+s)e^{i\psi}\), one observation contributes
\[
    W_{n,i}(\psi,s)
    =
    \frac1n
    \exp\left\{(\rho_n+s)e^{i\psi}\cdot Z_i
        -\frac{(\rho_n+s)^2}{2}\right\}.
\]
If \(\xi_i\), \(s\), and
\[
    y_i=\rho_n\operatorname{ang}(\Theta_i-\psi)
\]
are bounded, then
\[
    W_{n,i}(\psi,s)
    =
    \exp\{\xi_i-s^2/2-y_i^2/2+o(1)\}.
\]
Thus a radial edge point \((\xi_j,\Theta_j)\), viewed in the angular coordinate
\(\psi=\Theta_j+y/\rho_n\), creates the local spike
\[
    e^{\xi_j-s^2/2-y^2/2}.
\]
The spike has \(O(1)\) width in blown-up coordinates, hence
\(O(1/\rho_n)\) width on the original circle.

For fixed \(A,C<\infty\), the \(O(1)\)-radial edge spikes are separated:
\[
    \Pbb\left\{
    \exists i\ne j:\ \xi_i,\xi_j\ge-A,\
    \rho_n|\operatorname{ang}(\Theta_i-\Theta_j)|\le C
    \right\}\to0 ,
\]
because the expected number of such pairs is
\[
    O\left\{
        \left(\int_{-A}^{\infty}e^{-x}\,dx\right)^2
        \frac C{\rho_n}
    \right\}=o(1).
\]
Consequently, after the uniform compactness estimates are proved, at most one
radial extreme contributes at order one near any maximizing
direction. All other observations form the deterministic background. The
limiting event is therefore determined by the largest radial excess rather
than by a nontrivial random process on the entire circle.

The following calculation identifies the background. Let
\(t=(\rho_n+s)v\), \(v\in S^1\), and truncate observations at the
one-dimensional projection edge in direction \(v\):
\[
    H_n^{\mathrm{diff}}(v,s)
    =
    \frac1n\sum_{i=1}^n
    \exp\left\{(\rho_n+s)v\cdot Z_i-\frac{(\rho_n+s)^2}{2}\right\}
    \mathbf 1_{\{v\cdot Z_i\le\rho_n\}}
\]
satisfies
\[
    \E H_n^{\mathrm{diff}}(v,s)=\Phi(-s),
    \qquad
    \operatorname{Var}H_n^{\mathrm{diff}}(v,s)
    \le
    \frac1n e^{(\rho_n+s)^2}\Phi(-\rho_n-2s)
    \sim
    \frac{e^{-s^2}}{\sqrt{2\pi}\rho_n}
\]
uniformly for \(s\) in compact sets.  Thus the non-extreme cloud contributes \(\Phi(-s)\) to \(H_n\), or \(e^{-\lambda}\Phi(-s)\) after the damping factor in \(D_n\).

A radial edge point with
\[
    \xi_i=\frac{|Z_i|^2}{2}-\log n
\]
and direction \(U_i\) contributes, at \(t=(\rho_n+s)U_i\),
\[
    e^{-\lambda}\exp\{\xi_i-s^2/2+o(1)\}.
\]
Combining the two contributions, a radial excess \(x\) can produce a crossing
precisely when
\[
    g(x):=\sup_{s\in\R}\{\Phi(-s)+e^{x-s^2/2}\}
\]
exceeds the damped threshold \(e^\lambda\). This gives the one-atom limit
\[
    Q_2(\lambda)=\exp\{-e^{-x_\lambda}\},
\]
where \(x_\lambda\) is uniquely determined by
\[
    \sup_{s\in\R}\{\Phi(-s)+e^{x_\lambda-s^2/2}\}=e^\lambda .
\]
Equivalently, if \(a_\lambda>0\) solves
\[
    \Phi(a_\lambda)+\frac{\phi(a_\lambda)}{a_\lambda}=e^\lambda,
\]
then
\[
    x_\lambda=-\log(\sqrt{2\pi}a_\lambda),
    \qquad
    Q_2(\lambda)=\exp\{-\sqrt{2\pi}a_\lambda\}.
\]
The optimizer is \(s=-a_\lambda\), obtained from
\[
    -\phi(s)-s e^{x-s^2/2}=0.
\]
The map \(a\mapsto \Phi(a)+\phi(a)/a\) is strictly decreasing, since its derivative is \(-\phi(a)/a^2\), so \(a_\lambda\) is unique.

The optimized profile crosses every level above one exactly once. Its strict
monotonicity later converts a fixed radial gap into a fixed KKT gap.
\begin{lemma}
\label{lem:q-two-curve}
Let
\[
    g(x)=\sup_{s\in\R}\{\Phi(-s)+e^{x-s^2/2}\}.
\]
Then \(g\) is finite, continuous, and strictly increasing, with
\(\lim_{x\to-\infty}g(x)=1\) and \(\lim_{x\to\infty}g(x)=\infty\).  Hence,
for every \(\lambda>0\), there is a unique \(x_\lambda\) such that
\(g(x_\lambda)=e^\lambda\), and for every \(\eta>0\),
\[
    g(x_\lambda-\eta)<e^\lambda<g(x_\lambda+\eta).
\]
\end{lemma}

\begin{proof}
For every finite \(x\), the objective exceeds \(1\) at some finite negative
\(s\). Indeed, with \(s=-a\), Mills' ratio gives
\(1-\Phi(a)\sim\phi(a)/a\), whereas
\(e^{x-a^2/2}=\sqrt{2\pi}e^x\phi(a)\); the second term dominates the first
for sufficiently large \(a\). On a compact \(x\)-interval, the objective
converges uniformly to \(1\) as \(s\to-\infty\) and to \(0\) as
\(s\to+\infty\). Its supremum is therefore attained in a common compact
\(s\)-interval. Continuity follows from compact maximization, and strict
monotonicity follows because \(e^{x-s^2/2}\) is strictly increasing in \(x\)
at every finite maximizer. Finally,
\(1\le g(x)\le1+e^x\) implies \(g(x)\to1\) as \(x\to-\infty\), while
\(g(x)\ge e^x\to\infty\) as \(x\to\infty\).
\end{proof}

The explicit parametrization also gives the endpoint behavior
\[
    1-Q_2(\lambda)\sim e^{-\lambda}\quad(\lambda\to\infty),
    \qquad
    \log Q_2(\lambda)
    =
    -2\sqrt{\pi\log(1/\lambda)}\{1+o(1)\}
    \quad(\lambda\downarrow0).
\]
Indeed, as \(\lambda\to\infty\), the defining equation gives
\(a_\lambda\sim\phi(0)e^{-\lambda}\). Since
\(e^{-x_\lambda}=\sqrt{2\pi}a_\lambda\), this proves the first relation.
As \(\lambda\downarrow0\), the two-term Mills expansion gives
\[
 e^\lambda-1
 =\frac{\phi(a_\lambda)}{a_\lambda^3}
   \{1+O(a_\lambda^{-2})\}.
\]
Thus \(a_\lambda^2\sim2\log(1/\lambda)\), which proves the second relation.
For large \(\lambda\), the diffuse background is negligible and the radial
Gumbel tail determines the limit. By contrast,
\(Q_2(\lambda)\downarrow0\) as the shrinkage tends to zero.

\subsection{Radial compactness}

Radial compactness reduces all possible crossings to
\(|t|=\rho_n+O(1)\), without yet identifying the limiting field.

\begin{proposition}
\label{prop:d-two-radial-compactness}
Assume \(d=2\) and \(\lambda_n=\kappa_n\rho_n^2\to\lambda\in(0,\infty)\),
where \(\rho_n=\sqrt{2\log n}\).  Then
\[
    \lim_{A\to\infty}\limsup_{n\to\infty}
    \Pbb\left\{
        \sup_{\left||t|-\rho_n\right|>A}D_n(t)>1
    \right\}=0 .
\]
Thus, up to a probability which vanishes after \(n\to\infty\) and then
\(A\to\infty\), every crossing has the form
\[
    t=(\rho_n+s)v,\qquad |s|\le A,\quad v\in S^1 .
\]
\end{proposition}

Write
\[
    H_n^{\mathrm{diff}}(r,v)=\frac1n\sum_i e^{rv\cdot Z_i-r^2/2}
        \mathbf 1_{\{v\cdot Z_i\le \rho_n\}},
    \qquad
    T_n(r,v)=\frac1n\sum_i e^{rv\cdot Z_i-r^2/2}
        \mathbf 1_{\{v\cdot Z_i> \rho_n\}} .
\]
Thus \(H_n(rv)=H_n^{\mathrm{diff}}(r,v)+T_n(r,v)\).  The split is at the
one-dimensional edge in direction \(v\).

The diffuse part cannot create a crossing below the compact window or at its
upper boundary.
\begin{lemma}
\label{lem:d-two-diffuse-outside}
For every \(\delta>0\),
\[
\begin{aligned}
    &\lim_{A\to\infty}\limsup_{n\to\infty}
    \Pbb\left\{
        \sup_{\rho_n/2\le r\le \rho_n-A,\ v\in S^1}
        H_n^{\mathrm{diff}}(r,v)>1+\delta
    \right\}=0,\\
    &\lim_{A\to\infty}\limsup_{n\to\infty}
    \Pbb\left\{
        \sup_{v\in S^1}H_n^{\mathrm{diff}}(\rho_n+A,v)>\delta
    \right\}=0 .
\end{aligned}
\]
\end{lemma}

\begin{proof}
Put
\[
    \xi_i=\frac{|Z_i|^2}{2}-\log n,\qquad \ell_n^{\mathrm{far}}=8\log\rho_n,
\]
and let
\[
    \mathcal R_n(A)=\{r:\rho_n/2\le r\le \rho_n-A\}\cup\{\rho_n+A\}.
\]
Completing the square gives the radial identity
\begin{equation}
\label{eq:d2-diffuse-radial-identity}
    \frac1n\exp\{rv\cdot z-r^2/2\}
    =
    \exp\left\{x(z)-\frac{|z-rv|^2}{2}\right\},
    \qquad
    x(z)=\frac{|z|^2}{2}-\log n .
\end{equation}

First remove the near-edge layer. Let \(H_n^{\mathrm{diff},0}\) be
\(H_n^{\mathrm{diff}}\) with the additional restriction
\(\xi_i\le -\ell_n^{\mathrm{far}}\). Every summand retained in
\(H_n^{\mathrm{diff},0}\) is at most
\(e^{-\ell_n^{\mathrm{far}}}\) after normalization. Hence the unnormalized summand is bounded
by \(ne^{-\ell_n^{\mathrm{far}}}\), and its second moment is at most \(ne^{-\ell_n^{\mathrm{far}}}\) times its
first moment.  Since
\[
    \E H_n^{\mathrm{diff},0}(r,v)
    \le \E H_n^{\mathrm{diff}}(r,v)=\Phi(\rho_n-r)\le1
    \qquad (r\in\mathcal R_n(A)),
\]
Bernstein's inequality gives, at each fixed \((r,v)\),
\[
    \Pbb\{|H_n^{\mathrm{diff},0}(r,v)
                 -\E H_n^{\mathrm{diff},0}(r,v)|>\eta\}
    \le 2\exp\{-c_\eta e^{\ell_n^{\mathrm{far}}}\}.
\]
Let \(\mathcal F_n\) be the class of normalized summands
\[
    z\mapsto \frac1n e^{rv\cdot z-r^2/2}
        \mathbf 1_{\{v\cdot z\le\rho_n\}}
        \mathbf 1_{\{x(z)\le -\ell_n^{\mathrm{far}}\}},
    \qquad r\in\mathcal R_n(A),\ v\in S^1,
\]
so that \(H_n^{\mathrm{diff},0}(r,v)\) is the sum of \(n\) functions from
\(\mathcal F_n\). This is a VC-subgraph class of fixed dimension: after taking
logarithms in the subgraph coordinate, its subgraphs are intersections of a
fixed ball, a halfspace in \(z\), and a halfspace in \((z,\log y)\).
Standard closure properties therefore give polynomial covering numbers,
uniformly in \(n\); see
\cite[Sections~2.6.2 and~2.14.1]{vanderVaartWellner1996}. In the present
bounded-envelope setting, the Bernstein maximal inequality from that entropy
bound has the form
\[
 \Pbb\left\{\sup_{f\in\mathcal F_n}
   \left|\sum_{i=1}^n\{f(Z_i)-\E f(Z_i)\}\right|>\eta\right\}
 \le C\rho_n^C
 \exp\left\{-\frac{c\eta^2}
 {e^{-\ell_n^{\mathrm{far}}}+\eta e^{-\ell_n^{\mathrm{far}}}}\right\}.
\]
Moreover, identity~\eqref{eq:d2-diffuse-radial-identity}
bounds every member of \(\mathcal F_n\) by \(e^{-\ell_n^{\mathrm{far}}}\), and
the sum of its variances by \(e^{-\ell_n^{\mathrm{far}}}\). The Bernstein maximal inequality for
VC-subgraph classes now upgrades the pointwise estimate to the whole class.
Its entropy cost is \(O(\log\rho_n)\), whereas its deviation exponent is of
order \(e^{\ell_n^{\mathrm{far}}}=\rho_n^8\). Hence
\[
    \sup_{r\in\mathcal R_n(A),\,v\in S^1}
    |H_n^{\mathrm{diff},0}(r,v)-\E H_n^{\mathrm{diff},0}(r,v)|
    \xrightarrow{\Pbb}0 .
\]

It remains to control the removed layer.  Write \(Z_i=R_iU_i\) and
\(Y_i(v)=\rho_n\operatorname{ang}(U_i,v)\).  If \(\xi_i>0\) and \(v\cdot Z_i\le\rho_n\),
then
\[
    |Z_i-rv|^2
    \ge |Z_i|^2+r^2-2r\rho_n
    =2\xi_i+(r-\rho_n)^2,
\]
so every such positive-edge term is bounded by \(e^{-A^2/2}\), uniformly
over \(r\in\mathcal R_n(A)\).  Thus their total contribution is at most
\(e^{-A^2/2}N_n^+\), where \(N_n^+=\#\{i:\xi_i>0\}=O_{\Pbb}(1)\).

For \(-\ell_n^{\mathrm{far}}<\xi_i\le0\), \(R_i=\rho_n+O(\ell_n^{\mathrm{far}}/\rho_n)\).  Uniformly over
\(r\in\mathcal R_n(A)\), for all large \(n\),
\[
    |Z_i-rv|^2
    \ge cA^2+cY_i(v)^2 .
\]
Therefore the negative part of the removed layer is bounded by
\[
    e^{-cA^2}
    \sup_{v\in S^1}
    \sum_{-\ell_n^{\mathrm{far}}<\xi_i\le0} e^{\xi_i}e^{-cY_i(v)^2}.
\]
The last supremum is tight.  Indeed, partition \((-\ell_n^{\mathrm{far}},0]\) into the slabs
\(-k-1<\xi_i\le-k\), \(0\le k\le\lceil \ell_n^{\mathrm{far}}\rceil\), and cover \(S^1\) by
\(O(\rho_n)\) arcs of length \(O(\rho_n^{-1})\).  If \(N_{k\ell}\) denotes
the number of points in the \(k\)-th slab and in the \(\ell\)-th enlarged
arc, then the Gaussian angular kernel has summable tails and
\[
    \sup_v\sum_{-\ell_n^{\mathrm{far}}<\xi_i\le0} e^{\xi_i}e^{-cY_i(v)^2}
    \le
    C\sum_{k\le\lceil \ell_n^{\mathrm{far}}\rceil} e^{-k}\max_\ell N_{k\ell}.
\]
Here \(\E N_{k\ell}\le C e^k/\rho_n\).  Chernoff's bound and a union bound
over the \(O(\rho_n)\) arcs give, with probability tending to one,
\[
    \max_\ell N_{k\ell}
    \le C\left(\frac{e^k}{\rho_n}+k+1\right)
    \qquad\text{for all }k\le\lceil \ell_n^{\mathrm{far}}\rceil .
\]
Consequently the displayed shot-noise supremum is \(O_{\Pbb}(1)\).  We have
proved
\[
    \sup_{r\in\mathcal R_n(A),\,v\in S^1}
    \{H_n^{\mathrm{diff}}(r,v)-H_n^{\mathrm{diff},0}(r,v)\}
    \le e^{-A^2/2}O_{\Pbb}(1)+e^{-cA^2}O_{\Pbb}(1).
\]

On the lower side, \(\E H_n^{\mathrm{diff},0}(r,v)\le1\), so the first
assertion follows
from the concentration estimate and then \(A\to\infty\).  At the upper
boundary,
\(\E H_n^{\mathrm{diff},0}(\rho_n+A,v)
\le\E H_n^{\mathrm{diff}}(\rho_n+A,v)=\Phi(-A)\), and
\(\Phi(-A)\to0\). The concentration and removed-layer bounds above therefore
prove the second assertion as well.
\end{proof}

The tail part \(T_n\) is likewise negligible away from
\(|r-\rho_n|=O(1)\).
\begin{lemma}
\label{lem:d-two-tail-outside}
For every \(\delta>0\),
\[
    \lim_{A\to\infty}\limsup_{n\to\infty}
    \Pbb\left\{
        \sup_{\substack{r\ge\rho_n/2,\ v\in S^1\\ |r-\rho_n|\ge A}}
        T_n(r,v)>\delta
    \right\}=0 .
\]
\end{lemma}

\begin{proof}
Put \(\xi_i=|Z_i|^2/2-\log n\) and \(N_n^+=\#\{i:\xi_i>0\}\).  Since
\(|Z_i|^2/2\) is exponential in dimension two, \(N_n^+\Rightarrow
\operatorname{Pois}(1)\) and
\[
    \Pbb\{\max_i\xi_i>A^2/4\}\to 1-\exp\{-e^{-A^2/4}\}.
\]
Hence the event
\[
    \max_i\xi_i\le A^2/4,\qquad N_n^+\le e^{A^2/8}
\]
has probability tending to one as first \(n\to\infty\) and then
\(A\to\infty\).  If \(v\cdot Z_i>\rho_n\), then \(\xi_i>0\).  On the
displayed event, write \(R_i=|Z_i|\).  For \(r\le\rho_n-A\),
\[
    r\,v\cdot Z_i-\frac{r^2}{2}-\log n
    \le
    (\rho_n-A)R_i-\frac{(\rho_n-A)^2}{2}-\log n
    \le
    \xi_i-\frac{A^2}{2}+o(1).
\]
For \(r\ge\rho_n+A\), the maximum over \(r\) is attained either at
\(r=\rho_n+A\) or at \(r=v\cdot Z_i\).  Since \(\xi_i\le A^2/4\) implies
\(R_i=\rho_n+O(A^2/\rho_n)<\rho_n+A\) for large \(n\), the boundary case
applies and gives
\[
    r\,v\cdot Z_i-\frac{r^2}{2}-\log n
    \le
    (\rho_n+A)R_i-\frac{(\rho_n+A)^2}{2}-\log n
    \le
    \xi_i-\frac{A^2}{2}+o(1)
    \le -A^2/4+o(1).
\]
Therefore
\[
    \sup_{\substack{r\ge\rho_n/2,\ v\in S^1\\ |r-\rho_n|\ge A}}
    T_n(r,v)
    \le
    N_n^+e^{-A^2/4+o(1)}
    \le e^{-A^2/8+o(1)}
\]
on the displayed event, which proves the lemma.
\end{proof}

An ordinary bulk estimate removes the remaining range
\(0<|t|\le\rho_n/2\).
\begin{lemma}
\label{lem:d2-bulk-exclusion}
Assume \(d=2\) and \(\lambda_n=\kappa_n\rho_n^2\to\lambda\in(0,\infty)\).
Then
\[
    \Pbb\left\{\sup_{0<|t|\le\rho_n/2}D_n(t)>1\right\}\to0 .
\]
\end{lemma}

\begin{proof}
We repeat the three estimates behind
Proposition~\ref{prop:bulk-mgf-exclusion}, proved for fixed \(d\ge3\) in
Section~\ref{sec:sharp-high-dimensional-estimates}, and verify them at the
two-dimensional scale
\(\kappa_n\asymp\rho_n^{-2}\asymp1/\log n\).

First consider \(0<|t|\le c_0\kappa_n\). The sample covariance satisfies
\[
    \left\|\frac1n\sum_i
    (Z_i-\bar Z_n)(Z_i-\bar Z_n)^\top-I_2\right\|_{\mathrm{op}}
    =O_{\Pbb}(n^{-1/2})
    =o_{\Pbb}(\kappa_n).
\]
The Hessian argument in Lemma~\ref{lem:microscopic-bulk} therefore applies
unchanged: for a sufficiently small fixed \(c_0>0\), with probability tending
to one,
\[
    D_n(t)\le1-c\kappa_n|t|^2<1
    \qquad(0<|t|\le c_0\kappa_n).
\]

For \(c_0\kappa_n\le |t|\le1\), use the quadratic empirical-process
expansion from Lemma~\ref{lem:bulk-exponential-process}. The only scale
conditions needed there are
\[
    \frac{n^{-1/4}}{\kappa_n}\to0,
    \qquad
    \frac{|\bar Z_n|}{\kappa_n^2}
    =
    O_{\Pbb}(n^{-1/2}\rho_n^4)\to0 .
\]
They imply, uniformly on this range,
\[
    \frac{|H_n(t)-1|+|t\cdot\bar Z_n|}
         {1-e^{-\kappa_n|t|^2}}
    =o_{\Pbb}(1).
\]

For \(1\le |t|\le\rho_n/2\), repeat the dyadic decomposition and clipping
argument from Lemma~\ref{lem:bulk-exponential-process}. Explicitly, partition the squared radii into
dyadic intervals \([x_j,x_{j+1}]\), put
\[
 R_j=x_{j+1}^{1/2},\qquad
 \tau_n^2=20\log(1/\kappa_n)+100\log\rho_n,qquad
 B_j=\exp\{R_j^2/2+R_j\tau_n\},
\]
and clip each likelihood ratio at \(B_j\). Here
\(|t|^2\le(\log n)/2\), so the variance and envelope bounds are
\(e^{|t|^2}\le n^{1/2}\) and \(B_j\le n^{1/4+o(1)}\). The net entropy is
still \(O(\log n)\), while the
Bernstein exponent is at least
\[
    n^{1/2-o(1)}\kappa_n^2
    \gg \log n .
\]
Thus there is a deterministic \(\delta_n\downarrow0\) such that, with
probability tending to one,
\[
    |H_n(t)-1|+|t\cdot\bar Z_n|
    \le\delta_n g(t),
    \qquad
    g(t)=1-e^{-\kappa_n|t|^2},
\]
uniformly on \(c_0\kappa_n\le|t|\le\rho_n/2\). (In particular,
\(\rho_n|\bar Z_n|=o_{\Pbb}(1)\).) The algebra in the proof of
Proposition~\ref{prop:bulk-mgf-exclusion} then gives
\[
    D_n(t)
    \le e^{\delta_ng(t)}\{1-g(t)\}\{1+\delta_ng(t)\}
    <1
\]
uniformly on that range for all large \(n\). Together with the microscopic
estimate, this proves the lemma.
\end{proof}

\begin{proof}[Proof of Proposition~\ref{prop:d-two-radial-compactness}]
Lemma~\ref{lem:d2-bulk-exclusion} controls the bulk range:
\[
    \Pbb\left\{\sup_{0<|t|\le\rho_n/2}D_n(t)>1\right\}\to0 .
\]

For the upper exterior, use the centered representation
\[
    D_n(rv)=\frac1n\sum_i
    \exp\left\{rv\cdot (Z_i-\bar Z_n)-\frac{1+2\kappa_n}{2}r^2\right\}.
\]
With probability tending to one, \(\max_i|Z_i-\bar Z_n|<\rho_n+A/2\).
On this event every summand is decreasing in \(r\) for \(r\ge\rho_n+A\),
so the upper exterior reduces to the boundary \(r=\rho_n+A\).  Since
\(\rho_n|\bar Z_n|=o_{\Pbb}(1)\),
\[
    \sup_{v}D_n((\rho_n+A)v)
    \le
    e^{-\lambda+o_{\Pbb}(1)}
    \sup_v\{H_n^{\mathrm{diff}}(\rho_n+A,v)+T_n(\rho_n+A,v)\}.
\]
The upper-bound parts of Lemmas~\ref{lem:d-two-diffuse-outside} and
\ref{lem:d-two-tail-outside} make the right side smaller than one with
probability tending to one after \(A\to\infty\).

It remains to consider \(\rho_n/2\le r\le\rho_n-A\).  On this range
\[
    e^{-\kappa_nr^2}\le e^{-\lambda/4+o(1)},
    \qquad
    \sup_{r,v}e^{-rv\cdot\bar Z_n}=e^{o_{\Pbb}(1)} .
\]
Lemma~\ref{lem:d-two-diffuse-outside} gives
\(H_n^{\mathrm{diff}}(r,v)\le1+o_{\Pbb}(1)+o_A(1)\) uniformly, and
Lemma~\ref{lem:d-two-tail-outside} gives
\(T_n(r,v)=o_{\Pbb}(1)+o_A(1)\) uniformly.  Since \(e^{-\lambda/4}<1\),
the lower exterior also cannot contain a crossing once \(A\) is large.
\end{proof}

\subsection{Compact-window decomposition}

We now approximate the full field on the compact radial window. Observations
below a slowly receding radial cutoff form the
deterministic background; those above it produce narrow angular kernels.
\begin{proposition}
\label{prop:two-d-compact-decomposition}
Assume \(d=2\), put \(\rho_n=\sqrt{2\log n}\), and let \(K\subset\R\) be compact.  Let \(\operatorname{ang}(u,v)\in(-\pi,\pi]\) denote the principal signed angular displacement from \(v\) to \(u\).  For \(Z_i\ne0\), write \(Z_i=R_iU_i\), and set
\[
    \xi_i=\frac{R_i^2}{2}-\log n,
    \qquad
    Y_i(v)=\rho_n\,\operatorname{ang}(U_i,v),
\]
so that \(Y_i(v)\) is the blown-up signed angular displacement from \(v\) to
the observation direction \(U_i\). If \(I\subset S^1\) is a closed arc and
\(\ell_n^{\mathrm{edge}}=4\log\rho_n\), then
\[
    \sup_{v\in I,\ s\in K}
    \left|
    H_n((\rho_n+s)v)-\Phi(-s)
    -
    \sum_{i:\xi_i\ge -\ell_n^{\mathrm{edge}}}
    \exp\left\{\xi_i-\frac{s^2}{2}-\frac{Y_i(v)^2}{2}\right\}
    \right|
    \xrightarrow{\Pbb}0 .
\]
\end{proposition}

\begin{proof}
Write \(r_s=\rho_n+s\).  First separate the diffuse contribution below the
radial edge.  For \(L>0\), define
\[
    H^{\rm diff}_{n,L}(v,s)
    =
    \frac1n\sum_{i=1}^n
    \exp\left\{r_sv\cdot Z_i-\frac{r_s^2}{2}\right\}
    \mathbf 1_{\{\xi_i<-L\}} .
\]
We use \(L=\ell_n^{\mathrm{edge}}\).  If \(\xi_i<-L\), then \(R_i\le R_L=(\rho_n^2-2L)^{1/2}\), and uniformly for \(s\in K\),
\[
    \frac1n\exp\left\{r_sv\cdot Z_i-\frac{r_s^2}{2}\right\}
    \le
    \exp\left\{-L-\frac{(r_s-R_L)^2}{2}\right\}
    \le C_K e^{-L}.
\]
The sum of the variances is therefore at most \(C_Ke^{-L}\).  Bernstein's inequality gives, for each fixed \((v,s)\),
\[
    \Pbb\{|H^{\rm diff}_{n,L}(v,s)-\E H^{\rm diff}_{n,L}(v,s)|>\eta\}
    \le
    2\exp\{-c_{K,\eta}e^L\}.
\]
The field \(H^{\rm diff}_{n,L}\) is smooth in \((v,s)\). Each summand has
logarithmic derivatives bounded by \(C_K\rho_n^2\) in the angular coordinate
and by \(C_K\rho_n\) in \(s\). Fix \(\eta>0\), and choose an angular mesh of
order \(\rho_n^{-2}\) and an \(s\)-mesh of order \(\rho_n^{-1}\). The mesh
constants can be chosen so that neighboring values of both
\(H^{\rm diff}_{n,L}\) and \(\E H^{\rm diff}_{n,L}\) differ by a factor at
most \(1+\eta\). The grid has polynomial size in \(\rho_n\), so a union bound
upgrades the pointwise Bernstein estimate to the grid. Because
\(\E H^{\rm diff}_{n,L}\le1\), interpolation from the grid costs only an
additive \(O(\eta)\); letting \(\eta\downarrow0\) gives
\[
    \sup_{v\in I,\ s\in K}
    |H^{\rm diff}_{n,\ell_n^{\mathrm{edge}}}(v,s)-\E H^{\rm diff}_{n,\ell_n^{\mathrm{edge}}}(v,s)|
    \xrightarrow{\Pbb}0,
\]
because \(e^{\ell_n^{\mathrm{edge}}}=\rho_n^4\) dominates the grid entropy.

It remains to compute the mean.  By exponential tilting,
\[
    \E H^{\rm diff}_{n,L}(v,s)
    =
    \Pbb\{|G+r_sv|^2\le \rho_n^2-2L\},
    \qquad G\sim\cN(0,I_2).
\]
After rotating \(v=e_1\), this probability is
\[
    \Pbb\left\{(r_s+G_1)^2+G_2^2\le \rho_n^2-2L\right\}.
\]
Condition on \(G_2\).  On the event \(|G_2|\le\rho_n^{1/3}\), the boundary for \(G_1\) is
\[
    G_1
    \le
    -s-\frac{\ell_n^{\mathrm{edge}}+G_2^2/2}{\rho_n}
    +O_K\left(\frac{(\ell_n^{\mathrm{edge}}+G_2^2+1)^2}{\rho_n^3}\right),
\]
uniformly for \(s\in K\). The quadratic inequality also has the lower
boundary
\[
 G_1\ge-r_s-\sqrt{\rho_n^2-2\ell_n^{\mathrm{edge}}-G_2^2}
       =-2\rho_n+O_K(1)
\]
on this event, whose Gaussian probability is uniformly \(o(1)\). The
complement of \(\{|G_2|\le\rho_n^{1/3}\}\) also has probability \(o(1)\),
and the displayed upper boundary converges to \(-s\) uniformly on \(K\)
after conditioning on \(G_2\). Dominated convergence and the Lipschitz
continuity of the Gaussian distribution function give
\[
    \sup_{v\in I,\ s\in K}
    |\E H^{\rm diff}_{n,\ell_n^{\mathrm{edge}}}(v,s)-\Phi(-s)|\to0 .
\]
Thus the contribution from \(\{\xi_i<-\ell_n^{\mathrm{edge}}\}\) is \(\Phi(-s)+o_{\Pbb}(1)\), uniformly on \(I\times K\).

Now consider the edge contribution
\[
    A^{\rm ex}_n(v,s)
    =
    \sum_{i:\xi_i\ge -\ell_n^{\mathrm{edge}}}
    \frac1n
    \exp\left\{r_sv\cdot Z_i-\frac{r_s^2}{2}\right\}.
\]
We compare it with the displayed atom sum. Let
\[
    C_n=\log\rho_n,
    \qquad
    Y_n=(\log\rho_n)^{1/4}.
\]
The event \(\max_i\xi_i\le C_n\) has probability tending to one, since
\(n\Pbb\{\xi_i>C_n\}=e^{-C_n}\). Also
\[
    \sum_{i:-\ell_n^{\mathrm{edge}}\le \xi_i\le C_n} e^{\xi_i}
    =
    O_{\Pbb}(\ell_n^{\mathrm{edge}}+C_n),
\]
because \(\xi_i\) has shifted exponential intensity \(e^{-x}\,dx\) in dimension two.

On the event \(\max_i\xi_i\le C_n\), the exact contribution of an edge observation
with \(-\ell_n^{\mathrm{edge}}\le \xi_i\le C_n\) satisfies
\[
    \frac1n
    \exp\left\{r_sv\cdot Z_i-\frac{r_s^2}{2}\right\}
    \le
    C_K e^{\xi_i}
    \exp\{-c\min(Y_i(v)^2,\rho_n^2)\}.
\]
This follows from
\[
    r_sv\cdot Z_i-\frac{r_s^2}{2}-\log n
    =
    \xi_i-\frac{(R_i-r_s)^2}{2}-r_sR_i\{1-\cos\operatorname{ang}(U_i,v)\},
\]
and \(r_sR_i\asymp\rho_n^2\) in the displayed range.  Hence the total
contribution of terms with \(|Y_i(v)|>Y_n\), in either \(A^{\rm ex}_n\) or
the Gaussian atom sum, is
\[
    O_{\Pbb}(\ell_n^{\mathrm{edge}}+C_n)e^{-cY_n^2}=o_{\Pbb}(1)
\]
uniformly in \(v\in I\) and \(s\in K\), by the definition of \(Y_n\).
For the remaining terms
\(-\ell_n^{\mathrm{edge}}\le \xi_i\le C_n\) and \(|Y_i(v)|\le Y_n\), Taylor expansion gives,
uniformly in \(i,v,s\),
\[
\begin{aligned}
    &\log\left[
    \frac1n\exp\left\{r_sv\cdot Z_i-\frac{r_s^2}{2}\right\}
    \right]
    -\left\{\xi_i-\frac{s^2}{2}-\frac{Y_i(v)^2}{2}\right\} \\
    &\qquad=
    O_K\left(\frac{\ell_n^{\mathrm{edge}}+C_n}{\rho_n}
    +\frac{Y_n^2}{\rho_n}
    +\frac{Y_n^4}{\rho_n^2}\right)
    =O_K\left(\frac{\log\rho_n}{\rho_n}\right).
\end{aligned}
\]
The total atom weight is \(O_{\Pbb}(\log\rho_n)\), while the relative error is
\(O(\log\rho_n/\rho_n)\). Their product is \(o_{\Pbb}(1)\), and hence
\[
    \sup_{v\in I,\ s\in K}
    \left|
    A^{\rm ex}_n(v,s)
    -
    \sum_{i:\xi_i\ge -\ell_n^{\mathrm{edge}}}
    \exp\left\{\xi_i-\frac{s^2}{2}-\frac{Y_i(v)^2}{2}\right\}
    \right|
    \xrightarrow{\Pbb}0 .
\]
Combining this with \(H_n=H^{\rm diff}_{n,\ell_n^{\mathrm{edge}}}+A^{\rm ex}_n\) proves the proposition.
\end{proof}

Proposition~\ref{prop:two-d-compact-decomposition} uses the slowly receding
cutoff \(\ell_n^{\mathrm{edge}}=4\log\rho_n\), which makes the background concentrate uniformly.
For the final supremum bound, we replace it by a fixed cutoff. For every fixed
\(A<\infty\), the
edge points with \(\xi_i>-A\) are finite in the limit and have no
\(1/\rho_n\)-scale
angular collisions:
\[
    \Pbb\left\{
    \exists i\ne j:\ \xi_i,\xi_j>-A,\
    \rho_n|\operatorname{ang}(U_i,U_j)|\le L
    \right\}\le C_A L/\rho_n\to0 .
\]

Edge points in the lower radial strip \(-\ell_n^{\mathrm{edge}}\le \xi_i\le-A\) have negligible
total contribution after \(A\to\infty\), allowing a fixed cutoff.
\begin{lemma}
\label{lem:d2-low-atom-absorption}
Let \(\ell_n^{\mathrm{edge}}=4\log\rho_n\), and write
\(Y_i(v)=\rho_n\operatorname{ang}(U_i,v)\) for the principal angular
displacement between \(U_i\) and \(v\).  For every
\(\delta>0\),
\[
    \lim_{A\to\infty}\limsup_{n\to\infty}
    \Pbb\left\{
    \sup_{v\in S^1}
    \sum_{i:-\ell_n^{\mathrm{edge}}\le \xi_i\le -A}
    e^{\xi_i-Y_i(v)^2/2}>\delta
    \right\}=0 .
\]
\end{lemma}

\begin{proof}
Put \(m_n^\star=\lceil \ell_n^{\mathrm{edge}}\rceil\) and, for integers \(m\ge0\),
\[
    I_{m,n}=\{i:-m-1<\xi_i\le -m\}.
\]
These are unit radial-excess layers; on \(I_{m,n}\), the weight \(e^{\xi_i}\)
is of order \(e^{-m}\).  Since
\(|Z_i|^2/2\) is exponential in dimension two,
\[
    \Pbb\{i\in I_{m,n}\}\le C e^m/n,
    \qquad m\le m_n^\star,
\]
and the angle \(\Theta_i\) is independent and uniform on the circle.

For \(q=0,1,\ldots,\lfloor\pi\rho_n\rfloor\), let
\[
    N_{m,q,n}
    =
    \sup_{\phi\in\mathbb T}
    \#\left\{
        i\in I_{m,n}:
        \rho_n|\operatorname{ang}(\Theta_i-\phi)|\le q+1
    \right\}.
\]
Covering the circle by \(O(\rho_n)\) arcs shows that every interval appearing
in the supremum is contained in one of the covering arcs, with length enlarged
by only a fixed factor.  For one such enlarged arc the count is binomial with
mean
\[
    \mu_{m,q,n}\le C(q+1)e^m/\rho_n .
\]
Set
\[
    a_{m,q,n}
    =
    1+\frac{(q+1)e^m}{\rho_n}
    +
    \frac{\log\rho_n}
        {1+\log_+\!\left(\rho_n/((q+1)e^m)\right)},
    \qquad
    \log_+ x=\max\{\log x,0\}.
\]
The binomial Chernoff bound
\(\Pbb\{\operatorname{Bin}(N,p)\ge r\}\le (eNp/r)^r\) for
\(r\ge eNp\), together with the multiplicative Chernoff bound when
\(r\) is a fixed multiple of \(Np\), implies that for a sufficiently large
numerical constant \(B\),
\[
    \Pbb\left\{
        N_{m,q,n}>B a_{m,q,n}
        \text{ for some } A\le m\le m_n^\star,\ 0\le q\le \pi\rho_n
    \right\}\to0
\]
for each fixed \(A\).  Indeed, if \(\mu_{m,q,n}<1\), the last term in
\(a_{m,q,n}\) makes
\[
    a_{m,q,n}\log\{a_{m,q,n}/\mu_{m,q,n}\}\ge c\log\rho_n,
\]
while if \(\mu_{m,q,n}\ge1\), the term \(\log\rho_n\) gives the same
union-bound margin unless \(\mu_{m,q,n}\ge\log\rho_n\), in which case the
linear mean term does.  Choosing \(B\) large makes the tail summable over
the \(O(\rho_n^2\log\rho_n)\) enlarged arcs, layers, and angular scales.

On this high-probability occupancy event, uniformly in \(v\),
\[
\begin{aligned}
    \sum_{i:-\ell_n^{\mathrm{edge}}\le \xi_i\le-A}e^{\xi_i-Y_i(v)^2/2}
    &\le
    C\sum_{m=\lfloor A\rfloor}^{m_n^\star} e^{-m}
    \sum_{q=0}^{\lfloor\pi\rho_n\rfloor}
        e^{-q^2/2}N_{m,q,n} \\
    &\le
    C\sum_{q=0}^{\lfloor\pi\rho_n\rfloor}e^{-q^2/2}
    \sum_{m=\lfloor A\rfloor}^{m_n^\star} e^{-m}a_{m,q,n}.
\end{aligned}
\]
The first two terms in \(a_{m,q,n}\) contribute
\[
    O(e^{-A})+O(\ell_n^{\mathrm{edge}}/\rho_n).
\]
For the logarithmic occupancy term, fix \(A\) and split the \(m\)-sum at
\(\log\{\rho_n/(q+1)\}\).  This gives the deterministic bound
\[
    \sum_{m=\lfloor A\rfloor}^{m_n^\star}
    e^{-m}
    \frac{\log\rho_n}
        {1+\log_+\!\left(\rho_n/((q+1)e^m)\right)}
    \le
    C e^{-A}\{1+\log(q+2)\}
    +C\frac{(q+1)\log\rho_n}{\rho_n}
\]
for all \(n\ge n_0(A)\), uniformly in \(q\). This fixed-\(A\) qualification
is essential; the estimate is not asserted for a simultaneous choice
\(A=A_n\).
After summing in \(q\) against \(e^{-q^2/2}\), the right side is
\[
    O(e^{-A})+o(1).
\]
Thus the displayed supremum is at most \(C e^{-A}+o(1)\) with probability
tending to one, and the lemma follows by letting \(A\to\infty\).
\end{proof}

The moving-cutoff decomposition can therefore be replaced by the fixed-cutoff
form needed for the reverse inequality.
\begin{lemma}
\label{lem:d2-fixed-cutoff-decomposition}
For each compact \(K\subset\R\), each \(C<\infty\), and each \(\delta>0\),
\[
\begin{aligned}
&\lim_{A\to\infty}\limsup_{n\to\infty}
\Pbb\Biggl\{
\max_i\xi_i\le C,\
\sup_{\substack{v\in S^1\\ s\in K}}
\Biggl|
H_n((\rho_n+s)v)-\Phi(-s)\\
&\hspace{35mm}
-\sum_{i:-A<\xi_i\le C}
e^{\xi_i-s^2/2-Y_i(v)^2/2}
\Biggr|>\delta
\Biggr\}=0 .
\end{aligned}
\]
\end{lemma}

\begin{proof}
Apply Proposition~\ref{prop:two-d-compact-decomposition} on a finite closed-arc
cover of \(S^1\). On \(\{\max_i\xi_i\le C\}\), its fixed-cutoff sum omits only
the observations with
\(-\ell_n^{\mathrm{edge}}\le \xi_i\le-A\). Their
absolute contribution is bounded uniformly in \(s\in K\) by
\[
    \sup_{v\in S^1}
    \sum_{i:-\ell_n^{\mathrm{edge}}\le \xi_i\le-A} e^{\xi_i-Y_i(v)^2/2},
\]
because \(e^{-s^2/2}\le1\).  Lemma~\ref{lem:d2-low-atom-absorption} makes
this term vanish after first taking \(n\to\infty\) and then \(A\to\infty\),
proving the fixed-cutoff form.
\end{proof}


\section{The Spatial Edge in Dimensions \texorpdfstring{\(d_n\ge3\)}{at Least Three}}
\label{sec:radial-transition}
\label{sec:spatial-edge}

This section proves the spatial-exclusion estimates needed for the
\(d_n\ge3\) clause of Theorem~\ref{thm:main-support-law}. The exact Gamma
normalization gives the same radial Poisson limit in every growth regime, but
the argument showing that subthreshold observations cannot create additional
KKT contacts changes with dimension. We first derive the common radial
coordinate, then prove the required exclusion in bounded and growing
dimensions below the simplex scale. Section~\ref{sec:simplex} treats the
remaining high-dimensional regimes.

\subsection{The radial coordinate}

The exact coordinate \(s_n^\star\) appears in
Theorem~\ref{thm:main-support-law}. In the range
\(d_n\log n=o(n^2)\), the spatial estimates are cleaner in the following
\(\log n\)-based coordinate:
\begin{equation}
\label{eq:s-definition}
  s_n
  =
  \frac{\gamma_nL_n-b_n}{a_n}.
\end{equation}
The corresponding critical variance defect is
\begin{equation}
\label{eq:critical-epsilon}
  \eps_{c,n}=1-\frac{L_n}{b_n},
\end{equation}
and an \(O(1)\) change in \(s_n\) changes \(\eps_n\) on the scale
\begin{equation}
\label{eq:epsilon-window}
  \Delta\eps_n
  \asymp
  \frac{a_nL_n}{b_n^2}.
\end{equation}
Lemma~\ref{lem:exact-old-edge-coordinate} shows that \(s_n\) agrees with the
exact coordinate \(s_n^\star\) whenever \(d_n\log n=o(n^2)\). Define
\[
  \xi_{n,i}
  =
  \frac{|\widetilde Z_i|^2/2-b_n}{a_n},
  \qquad
  \delta_n=\frac{a_n}{\gamma_n}.
\]
At the normalization \eqref{eq:s-definition}, the bump height is exactly
\begin{equation}
\label{eq:bump-height-xi}
  h_i
  =
  \exp\{\delta_n(\xi_{n,i}-s_n)\}.
\end{equation}
The score radius satisfies the exact identity
\begin{equation}
\label{eq:universal-score-radius}
  |q_i|^2
  =
  \frac{|\widetilde Z_i|^2}{\gamma_n}
  =2L_n+2\delta_n(\xi_{n,i}-s_n).
\end{equation}
Consequently, $|q_i|^2=2L_n+O_{\Pbb}(\delta_n)$ for edge observations with
$\xi_{n,i}=O_{\Pbb}(1)$ whenever $s_n=O(1)$.
Thus the top bumps live on a sphere of radius \(\sqrt{2L_n}\). The dimension
affects their angular geometry and the height scale \(\delta_n\), but not
their leading score radius.

When \(d_n\gg L_n\), however,
\(\delta_n\to0\), so a radial gap of \(a_n\eta\) below threshold creates
only a height gap of order \(\delta_n\eta\). The remainder of the field must
therefore be \(o_{\Pbb}(\delta_n)\), not merely \(o_{\Pbb}(1)\).

\subsection{Consequences across dimension scales}
\label{sec:regime-map}

The exact definitions \eqref{eq:b-a-definition} are preferable in the main
theorem, but their asymptotic expansions reveal how the critical defect and
window width depend on dimension. In this subsection
\(\eps_{c,n}=1-L_n/b_n\) denotes a leading-order proxy for the critical
defect. Section~\ref{sec:full-growing-transition} restores the finite-sample
correction that becomes visible when \(d_n\asymp n^2\log n\).

\medskip
\noindent\emph{Sublogarithmic dimension.}

Suppose \(k_n=o(L_n)\). Then \(b_n/L_n\to1\), \(a_n\to1\), and
\begin{equation}
\label{eq:b-sublog}
  b_n
  =
  L_n+(k_n-1)\log L_n-\log\Gamma(k_n)
  +O\left(
    \frac{k_n}{L_n}
    \left\{1+(k_n-1)\log L_n-\log\Gamma(k_n)\right\}
  \right).
\end{equation}
When in addition \(k_n\to\infty\), Stirling's formula gives the simpler
leading term
\begin{equation}
\label{eq:b-sublog-stirling}
  b_n-L_n
  \sim
  k_n\log\frac{L_n}{k_n}.
\end{equation}
Consequently,
\begin{equation}
\label{eq:epsilon-sublog}
  \eps_{c,n}
  \sim
  \frac{k_n}{L_n}\log\frac{L_n}{k_n}
  =
  \frac{d_n}{2\log n}
  \log\frac{2\log n}{d_n},
  \qquad
  \Delta\eps_n\asymp\frac1{\log n}.
\end{equation}

For fixed \(d\ge3\), \eqref{eq:b-sublog} reduces to
\[
  b_n
  =
  \log n+\left(\frac d2-1\right)\log\log n
  -\log\Gamma(d/2)+o(1),
\]
recovering the fixed-dimensional normalization used in
Section~\ref{sec:fixed-d-patches}. If
\(d_n\to\infty\), then \(b_n-L_n\to\infty\); this makes the diffuse edge
background vanish, leaving the radial maximum as the limiting obstruction.

\begin{proof}[Justification of \eqref{eq:b-sublog}]
For \(y>k_n-1\), \eqref{eq:tail-density-ratio} and
\(1\le\overline F_n(y)/f_n(y)\le\{1-(k_n-1)/y\}^{-1}\) give
\[
  \log\overline F_n(y)
  =
  -y+(k_n-1)\log y-\log\Gamma(k_n)+O(k_n/y).
\]
Standard Gamma Chernoff bounds first give \(b_n/L_n\to1\). Substitute
\(y=b_n\), write \(b_n=L_n+r_n\), and expand
\(\log b_n=\log L_n+O(r_n/L_n)\). The resulting equation first gives
\(r_n=O((k_n-1)\log L_n-\log\Gamma(k_n)+1)\), and a second substitution
gives \eqref{eq:b-sublog}. Formula \eqref{eq:b-sublog-stirling} follows
because \(L_n/k_n\to\infty\).
\end{proof}

\medskip
\noindent\emph{Logarithmic dimension.}

Suppose \(k_n/L_n\to c\in(0,\infty)\). Let \(x_c>1\) be the unique
solution of
\begin{equation}
\label{eq:xc}
  x_c-1-\log x_c=\frac1c.
\end{equation}
Stirling's formula and the Gamma saddlepoint estimate yield
\begin{equation}
\label{eq:logarithmic-regime}
  \frac{b_n}{k_n}\to x_c,
  \qquad
  a_n\to\frac{x_c}{x_c-1},
  \qquad
  \eps_{c,n}\to1-\frac1{cx_c}.
\end{equation}
The transition therefore occurs at a nonzero constant variance defect, while
its width in \(\eps_n\) remains of order \(1/\log n\).

\medskip
\noindent\emph{Superlogarithmic dimension.}

Suppose \(k_n/L_n\to\infty\). The Gamma quantile is in its
moderate-deviation range:
\begin{equation}
\label{eq:superlog-regime}
  b_n
  =
  k_n+\sqrt{2k_nL_n}\{1+o(1)\},
  \qquad
  a_n
  =
  \sqrt{\frac{k_n}{2L_n}}\{1+o(1)\}.
\end{equation}
Hence
\begin{equation}
\label{eq:epsilon-superlog}
  \eps_{c,n}
  =
  1-\frac{L_n}{k_n}\{1+o(1)\},
  \qquad
  \Delta\eps_n
  \asymp
  \frac{\sqrt{L_n}}{k_n^{3/2}},
  \qquad
  \delta_n
  \asymp
  \sqrt{\frac{L_n}{k_n}}.
\end{equation}
The fitted variance must now be much larger than the data variance: the
critical defect tends to one. The shrinking factor \(\delta_n\) in
\eqref{eq:epsilon-superlog} is also the precision required of the
single-bump approximation.

In particular, the usual proportional-dimensional regime
\(d_n/n\to\tau\in(0,\infty)\) is a special case of the superlogarithmic
regime, not the logarithmic regime \(d_n\asymp\log n\).  In that case,
\[
  1-\eps_{c,n}
  \sim\frac{2\log n}{d_n}
  \sim\frac{2\log n}{\tau n},
  \qquad
  \Delta\eps_n\asymp\frac{\sqrt{\log n}}{n^{3/2}}.
\]
The raw squared radii are of order \(n\), but after the variance rescaling
the likelihood of one empirical bump still has to overcome its factor
\(1/n\); this is why the KKT threshold remains \(\log n\).

For completeness, all three regimes follow from the same saddlepoint formula.
If \(k_n\to\infty\), put \(x_n=b_n/k_n>1\). Uniformly at the extreme
quantile,
\begin{equation}
\label{eq:gamma-saddlepoint}
  \overline F_n(k_nx_n)
  =
  \frac{1+o(1)}{(x_n-1)\sqrt{2\pi k_n}}
  \exp\{-k_n(x_n-1-\log x_n)\}.
\end{equation}
Indeed, Stirling's formula gives the density at \(k_nx_n\), while
Lemma~\ref{lem:gamma-hazard} gives the tail-to-density ratio
\(x_n/(x_n-1)\{1+o(1)\}\). Equating the logarithm of
\eqref{eq:gamma-saddlepoint} to \(-L_n\) proves
\eqref{eq:logarithmic-regime}; when \(k_n/L_n\to\infty\), expanding
\(x-1-\log x=(x-1)^2/2+O((x-1)^3)\) proves
\eqref{eq:superlog-regime}.

Finally, \eqref{eq:centering-condition} is automatic when \(k_n=O(L_n)\).
If \(k_n/L_n\to\infty\), \eqref{eq:superlog-regime} reduces it to
\begin{equation}
\label{eq:centering-superlog}
  k_nL_n=o(n).
\end{equation}


\subsection{Bounded-dimensional radial patches}
\label{sec:fixed-d-patches}
\label{sec:sharp-high-dimensional-edge}

This subsection proves the fixed-dimensional spatial exclusion needed for the
\(d\ge3\) support law in Theorem~\ref{thm:main-support-law}. The moving radial
shell identifies the candidate KKT patches, and the uniform estimates in
Subsection~\ref{sec:sharp-high-dimensional-estimates} exclude every other
score location.

Let \(Z_1,\ldots,Z_n\) be i.i.d.\ \(\cN(0,I_d)\), \(d\ge3\), and set
\[
    H_n(t)
    =
    \frac1n\sum_{i=1}^n
    \exp\{t\cdot Z_i-|t|^2/2\},
    \qquad t\in\R^d .
\]

\subsubsection{The moving radial shell}

Let
\[
    Y_i=\frac{|Z_i|^2}{2},\qquad U_i=\frac{Z_i}{|Z_i|}\in S^{d-1},
    \qquad k=\frac d2 .
\]
Since \(Y_i\) has the \(\Gamma(k,1)\) law,
\[
    b_{n,d}
    =
    \log n+(k-1)\log\log n-\log\Gamma(k),
\]
the fixed-dimensional specialization of the Gamma extreme-value theory in
Section~\ref{sec:mv-gamma-normalization} gives the marked radial limit.
\begin{lemma}
\label{lem:radial-gamma-extremes}
The extremal process
\[
    \Xi_n=\sum_{i=1}^n\delta_{(Y_i-b_{n,d},\,U_i)}
\]
converges to a Poisson point process on \(\R\times S^{d-1}\) with intensity
\(e^{-x}dx\,\sigma_{d-1}(du)\), where \(\sigma_{d-1}\) is normalized surface
measure.  In particular, \(M_n=\max_i(Y_i-b_{n,d})\) satisfies
\[
    \Pbb\{M_n\le x\}\to\exp\{-e^{-x}\}.
\]
\end{lemma}

\begin{proof}
Proposition~\ref{prop:gamma-max} and
Lemma~\ref{lem:fixed-d-normalization-bridge} give the unmarked radial limit.
Radius and direction are independent, and \(U_i\) has law
\(\sigma_{d-1}\), so the standard marking theorem gives the displayed
point-process convergence.
\end{proof}

Put \(\rho_n=\sqrt{2b_{n,d}}\) and
\[
    \lambda_n=\kappa_n\rho_n^2 .
\]
In the critical window \(\lambda_n=O(\log\log n)\), and it is useful to write
\[
    s_n=\lambda_n-(k-1)\log\log n+\log\Gamma(k).
\]
Completing the square in one summand of \(D_n\), or equivalently using
Lemma~\ref{lem:bump-representation}, shows that an observation with
\(Y_i=b_{n,d}+x\) has optimized log contribution \(x-s_n+o_{\Pbb}(1)\),
uniformly for bounded \(x\). Thus an observation is locally supercritical
precisely when its radial excess \(x\) exceeds \(s_n\); these observations
index the candidate KKT crossings.

To resolve each candidate spatially, write the score and observation in normal
coordinates about a direction \(v\). Let
\(\exp_v:T_vS^{d-1}\to S^{d-1}\) be the Riemannian exponential map and put
\[
    t=\left(\rho_n+\frac q{\rho_n}\right)v,\qquad
    z=\left(\rho_n+\frac x{\rho_n}\right)\exp_v\left(\frac y{\rho_n}\right),
    \qquad y\in T_vS^{d-1}.
\]
Then, uniformly for \(x,q,|y|=O(\log\log n)\),
\[
    t\cdot z-\frac{|t|^2}{2}-\log n-\kappa_n|t|^2
    =
    (k-1)\log\log n-\log\Gamma(k)-\lambda_n
    +x-\frac{|y|^2}{2}
    +
    O\!\left(\frac{(\log\log n)^4}{\rho_n^2}\right).
\]
The radial score offset \(q\) cancels to first order. A radial excess \(x\)
therefore raises the optimized contribution by \(x\), whereas an angular
displacement \(y/\rho_n\) lowers it by \(|y|^2/2\); the variance defect
subtracts \(\lambda_n\). The estimates below show that the rest of the sample
contributes \(o_{\Pbb}(1)\) uniformly on each candidate patch.
Subsection~\ref{sec:higher-dimensional-support-counts} uses these patch
estimates to recover the fitted support.

For the uniform exclusion argument, use the centered radial maximum
\[
    M_n^c
    =
    \max_i\left(\frac{|Z_i-\bar Z_n|^2}{2}-b_{n,d}\right).
\]
It differs from \(M_n=\max_i(Y_i-b_{n,d})\) by \(o_{\Pbb}(1)\), since
\[
    |M_n^c-M_n|
    \le
    \max_i|Z_i|\,|\bar Z_n|+\frac{|\bar Z_n|^2}{2}
    =
    o_{\Pbb}(1).
\]

\subsection{Uniform moving-shell estimates}
\label{sec:sharp-high-dimensional-estimates}

This subsection supplies the estimates used to prune subthreshold radial
patches in Lemma~\ref{lem:support-dge3-deleted-field}. Bennett's inequality
controls the center of each angular patch, while scaled derivative bounds and
Morrey's inequality control oscillation across patches of diameter
\(1/\rho_n\). The following bulk estimates handle smaller score radii by a
clipped empirical-process argument.

For the moving-shell bound, replace the field by an independent truncated
field. On the radial-maximum event used in the comparison, every observation
lies below the deterministic truncation radius, so the two fields agree.

Fix a gap \(a>0\).  Set \(r=|t|\),
\[
    T_n(a)=\log n+\lambda_n-a,
    \qquad
    R_n(a)=\{2T_n(a)\}^{1/2},
\]
and write \(R_a=R_n(a)\).  Define
\[
    V_{i,t}^{(a)}
    =
    \frac1n
    e^{-\kappa_nr^2}
    \exp\left\{t\cdot Z_i-\frac{r^2}{2}\right\}
    \mathbf 1_{\{|Z_i|\le R_a\}},
    \qquad
    S_t^{(a)}=\sum_iV_{i,t}^{(a)}.
\]
Let
\[
    \mathcal T_n(B)=
    \left\{t:\ \left||t|^2/2-b_{n,d}\right|\le B\log\log n\right\}.
\]
The centered field is compared to this independent truncated field with a
small deterministic loss in the gap.  Put
\[
    \alpha_n=(\log\log n)^{-2},
    \qquad
    G_n=\{\rho_n|\bar Z_n|\le\alpha_n\}.
\]
Since \(\alpha_n\sqrt n/\rho_n\to\infty\), \(\Pbb(G_n^c)\to0\), and on
\(G_n\), for all large \(n\),
\[
    \sup_{t\in\mathcal T_n(B)}|t\cdot\bar Z_n|\le2\alpha_n .
\]
If \(M_n^c\le s_n-\eta\), then
\(|Z_i-\bar Z_n|^2/2\le\log n+\lambda_n-\eta\) for all \(i\).  On \(G_n\)
this implies \(|Z_i|^2/2\le\log n+\lambda_n-\eta/2\) for all large \(n\).
Consequently, uniformly over \(t\in\mathcal T_n(B)\),
\[
    D_n(t)\le e^{2\alpha_n}S_t^{(\eta/2)} .
\]
It is therefore enough to prove that, for fixed \(a>0\), the independent
truncated field stays below \(1-o(1)\) uniformly on \(\mathcal T_n(B)\).

The first input bounds the moments of one truncated summand. Exponential
tilting turns each moment into the probability that a noncentral Gaussian lies
in the truncation ball; a one-dimensional Gaussian tail estimate then gives
the required decay.
\begin{proposition}
\label{prop:noncentral-ball}
Fix \(d\ge3\), \(C_0<\infty\), \(B<\infty\), \(p>1\), and \(a>0\).  Suppose
the sharp-scale relation
\[
    s_n
    =
    \lambda_n-\left(\frac d2-1\right)\log\log n+\log\Gamma(d/2)
\]
satisfies \(|s_n|\le C_0\).  Uniformly over deterministic sequences satisfying
this relation, over directions \(t/|t|\), and over
\[
    \left|\frac{r^2}{2}-b_{n,d}\right|\le B\log\log n,
\]
one has
\[
    \sum_{i=1}^n\E \{V_{i,t}^{(a)}\}^p
    \le
    C_{p,a,B,C_0,d}\rho_n^{-1}
    \exp\{-\lambda_n-(p-1)a+o(1)\}.
\]
\end{proposition}

\begin{proof}
The exact identity is
\[
    \sum_{i=1}^n\E \{V_{i,t}^{(a)}\}^p
    =
    n^{1-p}
    e^{-p\kappa_nr^2+(p^2-p)r^2/2}
    \Pbb\{|G+pt|\le R_n(a)\},
    \]
where \(G\sim\cN(0,I_d)\).  Align \(t=re_1\), write \(G=We_1+U\), and put
\(\Delta=pr-R_n(a)\).  In the moving window,
\(\Delta=(p-1)\rho_n+O((\log\log n)/\rho_n)\), so
\(\Delta\ge c_p\rho_n\).  If \(|G+pt|\le R_n(a)\), then
\[
    W\le -\left(pr-\sqrt{R_n(a)^2-|U|^2}\right).
\]
Mills' bound and the inequality
\[
    pr-\sqrt{R_n(a)^2-|U|^2}
    \ge
    \Delta+\frac{|U|^2}{2R_n(a)}
\]
give
\[
    \Pbb\{|G+pt|\le R_n(a)\}
    \le
    C_{p,d}\rho_n^{-1}
    \exp\left\{-\frac{(pr-R_n(a))^2}{2}\right\}.
\]
Substitution gives
\[
    \sum_{i=1}^n\E \{V_{i,t}^{(a)}\}^p
    \le
    C_{p,d}\rho_n^{-1}
    \exp\left\{
        -p\kappa_nr^2
        +(p-1)(\lambda_n-a)
        -\frac p2(r-R_n(a))^2
    \right\}.
\]
Since the sharp-scale assumption gives \(\lambda_n=O(\log\log n)\), one has
\(\kappa_nr^2=\lambda_n+o(1)\) uniformly in the displayed shell. This proves
the stated estimate.
\end{proof}

Under the same sharp-scale assumptions, the truncation also gives the
pointwise jump bound
\[
    V_{i,t}^{(a)}\le e^{-a+o(1)}
\]
uniformly in the moving shell.  Indeed
\[
    \log V_{i,t}^{(a)}
    \le
    -\log n-\kappa_nr^2-\frac{r^2}{2}+rR_n(a)
    =
    \lambda_n-a-\kappa_nr^2-\frac12(r-R_n(a))^2.
\]
Thus \(V_{i,t}^{(a)}\le c_a<1\) for large \(n\), where \(c_a\) is deterministic.

The moment and jump bounds give a fixed-location tail strong enough for the
angular union bound after the Morrey step.
\begin{proposition}
\label{prop:fixed-t-bennett}
Fix \(d\ge3\), \(C_0<\infty\), \(B<\infty\), and \(a>0\), and suppose the
sharp-scale relation in Proposition~\ref{prop:noncentral-ball} satisfies
\(|s_n|\le C_0\).  Uniformly over deterministic sequences satisfying this
relation and over the moving shell, for every deterministic
\(\delta_n\to0\) with \(0\le\delta_n\le1/2\) eventually,
\[
    \Pbb\{S_t^{(a)}\ge1-\delta_n\}
    \le
    C_{a,B,C_0,d}
    \left(\rho_n^{-1}e^{-\lambda_n}\right)^{
        e^a(1-\delta_n)-o(1)} .
\]
In particular, at level \(1\) the exponent is \(e^a-o(1)\).
\end{proposition}

\begin{proof}
Let \(\mu_t=\E S_t^{(a)}\), \(v_t=\sum_i\operatorname{Var}(V_{i,t}^{(a)})\),
\(A_n=\rho_n^{-1}e^{-\lambda_n}\), and \(x_n=1-\delta_n\).
Proposition~\ref{prop:noncentral-ball} with \(p=2\) gives
\[
    v_t\le C\rho_n^{-1}e^{-\lambda_n-a+o(1)},
    \qquad
    \mu_t\le e^{-\lambda_n+o(1)}=o(1)
\]
under the displayed \(d\ge3\) sharp-scale assumption.  Bennett's inequality for
\(0\le V_{i,t}^{(a)}\le c_a<1\) gives
\[
    \Pbb\{S_t^{(a)}\ge x_n\}
    \le
    \exp\left[
        -\frac{v_t}{c_a^2}
        h\left(\frac{c_a(x_n-\mu_t)}{v_t}\right)
    \right],
    \qquad h(u)=(1+u)\log(1+u)-u .
\]
Using the sharp jump bound \(V_{i,t}^{(a)}\le e^{-a+o(1)}\) and
\(\log A_n^{-1}=(d-1)\log\rho_n+O(1)\), the Bennett rate is at least
\[
    \{e^a(1-\delta_n)-o(1)\}\log A_n^{-1}.
\]
This proves the displayed estimate. Only the second moment is needed at a
fixed location; the derivative estimates used for the Morrey transfer require
higher moments.
\end{proof}

At the sharp scale \(s_n=O(1)\), equivalently
\(\lambda_n=(d-2)\log\rho_n+O(1)\), the quantity
\(\rho_n^{-1}e^{-\lambda_n}\) is \(\rho_n^{-(d-1)+o(1)}\).  The strict
near-unit exponent \(e^a(1-o(1))>1\) is therefore exactly what is needed to
union bound over \(\rho_n^{d-1+o(1)}\) angular locations.
Throughout the remaining \(d\ge3\) estimates, ``at the sharp scale
\(s_n=O(1)\)'' means uniformly over deterministic sequences satisfying
\(|s_n|\le C\), for each fixed \(C<\infty\); constants may depend on \(C\).

Scaled derivatives control oscillation on angular boxes at scale
\(1/\rho_n\), reducing the continuum to a net of size
\(\rho_n^{d-1+o(1)}\).
\begin{lemma}
\label{lem:scaled-derivative-morrey}
Fix \(d\ge3\), \(p>d\), \(a>0\), and suppose \(s_n=O(1)\).
Let \(Q\subset\R^{d-1}\) be a fixed cube, and let
\(\{u_\alpha\}\) be a bounded-overlap family of smooth charts such that
\(\{u_\alpha(Q/\rho_n)\}\) covers \(S^{d-1}\), the number of charts is
\(O(\rho_n^{d-1})\), and all first two derivatives of the charts are
bounded uniformly in \(\alpha\) and \(n\).  For each fixed \(B<\infty\), in
moving-shell coordinates
\[
    t=t(q,y)=\sqrt{2(b_{n,d}+q)}\,u_\alpha(y/\rho_n),
    \qquad |q|\le B\log\log n,\quad y\in Q,
\]
one has
\[
    \E\sum_\alpha\int_{|q|\le B\log\log n}\int_Q
    |\nabla_{q,y}S_{t(q,y)}^{(a)}|^p\,dy\,dq
    \le(\log\log n)^C .
\]
The displayed derivative-moment bound also holds on every fixed physical-offset annulus
\[
    t=t(h,y)=(R_a-h)u_\alpha(y/\rho_n),
    \qquad |h|\le A,\quad y\in Q,
\]
with \(C=C(a,A,B,d,p)\). Consequently, in either region and for every fixed
\(\gamma>0\), a net of cardinality
\[
    \rho_n^{d-1}(\log\log n)^{C_\gamma}=\rho_n^{d-1+o(1)}
\]
suffices, with probability tending to one, to approximate \(S_t^{(a)}\)
up to the prescribed error
\[
    \varepsilon_{n,\gamma}=(\log\log n)^{-\gamma}.
\]
\end{lemma}

\begin{proof}
The truncation set \(\{|Z_i|\le R_a\}\) is independent of \(t\), so
\(S_t^{(a)}\) is \(C^1\) in the displayed coordinates.  For any coordinate
\(\xi\in\{q,y_1,\ldots,y_{d-1}\}\) or
\(\xi\in\{h,y_1,\ldots,y_{d-1}\}\),
\[
    \partial_\xi V_{i,t}^{(a)}
    =
    V_{i,t}^{(a)}M_\xi(Z_i,t),
    \qquad
    M_\xi(z,t)=\{z-(1+2\kappa_n)t\}\cdot\partial_\xi t .
\]
The coordinate choices give
\[
    |\partial_qt|=O(\rho_n^{-1}),\qquad
    |\partial_ht|=1,\qquad
    |\partial_{y_j}t|=O(1),\qquad
    t\cdot\partial_{y_j}t=0 .
\]
For \(m\in\{2,p\}\), exponential tilting gives the exact identity
\[
    \sum_i\E|\partial_\xi V_{i,t}^{(a)}|^m
    =
    n^{1-m}e^{-m\kappa_nr^2+(m^2-m)r^2/2}
    \E\!\left[
        |M_\xi(G+mt,t)|^m\mathbf 1_{\{|G+mt|\le R_a\}}
    \right].
\]
We next bound the tilted expectation.  Align \(t=re_1\), write
\(z=G+mt\), and decompose
\[
    z=z_1e_1+z_\perp,\qquad
    \zeta=R_a-z_1,\qquad \Delta_m=mr-R_a .
\]
In both the moving shell and fixed physical-offset annuli,
\(\Delta_m=(m-1)\rho_n+O(\log\log n)\), and hence
\(\Delta_m\ge c_m\rho_n\) for large \(n\).  We claim that, for every fixed
integer \(L\),
\[
    \E\!\left[
        (1+\zeta_+ + |z_\perp|)^L\mathbf 1_{\{|z|\le R_a\}}
    \right]
    \le
    C_{L,m}\rho_n^{-1}e^{-\Delta_m^2/2}.
\]
Indeed, write \(G=We_1+U\).  If \(|G+mt|\le R_a\), then
\[
    W\le -\Delta_m-\frac{|U|^2}{2R_a},
\]
because \(\sqrt{R_a^2-|U|^2}\le R_a-|U|^2/(2R_a)\).  The one-dimensional
Gaussian tail-moment bound
\[
    \E\{(1+(-x-W)_+)^L\mathbf 1_{W\le -x}\}
    \le C_Lx^{-1}e^{-x^2/2},\qquad x\ge1,
\]
with \(x=\Delta_m+|U|^2/(2R_a)\), gives the claim after integration in
\(U\): since \(\Delta_m/R_a\) stays bounded away from zero and infinity,
the factor \(e^{-x^2/2}\) is at most
\(e^{-\Delta_m^2/2}e^{-c|U|^2}\).

In these coordinates the derivative marks have the following deterministic
bounds on \(\{|z|\le R_a\}\):
\[
\begin{aligned}
    M_q(z,t)
    &=
    \{z-(1+2\kappa_n)t\}\cdot\partial_qt
      =O\{(1+\zeta_+)/\rho_n+\kappa_n\rho_n\},\\
    M_h(z,t)
    &=
    \{z-(1+2\kappa_n)t\}\cdot\partial_ht
      =O_A(1+\zeta_++\kappa_n\rho_n),\\
    M_{y_j}(z,t)
    &=
    \{z-(1+2\kappa_n)t\}\cdot\partial_{y_j}t
      =O\{1+|z_\perp|+(\log\log n)(1+\zeta_+)/\rho_n\}.
\end{aligned}
\]
For \(M_q\), use \(R_a-r=O(\log\log n/\rho_n)\) in the moving shell; for
\(M_h\), use \(|h|\le A\); for \(M_{y_j}\), use
\(t\cdot\partial_{y_j}t=0\) and the uniform chart derivative bounds.  Since
\(\kappa_n\rho_n=O(\log\log n/\rho_n)=o(1)\), the tilted tail-moment estimate
from Proposition~\ref{prop:noncentral-ball} shows that the derivative marks cost only fixed powers
of \(\log\log n\).  Combining this with the same algebra as in
Proposition~\ref{prop:noncentral-ball} yields, uniformly in the moving shell
and in fixed physical-offset annuli,
\[
    \sum_i\E|\partial_\xi V_{i,t}^{(a)}|^m
    \le
    (\log\log n)^C\rho_n^{-1}e^{-\lambda_n}
    \le
    (\log\log n)^C\rho_n^{-(d-1)} .
\]
Here we used \(e^{-\lambda_n}\le C\rho_n^{-(d-2)}\), which follows from
\(s_n=O(1)\).

Apply Rosenthal's inequality to
\(\partial_\xi S_t^{(a)}=\sum_i\partial_\xi V_{i,t}^{(a)}\).  The centered
\(p\)th-moment term is bounded by the display with \(m=p\), and the variance
term is bounded by the display with \(m=2\) raised to the power \(p/2\),
which is smaller.  It remains only to bound the mean.  Rotational symmetry
makes the angular means zero.  For radial derivatives, write
\[
    \E S_t^{(a)}=e^{-\kappa_nr^2}P_r,\qquad
    P_r=\Pbb\{|G+re_1|\le R_a\}.
\]
The derivative \(P_r'\) is uniformly bounded, because
\[
    P_r'=\int_{|z|\le R_a}(z_1-r)\phi_d(z-re_1)\,dz
\]
and the absolute value of this integral is at most \(\E|G_1|\).  Hence
\(\partial_q\E S_t^{(a)}=O(e^{-\lambda_n}/\rho_n)\), since
\(\partial_q r=1/r=O(\rho_n^{-1})\) and
\(\kappa_n=O(\log\log n/\rho_n^2)\).  On a fixed physical-offset annulus,
\(\partial_h\E S_t^{(a)}=O_A(e^{-\lambda_n})\).  Their \(p\)th powers are
also \(O(\rho_n^{-(d-1)})\), because \(p>d\) and \(d\ge3\).  Therefore
\[
    \E|\partial_\xi S_t^{(a)}|^p
    \le(\log\log n)^C\rho_n^{-(d-1)}
\]
uniformly in all coordinates and all displayed parameter values.

Integrating this pointwise bound over \(O(\rho_n^{d-1})\) angular charts and
over \(q\)-length \(O(\log\log n)\), or fixed \(h\)-length \(O_A(1)\), gives
the displayed integral bound.  By Markov, the total gradient integral is at
most \((\log\log n)^{C'}\) with probability tending to one.  Fix
\(\gamma>0\), choose
\[
    K>\frac{\gamma+C'/p}{1-d/p},
\]
and, on this event, partition every coordinate patch into \(d\)-dimensional
cubes.  Take the side length to be
\[
    \ell_n=(\log\log n)^{-K}.
\]
The Morrey--Poincare inequality on each cube gives
\[
    O\!\left(\ell_n^{1-d/p}(\log\log n)^{C'/p}\right)
    =
    O\!\left((\log\log n)^{-K(1-d/p)+C'/p}\right),
\]
as a uniform bound for the oscillation of \(S_t^{(a)}\) on that cube.  For
the displayed choice of \(K\), this is at most
\((\log\log n)^{-\gamma}\).  Taking one point from each cube gives a net of
size \(\rho_n^{d-1}(\log\log n)^C\) in either coordinate region.
\end{proof}

The Morrey estimate is used below only to transfer fixed-location bounds to
the angular continuum. Its net has size \(\rho_n^{d-1+o(1)}\), which is small
enough for the Bennett exponents in Proposition~\ref{prop:fixed-t-bennett}.

For radii below the moving shell, use physical distance from the truncated
sample edge. Put
\[
    R_a=R_n(a),\qquad h=R_a-|t|.
\]
A bounded interval of \(h\) corresponds to a radial-energy interval of length
\(O(\rho_n)\), so a direct energy-coordinate net would be unnecessarily large.
The fixed-location estimates retain the Gaussian penalty in \(h\).
\begin{lemma}
\label{lem:physical-offset-fixed-t}
Fix \(d\ge3\), \(a>0\), and \(A<\infty\), and suppose \(s_n=O(1)\).
Uniformly over directions and \(|h|\le A\), for each fixed \(p>1\),
\[
 V_{i,t}^{(a)}\le e^{-a-h^2/2+o(1)},\qquad
 \sum_i\E\{V_{i,t}^{(a)}\}^p
 \le C_{p,a,A,d}\rho_n^{-1}
 e^{-\lambda_n-(p-1)a-ph^2/2+o(1)}.
\]
Consequently, for every fixed \(\zeta\in(0,1)\),
\[
 \Pbb\{S_t^{(a)}\ge\zeta\}
 \le \left(\rho_n^{-1}e^{-\lambda_n}\right)^{
 e^{a+h^2/2}\zeta-o(1)}.
\]
\end{lemma}

\begin{proof}
The completing-square and Mills-ratio calculation used in
Proposition~\ref{prop:noncentral-ball} sharpens in the physical coordinate
\(h=R_a-r\).  Uniformly for bounded \(h\) and every fixed \(p>1\),
\[
    V_{i,t}^{(a)}\le \exp\{-a-h^2/2+o(1)\},
    \qquad
    \sum_i\E\{V_{i,t}^{(a)}\}^p
    \le
    C\rho_n^{-1}\exp\{-\lambda_n-(p-1)a-ph^2/2+o(1)\}.
\]
The first bound is the completing-square inequality
\[
    \log V_{i,t}^{(a)}
    \le
    \lambda_n-a-\kappa_nr^2-\frac12(R_a-r)^2
    =
    -a-h^2/2+o(1).
\]
For the second bound, exponential tilting reduces
\(\sum_i\E\{V_{i,t}^{(a)}\}^p\) to a Gaussian probability of the ball
\(\{|G|\le R_a\}\) under a mean shifted by \(pr\) in the \(t\)-direction.
The nearest boundary distance is
\[
    pr-R_a=(p-1)R_a-ph.
\]
Using the Mills bound as in Proposition~\ref{prop:noncentral-ball} gives
\[
    \Pbb\{|G+pt|\le R_a\}
    \le
    C_{p,A,d}\rho_n^{-1}
    \exp\left\{-\frac{(pr-R_a)^2}{2}\right\}.
\]
Substituting \(r=R_a-h\), \(R_a^2/2=\log n+\lambda_n-a\), and
\(\kappa_nr^2=\lambda_n+o(1)\) into the exact tilted identity gives
\[
\begin{aligned}
&n^{1-p}e^{-p\kappa_nr^2+(p^2-p)r^2/2}
    \exp\left\{-\frac{(pr-R_a)^2}{2}\right\}  \\
&\qquad\le
    \exp\{-\lambda_n-(p-1)a-ph^2/2+o(1)\},
\end{aligned}
\]
which is the displayed moment estimate for \(p>1\).  Separately,
\[
    \mu_t=\E S_t^{(a)}
    =
    e^{-\kappa_nr^2}\Pbb\{|G+t|\le R_a\}
    \le e^{-\kappa_nr^2}
    =
    e^{-\lambda_n+o(1)}
    =
    o(1).
\]
The \(p=2\) case gives
\[
    v_t:=\sum_i\operatorname{Var}(V_{i,t}^{(a)})
    \le C\rho_n^{-1}e^{-\lambda_n-a-h^2+o(1)},
\]
and the jump bound is \(V_{i,t}^{(a)}\le e^{-a-h^2/2+o(1)}\). Bennett's
inequality gives
\[
    \Pbb\{S_t^{(a)}\ge\zeta\}
    \le
    \left(\rho_n^{-1}e^{-\lambda_n}\right)^{e^{a+h^2/2}\zeta-o(1)}.
\]
\end{proof}

For growing physical offset \(h=R_a-|t|\), the fixed-location tail becomes
much smaller.
\begin{lemma}
\label{lem:growing-offset-fixed-t}
Fix \(d\ge3\), \(a>0\), \(A_0>0\), and \(\tau\in(0,1)\).  At the sharp
scale \(s_n=O(1)\), uniformly over directions and over
\[
    h=R_a-|t|\in[A_0,\rho_n/2+2],
\]
one has
\[
    \mu_t:=\E S_t^{(a)}=o(1),\qquad
    \max_i V_{i,t}^{(a)}\le \beta_h:=\exp\{-a-h^2/4\},
\]
and
\[
    v_t:=\sum_i\operatorname{Var}(V_{i,t}^{(a)})
    \le
    \beta_h\,\rho_n^{-(d-2)+o(1)} .
\]
Consequently,
\[
    \Pbb\{S_t^{(a)}\ge\tau\}
    \le
    \exp\{-c_{\tau,d}e^{a+h^2/4}\log\rho_n\}
\]
for all large \(n\), with \(c_{\tau,d}>0\).
\end{lemma}

\begin{proof}
Write \(R=R_a\), \(r=|t|=R-h\), and use \(R=\rho_n+O((\log\log n)/\rho_n)\),
\(\lambda_n=\kappa_n\rho_n^2=O(\log\log n)\).  Since \(h\ge A_0\),
\[
    \kappa_n(\rho_n^2-r^2)
    \le
    C\lambda_n h/\rho_n
    \le h^2/4
\]
uniformly on the displayed range, for all large \(n\).  The pointwise
completing-square bound therefore gives
\[
    \log V_{i,t}^{(a)}
    \le
    \lambda_n-a-\kappa_nr^2-h^2/2
    =
    -a-h^2/2+\kappa_n(\rho_n^2-r^2)
    \le -a-h^2/4 .
\]
This proves the jump bound.

The mean satisfies
\[
    \mu_t
    =
    e^{-\kappa_nr^2}\Pbb\{|G+t|\le R\}
    \le
    e^{-\kappa_nr^2}
    \le
    e^{-\lambda_n/10}=o(1),
\]
because \(r\ge R-\rho_n/2-2\ge2\rho_n/5\) for large \(n\), and
\(\lambda_n=(d-2)\log\rho_n+O(1)\).

It remains to bound the variance.  The exact second-moment identity is
\[
    \sum_i\E\{V_{i,t}^{(a)}\}^2
    =
    n^{-1}e^{-2\kappa_nr^2+r^2}
    \Pbb\{|G+2t|\le R\}.
\]
If \(h\le R/2-1\), then the ball event implies a one-dimensional
left-tail event of size \(R-2h\).  Hence
\[
    \Pbb\{|G+2t|\le R\}
    \le
    C\exp\{-(R-2h)^2/2\},
\]
and the algebra
\[
    -\log n+r^2-(R-2h)^2/2
    =
    \lambda_n-a-h^2
\]
gives
\[
    \sum_i\E\{V_{i,t}^{(a)}\}^2
    \le
    C\exp\{-a-h^2-\lambda_n
        +2\kappa_n(\rho_n^2-r^2)\}
    \le
    C e^{-a-h^2/2}\rho_n^{-(d-2)+o(1)} .
\]
This is bounded by \(\beta_h\rho_n^{-(d-2)+o(1)}\).  If
\(h>R/2-1\), then \(\Pbb\{|G+2t|\le R\}\le1\) and
\[
    n^{-1}e^{-2\kappa_nr^2+r^2}
    \le
    \exp\{-(1/2+o(1))\log n\},
\]
because \(r\le R/2+1\).  Since \(h\le\rho_n/2+2\), this is again
\(\le\beta_h\rho_n^{-(d-2)+o(1)}\): here the left side is
\(\exp\{-(1/2+o(1))\log n\}\), whereas \(\beta_h^{-1}\) is at most
\(\exp\{(1/8+o(1))\log n\}\).  The variance bound follows.

For large \(n\), \(\mu_t\le\tau/2\).  Bennett's inequality for centered
summands bounded by \(\beta_h\) and total variance \(v_t\) yields
\[
    \Pbb\{S_t^{(a)}\ge\tau\}
    \le
    \exp\left[
        -\frac{v_t}{\beta_h^2}
        H\!\left(\frac{\beta_h(\tau-\mu_t)}{v_t}\right)
    \right],
    \qquad
    H(u)=(1+u)\log(1+u)-u .
\]
Since \(\beta_h(\tau-\mu_t)/v_t\ge c_\tau\rho_n^{d-2-o(1)}\), the
exponent is at least \(c_{\tau,d}\beta_h^{-1}\log\rho_n\).  This is the
claimed tail bound.
\end{proof}

The exponential decay in \(h^2\) permits a union bound over a polynomial-size
Euclidean net. The next lemmas handle the remaining bulk
\(|t|\le\rho_n/2\).

Very close to the origin, negative sample curvature rules out a KKT crossing
without any edge analysis. The next lemma uses that curvature to prove the
small-radius part of Proposition~\ref{prop:bulk-mgf-exclusion}.
\begin{lemma}
\label{lem:microscopic-bulk}
Fix \(d\ge3\). At the sharp scale \(s_n=O(1)\),
\[
    \Pbb\left\{\sup_{0<|t|\le c_0\kappa_n}D_n(t)>1\right\}\to0
\]
for a sufficiently small constant \(c_0>0\).
\end{lemma}

\begin{proof}
Write
\[
    D_n(t)=\frac1n\sum_{i=1}^n
    \exp\{t\cdot \widetilde Z_i-(1/2+\kappa_n)|t|^2\},
    \qquad \widetilde Z_i=Z_i-\bar Z_n .
\]
On an event whose probability tends to one,
\[
    \left\|\frac1n\sum_i\widetilde Z_i\widetilde Z_i^\top-I_d\right\|_{\mathrm{op}}
    \le \kappa_n/4,\qquad
    \max_i|\widetilde Z_i|\le3\rho_n,\qquad
    \frac1n\sum_i|\widetilde Z_i|^3\le C .
\]
At \(t=0\), \(D_n(0)=1\), \(\nabla D_n(0)=0\), and
\[
    \nabla^2D_n(0)
    =
    \frac1n\sum_i\widetilde Z_i\widetilde Z_i^\top-(1+2\kappa_n)I_d
    \preceq-\frac74\kappa_nI_d .
\]
For \(|t|\le c_0\kappa_n\), uniformly over unit vectors \(v\),
\[
\begin{aligned}
    &\left|
    v^\top\{\nabla^2D_n(t)-\nabla^2D_n(0)\}v
    \right|  \\
    &\qquad\le
    C|t|\frac1n\sum_i(1+|\widetilde Z_i|^3)+C|t|^2
    \le Cc_0\kappa_n .
\end{aligned}
\]
The first inequality follows from the explicit Hessian formula and
\(\max_i|t\cdot \widetilde Z_i|=o(1)\).  Choosing \(c_0\) small enough makes the
Hessian at most \(-\kappa_n I_d\) throughout the ball.  Taylor's formula
then gives \(D_n(t)\le1-c\kappa_n|t|^2<1\) for \(0<|t|\le c_0\kappa_n\).
\end{proof}

Empirical centering is negligible relative to the deterministic bulk damping.
\begin{lemma}
\label{lem:bulk-empirical-centering}
Fix \(d\ge3\). At the sharp scale \(s_n=O(1)\), for every fixed \(c_0>0\),
\[
    \sup_{c_0\kappa_n\le |t|\le\rho_n/2}
    \frac{|t\cdot\bar Z_n|}
         {1-\exp\{-\kappa_n|t|^2\}}
    \xrightarrow{\Pbb}0 .
\]
\end{lemma}

\begin{proof}
Since the supremum over directions is \(|t|\,|\bar Z_n|\), it remains to
bound
\[
    \sup_{c_0\kappa_n\le r\le\rho_n/2}
    \frac{r|\bar Z_n|}{1-e^{-\kappa_nr^2}} .
\]
For \(r\le\kappa_n^{-1/2}\), the denominator is at least
\(c\kappa_nr^2\), and this part is bounded by
\(C|\bar Z_n|\kappa_n^{-2}\).  For \(r\ge\kappa_n^{-1/2}\), the
denominator is bounded below by a positive constant, and this part is
bounded by \(C\rho_n|\bar Z_n|\).  At the sharp scale
\(\kappa_n\asymp(\log\log n)/\log n\), while
\(|\bar Z_n|=O_{\Pbb}(n^{-1/2})\).  Hence both bounds are
\(o_{\Pbb}(1)\).
\end{proof}

The full empirical moment-generating function is likewise negligible relative
to the bulk damping.
\begin{lemma}
\label{lem:bulk-exponential-process}
Fix \(d\ge3\). At the sharp scale \(s_n=O(1)\), for every fixed \(c_0>0\),
\[
    \sup_{c_0\kappa_n\le |t|\le\rho_n/2}
    \frac{|H_n(t)-1|+|t\cdot\bar Z_n|}
         {1-\exp\{-\kappa_n|t|^2\}}
    \xrightarrow{\Pbb}0 .
\]
\end{lemma}

\begin{proof}
Lemma~\ref{lem:bulk-empirical-centering} proves the empirical-centering
part, so it remains to prove the same display with \(|H_n(t)-1|\) in the
numerator.  Write
\[
    m(r)=1-\exp\{-\kappa_nr^2\}.
\]

We first remove the tiny radii.  Put \(t=ru\), \(u\in S^{d-1}\), and
define, continuously at \(r=0\),
\[
    \Psi_{r,u}(z)
    =
    \frac{\exp\{ru\cdot z-r^2/2\}-1-ru\cdot z}{r^2}.
\]
Then \(\E\Psi_{r,u}(Z)=0\).  On the compact parameter set
\([0,1]\times S^{d-1}\), the first derivatives of \(\Psi_{r,u}\) are
bounded by \(A(z)=C_de^{|z|}(1+|z|^3)\), which has finite moments of all
orders.  Take an \(n^{-1/4}\)-net, whose size is \(O(n^{d/4})\).  Rosenthal's
inequality gives an \(O(n^{-q/2})\) \(q\)-moment bound for any fixed \(q>d\)
large enough, hence an \(O(n^{-q/4})\) tail at threshold \(n^{-1/4}\).  The
union bound is \(O(n^{(d-q)/4})=o(1)\).  The bound
\(n^{-1}\sum_iA(Z_i)=O_{\Pbb}(1)\) then extends the estimate from the net
to the continuum and gives
\[
    \sup_{0<|t|\le1}
    \frac{|H_n(t)-1-t\cdot\bar Z_n|}{|t|^2}
    =O_{\Pbb}(n^{-1/4}).
\]
Since \(m(r)\ge c\kappa_nr^2\) for \(r\le1\), uniformly on
\(c_0\kappa_n\le |t|\le1\),
\[
    \frac{|H_n(t)-1|}{m(|t|)}
    \le
    C\frac{|\bar Z_n|}{\kappa_n|t|}
    +C\frac{n^{-1/4}}{\kappa_n}
    =
    o_{\Pbb}(1).
\]
Here \(|t|\ge c_0\kappa_n\), \(|\bar Z_n|=O_{\Pbb}(n^{-1/2})\), and
\(\kappa_n\asymp(\log\log n)/\log n\).

It remains to control \(1\le |t|\le\rho_n/2\).  Let
\[
    1=x_0<x_1<\cdots <x_J=b_{n,d}/2,\qquad x_{j+1}\le2x_j,
\]
be the dyadic partition in squared radius, and set
\[
    \mathcal A_j=\{t:x_j\le |t|^2\le x_{j+1}\},\quad
    R_j=x_{j+1}^{1/2},\quad
    m_j=1-e^{-\kappa_nx_j}.
\]
Since \(m(|t|)\ge m_j\) on \(\mathcal A_j\), it is enough to show that
\[
    \max_j\frac1{m_j}\sup_{t\in\mathcal A_j}|H_n(t)-1|
    \xrightarrow{\Pbb}0 .
\]
Fix \(\eta>0\).  Put
\[
    \tau_n^2=20\log(1/\kappa_n)+20(d+3)\log\rho_n,\qquad
    B_j=\exp\{R_j^2/2+R_j\tau_n\},
\]
and write \(f_t(z)=\exp\{t\cdot z-|t|^2/2\}\), \(f_{t,j}^B=f_t\wedge B_j\).

The discarded lognormal tail is uniformly negligible.  For \(t\in
\mathcal A_j\),
\[
    \E f_t(Z)\mathbf 1_{\{f_t(Z)>B_j\}}
    \le \Pbb\{G>\tau_n\}
    =o(m_j),\qquad G\sim\cN(0,1).
\]
For the empirical tail, use
\[
    \sup_{|t|\le R}f_t(z)
    =
    \begin{cases}
        e^{|z|^2/2},& |z|\le R,\\
        e^{R|z|-R^2/2},& |z|>R.
    \end{cases}
\]
Because \(B_j>e^{R_j^2/2}\), the event
\(\sup_{|t|\le R_j}f_t(Z)>B_j\) implies \(|Z|>R_j+\tau_n\). Therefore
\[
    \E\sup_{|t|\le R_j}
    f_t(Z)\mathbf 1_{\{f_t(Z)>B_j\}}
    \le
    C_d(R_j+\tau_n)^{d-1}e^{-\tau_n^2/2}
    =o(m_j/J).
\]
Markov's inequality and a union bound over \(j\) give
\[
    \max_j\frac1{m_j}
    \sup_{t\in\mathcal A_j}
    \frac1n\sum_{i=1}^n
    f_t(Z_i)\mathbf 1_{\{f_t(Z_i)>B_j\}}
    \xrightarrow{\Pbb}0 .
\]

It remains to control the clipped process.  On the event
\(\max_i|Z_i|\le3\rho_n\), whose probability tends to one,
\[
    t\mapsto \frac1n\sum_i f_{t,j}^B(Z_i)
\]
is \(C\rho_nB_j\)-Lipschitz on \(\mathcal A_j\), and its expectation is
no less regular: a.e.\ in \(t\),
\[
    |\nabla_t \E f_{t,j}^B(Z)|
    \le
    \E\{|Z-t|f_t(Z)\}
    \le C_d .
\]
Take a Euclidean net of \(\mathcal A_j\) with mesh
\[
    \delta_j=\frac{\eta m_j}{16C\rho_nB_j}.
\]
Its cardinality satisfies
\[
    \log N_j
    \le C_d\{\log(\rho_n/\eta m_j)+\log B_j\}
    \le C_d\log n.
\]
At a fixed net point, Bernstein's inequality gives
\[
    \Pbb\left\{
        \left|
        \frac1n\sum_i\{f_{t,j}^B(Z_i)-\E f_{t,j}^B(Z)\}
        \right|>\eta m_j/4
    \right\}
    \le
    2\exp\left[
        -c\frac{n m_j^2}{e^{R_j^2}+B_jm_j}
    \right].
\]
Here
\[
    \operatorname{Var}\{f_{t,j}^B(Z)\}
    \le
    \E f_t(Z)^2
    =
    e^{|t|^2}
    \le e^{R_j^2},
    \qquad
    |f_{t,j}^B|\le B_j,
\]
which gives the displayed Bernstein denominator.
Here \(R_j^2\le b_{n,d}/2\), \(B_j\le n^{1/4+o(1)}\), and
\(m_j\ge c\kappa_n\).  Hence the exponent is at least
\[
    n^{1/2-o(1)}\kappa_n^2,
\]
which dominates \(\log N_j+\log J=O(\log n)\).  A union bound over all
net points and all \(j\), followed by the Lipschitz extension from the
net, proves
\[
    \max_j\frac1{m_j}
    \sup_{t\in\mathcal A_j}
    \left|
        \frac1n\sum_i\{f_{t,j}^B(Z_i)-\E f_{t,j}^B(Z)\}
    \right|
    \xrightarrow{\Pbb}0 .
\]
Combining the clipped concentration, the truncation bias, and the
empirical tail bound proves the dyadic estimate and the lemma.
\end{proof}

The two bulk estimates exclude all crossings with \(|t|\le\rho_n/2\).
\begin{proposition}
\label{prop:bulk-mgf-exclusion}
Fix \(d\ge3\). At the sharp scale \(s_n=O(1)\),
\[
    \Pbb\left\{\sup_{|t|\le\rho_n/2}D_n(t)>1\right\}\to0 .
\]
\end{proposition}

\begin{proof}
Lemma~\ref{lem:microscopic-bulk} handles \(0<|t|\le c_0\kappa_n\).
On the event in Lemma~\ref{lem:bulk-exponential-process}, put
\(g(t)=1-\exp\{-\kappa_n|t|^2\}\).  Uniformly on
\(c_0\kappa_n\le |t|\le\rho_n/2\), for a deterministic \(\delta_n\downarrow0\),
\[
    |H_n(t)-1|+|t\cdot\bar Z_n|\le\delta_n g(t).
\]
Therefore
\[
    D_n(t)
    =
    e^{-t\cdot\bar Z_n}e^{-\kappa_n|t|^2}H_n(t)
    \le
    e^{\delta_ng(t)}\{1-g(t)\}\{1+\delta_ng(t)\}.
\]
The right side is at most \(1-cg(t)\) for all large \(n\), uniformly in
this range.  Hence \(D_n(t)<1\) there with probability tending to one,
and the microscopic lemma completes the proof.
\end{proof}

\subsection{Sparse edges in growing dimension}
\label{sec:proof-program}

This subsection extends the spatial exclusion required by
Theorem~\ref{thm:main-support-law} from fixed to growing dimension below the
simplex regime. We split the sample into a sparse upper radial cloud and a
truncated lower cloud. Angular separation controls the upper cloud, while
dimension-uniform concentration controls the lower cloud.

Fix a sequence \(A_n\to\infty\) and define the upper radial cloud
\begin{equation}
\label{eq:upper-cloud}
  \mathcal I_n(A_n)
  =
  \{i:\xi_{n,i}\ge-A_n\}.
\end{equation}
The expected size of this set is \(e^{A_n+o(A_n)}\) whenever
\(A_n=o(L_n)\). Write \(B_n=B_n^{\rm hi}+B_n^{\rm lo}\) for the two sums
in \eqref{eq:bump-representation}, split according to
\(\mathcal I_n(A_n)\).

\subsubsection{Separation of the upper radial cloud}

Conditional on their radii, the uncentered Gaussian directions are
independent and uniform on \(S^{d_n-1}\). If
\(\#\mathcal I_n(A_n)=\exp\{O_{\Pbb}(A_n)\}\) and \(A_n=o(d_n)\), spherical
concentration gives
\begin{equation}
\label{eq:direction-separation}
  \max_{i\ne j\in\mathcal I_n(A_n)}
  \left|
    \frac{q_i}{|q_i|}\cdot\frac{q_j}{|q_j|}
  \right|
  =o_{\Pbb}(1).
\end{equation}
Since the upper centers have radius
\(\sqrt{2L_n}\{1+o_{\Pbb}(1)\}\),
\eqref{eq:direction-separation} implies
\[
  \min_{i\ne j\in\mathcal I_n(A_n)}|q_i-q_j|^2
  =
  4L_n\{1+o_{\Pbb}(1)\}.
\]
Thus no location \(w\) can lie within \(o(\sqrt{L_n})\) of two upper
centers. The contribution of all nonnearest upper bumps is exponentially
small in \(L_n\).

Centering by \(\bar Z_n\) perturbs every direction by
\(O_{\Pbb}(\sqrt{d_n/n})\) in Euclidean norm. Under
\eqref{eq:centering-condition} and the rate ranges considered below, this
does not affect \eqref{eq:direction-separation}.

\subsubsection{The lower cloud and the inner field}

At the critical normalization, every low bump has height at most
\begin{equation}
\label{eq:low-jump}
  \max_{i\notin\mathcal I_n(A_n)}h_i
  \le
  \exp\{-\delta_n(A_n+s_n)\}.
\end{equation}
Before empirical centering, let
\[
  B_n^\circ(w)
  =
  \frac1n\sum_{i=1}^n
  \exp\left\{
    w\cdot\frac{Z_i}{\sqrt{\gamma_n}}-\frac{|w|^2}{2}
  \right\}.
\]
For a fixed \(w\), its population mean is
\begin{equation}
\label{eq:population-field}
  \E B_n^\circ(w)
  =
  \exp\left\{-\frac{\eps_n|w|^2}{2}\right\}.
\end{equation}
Thus the deterministic background is already negligible near
\(|w|=\sqrt{2L_n}\) whenever \(d_n\to\infty\), because then
\(\eps_nL_n\to\infty\) at the critical scale.

The proof below establishes the following low-cloud estimate.  Its precision
\(o_{\Pbb}(\delta_n)\) is essential in superlogarithmic dimension: for a
deterministic radius \(r_n=o(\sqrt{L_n})\),
\begin{equation}
\label{eq:outer-low-target}
  \sup_{|w|\ge r_n}B_n^{\rm lo}(w)
  =o_{\Pbb}(\delta_n),
\end{equation}
provided \(A_n\) is chosen so that
\begin{equation}
\label{eq:A-conditions}
  A_n=o(L_n\wedge d_n),
  \qquad
  \delta_nA_n-\log(d_nL_n)\longrightarrow\infty.
\end{equation}

For fixed \(w\), \eqref{eq:low-jump}, \eqref{eq:population-field}, and
Bernstein's inequality give much more than \eqref{eq:outer-low-target}. The
work is to transfer that bound over \(w\). A direct Euclidean net of the
relevant ball has logarithmic size \(O(d_n\log L_n)\); the second condition
in \eqref{eq:A-conditions} is designed to make the bounded-jump exponent
dominate this entropy. Radial peeling avoids an unnecessarily fine net
near the origin.

The inner region contains the deterministic equality \(B_n(0)=1\), so it
cannot satisfy \eqref{eq:outer-low-target}.  The corresponding target is
strict concavity away from zero:
\begin{equation}
\label{eq:inner-target}
  \Pbb\left\{
    \sup_{0<|w|\le r_n}B_n(w)<1
  \right\}
  \longrightarrow1.
\end{equation}

At the origin,
\[
  \nabla B_n(0)=0,
  \qquad
  \nabla^2B_n(0)
  =
  \frac1n\sum_i
  \frac{\widetilde Z_i\widetilde Z_i^\top}{\gamma_n}
  -I_{d_n}.
\]
Its expectation is \(-\eps_nI_{d_n}+o(1)\). Standard covariance concentration
therefore gives a negative Hessian whenever
\(\sqrt{d_n/n}=o(\eps_n)\). The remaining task is to propagate this local
curvature through the interval up to \(r_n\) using a clipped exponential
process. This is easier than the outer-edge estimate because the population
gap in \eqref{eq:population-field} grows monotonically with \(|w|\).

Equations~\eqref{eq:outer-low-target} and \eqref{eq:inner-target} give the
field exclusion used later in
Proposition~\ref{prop:buffered-support-below-simplex}. On the event
\(\max_i\xi_{n,i}\le s_n-\eta\), every upper bump has height at most
\[
  e^{-\delta_n\eta},
\]
and hence lies below one by at least \(c_\eta\delta_n\) whenever
\(0<\delta_n\le1\).
At the center of any upper bump, the separation estimate makes all other
upper bumps \(o_{\Pbb}(\delta_n)\), and \eqref{eq:outer-low-target} makes the
low-cloud contribution \(o_{\Pbb}(\delta_n)\). Hence the outer field is below one, while
\eqref{eq:inner-target} handles the inner field.

\subsubsection{The direct truncation range}

The truncation and angular-separation conditions can be imposed
simultaneously in the range
\begin{equation}
\label{eq:direct-truncation-range}
  \log\log n\ll d_n,
  \qquad
  d_n\{\log d_n\}^2=o\{(\log n)^3\},
  \qquad
  d_n\log n=o(n).
\end{equation}
Indeed, in the worst part of this range,
\(\delta_n\asymp\sqrt{L_n/d_n}\). Conditions
\eqref{eq:A-conditions} admit a choice of \(A_n\) precisely when
\[
  \delta_n^{-1}\log(d_nL_n)=o(L_n\wedge d_n),
\]
which is implied by \eqref{eq:direct-truncation-range}. The lower bound
\(d_n\gg\log\log n\) ensures \(A_n=o(d_n)\), so the upper radial points are
uniformly separated. Smaller dimensions are handled below by constant-mesh
estimates. At the upper end, where \(\delta_n\) is too small for
\eqref{eq:A-conditions}, Section~\ref{sec:explicit-regime} replaces spatial
truncation by the entropy dual.

\subsubsection{Slowly growing dimension}
\label{sec:slow-dimension}

We first treat growing dimensions below \((\log L_n)^2\). A constant-mesh
argument is enough; its key input is an exact semiconvexity property of every
Gaussian-bump field. Fixed dimensions use the moving-shell estimates in
Section~\ref{sec:fixed-d-patches}.

\begin{lemma}
\label{lem:slow-truncated-shell}
Let
\[
 F(w)=\sum_{j=1}^m c_j
 \exp\left\{w\cdot q_j-\frac{|w|^2}{2}\right\},
 \qquad c_j>0.
\]
Then
\begin{equation}
\label{eq:semiconvex-log-field}
 \nabla^2\log F(w)=\operatorname{Cov}_{p_w}(q_j)-I\succeq-I,
\end{equation}
where \(p_w\) is proportional to
\(c_j e^{w\cdot q_j-|w|^2/2}\).  Every global maximizer \(w_*\) belongs
to the convex hull of \(q_1,\ldots,q_m\), and
\begin{equation}
\label{eq:semiconvex-net-transfer}
 F(w)\ge F(w_*)e^{-|w-w_*|^2/2}
 \qquad(w\in\R^d).
\end{equation}
\end{lemma}

\begin{proof}
Differentiating \(\log F\) gives \eqref{eq:semiconvex-log-field}.  At a
maximizer, \(0=\nabla\log F(w_*)=\sum_jp_{w_*}(j)q_j-w_*\), which proves
the convex-hull assertion.  Integrating the lower Hessian bound along the
segment from \(w_*\) to \(w\) proves
\eqref{eq:semiconvex-net-transfer}.
\end{proof}

Fix \(\alpha>0\), put
\[
 Q_i=\frac{Z_i}{\sqrt{\gamma_n}},\qquad
 R_{n,\alpha}=\sqrt{2(L_n-\alpha)},
\]
and define the independent truncated field
\[
 S_{n,\alpha}(w)
 =\frac1n\sum_{i=1}^n
 e^{w\cdot Q_i-|w|^2/2}
 \mathbf1_{\{|Q_i|\le R_{n,\alpha}\}}.
\]
The next lemma gives the only outer-field estimate needed below.  It uses
only bounded summands, their population mean, and a constant Euclidean net.

\begin{lemma}
\label{lem:slow-outside-shell}
Suppose
\[
 d_n\longrightarrow\infty,\qquad
 d_n\le(\log L_n)^2,\qquad s_n=O(1).
\]
For every fixed \(\alpha>0\), choose a fixed \(\zeta>0\) with
\(\zeta^2/2<\alpha\), let \(\mathcal G_n\) be a \(\zeta\)-net of
\(B(0,R_{n,\alpha})\), and fix
\[
 0<\rho<e^{-\zeta^2/2}-e^{-\alpha}.
\]
Then
\begin{equation}
\label{eq:slow-outside-shell}
 \Pbb\left\{
  \max_{\substack{w\in\mathcal G_n\\
       |w|\ge R_{n,\alpha}/2-\zeta}}
  S_{n,\alpha}(w)
  \ge e^{-\zeta^2/2}-\rho
 \right\}\longrightarrow0.
\end{equation}
\end{lemma}

\begin{proof}
For fixed \(w\), every summand of \(S_{n,\alpha}(w)\) is at most
\begin{equation}
\label{eq:slow-jump-simple}
 M_n(w)=\exp\left\{-\alpha
 -\frac{(R_{n,\alpha}-|w|)^2}{2}\right\},
\end{equation}
while
\begin{equation}
\label{eq:slow-mean-simple}
 \E S_{n,\alpha}(w)\le e^{-\eps_n|w|^2/2}.
\end{equation}
If independent nonnegative summands have sum of means \(\mu\), are bounded
by \(M\), and \(t>\mu\), the exponential-Chernoff bound is
\begin{equation}
\label{eq:bounded-poisson-chernoff}
 \Pbb\left\{\sum_jX_j\ge t\right\}
 \le
 \exp\left[-\frac1M
 \left\{t\log\frac t\mu-t+\mu\right\}\right].
\end{equation}
Apply this with \(t=e^{-\zeta^2/2}-\rho\).  Uniformly for
\(R_{n,\alpha}/2-\zeta\le|w|\le R_{n,\alpha}\),
\[
 -\log\Pbb\{S_{n,\alpha}(w)\ge t\}
 \ge(1-o(1))e^\alpha t\,\eps_nL_n.
\]
Indeed, the infimum of
\(e^{(R_{n,\alpha}-r)^2/2}r^2\) over this interval is
\((1-o(1))R_{n,\alpha}^2\).

In the present range, \(k_n=d_n/2\to\infty\) and
\begin{equation}
\label{eq:slow-entropy-match}
 \eps_nL_n
 =(1+o(1))k_n\log\frac{L_n}{k_n}
 =(1-o(1))k_n\log L_n.
\end{equation}
On the other hand,
\[
 \log|\mathcal G_n|
 \le d_n\log\frac{CR_{n,\alpha}}\zeta
 =k_n\log L_n+O(d_n).
\]
Since \(e^\alpha t>1\), the union bound proves
\eqref{eq:slow-outside-shell}.
\end{proof}

The population gap in \eqref{eq:population-field} controls the entire inner
half-ball once the
microscopic curvature estimate is combined with a relative empirical-process
bound.
\begin{lemma}
\label{lem:slow-bulk}
Suppose \(3\le d_n\le(\log L_n)^2\) and \(s_n=O(1)\).  Then
\begin{equation}
\label{eq:slow-bulk}
 \Pbb\left\{
  \sup_{0<|w|\le R_{n,\alpha}/2}B_n(w)>1
 \right\}\longrightarrow0.
\end{equation}
\end{lemma}

\begin{proof}
At the origin, \(B_n(0)=1\), \(\nabla B_n(0)=0\), and Gaussian
sample-covariance concentration gives
\[
 \nabla^2B_n(0)
 =\frac1n\sum_i\frac{\widetilde Z_i\widetilde Z_i^\top}{\gamma_n}-I
 \preceq-\frac{\eps_n}{2}I
\]
with probability tending to one.  Since
\(\max_i|\widetilde Z_i|=O_{\Pbb}(\sqrt{L_n})\), Taylor's theorem makes
the inequality strict on
\(0<|w|\le r_{0,n}:=\eps_n/(16L_n^{3/2})\).

For completeness, we give the empirical-process estimate used beyond this
microscopic ball.  Before centering, put
\[
 H_n(w)=\frac1n\sum_i
 \exp\left\{w\cdot Q_i-\frac{|w|^2}{2\gamma_n}\right\},
 \qquad \E H_n(w)=1,
\]
and \(g_n(r)=1-e^{-\eps_nr^2/2}\).  On
\(r_{0,n}\le|w|\le R_{n,\alpha}/2\),
\begin{equation}
\label{eq:slow-bulk-process}
 \sup_w\frac{|H_n(w)-1|}{g_n(|w|)}\longrightarrow_{\Pbb}0.
\end{equation}
To verify this, split the range into squared-radius annuli
\(\mathcal A_j=\{r_j\le |w|\le R_j\}\) with \(R_j^2\le2r_j^2\), and put
\(g_j=g_n(r_j)\).  There are \(O(\log L_n)\) annuli.  On an annulus with
outer radius \(R=R_j\), clip the lognormal summands at
\[
 B_R=\exp\{R^2/(2\gamma_n)+RT_n\},
 \qquad
 T_n^2=C\{d_n\log L_n+\log g_n(r_{0,n})^{-1}\}.
\]
We first control both discarded tails.  For a fixed \(w\), exponential
tilting gives
\[
 \sup_{|w|\le R}
 \E\!\left[
  e^{w\cdot Q-|w|^2/(2\gamma_n)}
  \mathbf1_{\{e^{w\cdot Q-|w|^2/(2\gamma_n)}>B_R\}}
 \right]
 \le e^{-cT_n^2}.
\]
Indeed, tilt \(Q\) from
\(\mathcal N(0,\gamma_n^{-1}I)\) to
\(\mathcal N(w/\gamma_n,\gamma_n^{-1}I)\); the threshold is then at least
\(\sqrt{\gamma_n}T_n\) standard deviations in the \(w\)-direction.  For the
empirical envelope, use
\[
 \sup_{|w|\le R}
 e^{w\cdot q-|w|^2/(2\gamma_n)}
 =
 \begin{cases}
  e^{\gamma_n|q|^2/2},&\gamma_n|q|\le R,\\
  e^{R|q|-R^2/(2\gamma_n)},&\gamma_n|q|>R.
 \end{cases}
\]
Thus exceeding \(B_R\) forces \(|q|>R/\gamma_n+T_n\).  Completing the
square against the \(\mathcal N(0,\gamma_n^{-1}I)\) density gives
\[
 \E\left[
  \sup_{|w|\le R}
  e^{w\cdot Q-|w|^2/(2\gamma_n)}
  \mathbf1_{\{\sup_{|v|\le R}
    e^{v\cdot Q-|v|^2/(2\gamma_n)}>B_R\}}
 \right]
 \le
 C e^{C d_n\log L_n}
 \int_{T_n}^{\infty}e^{-cu^2}\,du.
\]
Taking the constant in \(T_n^2\) sufficiently large makes this
\(o(g_j/\log L_n)\).  Markov's inequality and a union bound over the
annuli control the empirical discarded tails at the same order.

It remains to control the clipped fields.  Fix \(\rho>0\).  On the event
\(\max_i|Q_i|\le C\sqrt{L_n}\), each is
\(C\sqrt{L_n}B_R\)-Lipschitz.  Use a net on \(\mathcal A_j\) of mesh
\[
 \frac{\rho g_j}{C\sqrt{L_n}B_R}.
\]
Its logarithmic cardinality is \(O(d_nL_n)\).  Because
\(R^2\le L_n/2+O(1)\), one has \(B_R\le n^{1/4+o(1)}\), while Bernstein's
inequality at deviation \(\rho g_j/4\) has exponent at least
\[
 \frac{c n g_n(r)^2}{e^{R^2/\gamma_n}+B_Rg_n(r)}
 \ge n^{1/2-o(1)}g_n(r_{0,n})^2.
\]
This dominates \(d_nL_n\).  The union bound over all net points and all
annuli shows that the probability that the left side of
\eqref{eq:slow-bulk-process} exceeds \(\rho\) tends to zero.  Since
\(\rho>0\) was arbitrary, \eqref{eq:slow-bulk-process} follows.

Now
\[
 B_n(w)
 =e^{-w\cdot\bar Z_n/\sqrt{\gamma_n}}
   e^{-\eps_n|w|^2/2}H_n(w).
\]
Moreover,
\[
 \sup_{r_{0,n}\le r\le R_{n,\alpha}/2}
 \frac{r|\bar Z_n|/\sqrt{\gamma_n}}{g_n(r)}
 \le
 C\frac{|\bar Z_n|}{\eps_nr_{0,n}}
 =o_{\Pbb}(1),
\]
because \(d_n\le(\log L_n)^2\),
\(|\bar Z_n|=O_{\Pbb}(\sqrt{d_n/n})\), and
\(\eps_n\ge c\log L_n/L_n\).  Thus
\eqref{eq:slow-bulk-process} preserves a fixed fraction of the population
gap \(g_n(|w|)\), proving \eqref{eq:slow-bulk}.
\end{proof}


\subsubsection{Moderately growing dimension}
\label{sec:moderate-dimensional-theorem}

We now prove \eqref{eq:outer-low-target} and \eqref{eq:inner-target} under
\eqref{eq:direct-truncation-range}. The
argument deliberately uses a deeper truncation than
\eqref{eq:A-conditions}: this makes every remaining low-cloud summand so
small that elementary Bernstein bounds survive a Euclidean net.

Put
\[
  \mathfrak h_n=\log\frac{d_nL_n}{\delta_n\eps_n}.
\]
The expansions in Section~\ref{sec:regime-map} and
\eqref{eq:direct-truncation-range} imply
\begin{equation}
\label{eq:moderate-scale-relations}
  \mathfrak h_n\asymp\log L_n,
  \qquad
  \frac{\eps_nL_n}{\mathfrak h_n}\longrightarrow\infty,
  \qquad
  \log\frac{L_n^3}{\eps_n^3}=O(\log L_n).
\end{equation}
For \(d_n=o(L_n)\), these statements follow from
\(\delta_n=1+o(1)\) and
\(\eps_nL_n\sim(d_n/2)\log(2L_n/d_n)\); in particular,
\[
 \mathfrak h_n=2\log L_n-\log\log(2L_n/d_n)+O(1).
\]
When \(d_n\asymp L_n\), both \(\eps_n\) and \(\delta_n\) stay bounded away
from zero and infinity.  When \(d_n/L_n\to\infty\),
\(\eps_n\to1\) and \(\delta_n\asymp\sqrt{L_n/d_n}\).  The last two
conditions in \eqref{eq:direct-truncation-range} then give all three relations in
\eqref{eq:moderate-scale-relations}.
They also permit deterministic sequences \(\omega_n\to\infty\) and
\begin{equation}
\label{eq:moderate-cutoff}
  A_n=\frac{\omega_n\mathfrak h_n}{\delta_n}
  =o(L_n\wedge d_n).
\end{equation}
Indeed, this is immediate when \(d_n=O(L_n)\); when \(d_n/L_n\to\infty\),
it is equivalent, up to logarithmic factors, to
\(d_n(\log d_n)^2=o(L_n^3)\). Finally set
\begin{equation}
\label{eq:moderate-inner-radius}
  r_n^2=\frac{40\mathfrak h_n}{\eps_n}=o(L_n).
\end{equation}

We shall also use a uniform count for the upper radial cloud.  It is a
moderate-deviation extension of Lemma~\ref{lem:gamma-hazard}.

\begin{lemma}
\label{lem:moderate-upper-count}
Under \eqref{eq:direct-truncation-range}, if \(A_n\to\infty\) and
\(A_n=o(L_n\wedge d_n)\), then
\[
 n\Pbb\{Y_n\ge b_n-a_nA_n\}
 =\exp\{(1+o(1))A_n\}.
\]
Consequently, the number of indices with
\((|Z_i|^2/2-b_n)/a_n\ge-A_n\) is at most \(e^{2A_n}\) with probability
tending to one.
\end{lemma}

\begin{proof}
On \([b_n-a_nA_n,b_n]\), the reciprocal Gamma hazard equals
\(a_n\{1+o(1)\}\) uniformly.  Indeed,
\eqref{eq:tail-density-ratio} gives this once
\(a_nA_n=o(b_n-k_n+1)\).  For \(k_n=O(L_n)\), that relation follows from
\(A_n=o(d_n\wedge L_n)\); for \(k_n/L_n\to\infty\), it follows from
\(a_n\asymp\sqrt{k_n/L_n}\) and \(A_n=o(L_n)\).  Integrating the
hazard gives the displayed tail ratio.  The count is binomial with mean
\(e^{(1+o(1))A_n}\), so Markov's inequality proves the second claim.
\end{proof}

The next proposition proves the low-cloud estimate with more precision than
is needed for the Gumbel law. The stronger factor \(L_n^{-10}\) leaves enough
slack to differentiate the field up to three times in the support-count
argument.

\begin{proposition}
\label{prop:moderate-low-cloud}
Suppose \eqref{eq:direct-truncation-range} holds and \(s_n=O(1)\). Choose \(A_n,r_n\)
as in \eqref{eq:moderate-cutoff}--\eqref{eq:moderate-inner-radius}. Then
\begin{equation}
\label{eq:moderate-low-cloud}
  \sup_{|w|\ge r_n}B_n^{\rm lo}(w)
  =O_{\Pbb}(\delta_nL_n^{-10}).
\end{equation}
The same conclusion holds after differentiating the bump field up to three
times, with the right side replaced by \(o_{\Pbb}(\delta_nL_n^{-7})\).
\end{proposition}

\begin{proof}
We first work before empirical centering. Put
\(\xi_{n,i}^{\circ}=(|Z_i|^2/2-b_n)/a_n\) and define
\[
 S_n^{\rm lo}(w)
 =\frac1n\sum_{i=1}^n
   \exp\left\{w\cdot\frac{Z_i}{\sqrt{\gamma_n}}-\frac{|w|^2}{2}\right\}
   \mathbf1_{\{\xi_{n,i}^{\circ}<-A_n+1\}}.
\]
Completing the square and using
\(\gamma_nL_n=b_n+a_ns_n\) show that every summand is bounded by
\begin{equation}
\label{eq:moderate-jump-bound}
  \mathfrak j_n
  =\exp\{-\delta_n(A_n-1+s_n)\}
  =\exp\{-(1+o(1))\omega_n\mathfrak h_n\}.
\end{equation}
For fixed \(w\),
\begin{equation}
\label{eq:moderate-low-mean}
  \E S_n^{\rm lo}(w)
  \le \exp\{-\eps_n|w|^2/2\}.
\end{equation}
Since a nonnegative summand bounded by \(\mathfrak j_n\) has variance at most
\(\mathfrak j_n\) times its mean, Bernstein's inequality gives
\begin{equation}
\label{eq:moderate-bernstein}
 \Pbb\{S_n^{\rm lo}(w)-\E S_n^{\rm lo}(w)>t\}
 \le
 \exp\left\{-c\frac{t^2}
   {\mathfrak j_n(\E S_n^{\rm lo}(w)+t)}\right\}.
\end{equation}

With probability tending to one, all score centers have norm at most
\(C\sqrt{L_n}\). Every local maximum of the low bump field lies in their
convex hull, so it is enough to consider \(|w|\le C\sqrt{L_n}\). On this
ball the field is multiplicatively \(C\sqrt{L_n}\)-Lipschitz:
\begin{equation}
\label{eq:multiplicative-lipschitz}
  S_n^{\rm lo}(w')
  \ge e^{-C\sqrt{L_n}|w-w'|}S_n^{\rm lo}(w).
\end{equation}
Take a net of mesh \(c/\sqrt{L_n}\); its logarithmic cardinality is at most
\(C d_n\log L_n\). For \(|w|\ge r_n\),
\eqref{eq:moderate-low-mean} is at most \(e^{-20\mathfrak h_n}\). Apply
\eqref{eq:moderate-bernstein} with
\(t=\delta_nL_n^{-10}\). Equations
\eqref{eq:moderate-scale-relations} and \eqref{eq:moderate-jump-bound} give
\[
  \frac{t}{\mathfrak j_n}\gg d_n\log L_n,
  \qquad
  \frac{t^2}{\mathfrak j_ne^{-20\mathfrak h_n}}\gg d_n\log L_n.
\]
The union bound and \eqref{eq:multiplicative-lipschitz} prove the uncentered
version of \eqref{eq:moderate-low-cloud}.

The proof of Lemma~\ref{lem:centered-radial-max} gives
\[
 \max_i|\xi_{n,i}-\xi_{n,i}^{\circ}|=o_{\Pbb}(1).
\]
Moreover,
\[
 B_n(w)
 =\exp\left\{-w\cdot\frac{\bar Z_n}{\sqrt{\gamma_n}}\right\}
   B_n^{\circ}(w),
 \qquad
 \sup_{|w|\le C\sqrt{L_n}}
 \left|w\cdot\frac{\bar Z_n}{\sqrt{\gamma_n}}\right|
 =O_{\Pbb}(L_n/\sqrt n)=o_{\Pbb}(\delta_n).
\]
Replacing \(A_n-1\) by \(A_n\) therefore proves the centered statement.
For a multi-index \(\beta\) of order \(j\le3\), differentiation gives
\[
  \left|\partial_w^\beta e^{-|w-q_i|^2/2}\right|
  \le C_j(1+|w-q_i|^j)e^{-|w-q_i|^2/2}.
\]
On a fixed enlargement of the convex hull, both \(|w|\) and
\(\max_i|q_i|\) are \(O_{\Pbb}(\sqrt{L_n})\). Repeating the net-and-Bernstein
argument used for the order-zero field with this polynomial envelope costs at most
\(L_n^{j/2}\); the base estimate
\(O_{\Pbb}(\delta_nL_n^{-10})\) is therefore
\(o_{\Pbb}(\delta_nL_n^{-7})\) for every \(j\le3\). Outside that enlargement,
the polynomially weighted Gaussian envelope is radially decreasing toward the convex hull
once the enlargement radius exceeds \(\sqrt3\). Its supremum is consequently
attained on the controlled boundary, proving the derivative estimate globally.
\end{proof}

The inner field requires a different argument because it equals one at the
origin. Sample-covariance curvature handles a microscopic ball, and the
bounded-jump estimate \eqref{eq:moderate-bernstein} handles the remaining
annuli.

\begin{proposition}
\label{prop:moderate-inner-field}
Under the assumptions of Proposition~\ref{prop:moderate-low-cloud},
\begin{equation}
\label{eq:moderate-inner-field}
  \Pbb\left\{\sup_{0<|w|\le r_n}B_n(w)<1\right\}
  \longrightarrow1.
\end{equation}
\end{proposition}

\begin{proof}
At the origin,
\[
  \nabla B_n(0)=0,
  \qquad
  \nabla^2B_n(0)=\frac1n\sum_iq_iq_i^\top-I.
\]
The Gaussian sample-covariance bound and
\eqref{eq:direct-truncation-range} imply
\begin{equation}
\label{eq:moderate-origin-curvature}
  \left\|\frac1n\sum_iq_iq_i^\top-\frac1{\gamma_n}I\right\|_{\rm op}
  =o_{\Pbb}(\eps_n),
  \qquad
  \nabla^2B_n(0)\preceq-\frac{\eps_n}{2}I
\end{equation}
with probability tending to one. On the same event
\(\max_i|q_i|=O(\sqrt{L_n})\), so the Hessian is
\(O(L_n^{3/2})\)-Lipschitz near zero. It follows that
\begin{equation}
\label{eq:moderate-microscopic-ball}
  B_n(w)<1
  \quad\text{for}\quad
  0<|w|\le r_{0,n}:=\frac{\eps_n}{16L_n^{3/2}}.
\end{equation}

It remains to consider \(r_{0,n}\le|w|\le r_n\). Put
\(g(r)=1-e^{-\eps_nr^2/2}\). In this range
\[
  g(r)\ge c\eps_nr_{0,n}^2
  \ge c\frac{\eps_n^3}{L_n^3}=:g_{0,n}.
\]
The scale relations above give
\(\log g_{0,n}^{-1}=O(\log L_n)\). Cover each dyadic radial annulus by a
net of mesh \(c g(r)/\sqrt{L_n}\). The total logarithmic cardinality is
\(O(d_n\log L_n)\), and
\eqref{eq:multiplicative-lipschitz} shows that a crossing of level one would
produce a net point at level at least \(1-g(r)/8\).

For the low cloud, apply \eqref{eq:moderate-bernstein} with
\(t=g(r)/4\). Uniformly over all annuli, its exponent is bounded below by
\[
  c\frac{g_{0,n}^2}{\mathfrak j_n}\gg d_n\log L_n.
\]
Its mean is at most \(1-g(r)\) by
\eqref{eq:moderate-low-mean}. Lemma~\ref{lem:moderate-upper-count}, with a
one-unit buffer, shows that the upper cloud contains at most
\(e^{3A_n}\) points with probability tending to one. On
\(|w|\le r_n=o(\sqrt{L_n})\), each corresponding bump is at most
\(e^{-L_n+o(L_n)}\), so their total contribution is
\(e^{-L_n+o(L_n)}=o(g_{0,n})\). Finally, empirical centering changes the
logarithm of the field by at most
\[
  |w|\frac{|\bar Z_n|}{\sqrt{\gamma_n}},
  \qquad
  \sup_{r\ge r_{0,n}}
  \frac{|\bar Z_n|/\sqrt{\gamma_n}}{\eps_nr}
  =o_{\Pbb}(1),
\]
where the last relation follows from
\(|\bar Z_n|/\sqrt{\gamma_n}=O_{\Pbb}(\sqrt{L_n/n})\) and
\[
 \frac{L_n^2}{\sqrt n\,\eps_n^2}=o(1),
\]
which follows from \eqref{eq:direct-truncation-range}. The union bound now excludes every crossing outside
the microscopic ball. Together with
\eqref{eq:moderate-microscopic-ball}, this proves
\eqref{eq:moderate-inner-field}.
\end{proof}

Together with the upper-cloud separation established at the start of this
subsection, Propositions~\ref{prop:moderate-low-cloud} and
\ref{prop:moderate-inner-field} exclude all subthreshold contacts in the
direct-truncation range.

\subsubsection{The logarithmic crossover}
\label{sec:constant-crossover}

The slowly growing and direct-truncation arguments cover the sparse-edge
regimes except the crossover scale
\(d_n\asymp L_n^3/(\log L_n)^2\); the larger entropy regime is handled
in Section~\ref{sec:simplex}. Here a direct entropy comparison is
unnecessarily delicate.  A spatial argument is more effective: distinct
sample centers are separated by order \(\sqrt{L_n}\), whereas a neighborhood
of radius \(\sqrt{\log L_n}\) already makes every nonnearest Gaussian bump
negligible.

\begin{lemma}
\label{lem:crossover-geometry}
Fix \(0<c_-<c_+<\infty\), put \(\ell_n=\log L_n\), and suppose
\[
 c_-\frac{L_n^3}{\ell_n^2}
 \le d_n\le
 c_+\frac{L_n^3}{\ell_n^2},
 \qquad s_n=O(1).
\]
Then
\begin{equation}
\label{eq:crossover-scales}
 \gamma_n\asymp\frac{L_n^2}{\ell_n^2},
 \qquad
 \delta_n:=\frac{a_n}{\gamma_n}\asymp\frac{\ell_n}{L_n},
 \qquad
 \eps_n\longrightarrow1.
\end{equation}
Moreover, with probability tending to one,
\begin{equation}
\label{eq:crossover-separation}
 \max_i\frac{|Z_i|}{\sqrt{\gamma_n}}\le C\sqrt{L_n},
 \qquad
 \min_{i\ne j}|q_i-q_j|^2\ge3L_n,
 \qquad
 \frac{|\bar Z_n|}{\sqrt{\gamma_n}}
 =O_{\Pbb}\!\left(\sqrt{\frac{L_n}{n}}\right).
\end{equation}
\end{lemma}

\begin{proof}
The scale relations follow from \eqref{eq:superlog-regime} and
\eqref{eq:epsilon-superlog}.  Also
\(|\bar Z_n|=O_{\Pbb}(\sqrt{d_n/n})\), which gives the last estimate in
\eqref{eq:crossover-separation}.  A Gamma Chernoff bound and a union bound
over \(i\) give the first estimate there.  Finally,
\[
 q_i-q_j=\frac{Z_i-Z_j}{\sqrt{\gamma_n}},
 \qquad
 \frac{|Z_i-Z_j|^2}{2}\sim\chi^2_{d_n}.
\]
Since \(d_n/\gamma_n\sim2L_n\), the lower-tail Chernoff bound for
\(\chi^2_{d_n}\) is at most \(e^{-c d_n}\) at the threshold
\(|q_i-q_j|^2=3L_n\).  This remains summable over the fewer than \(n^2\)
pairs because \(d_n/L_n\to\infty\).
\end{proof}

We next use this separation together with truncation of the independent,
uncentered field.  The truncation is only needed away from the sample
centers, where each individual bump is already small.

\begin{proposition}
\label{prop:crossover-exclusion}
Under the assumptions of Lemma~\ref{lem:crossover-geometry}, for
every fixed \(\eta>0\),
\[
 \Pbb\left\{
  \max_i\xi_{n,i}\le s_n-\eta,\quad
  \sup_{w\in\mathbb R^{d_n}}B_n(w)>1
 \right\}\longrightarrow0.
\]
\end{proposition}

\begin{proof}
Work on the high-probability events in Lemma~\ref{lem:crossover-geometry}
and on
\(\{\max_i\xi_{n,i}\le s_n-\eta\}\).  Then
\[
 \max_i\frac{|q_i|^2}{2}\le L_n-g_n,
 \qquad
 g_n=\eta\frac{a_n}{\gamma_n}\asymp_\eta\frac{\ell_n}{L_n}.
\]
 Put
 \[
  Q_i=\frac{Z_i}{\sqrt{\gamma_n}},
  \qquad c_n=\frac{\bar Z_n}{\sqrt{\gamma_n}},
  \qquad
  S_n(w)=\frac1n\sum_i e^{w\cdot Q_i-|w|^2/2}.
 \]
 Then \(q_i=Q_i-c_n\) and, exactly,
 \begin{equation}
 \label{eq:crossover-centering-identity}
  B_n(w)=e^{-w\cdot c_n}S_n(w),
  \qquad
  \E S_n(w)=e^{-\eps_n|w|^2/2}.
 \end{equation}
 Lemma~\ref{lem:crossover-geometry} gives
 \begin{equation}
 \label{eq:crossover-centering-small}
  \sup_{|w|\le\sqrt{2L_n}}|w\cdot c_n|
 =O_{\Pbb}(L_n/\sqrt n)=o_{\Pbb}(g_n).
\end{equation}
In particular,
\[
 |c_n|=O_{\Pbb}\!\left(\sqrt{L_n/n}\right)
 =o_{\Pbb}(L_n^{-3/2}),
\]
which will make centering negligible even in the smallest annulus below.

 We first exclude neighborhoods of the sample centers.  Write
 \(R_n=\sqrt{8\ell_n}\).  The bump representation and the radial buffer
 give
 \[
  h_i=\frac1n e^{|q_i|^2/2}\le e^{-g_n}\le1-\frac{g_n}{2}
 \]
 for all large \(n\).  If \(|w-q_i|\le R_n\), pairwise separation gives
 \[
  B_n(w)
  \le e^{-g_n}
  +n\exp\left\{-\frac12(\sqrt{3L_n}-R_n)^2\right\}
  =1-\frac{g_n}{2}+e^{-L_n/2+o(L_n)}<1.
 \]
 Thus a crossing cannot occur within distance \(R_n\) of any \(q_i\).

We next exclude a fixed ball around the origin.  Standard Gaussian
sample-covariance concentration and \(d_n/n\to0\) imply
\[
 \frac1n\sum_iq_iq_i^\top\preceq\frac14I_{d_n}
\]
with probability tending to one.  Hence
\(\nabla^2B_n(0)\preceq-3I_{d_n}/4\).  On the radial-buffer event,
\(\max_i|q_i|\le\sqrt{2L_n}\), so the operator norm of the third derivative
of \(B_n\) is at most \(CL_n^{3/2}\) throughout a ball of radius
\(c_0L_n^{-3/2}\), where \(c_0>0\) is a sufficiently small universal
constant.  Taylor's theorem therefore gives
\begin{equation}
\label{eq:crossover-microscopic}
 B_n(w)\le1-\frac{|w|^2}{4},
 \qquad 0<|w|\le r_{0,n}:=c_0L_n^{-3/2}.
\end{equation}

It remains to pass from this microscopic ball to a fixed radius.  On the
event \(\max_i|Q_i|\le C\sqrt{L_n}\),
\begin{equation}
\label{eq:crossover-log-lipschitz}
 \sup_{|w|\le3}|\nabla\log S_n(w)|\le C\sqrt{L_n}.
\end{equation}
Partition \(\{r_{0,n}\le|w|\le3\}\) into the \(O(\ell_n)\) annuli
\(\mathcal A_j=\{r_j\le|w|\le2r_j\}\), where \(r_j=2^jr_{0,n}\),
truncating the last annulus at radius \(3\).  For each \(j\), take a
deterministic net of \(\mathcal A_j\) with mesh
\[
 \rho_j=c_1\frac{r_j^2}{\sqrt{L_n}},
\]
where \(c_1>0\) is a sufficiently small constant.  The logarithm of the
total number of net points is \(O(d_n\ell_n)=O(L_n^3/\ell_n)\).

For a fixed net point \(v\), define the independent variables
\[
 Y_i(v)=\min\left\{
  \frac1n e^{v\cdot Q_i-|v|^2/2},
  U_n
 \right\},
 \qquad
 U_n=n^{-1}e^{C_1\sqrt{L_n}},
\]
with \(C_1\) chosen so that truncation is inactive on
\(\{\max_i|Q_i|\le C\sqrt{L_n}\}\) for \(|v|\le3\).  Since
\(\eps_n\ge1/2\) for all large \(n\),
\[
 \sum_i\E Y_i(v)
 \le e^{-\eps_n|v|^2/2}
 \le1-c_2|v|^2,
 \qquad |v|\le3,
\]
for a numerical \(c_2>0\), while
\(\sum_i\operatorname{Var}Y_i(v)\le U_n\).  Bernstein's inequality gives
\[
 \Pbb\left\{\sum_iY_i(v)>
  e^{-\eps_n|v|^2/2}+c_3r_j^2\right\}
 \le\exp\{-c r_j^4/U_n\}.
\]
The exponent is at least \(n^{1-o(1)}L_n^{-6}\), uniformly in \(j\), and
therefore dominates the total net entropy.  Thus, with probability tending
to one, \(S_n(v)\le1-c_4r_j^2\) at every such net point.

If \(B_n(w)\ge1\) at some \(w\in\mathcal A_j\), then
\eqref{eq:crossover-centering-identity} and
\(|c_n|=o_{\Pbb}(r_{0,n})\) give \(S_n(w)\ge1-o(r_j^2)\).
At the nearest net point \(v\), \eqref{eq:crossover-log-lipschitz} and the
choice of \(c_1\) then give \(S_n(v)>1-c_4r_j^2\), a contradiction.  In
combination with \eqref{eq:crossover-microscopic}, this proves
\begin{equation}
\label{eq:crossover-inner}
 \sup_{0<|w|\le3}B_n(w)<1
\end{equation}
with probability tending to one on the radial-buffer event.

We finally exclude the part of the field outside the fixed ball and the
center neighborhoods.  Let \(\mathcal G_n\) be a deterministic
\(1/4\)-net of \(B(0,\sqrt{2L_n})\).  Then
\[
 \log|\mathcal G_n|
 \le d_n\log(C\sqrt{L_n})
 =O(L_n^3/\ell_n).
\]
For \(v\in\mathcal G_n\) with \(|v|\ge11/4\), set
\[
 \widetilde Y_i(v)
 =
 \min\left\{
  \frac1n e^{v\cdot Q_i-|v|^2/2},
  \mathfrak b_n
 \right\},
 \qquad \mathfrak b_n=2L_n^{-15/4}.
\]
These variables are independent, and
\[
 \sum_i\E\widetilde Y_i(v)
 \le e^{-\eps_n|v|^2/2}<0.05
\]
for all large \(n\).  Since
\(\sum_i\operatorname{Var}\widetilde Y_i(v)\le0.05\mathfrak b_n\), Bernstein's
inequality and a union bound give
\begin{equation}
\label{eq:crossover-outer-net}
 \Pbb\left\{
  \max_{\substack{v\in\mathcal G_n\\|v|\ge11/4}}
  \sum_i\widetilde Y_i(v)>0.9
 \right\}
 \le
 \exp\{O(L_n^3/\ell_n)-cL_n^{15/4}\}
 =o(1).
\end{equation}

Suppose now that \(\sup_wB_n(w)>1\), and let \(w_*\) be a global
maximizer.  Lemma~\ref{lem:slow-truncated-shell} places \(w_*\) in the
convex hull of the \(q_i\)'s, hence in \(B(0,\sqrt{2L_n})\).  The preceding
two parts of the proof imply
\[
 |w_*|>3,\qquad \min_i|w_*-q_i|>R_n.
\]
Let \(v\in\mathcal G_n\) satisfy \(|v-w_*|\le1/4\).  The semiconvex
estimate \eqref{eq:semiconvex-net-transfer} gives
\[
 B_n(v)\ge e^{-1/32}B_n(w_*)>e^{-1/32}.
\]
Moreover, \(|v|\ge11/4\) and
\(\min_i|v-q_i|\ge R_n-1/4\).  The exact identity
\[
 \frac1n e^{v\cdot Q_i-|v|^2/2}
 =e^{v\cdot c_n}h_i e^{-|v-q_i|^2/2},
\]
together with the radial-buffer bound on \(h_i\) and
\eqref{eq:crossover-centering-small}, shows that every uncentered summand
at \(v\) is at most
\[
 e^{|v\cdot c_n|}e^{-(R_n-1/4)^2/2}
 \le \mathfrak b_n.
\]
Thus truncation is inactive and, by
\eqref{eq:crossover-centering-identity},
\[
 \sum_i\widetilde Y_i(v)=S_n(v)
 =e^{v\cdot c_n}B_n(v)>0.9
\]
for all large \(n\), contradicting \eqref{eq:crossover-outer-net}.
\end{proof}

The exclusion propositions in this section cover, respectively, slowly
growing dimensions, the direct truncation range, and the remaining
crossover band. They are combined with the entropy estimate in
Proposition~\ref{prop:buffered-support-below-simplex}; no separate
regime-by-regime limit theorem is needed.


\section{Entropy and the Simplex Regime}
\label{sec:simplex}

This section completes the spatial-exclusion part of the
\(d_n\ge3\) analysis in Theorem~\ref{thm:main-support-law}. Once Euclidean
covering becomes inefficient, an entropy dual reduces the KKT optimization to
the nearly regular simplex formed by the centered sample. The resulting
estimates both rule out subthreshold contacts and identify the finite-sample
correction in the exact Gumbel coordinate \(s_n^\star\).

\subsection{Entropy reduction above the sparse-edge range}
\label{sec:explicit-regime}

The spatial truncation argument is no longer efficient once the dimension is
much larger than \((\log n)^3\). We use the entropy dual under
\begin{equation}
\label{eq:entropy-regime}
  \frac{d_n\{\log(d_n/L_n)\}^2}{L_n^3}\longrightarrow\infty,
  \qquad
  d_nL_n=o(n^2).
\end{equation}
This range includes
\(d_n=\lfloor(\log n)^4\rfloor\) and \(d_n\asymp n\). The entropy dual
\eqref{eq:entropy-dual}
provides a different route: away from the vertices of the probability
simplex, entropy beats the quadratic likelihood gain by a quantitative
margin; near a vertex, the corresponding radial observation controls the
optimization.

\subsubsection{Deterministic entropy inequalities}

We first record the entropy estimate used to separate the diffuse and
near-vertex regions of the simplex. Write \(H(p)=-\sum_i p_i\log p_i\), so that
\(K(p\Vert \mathbf u_n)=L_n-H(p)\), and let
\[
  h_{\rm bin}(y)=-y\log y-(1-y)\log(1-y)
\]
denote binary entropy, with the usual continuous values at zero and one.

\begin{lemma}
\label{lem:simplex-entropy}
There are universal constants \(c_0>0\) and \(n_0<\infty\) such that, for
all \(n\ge n_0\) and \(p\in\Delta_n\),
\begin{equation}
\label{eq:entropy-collision}
  K(p\Vert \mathbf u_n)
  \ge
  \frac{n-1}{n-2}\log(n-1)|p-\mathbf u_n|_2^2
  \ge
  L_n|p-\mathbf u_n|_2^2.
\end{equation}
Moreover, if
\(\max_i p_i\le1-L_n^{-2}\), then
\begin{equation}
\label{eq:entropy-stability}
  K(p\Vert \mathbf u_n)-L_n|p-\mathbf u_n|_2^2
  \ge
  \frac{c_0}{L_n}|p-\mathbf u_n|_1^2.
\end{equation}
\end{lemma}

\begin{proof}
Write
\[
  \mathfrak c_n=\frac{n-1}{n-2}\log(n-1),
  \qquad \beta_n=2\mathfrak c_n,
\]
and first consider
\[
  G_n(p)=K(p\Vert \mathbf u_n)-\mathfrak c_n|p-\mathbf u_n|_2^2.
\]
We identify the global minima of \(G_n\). A minimum cannot lie on the
boundary of the simplex: if \(p_j=0<p_i\), transferring mass \(t\) from
coordinate \(i\) to coordinate \(j\) changes the entropy part by
\(t\log t+O(t)\), whereas the quadratic part changes by \(O(t)\). The
total change is therefore negative for all sufficiently small \(t>0\).

At an interior minimum, the Lagrange equations are
\begin{equation}
\label{eq:potts-stationarity}
  \log p_i-\beta_np_i=\text{constant}.
\end{equation}
The function \(x\mapsto\log x-\beta_nx\) increases up to
\(1/\beta_n\) and decreases thereafter, so the coordinates of a stationary
point take at most two values. If these values are \(a>b\), then
\(b<1/\beta_n<a\). Moreover, the larger value has multiplicity one:
if two coordinates were equal to \(a\), the second variation in the
direction that transfers mass between them would be
\(2(1/a-\beta_n)<0\). Thus every nonuniform local minimum has the form
\[
  p(a)=\left(a,\frac{1-a}{n-1},\ldots,
                  \frac{1-a}{n-1}\right),
  \qquad \frac1n<a<1.
\]

Put \(x=(n-1)a/(1-a)\). Direct differentiation gives
\[
  \frac{d}{da}G_n(p(a))
  =g_n(x),
  \qquad
  g_n(x)=\log x-\beta_n\frac{x-1}{x+n-1}.
\]
The three zeros
\[
  x=1,\qquad x=n-1,\qquad x=(n-1)^2
\]
are immediate from the definition of \(\beta_n\). They are the only
zeros: indeed,
\[
  g_n'(x)=\frac1x-\frac{n\beta_n}{(x+n-1)^2}
\]
has at most two zeros. Also \(g_n'(1)>0\) and
\(g_n'((n-1)^2)>0\), because \(\beta_n<n\). For \(n=3\) the second
inequality is immediate; for \(n\ge4\) it follows from
\(\log m<m/2-1/(2m)\), \(m=n-1\): this holds at \(m=3\), and the
derivative of the right side minus the left side is
\((m-1)^2/(2m^2)>0\). Hence \(g_n\) is positive on \((1,n-1)\), negative on
\((n-1,(n-1)^2)\), and positive on \(((n-1)^2,\infty)\).
Now \(G_n(\mathbf u_n)=0\), while at \(x=(n-1)^2\),
\[
  p(a)=\left(\frac{n-1}{n},
      \frac1{n(n-1)},\ldots,\frac1{n(n-1)}\right)
  \quad\text{and}\quad G_n(p(a))=0.
\]
The preceding sign pattern therefore gives \(G_n(p(a))\ge0\) for every
\(a\in[1/n,1]\). Since every global minimum of \(G_n\) is among the
points just considered, this proves
\(K(p\Vert\mathbf u_n)\ge \mathfrak c_n|p-\mathbf u_n|_2^2\). Finally,
\(\mathfrak c_n\ge\log n=L_n\) for \(n\ge3\). For example, after writing
\(m=n-1\), this is equivalent to
\(r(m):=m\log m-(m-1)\log(m+1)\ge0\). Here \(r(2)>0\), and
\[
  r'(m)=\log\frac{m}{m+1}+\frac2{m+1}>0
  \qquad(m\ge2),
\]
because \(-\log(1-t)<t/(1-t)<2t\) for
\(t=1/(m+1)\). This proves \eqref{eq:entropy-collision}.

We next prove the quantitative gap from the vertices. In the remainder of
the proof take \(L=L_n\ge16\) and set
\[
  F_n(p)=K(p\Vert\mathbf u_n)-L|p-\mathbf u_n|_2^2,
  \qquad T(p)=|p-\mathbf u_n|_1.
\]
We shall prove
\begin{equation}
\label{eq:entropy-stability-explicit}
  F_n(p)\ge\frac1{32L}T(p)^2
  \qquad\text{if}\qquad
  \max_i p_i\le1-L^{-2}.
\end{equation}

We first dispose of diffuse vectors. The elementary scalar inequality
\[
  x\log x-x+1\ge\frac{(x-1)^2}{2(x+1)},
  \qquad x\ge0,
\]
follows by differentiation. Applied coordinatewise, it gives
\[
  K(p\Vert\mathbf u_n)\ge\frac12 A(p),
  \qquad
  A(p)=\sum_i\frac{(p_i-1/n)^2}{p_i+1/n}.
\]
If \(\max_i p_i\le(8L)^{-1}\), then
\(p_i+1/n\le(4L)^{-1}\) for all \(i\), once \(L\ge16\). Consequently,
\[
  L|p-\mathbf u_n|_2^2\le\frac14A(p),
  \qquad
  T(p)^2\le A(p)\sum_i(p_i+1/n)=2A(p),
\]
and hence \(F_n(p)\ge T(p)^2/8\), which is stronger than
\eqref{eq:entropy-stability-explicit}.

It remains to consider \(m=\max_i p_i\ge(8L)^{-1}\). Relabel so that
\(p_1=m\), put \(y=1-m\), and write
\[
  p=(m,yq_1,\ldots,yq_{n-1}),
  \qquad q\in\Delta_{n-1}.
\]
Let \(\bar p_m=(m,y/(n-1),\ldots,y/(n-1))\), and define
\[
  f_n(m)=F_n(\bar p_m),
  \qquad t_0=|\bar p_m-\mathbf u_n|_1=2(m-1/n).
\]
The entropy chain rule and the orthogonal decomposition of the squared
Euclidean distance give the exact identity
\begin{equation}
\label{eq:largest-coordinate-decomposition}
  F_n(p)
  =f_n(m)+y\left\{
     K(q\Vert\mathbf u_{n-1})-Ly|q-\mathbf u_{n-1}|_2^2
   \right\}.
\end{equation}
The first part of the lemma in dimension \(n-1\) gives
\[
  K(q\Vert\mathbf u_{n-1})\ge C_{n-1}|q-\mathbf u_{n-1}|_2^2,
  \qquad
  C_{n-1}=\frac{n-2}{n-3}\log(n-2)\ge L.
\]
The last inequality holds for \(n\ge4\): with \(x=n-2\), it is equivalent
to \(x\log x\ge(x-1)\log(x+2)\); the derivative of the difference is
\(\log(1-t)+3t/2\), where \(t=2/(x+2)\in(0,1/2]\). This derivative is
nonnegative because \(-\log(1-t)/t\) is increasing and its value at
\(t=1/2\) is \(2\log2<3/2\).
Pinsker's inequality then yields
\begin{equation}
\label{eq:conditional-tail-stability}
  K(q\Vert\mathbf u_{n-1})-Ly|q-\mathbf u_{n-1}|_2^2
  \ge
  \left(1-\frac{Ly}{C_{n-1}}\right)K(q\Vert\mathbf u_{n-1})
  \ge\frac m2|q-\mathbf u_{n-1}|_1^2.
\end{equation}
Here Pinsker's inequality with natural logarithms follows by grouping the
positive and negative coordinate deviations and applying the binary bound
whose second derivative is at least four.

It remains only to bound \(f_n\). Put
\[
  x=\frac{(n-1)m}{1-m},
  \qquad
  \widetilde g_n(x)=\log x-2L\frac{x-1}{x+n-1}.
\]
Then \(f_n'(m)=\widetilde g_n(x)\), while
\[
  \widetilde g_n'(x)=\frac1x-\frac{2nL}{(x+n-1)^2}
\]
has at most two zeros. Thus \(\widetilde g_n\) has at most three zeros on
\((0,\infty)\), one of which is \(x=1\). Moreover,
\(\widetilde g_n'(1)=1-2L/n>0\). At the left endpoint
\(m=(8L)^{-1}\), we have \(x_0=(n-1)/(8L-1)\) and, since
\[
  2L\frac{x_0-1}{x_0+n-1}=\frac{n-8L}{4(n-1)}<\frac14,
\]
\[
  \widetilde g_n(x_0)
  >L-\log(16L)-\frac14>0.
\]
At the right endpoint
\(M=1-L^{-2}\), where \(x=(n-1)(L^2-1)\),
\begin{align*}
  \widetilde g_n(x)
  &=-L+2\log L+\log(1-1/n)+\log(1-L^{-2})
       +\frac{2n}{(n-1)L}\\
  &\le-L+2\log L+\frac4L<0.
\end{align*}
Thus the derivative of \(f_n\) is positive at the left endpoint and
negative at the right endpoint. If \(f_n\) had an interior local minimum,
its derivative would have to cross first from positive to negative, then
from negative to positive at that minimum, and finally back to negative
before \(M\). Together with the zero at \(x=1\), this would give at least
four zeros of \(\widetilde g_n\), a contradiction. Therefore the minimum
of \(f_n\) on this interval occurs at an endpoint.

At \(m=(8L)^{-1}\), using
\(-(1-m)\log(1-m)\le m\) and
\(|\bar p_m-\mathbf u_n|_2^2\le2m^2\), we obtain
\[
  f_n(m)
  \ge m\{L-\log(8L)-1\}-2Lm^2
  =m\{L-\log(8L)-5/4\}
  \ge\frac1{16}.
\]
At \(m=M\), put \(y=L^{-2}\). Since
\(|\bar p_m-\mathbf u_n|_2^2\le m^2\) and
\(h_{\rm bin}(y)\le y\log(e/y)\),
\begin{align*}
  f_n(M)
  &\ge Ly(1-y)-h_{\rm bin}(y)\\
  &\ge \frac1L-\frac1{L^3}-\frac{1+2\log L}{L^2}
  \ge\frac1{2L}.
\end{align*}
All numerical inequalities in the last three displays hold for \(L\ge16\).
It follows that, throughout \([(8L)^{-1},M]\),
\begin{equation}
\label{eq:scalar-endpoint-stability}
  f_n(m)\ge\frac1{2L}\ge\frac1{8L}t_0^2.
\end{equation}

Combining \eqref{eq:largest-coordinate-decomposition},
\eqref{eq:conditional-tail-stability}, and
\eqref{eq:scalar-endpoint-stability}, and using
\(m\ge(8L)^{-1}\), gives
\[
  F_n(p)
  \ge\frac{t_0^2}{8L}
      +\frac{y^2}{16L}|q-\mathbf u_{n-1}|_1^2.
\]
Finally,
\(T(p)\le t_0+y|q-\mathbf u_{n-1}|_1\), so the last display implies
\eqref{eq:entropy-stability-explicit}. Thus
\eqref{eq:entropy-stability} holds with \(c_0=1/32\) and, for example,
\(n_0=\lceil e^{16}\rceil\).
\end{proof}

\begin{remark}
The sharp constant in \eqref{eq:entropy-collision} is the mean-field Potts
transition coefficient. The proof also identifies all equality cases:
the uniform distribution and the \(n\) permutations of
\(( (n-1)/n,1/[n(n-1)],\ldots,1/[n(n-1)] )\).
\end{remark}

The next proposition converts the entropy estimate into a deterministic
one-atom criterion. It is tailored to a radial gap much smaller than the
largest off-diagonal inner product.

\begin{proposition}
\label{prop:deterministic-gram}
Let \(q_1,\ldots,q_n\in\R^d\) satisfy \(\sum_iq_i=0\). Put
\[
  \mu_n^{\rm off}=\max_{i\ne j}|q_i\cdot q_j|.
\]
Suppose, for some \(0<g_n\le1\),
\begin{equation}
\label{eq:deterministic-gram-assumptions}
  \max_i\frac{|q_i|^2}{2}\le L_n-g_n,
  \qquad
  \mu_n^{\rm off}=o(L_n^{-1}),
  \qquad
  ng_n\longrightarrow\infty.
\end{equation}
Then, for all sufficiently large \(n\),
\[
  \sup_{w\in\R^d}B_n(w)\le1.
\]
\end{proposition}

\begin{proof}
Fix \(p\in\Delta_n\). Since \(\sum_iq_i=0\), writing
\(r=p-\mathbf u_n\) gives
\[
  \frac12\left|\sum_i p_iq_i\right|^2
  =
  \frac12\left|\sum_i r_iq_i\right|^2
  \le
  L_n|r|_2^2+\frac{\mu_n^{\rm off}}{2}|r|_1^2.
\]
If \(\max_i p_i\le1-L_n^{-2}\), Lemma~\ref{lem:simplex-entropy} and
\(\mu_n^{\rm off}=o(L_n^{-1})\) show that this is at most \(K(p\Vert \mathbf u_n)\).

It remains to consider a distribution with
\(\max_i p_i=1-y>1-L_n^{-2}\). Let \(i\) be its largest coordinate and
write
\[
  \sum_jp_jq_j=(1-y)q_i+y\bar q,
\]
where \(\bar q\) is a convex combination of the remaining \(q_j\)'s.
Convexity of the squared norm and
\eqref{eq:deterministic-gram-assumptions} give
\[
  \frac{|\bar q|^2}{2}\le L_n-g_n,
  \qquad
  q_i\cdot\bar q\le \mu_n^{\rm off}.
\]
Consequently,
\begin{equation}
\label{eq:near-vertex-energy}
  \frac12\left|\sum_jp_jq_j\right|^2
  \le
  (L_n-g_n)\{(1-y)^2+y^2\}
  +\mu_n^{\rm off}y(1-y).
\end{equation}
On the other hand,
\[
  H(p)\le h_{\rm bin}(y)+y\log(n-1)\le h_{\rm bin}(y)+yL_n.
\]
Subtracting \eqref{eq:near-vertex-energy} from
\(K(p\Vert \mathbf u_n)=L_n-H(p)\), and using \(y\le L_n^{-2}\), gives
\begin{align*}
  K(p\Vert \mathbf u_n)
  -\frac12\left|\sum_jp_jq_j\right|^2
  &\ge
  g_n+\{2(L_n-g_n)-\mu_n^{\rm off}\}y(1-y)-h_{\rm bin}(y)-yL_n\\
  &\ge
  g_n+y\{L_n-4+\log y\}
\end{align*}
for all large \(n\), because \(h_{\rm bin}(y)\le y\log(e/y)\).
The minimum of \(y\mapsto y(L_n-4+\log y)\) over \(y\ge0\) is
\(-e^3/n\). The last display is therefore positive by \(ng_n\to\infty\).
The entropy dual in Lemma~\ref{lem:entropy-dual} completes the proof.
\end{proof}

In very high dimension, the absolute off-diagonal bound in
Proposition~\ref{prop:deterministic-gram} discards the regular-simplex
geometry created by centering.  The next criterion retains it.  The
quantity \(C_n\) measures the deviation of
each row of the Gram matrix from the identity
\(q_i\cdot\{(n-1)^{-1}\sum_{j\ne i}q_j\}=-|q_i|^2/(n-1)\).

\begin{proposition}
\label{prop:local-simplex-criterion}
Let \(q_1,\ldots,q_n\in\mathbb R^d\) satisfy \(\sum_iq_i=0\), and put
\[
 G_n=(q_i\cdot q_j)_{i,j\le n},
 \qquad
 \mu_n^{\rm off}=\max_{i\ne j}|q_i\cdot q_j|,
 \qquad
 C_n=\max_{i\ne j}
 \left|q_i\cdot q_j+\frac{|q_i|^2}{n-1}\right|.
\]
Suppose, for some \(g_n>0\), that
\[
 \max_i\frac{|q_i|^2}{2}\le R_n^\star-g_n,
 \qquad
 \mu_n^{\rm off}=o(L_n^{-1}),
 \qquad
 \|G_n\|_{\rm op}=O(L_n),
 \qquad
 C_n=o(g_nL_n^2).
\]
Then, for all sufficiently large \(n\),
\[
 \sup_{w\in\mathbb R^d}B_n(w)\le1.
\]
\end{proposition}

\begin{proof}
Fix \(p\in\Delta_n\).  If
\(\max_i p_i\le1-L_n^{-2}\), the argument away from the vertices in
Proposition~\ref{prop:deterministic-gram} applies: since
\(R_n^\star<L_n\), Lemma~\ref{lem:simplex-entropy} and
\(\mu_n^{\rm off}=o(L_n^{-1})\) give
\[
 \frac12\left|\sum_i p_iq_i\right|^2\le K(p\Vert \mathbf u_n).
\]

Suppose instead that \(p_i=1-y>1-L_n^{-2}\).  Write
\(p_j=yr_j\) for \(j\ne i\), where
\(r=(r_j)_{j\ne i}\in\Delta_{n-1}\), and let \(\mathbf u_{n-1}\) denote the
uniform distribution on these \(n-1\) coordinates.  Set
\[
 c_y=1-\frac{ny}{n-1},
 \qquad
 v_r=\sum_{j\ne i}\left(r_j-\frac1{n-1}\right)q_j.
\]
Since \(\sum_{j\ne i}q_j=-q_i\),
\[
 \sum_jp_jq_j=c_yq_i+yv_r.
\]
The relative entropy separates as
\[
 K(p\Vert \mathbf u_n)
 =
 k_n(y)+yK(r\Vert \mathbf u_{n-1}),
\]
where
\[
 k_n(y)
 =(1-y)\log\{n(1-y)\}
 +y\log\frac{ny}{n-1}.
\]
Apply the sharp inequality \eqref{eq:entropy-collision} to the vector
\((1-y,y/(n-1),\ldots,y/(n-1))\).  Since its squared distance from
\(\mathbf u_n\) is \((n-1)c_y^2/n\), we obtain
\begin{equation}
\label{eq:potts-scalar-profile}
 k_n(y)\ge R_n^\star c_y^2.
\end{equation}

The definition of \(C_n\) gives
\[
 |q_i\cdot v_r|
 =
 \left|
  \sum_{j\ne i}r_j
  \left(q_i\cdot q_j+\frac{|q_i|^2}{n-1}\right)
 \right|
 \le C_n.
\]
Moreover,
\[
 |v_r|^2
 \le\|G_n\|_{\rm op}|r-\mathbf u_{n-1}|_2^2
 \le\frac{\|G_n\|_{\rm op}}{\log(n-1)}
       K(r\Vert \mathbf u_{n-1}),
\]
where the last inequality is \eqref{eq:entropy-collision} with \(n-1\)
in place of \(n\).  Because \(y\le L_n^{-2}\) and
\(\|G_n\|_{\rm op}=O(L_n)\), it follows that
\[
 K(r\Vert \mathbf u_{n-1})-\frac y2|v_r|^2\ge0
\]
for all large \(n\).  Combining this with
\eqref{eq:potts-scalar-profile} and
\(|q_i|^2/2\le R_n^\star-g_n\) gives
\begin{align*}
 K(p\Vert \mathbf u_n)
 -\frac12\left|\sum_jp_jq_j\right|^2
 &\ge
 g_nc_y^2-y|c_y|C_n\\
 &\ge
 g_n\{1-O(L_n^{-2})\}-\frac{C_n}{L_n^2}>0.
\end{align*}
The entropy dual in Lemma~\ref{lem:entropy-dual} completes the proof.
\end{proof}

The local-simplex criterion applies when \(d_nL_n/n^2\) is bounded below.
When \(d_nL_n=o(n^2)\), the following criterion is more convenient: it
separates diffuse probability vectors, controlled spectrally, from their
heavy coordinates, controlled by entropy and \(\mu_n^{\rm off}\).

\begin{proposition}
\label{prop:refined-deterministic-gram}
Let \(q_1,\ldots,q_n\in\mathbb R^d\) satisfy \(\sum_iq_i=0\), and put
\[
 G=(q_i\cdot q_j)_{i,j\le n},
 \qquad
 \upsilon_n=\frac{\|G\|_{\mathrm{op}}}{2n},
 \qquad
 \mu_n^{\rm off}=\max_{i\ne j}|q_i\cdot q_j|.
\]
Suppose \(\max_i|q_i|^2/2\le L_n-g_n\), where \(0<g_n\le1\).  Assume that
there are deterministic \(a_n^\circ\to\infty\) and \(y_n\downarrow0\),
with \(T_n=e^{a_n^\circ}\), such that
\begin{equation}
\label{eq:refined-gram-conditions}
 a_n^\circ=o(L_n),\qquad
 \frac{\upsilon_nT_n}{a_n^\circ}=o(1),\qquad
 \frac{\mu_n^{\rm off}+L_nT_n/n}{a_n^\circ}=o(1),\qquad
 \frac{L_ny_n}{a_n^\circ}\longrightarrow\infty,
\end{equation}
and
\begin{equation}
\label{eq:refined-vertex-condition}
 e^{2L_ny_n+\mu_n^{\rm off}}=o(ng_n).
\end{equation}
Then, for all sufficiently large \(n\),
\[
 \sup_{w\in\mathbb R^d}B_n(w)\le1.
\]
\end{proposition}

\begin{proof}
Fix \(p\in\Delta_n\), write \(f_i=np_i\), and let
\(\phi(x)=x\log x-x+1\).  Then
\[
 K(p\Vert \mathbf u_n)=\frac1n\sum_i\phi(f_i).
\]
Set
\[
 \ell_i=f_i\wedge T_n,\qquad h_i=(f_i-T_n)_+,\qquad
 \bar\ell=\frac1n\sum_i\ell_i.
\]
Since \(\sum_iq_i=0\),
\[
 \sum_i p_iq_i
 =
 \frac1n\sum_i(\ell_i-\bar\ell)q_i
 +\frac1n\sum_i h_iq_i
 =:v_0+v_1.
\]
The elementary inequality
\[
 (x-1)^2\le C\frac{T_n}{a_n^\circ}\phi(x),
 \qquad 0\le x\le T_n,
\]
and the definition of \(\upsilon_n\) give, with
\(K_0=n^{-1}\sum_i\phi(\ell_i)\),
\begin{equation}
\label{eq:light-energy}
 \frac{|v_0|^2}{2}
 \le C\frac{\upsilon_nT_n}{a_n^\circ}K_0=o(K_0).
\end{equation}

Put \(b_i=h_i/n\) and \(z=\sum_i b_i\).  The diagonal and off-diagonal
parts of the Gram matrix give
\begin{equation}
\label{eq:heavy-energy}
 \frac{|v_1|^2}{2}
 \le L_n\sum_i b_i^2+\frac{\mu_n^{\rm off}}{2}z^2,
 \qquad
 |v_0\cdot v_1|
 \le\left(2\mu_n^{\rm off}+\frac{2L_nT_n}{n}\right)z.
\end{equation}
Suppose first that \(\max_i p_i\le1-y_n\).  Uniformly for
\(T_n\le x\le n(1-y_n)\), one-variable calculus gives
\begin{equation}
\label{eq:heavy-entropy-scalar}
 \phi(x)-\phi(T_n)
 \ge\frac{L_n}{n}(x-T_n)^2
 +\frac{a_n^\circ}{4}(x-T_n).
\end{equation}
Indeed, after subtracting the right side, the derivative first increases
and then decreases, and is positive at \(T_n\).  The minimum is therefore
at an endpoint. At \(x_n=n(1-y_n)\), the difference between the two sides is
\begin{align*}
 &x_n\left\{L_ny_n+\log(1-y_n)-1-\frac{a_n^\circ}{4}\right\}\\
 &\quad+2L_n(1-y_n)T_n-\frac{3a_n^\circ T_n}{4}
       +T_n-\frac{L_nT_n^2}{n}.
\end{align*}
The first line is positive and of order \(nL_ny_n\), while the absolute
value of the second line is \(o(nL_ny_n)\) by
\eqref{eq:refined-gram-conditions}. Thus the upper-endpoint value is
positive. Summing
\eqref{eq:heavy-entropy-scalar} yields
\[
 K(p\Vert \mathbf u_n)-K_0
 \ge L_n\sum_i b_i^2+\frac{a_n^\circ}{4}z.
\]
Equations \eqref{eq:refined-gram-conditions},
\eqref{eq:light-energy}, and \eqref{eq:heavy-energy} now show that
\(|v_0+v_1|^2/2\le K(p\Vert \mathbf u_n)\): indeed, \(z\le1\), and the last two
conditions in \eqref{eq:refined-gram-conditions} make the off-diagonal and
cross terms in \eqref{eq:heavy-energy} at most
\(a_n^\circ z/4\), while \eqref{eq:light-energy} is at most \(K_0\).

It remains to consider \(p_i=1-y>1-y_n\) for some \(i\).  If \(\bar q\)
is the barycenter of the remaining \(q_j\)'s under their normalized
weights, convexity gives
\[
 \frac12|(1-y)q_i+y\bar q|^2
 \le(L_n-g_n)\{(1-y)^2+y^2\}+\mu_n^{\rm off}y(1-y).
\]
On the other hand,
\(K(p\Vert \mathbf u_n)\ge(1-y)L_n-h_{\rm bin}(y)\). Their difference is at least
\[
 \frac{g_n}{2}
 +y\{L_n-2L_ny-\mu_n^{\rm off}+\log y-1\}.
\]
The negative part of the second term is at most
\(\exp\{-L_n+2L_ny_n+\mu_n^{\rm off}+1\}=o(g_n)\) by
\eqref{eq:refined-vertex-condition}.  Thus the difference is positive.
The entropy dual in Lemma~\ref{lem:entropy-dual} completes the proof.
\end{proof}

\subsubsection{The random Gram matrix}

Under \eqref{eq:entropy-regime}, the next lemma verifies the off-diagonal and
spectral Gram assumptions of
Proposition~\ref{prop:refined-deterministic-gram}.

\begin{lemma}
\label{lem:random-off-diagonal}
Suppose the first condition in \eqref{eq:entropy-regime} holds and
\[
 \frac{\gamma_nL_n}{\sigma_n^2b_n}\longrightarrow1.
\]
For
\[
  q_i=\frac{Z_i-\bar Z_n}{\sqrt{\gamma_n}},
  \qquad
  \mu_n^{\rm off}=\max_{i\ne j}|q_i\cdot q_j|,
  \qquad
  C_n=\max_{i\ne j}
  \left|q_i\cdot q_j+\frac{|q_i|^2}{n-1}\right|,
  \qquad
  \upsilon_n=\frac{\|(q_i\cdot q_j)_{i,j\le n}\|_{\rm op}}{2n},
\]
one has
\begin{equation}
\label{eq:B-rate}
  \mu_n^{\rm off}
  =
  O_{\Pbb}\left\{
    \frac{L_n}{n}
    +\frac{L_n^{3/2}}{\sqrt{d_n}}
    +\frac{L_n^2}{d_n}
  \right\},
  \qquad
  \upsilon_n
  =
  O_{\Pbb}\left(\frac{L_n}{d_n}+\frac{L_n}{n}\right).
\end{equation}
Moreover,
\begin{equation}
\label{eq:delta-entropy-regime}
  \delta_n^\star:=\frac{\sigma_n^2a_n}{\gamma_n}
  =
  \sqrt{\frac{L_n}{d_n}}\{1+o(1)\}.
\end{equation}
In addition,
\begin{equation}
\label{eq:centered-row-gram-rate}
 C_n=O_{\Pbb}\left(\frac{L_n^{3/2}}{\sqrt{d_n}}\right).
\end{equation}
If \(d_nL_n/n^2\) is bounded below by a positive constant, then
\begin{equation}
\label{eq:high-simplex-gram-rates}
 \mu_n^{\rm off}=o_{\Pbb}(L_n^{-1}),
 \qquad
 \|(q_i\cdot q_j)_{i,j\le n}\|_{\rm op}=O_{\Pbb}(L_n).
\end{equation}
\end{lemma}

\begin{proof}
Section~\ref{sec:regime-map}, with \(k_n=d_n/2\gg L_n\), gives
\[
  b_n\sim k_n,\qquad
  a_n\sim\sqrt{\frac{k_n}{2L_n}},\qquad
  \gamma_n\sim\frac{\sigma_n^2b_n}{L_n}
  \sim\frac{d_n}{2L_n}.
\]
This proves \eqref{eq:delta-entropy-regime}.

For the Gram bound, let \(P_n=I_n-\mathbf1\mathbf1^\top/n\). For every
coordinate \(\ell\), the vector
\[
  (\widetilde Z_{1\ell},\ldots,\widetilde Z_{n\ell})
\]
is centered Gaussian with covariance \(P_n\), and these vectors are
independent over \(\ell\). Hence, for \(i\ne j\),
\[
  \E(\widetilde Z_i\cdot\widetilde Z_j)=-\frac{d_n}{n}.
\]
Products of two standard Gaussian variables have uniformly bounded
subexponential norm. Bernstein's inequality and a union bound over fewer than
\(n^2\) pairs therefore give
\[
  \max_{i\ne j}
  \left|
    \widetilde Z_i\cdot\widetilde Z_j+\frac{d_n}{n}
  \right|
  =
  O_{\Pbb}\{\sqrt{d_nL_n}+L_n\}.
\]
Dividing by \(\gamma_n\asymp d_n/L_n\) yields
\[
  \mu_n^{\rm off}
  =
  O_{\Pbb}\left\{
    \frac{L_n}{n}
    +\frac{L_n^{3/2}}{\sqrt{d_n}}
    +\frac{L_n^2}{d_n}
  \right\},
\]
as claimed.  The nonzero eigenvalues of the Gram matrix agree with those of
\(\gamma_n^{-1}\widetilde Z^\top\widetilde Z\).  Hence the standard
Gaussian operator-norm bound gives
\[
 \upsilon_n
 \le
 \frac{C}{\gamma_n}
 \left(1+\sqrt{\frac{d_n}{n}}\right)^2
 =
 O_{\Pbb}\left(\frac{L_n}{d_n}+\frac{L_n}{n}\right).
\]

For the centered row bound, Gaussian regression gives, conditionally on
\(\widetilde Z_i\),
\[
 \widetilde Z_j+\frac{\widetilde Z_i}{n-1}
 \sim
 \mathcal N\left(0,\frac{n-2}{n-1}I_{d_n}\right)
 \qquad (j\ne i).
\]
Thus
\[
 \widetilde Z_i\cdot
 \left(\widetilde Z_j+\frac{\widetilde Z_i}{n-1}\right)
\]
is conditionally centered Gaussian with variance at most
\(|\widetilde Z_i|^2\).  Since
\(\max_i|\widetilde Z_i|^2=O_{\Pbb}(d_n)\) in the present regime, a
conditional Gaussian tail bound and a union bound over the \(n(n-1)\)
pairs give
\[
 \max_{i\ne j}
 \left|
 \widetilde Z_i\cdot
 \left(\widetilde Z_j+\frac{\widetilde Z_i}{n-1}\right)
 \right|
 =O_{\Pbb}(\sqrt{d_nL_n}).
\]
Division by \(\gamma_n\asymp d_n/L_n\) proves
\eqref{eq:centered-row-gram-rate}.

Finally, suppose \(d_nL_n/n^2\) is bounded below.  Since
\[
 \mu_n^{\rm off}\le C_n+\frac{\max_i|q_i|^2}{n-1}
\]
and \(\max_i|q_i|^2=O_{\Pbb}(L_n)\),
\eqref{eq:centered-row-gram-rate} gives
\(\mu_n^{\rm off}=o_{\Pbb}(L_n^{-1})\).  Applying the standard Gaussian
operator-norm bound once more gives
\[
 \|(q_i\cdot q_j)_{i,j\le n}\|_{\rm op}
 =
 O_{\Pbb}\left\{\frac{(\sqrt{d_n}+\sqrt n)^2}{\gamma_n}\right\}
 =
 O_{\Pbb}(L_n).
\]
This proves \eqref{eq:high-simplex-gram-rates}.
\end{proof}


\subsection{Exact centering and the high-dimensional correction}
\label{sec:full-growing-transition}

The entropy reduction in this section uses \(\log n\) as the radial
threshold. When \(d_n\asymp n^2\log n\), empirical centering and the
regular-simplex geometry shift that threshold by an amount visible in the
Gumbel coordinate. The exact coordinate \(s_n^\star\) in
\eqref{eq:exact-transition-coordinate} includes both effects.

The following calculation translates the exact coordinate into a critical
variance defect and records when this finite-sample correction is visible.
\begin{lemma}
\label{lem:exact-critical-coordinate}
Define
\begin{equation}
\label{eq:exact-critical-defect-window}
 \eps_{c,n}^\star
 =1-\frac{R_n^\star}{\sigma_n^2b_n},
 \qquad
 w_n^\star
 =\frac{a_nR_n^\star}{\sigma_n^2b_n^2}.
\end{equation}
If \(x_n=(\eps_n-\eps_{c,n}^\star)/w_n^\star=O(1)\), then
\[
 s_n^\star=\frac{x_n}{1-a_nx_n/b_n}=x_n+o(1).
\]
Let
\[
 s_n^{\rm cen}
 =\frac{\gamma_n\log n/\sigma_n^2-b_n}{a_n}.
\]
If \(s_n^\star=O(1)\), then
\[
 s_n^\star=s_n^{\rm cen}+o(1)
 \quad\text{when }d_n=o(n^2\log n),
\]
whereas, if \(d_n/(n^2\log n)\to\tau\in(0,\infty)\),
\begin{equation}
\label{eq:terminal-coordinate-shift}
 s_n^\star=s_n^{\rm cen}-\sqrt{\tau}+o(1).
\end{equation}
\end{lemma}

\begin{proof}
The first identity follows by substituting
\(\eps_n=\eps_{c,n}^\star+x_nw_n^\star\) into
\(\gamma_n=(1-\eps_n)^{-1}\); Lemma~\ref{lem:gamma-hazard} gives
\(a_n/b_n\to0\).  Also
\[
 s_n^{\rm cen}-s_n^\star
 =\frac{\gamma_n(\log n-R_n^\star)}
        {\sigma_n^2a_n}.
\]
Under \(s_n^\star=O(1)\), the expansion
\(R_n^\star=\log n-n^{-1}+O(\log n/n^2)\), together with the Gamma
asymptotics in Section~\ref{sec:regime-map}, gives
\[
 s_n^{\rm cen}-s_n^\star
 =\frac{\sqrt{d_n}}{n\sqrt{\log n}}+o(1)
\]
in the superlogarithmic range; the error is smaller when
\(d_n=O(\log n)\).  The two assertions follow.
\end{proof}


\section{From Radial Edge Points to NPMLE Atoms}
\label{sec:support}

This section completes the proofs of Theorems~\ref{thm:main-support-law} and
\ref{thm:regularization-path}. Sections~\ref{sec:d2}--\ref{sec:simplex}
locate every possible noncentral KKT contact; the propositions below convert
those field crossings into an exact support identity. In both fixed and
growing dimensions, we show that each superthreshold radial observation
creates exactly one fitted atom and that quantitative gaps away from the edge
exclude all others. Subsection~\ref{sec:edge-count-reduction} combines these
identities with the radial Poisson limit, and
Subsection~\ref{sec:regularization-path} proves that the support identities
hold uniformly along the variance path.

\subsection{Fixed-dimensional likelihood localization}
\label{sec:support-likelihood-localization}

Fixed-dimensional likelihood localization is the first step in converting
the field estimates into the exact support identity of
Theorem~\ref{thm:main-support-law}. We apply the one-atom
geometry around a fit that already carries a small amount of mass away from
its central atom. The resulting comparison shows that edge mass cannot create
new KKT contacts where the central directional field has a quantitative gap
below one.

It is convenient to work in score coordinates. Replace a component location
\(u\) by \(t=\sqrt{1-\eps_n}\,u\); this bijection preserves support size.
We reuse \(G\) for the transformed mixing measure. Recall that
\(\gamma_n=(1-\eps_n)^{-1}=1+2\kappa_n\). If the central atom is at
\(\tau\), the likelihood ratio of the atom at \(\tau+a\) to the atom at
\(\tau\), evaluated at \(Z_i\), is
\[
    q_{n,\tau,a}(i)
    =\exp\left\{a\cdot(Z_i-\gamma_n\tau)
                    -\frac{\gamma_n|a|^2}{2}\right\}.
\]
For a perturbation \(G=(1-W)\delta_\tau+G^{\rm out}\), write
\[
    m_G(i)
    =1-W+\int q_{n,\tau,u-\tau}(i)\,G^{\rm out}(du)
\]
for its likelihood ratio to \(\delta_\tau\), and define
\[
    D_{n,\tau}(a)=\frac1n\sum_{i=1}^nq_{n,\tau,a}(i),
    \qquad
    \Psi_G(\tau+a)=\frac1n\sum_{i=1}^n
        \frac{q_{n,\tau,a}(i)}{m_G(i)}.
\]
The KKT conditions are \(\Psi_G\le1\), with equality at every support point.

The deterministic no-propagation estimate below separates the total mass of
the edge atoms from the geometry of their locations.
\begin{lemma}
\label{lem:support-no-percolation}
For every measure \(G=(1-W)\delta_\tau+G^{\rm out}\) defined above,
\[
    \Psi_G(\tau+a)\le \frac{D_{n,\tau}(a)}{1-W},
    \qquad a\in\mathbb R^d.
\]
Consequently, if \(D_{n,\tau}\le1-\eta\) on a set \(A\) and
\(W\le\eta/2\), then \(G\) has no support point in \(\tau+A\).
\end{lemma}

\begin{proof}
Every likelihood ratio is nonnegative, so \(m_G(i)\ge1-W\). Substitution in
the definition of \(\Psi_G\) gives the first inequality, and
\((1-\eta)/(1-\eta/2)<1\) gives the second. \qedhere
\end{proof}

An NPMLE places almost all its mass near the true score origin. The following
coarse estimate isolates a central component before the dimension-specific
edge analysis. For a mixing measure \(G\) in score coordinates, put
\[
    m_G(z)=\int
       \exp\left\{t\cdot z-\frac{\gamma_n|t|^2}{2}\right\}G(dt),
    \qquad
    \mathfrak D_n(G)=\int\{1-e^{-\kappa_n|t|^2}\}\,G(dt).
\]
If \(P\) is standard Gaussian measure, then
\(Pm_G=1-\mathfrak D_n(G)\), and hence
\begin{equation}
\label{eq:support-population-defect}
    P\log m_G\le\log Pm_G
       =\log\{1-\mathfrak D_n(G)\}\le-\mathfrak D_n(G).
\end{equation}

The population loss yields a uniform empirical bound whose coarse
polylogarithmic precision is already sharper than the critical KKT gaps used
in the fixed-dimensional support argument.
\begin{lemma}
\label{lem:support-coarse-defect}
For each fixed \(d\), there is \(C_d<\infty\) such that, uniformly over
\(0\le\eps_n<1\),
\[
    \mathfrak D_n(\widehat G_n)
    =O_{\Pbb}\left\{\frac{(\log n)^{C_d}}{n}\right\}
\]
for every NPMLE \(\widehat G_n\).
\end{lemma}

\begin{proof}
It suffices to optimize over mixing measures supported in
\(\operatorname{conv}\{Z_1/\gamma_n,\ldots,Z_n/\gamma_n\}\). Indeed, if
\(p\) is the Euclidean projection of \(t\) onto this convex hull, then
\[
 (t-p)\cdot Z_i-\frac{\gamma_n}{2}(|t|^2-|p|^2)
 =(t-p)\cdot(Z_i-\gamma_np)-\frac{\gamma_n}{2}|t-p|^2\le0,
\]
so moving \(t\) to \(p\) weakly increases every fitted density. With
probability tending to one this hull lies in
\(B(0,C\sqrt{\log n})\).

For \(\mathcal M(R)=\{m_G:\supp(G)\subset B(0,R)\}\), let
\(N_{[]}(\eta,\mathcal M(R),L^1(P))\) be the least number of brackets
\([l,u]=\{f:l\le f\le u\}\), with \(P(u-l)\le\eta\), needed to cover
\(\mathcal M(R)\). We use the compact Gaussian-mixture entropy bound from
\cite[Section~2]{supportDivergenceCompanion}; see also
\cite[Section~3]{ghosal2001entropies}:
\begin{equation}
\label{eq:support-coarse-entropy}
 \log N_{[]}\{\eta,\mathcal M(R),L^1(P)\}
 \le C_d(1+R)^d\{1+\log(1/\eta)\}^{d+1}.
\end{equation}
The constant is uniform in \(\kappa_n\), since
\[
    m_G(z)\phi_d(z)
    =\int e^{-\kappa_n|t|^2}\phi_d(z-t)\,G(dt),
\]
and an \(L^1(P)\) bracket for \(m_G\) is an ordinary \(L^1\) bracket for
a Gaussian location mixture of total mass at most one.

Write \(\mathbb P_nf=n^{-1}\sum_{i=1}^nf(Z_i)\). Fix a dyadic shell
\(\delta<\mathfrak D_n(G)\le2\delta\), and cover it by upper
brackets \(u\) of \(L^1(P)\)-width \(\delta/4\). If \(m_G\le u\), then
\(Pu\le1-3\delta/4\), and Markov's inequality gives
\[
 \Pbb\{\mathbb P_n\log m_G\ge-\delta/4,\ m_G\le u\}
 \le e^{n\delta/4}(Pu)^n\le e^{-n\delta/2}.
\]
For \(R=O(\{\log n\}^{1/2})\), the logarithm of the number of brackets is
at most \((\log n)^{C_d}\) whenever \(\delta\ge n^{-2}\). A union bound over
brackets and dyadic shells therefore shows that every measure with
\(\mathfrak D_n(G)>C(\log n)^{C_d}/n\) has negative likelihood gain relative to
\(\delta_0\), with probability tending to one. Such a measure cannot be an
NPMLE, proving the claim.
\end{proof}

The defect bound does not rule out several nearby atoms. The local collapse
lemma shows that every maximizer carries all central mass on one atom, with
the remaining edge mass smaller than the KKT gaps used later.
\begin{lemma}
\label{lem:support-fixed-core-collapse}
Assume either \(d=2\) and \(\lambda_n\to\lambda\in(0,\infty)\), or fixed
\(d\ge3\) with \(s_n^\star\) converging in \(\R\). For a
sufficiently small fixed $r_0>0$, every NPMLE can, with
probability tending to one, be written
\begin{equation}
\label{eq:support-central-decomposition}
    \widehat G_n=(1-W_n)\delta_{\tau_n}+G_n^{\rm out},
    \qquad
    \supp(G_n^{\rm out})\subset B(0,r_0)^c,
\end{equation}
where
\[
    W_n+|\tau_n|^2
    =O_{\Pbb}\left\{\frac{(\log n)^{C_d}}{n\kappa_n}\right\}.
\]
In all these regimes, \(W_n=o_{\Pbb}(\kappa_n)\).
\end{lemma}

\begin{proof}
Lemma~\ref{lem:support-coarse-defect} and
\(1-e^{-\kappa_n|t|^2}\ge c\kappa_n\min\{|t|^2,1\}\) imply the displayed
bounds for the mass outside \(B(0,r_0)\) and for the second moment of the
restriction to that ball. Normalize the central restriction to a probability
measure \(G^{\rm cen}\), and let \(\tau\) and \(V\) be its mean and
covariance. Cauchy--Schwarz gives the asserted bound for \(|\tau|^2\).

We show that replacing \(G^{\rm cen}\) by \(\delta_\tau\) strictly increases
the likelihood unless \(V=0\). Keep the outside part fixed and interpolate
\[
 G_u=(1-W)\{(1-u)\delta_\tau+uG^{\rm cen}\}+G^{\rm out},
 \qquad 0\le u\le1.
\]
By concavity it is enough to show that the derivative at \(u=0\) is negative.
With \(h=t-\tau\), put
\[
 g_i(h)=\exp\left\{h\cdot(Z_i-\gamma_n\tau)
                  -\frac{\gamma_n|h|^2}{2}\right\},
 \qquad
 Q_i^{\rm out}=\int q_{n,\tau,t-\tau}(i)\,G^{\rm out}(dt).
\]
The likelihood ratio under \(G_0\) is
\(m_i^0=1-W+Q_i^{\rm out}\). Define
\(\omega_i=(1-W)/m_i^0\). The derivative is
\begin{equation}
\label{eq:support-collapse-derivative}
 \sum_{i=1}^n\omega_i
 \int\{g_i(h)-1\}\,G^{\rm cen}(dt).
\end{equation}

First replace \(\omega_i\) by one. Let
\(\widehat F_\tau(h)=n^{-1}\sum_i g_i(h)\). Uniformly for
\(|\tau|=o_{\Pbb}(\sqrt{\kappa_n})\) and \(|h|\le2r_0\), a compact
finite-dimensional empirical-process bound after two differentiations gives
\[
 \sup_{|h|\le2r_0}
 \|\nabla^2\widehat F_\tau(h)-\nabla^2F_\tau(h)\|_{\rm op}
 =O_{\Pbb}(n^{-1/2})=o_{\Pbb}(\kappa_n).
\]
Here the differentiated class has an integrable envelope of the form
\(C(1+|Z|^2)e^{2r_0|Z|}\), and the population average is
\[
    F_\tau(h)=\exp\{-\kappa_n|h|^2-\gamma_n\tau\cdot h\}.
\]
Indeed,
\[
 \nabla^2F_\tau(h)
 =F_\tau(h)\left[
   (2\kappa_nh+\gamma_n\tau)(2\kappa_nh+\gamma_n\tau)^\top
   -2\kappa_n I_d\right].
\]
The bound on \(|\tau|^2\), together with
\(\kappa_n\gtrsim1/\log n\), gives
\(|\tau|^2=o_{\Pbb}(\kappa_n)\). Thus, after decreasing the fixed \(r_0\) if
needed, \(\nabla^2\widehat F_\tau(h)\preceq-c\kappa_nI_d\) throughout the
ball, with probability tending to one. Taylor's theorem and
\(\int h\,G^{\rm cen}(dt)=0\) now show that the unweighted version of
\eqref{eq:support-collapse-derivative} is at most
\(-cn\kappa_n\operatorname{tr}V\).

It remains to restore the weights. Let \(m_i^1\) be the likelihood ratio
under the original NPMLE \(G_1\), and let
\[
 r_i^1=\frac{\int q_{n,\tau,t-\tau}(i)\,G^{\rm out}(dt)}{m_i^1}
\]
be its posterior responsibility for the outside component. Integrating the
componentwise KKT equality over \(G^{\rm out}\) gives
\(\sum_i r_i^1=nW\).

We next compare the two sets of weights explicitly. Put
\(\bar g_i=\int g_i(h)\,G^{\rm cen}(dt)\), so that
\(m_i^1=(1-W)\bar g_i+Q_i^{\rm out}\). On the event
\(R_n:=\max_i|Z_i|\le C\sqrt{\log n}\),
\[
 e^{-K_n}\le \bar g_i\le e^{K_n},
 \qquad K_n=C_{r_0}(1+R_n),
\]
uniformly in \(i\). Consequently \(m_i^1/m_i^0\le e^{K_n}\), and hence
\[
 1-\omega_i=\frac{Q_i^{\rm out}}{m_i^0}
 \le e^{K_n}r_i^1,
 \qquad
 \sum_i(1-\omega_i)\le e^{K_n}nW.
\]
For each \(i\), another second-order expansion, again using
\(\int h\,G^{\rm cen}(dt)=0\), gives
\[
 \left|\int\{g_i(h)-1\}\,G^{\rm cen}(dt)\right|
 \le C_{r_0}(1+|Z_i|^2)e^{C_{r_0}|Z_i|}\operatorname{tr}V.
\]
It follows that restoring the weights changes
\eqref{eq:support-collapse-derivative} by at most
\[
 C_{r_0}nW(1+\log n)e^{C_{r_0}\sqrt{\log n}}\operatorname{tr}V.
\]
Relative to \(n\kappa_n\operatorname{tr}V\), the last expression is
\[
 O_{\Pbb}\left\{
   \frac{(\log n)^{C_d+1}e^{C_{r_0}\sqrt{\log n}}}
        {n\kappa_n^2}
 \right\}=o_{\Pbb}(1),
\]
since \(\kappa_n\gtrsim1/\log n\) in both stated critical windows. The
derivative is therefore negative whenever \(V\ne0\), contradicting the
maximality of \(G_1\). Hence \(V=0\), which proves
\eqref{eq:support-central-decomposition}. The same rate and
\(\kappa_n\gtrsim1/\log n\) give \(W=o_{\Pbb}(\kappa_n)\).
\end{proof}

\subsection{Fixed-dimensional support recovery}
\label{sec:higher-dimensional-support-counts}

This subsection proves the fixed-dimensional support identity used in both
parts of Theorem~\ref{thm:main-support-law}. Exposed radial points are
separated by much more than the \(1/\rho_n\) width of a likelihood patch, so
their local optimizations decouple: a superthreshold patch contains one fitted
atom and a subthreshold patch contains none. Here
\(\rho_n=\sqrt{2\log n}\) for \(d=2\) and
\(\rho_n=\sqrt{2b_{n,d}}\) for fixed \(d\ge3\), while \(\tau\) denotes the
central score location from Lemma~\ref{lem:support-fixed-core-collapse}. Put
\(\Delta_n^{\rm anc}=\gamma_n\tau-\bar Z_n\). The exact anchor identity is
\begin{equation}
\label{eq:support-high-d-anchor-translation}
 D_{n,\tau}(a)
 =\frac1n\sum_iq_{n,\tau,a}(i)
 =e^{-a\cdot\Delta_n^{\rm anc}}D_{n,*}(a),
 \qquad
 D_{n,*}(a)=\frac1n\sum_i
 e^{a\cdot(Z_i-\bar Z_n)-\gamma_n|a|^2/2}.
\end{equation}
The central location is tied to the sample mean by an exact barycenter
identity. For a score-coordinate NPMLE \(G\), let
\[
 \Pi_i(dt)=
 \frac{\exp\{t\cdot Z_i-\gamma_n|t|^2/2\}\,G(dt)}{m_G(Z_i)}
\]
be the posterior mixing distribution for observation \(i\). KKT equality
\(G\)-almost everywhere gives
\[
 \frac1n\sum_{i=1}^n\Pi_i(dt)=G(dt).
\]
On the other hand, differentiating the likelihood when every score location is
translated by the same vector gives
\[
 0=\frac1n\sum_{i=1}^n
 \left\{Z_i-\gamma_n\int t\,\Pi_i(dt)\right\}.
\]
Consequently,
\begin{equation}
\label{eq:support-npmle-barycenter}
    \bar Z_n=\gamma_n\int t\,G(dt).
\end{equation}
For the decomposition in Lemma~\ref{lem:support-fixed-core-collapse}, this
implies
\[
 \Delta_n^{\rm anc}
 =\gamma_n\int(\tau-t)\,G_n^{\rm out}(dt).
\]
Support reduction places \(G\) in
\(\operatorname{conv}\{Z_1/\gamma_n,\ldots,Z_n/\gamma_n\}\). Since
\(\max_i|Z_i|=O_{\Pbb}(\rho_n)\) in fixed dimension, the rate in
Lemma~\ref{lem:support-fixed-core-collapse} therefore gives
\begin{equation}
\label{eq:support-high-d-anchor-error}
    \rho_n|\Delta_n^{\rm anc}|
    =O_{\Pbb}(W_n\rho_n^2)
    =o_{\Pbb}(\kappa_n)
\end{equation}
in both fixed-dimensional critical windows. For the last equality, use
\[
 \rho_n^2=O(\log n),\qquad
 W_n=O_{\Pbb}\left(\frac{(\log n)^{C_d}}{n\kappa_n}\right),
 \qquad \kappa_n\gtrsim\frac1{\log n}.
\]

We use one symbol for the dimension-specific radial excess:
\begin{equation}
\label{eq:fixed-support-radial-excess}
 x_{n,i}=
 \begin{cases}
   |Z_i|^2/2-\log n,&d=2,\\[2mm]
   |Z_i-\bar Z_n|^2/2-b_{n,d},&d\ge3.
 \end{cases}
\end{equation}
In fixed dimension, replacing the centered radius in the second line by
\(|Z_i|\) changes all excesses by \(o_{\Pbb}(1)\), uniformly in \(i\).

We now fix the local coordinate maps used below. In \(d=2\), write
\(U_j=Z_j/|Z_j|\), let \(\mathsf R_\theta\) denote planar rotation through
angle \(\theta\), and set
\begin{equation}
\label{eq:support-d2-patch-map}
 a_{n,j}^{(2)}(r,z)
 = (\rho_n+r)\mathsf R_{z/\rho_n}U_j,
 \qquad (r,z)\in\mathbb R^2.
\end{equation}
For fixed \(d\ge3\), set
\begin{equation}
\label{eq:support-dge3-patch-map}
 a_{n,j}^{(\ge3)}(h)
 =a_{n,j}^*+h,
 \qquad
 a_{n,j}^*=\frac{Z_j-\gamma_n\tau}{\gamma_n},
 \qquad h\in\mathbb R^d.
\end{equation}
If \(\nu_j\) is a measure in the corresponding local coordinates, its
pushforward under \(a_{n,j}^{(d)}\) is a measure on score displacements. Below
we write
\begin{equation}
\label{eq:support-patch-pushforward}
    \nu=\sum_j(a_{n,j}^{(d)})_\#\nu_j
\end{equation}
for the combined displacement measure.

For a finite positive measure \(\nu\) on displacement space, with
\(m=\nu(\mathbb R^d)\le n\), consider the anchored mixture
\[
    G_{n,\tau,\nu}
    =\left(1-\frac mn\right)\delta_\tau
      +\frac1n(\tau+\cdot)_\#\nu .
\]
Its exact likelihood gain over \(\delta_\tau\) is
\begin{equation}
\label{eq:support-anchored-likelihood}
    \mathcal L_{n,\tau}(\nu)
    =\sum_{i=1}^n
      \log\left\{1+F_{n,i}(\nu)-\frac mn\right\},
    \qquad
    F_{n,i}(\nu)=\frac1n\int q_{n,\tau,a}(i)\,\nu(da).
\end{equation}
Thus \(m=nW\) is the natural rescaled edge mass.

\subsubsection{Local likelihood programs}

The field limits in Sections~\ref{sec:d2} and
\ref{sec:fixed-d-patches} determine the likelihood program associated with one
radial extreme. Optimizing this program shows that an active patch contains one
atom rather than a split mixing measure.

Suppose first that \(d\ge3\), along a sequence on which \(s_n\to s\). For an
observation of centered radial excess \(x\), completing the square at the
displacement \(a=a_{n,i}^{(\ge3)}(h)\) gives
\[
    \frac1nq_{n,\tau,a}(i)\longrightarrow
    A_x(h):=e^{x-s-|h|^2/2}.
\]
If a finite positive measure \(\nu\) of rescaled mass is placed in this
patch, its limiting gain is
\begin{equation}
\label{eq:support-dge3-local-program}
    \mathcal J_x^{(\ge3)}(\nu)
    =\log\left\{1+\int A_x(h)\,\nu(dh)\right\}
      -\nu(\mathbb R^d).
\end{equation}
For fixed total mass, the integral is uniquely maximized by putting all mass
at \(h=0\). The remaining scalar problem is
\[
    \max_{y\ge0}\{\log(1+e^{x-s}y)-y\},
\]
whose optimizer is \(y_x^*=(1-e^{-(x-s)})_+\). Hence the patch is active
exactly when \(x>s\), and every nonzero limiting optimizer is the one-atom
measure \(y_x^*\delta_0\).

In dimension two, use radial score offset \(r\) and blown-up angular
displacement \(z\), and write
\[
    B_\lambda(r)=e^{-\lambda}\Phi(-r),
    \qquad
    C_{\lambda,x}(r,z)=e^{-\lambda}e^{x-r^2/2-z^2/2}.
\]
The corresponding local program is
\begin{equation}
\label{eq:support-d2-local-program}
    \mathcal J_x^{(2)}(\nu)
    =\log\left\{1+\int C_{\lambda,x}\,d\nu\right\}
      -\int\{1-B_\lambda(r)\}\,d\nu(r,z).
\end{equation}
Put
\[
    R_x(r,z)=\frac{C_{\lambda,x}(r,z)}{1-B_\lambda(r)}.
\]
If \(u=\int C_{\lambda,x}\,d\nu\), then
\[
    \int(1-B_\lambda)\,d\nu\ge
    \frac{u}{\sup_{r,z}R_x(r,z)},
\]
with equality only when \(\nu\) is supported on the maximizers of \(R_x\).
The angular maximizer is uniquely \(z=0\). The radial maximizer is uniquely
\(r=-a_\lambda\), where
\[
    e^\lambda=\Phi(a_\lambda)+\frac{\phi(a_\lambda)}{a_\lambda},
\]
and the calculation in Section~\ref{sec:d-two-sharp-window} gives
\(\sup R_x=e^{x-x_\lambda}\). Thus the patch is active exactly when
\(x>x_\lambda\), and every nonzero limiting optimizer is a single atom at
\((-a_\lambda,0)\). Its rescaled mass is
\[
    q_x^*=
    \frac{1-e^{-(x-x_\lambda)}}
         {1-e^{-\lambda}\Phi(a_\lambda)}.
\]

For fixed \(d\ge3\), the contact argument requires differentiated control near
a random radial extreme.  The next lemma supplies that control without
conditioning on a sample-centered exceedance event.  It first fixes the raw
extreme, so the punctured sample remains independent, and then sums over the
possible distinguished observations by the Palm formula.
\begin{lemma}
\label{lem:support-dge3-palm-c2}
Fix \(d\ge3\) and suppose \(s_n=O(1)\).  Put
\[
 \widetilde x_{n,i}=\frac{|Z_i|^2}{2}-b_{n,d},
 \qquad
 \mathcal E_n(C)=\{j:|\widetilde x_{n,j}|\le C\}.
\]
For fixed \(A,C,R<\infty\), define, for \(j\in\mathcal E_n(C)\),
\[
 \mathcal R_{n,j,A}(h)
 =\frac1n\sum_{i\ne j}
   \exp\left\{
      \left(\frac{Z_j}{\gamma_n}+h\right)\!\cdot Z_i
      -\frac{\gamma_n}{2}
       \left|\frac{Z_j}{\gamma_n}+h\right|^2
   \right\}
   \mathbf 1_{\{\widetilde x_{n,i}\le-A\}},
 \qquad h\in B_R.
\]
Then
\begin{equation}
\label{eq:support-dge3-palm-c2}
 \max_{j\in\mathcal E_n(C)}
 \|\mathcal R_{n,j,A}\|_{C^2(B_R)}\xrightarrow{\Pbb}0.
\end{equation}
Moreover, if \(K_{n,j,i}(h)\) denotes the exponential summand in the
definition of \(\mathcal R_{n,j,A}\), then
\begin{equation}
\label{eq:support-dge3-palm-absolute}
 \max_{j\in\mathcal E_n(C)}\max_{|\beta|\le2}\sup_{h\in B_R}
 \frac1n\sum_{i\ne j}
 |\partial_h^\beta K_{n,j,i}(h)|
 \mathbf1_{\{\widetilde x_{n,i}\le-A\}}
 \xrightarrow{\Pbb}0.
\end{equation}
The same conclusions hold with the candidate condition
\(|x_{n,j}|\le C\) and lower-cloud cutoff \(x_{n,i}\le-A\), where \(x_{n,i}\)
is the sample-centered excess in \eqref{eq:fixed-support-radial-excess}.
\end{lemma}

\begin{proof}
Let \(R_{n,A}^2/2=b_{n,d}-A\), choose \(p>d\), and fix a deterministic
\(z\) with \(\bigl||z|^2/2-b_{n,d}\bigr|\le C\).  For
\(t_z(h)=z/\gamma_n+h\), set
\[
 V_{n,z,i}(h)=\frac1n
 e^{t_z(h)\cdot Z_i-\gamma_n|t_z(h)|^2/2}
 \mathbf 1_{\{|Z_i|\le R_{n,A}\}}.
\]
We first establish, uniformly over these \(z\), \(h\in B_R\),
\(q\in\{2,p\}\), and \(|\beta|\le3\),
\begin{equation}
\label{eq:support-palm-derivative-moments}
 \sum_{i=2}^n
 \E|\partial_h^\beta V_{n,z,i}(h)|^q
 \le C_{A,C,R,p,d}\rho_n^{-1}e^{-\lambda_n}.
\end{equation}
Indeed, differentiation inserts a polynomial of degree at most \(|\beta|\)
in \(Z_i-\gamma_nt_z(h)\).  Exponential tilting therefore gives
\begin{align*}
 &\sum_{i=2}^n
 \E|\partial_h^\beta V_{n,z,i}(h)|^q\\
 &\quad\le
 C n^{1-q}e^{-q\kappa_n|t_z(h)|^2+(q^2-q)|t_z(h)|^2/2}
 \E\left[
  \{1+|G+(q-\gamma_n)t_z(h)|\}^{q|\beta|}
  \mathbf1_{\{|G+qt_z(h)|\le R_{n,A}\}}
 \right].
\end{align*}
Aligning \(t_z(h)=re_1\), conditioning on \(G_\perp\), and applying Mills'
bound to \(G_1\) show that the last expectation is at most
\[
 C_{A,C,R,p,d}\rho_n^{-1}
   \exp\{-(qr-R_{n,A})^2/2\}.
\]
Here the polynomial causes only a constant loss: on the truncation event its
longitudinal argument is
\(R_{n,A}-\gamma_nr+O(1/\rho_n)\), which is bounded for
\(|r-|z|/\gamma_n|\le R\), while the transverse Gaussian moments are fixed.
Equivalently, this weighted Mills estimate follows by integrating
\((1+|u|)^m e^{-u^2/2}\) beyond the same one-dimensional boundary.

The exponent outside \(\rho_n^{-1}\) equals
\begin{align*}
 &(1-q)\log n-q\kappa_nr^2
   +\frac{q^2-q}{2}r^2-\frac12(qr-R_{n,A})^2\\
 &\qquad=(q-1)\{b_{n,d}-\log n-A\}
   -q\kappa_nr^2-\frac q2(r-R_{n,A})^2.
\end{align*}
Since \(b_{n,d}-\log n=\lambda_n-s_n\),
\(\kappa_nr^2=\lambda_n+o(1)\), and \(s_n=O(1)\), it is at most
\(-\lambda_n+O_{A,C,R,p,d}(1)\).  This proves
\eqref{eq:support-palm-derivative-moments}.  The same tilting calculation with
\(q=1\) gives
\begin{equation}
\label{eq:support-palm-derivative-means}
 \sup_{z,h,|\beta|\le3}
 \left|\sum_{i=2}^n\E\partial_h^\beta V_{n,z,i}(h)\right|
 \le C_{A,C,R,d}e^{-\lambda_n}=o(1).
\end{equation}
The same \(q=1\) calculation applies to the absolute values of the
derivatives.  Using the derivatives of order three and the fundamental theorem
of calculus on a fixed mesh of \(B_R\) gives
\[
 \sup_z\E\left[
  \max_{|\beta|\le2}\sup_{h\in B_R}
  \frac1n\sum_{i=2}^n
  |\partial_h^\beta K_{n,z,i}(h)|
  \mathbf1_{\{|Z_i|\le R_{n,A}\}}
 \right]
 \le C_{A,C,R,d}e^{-\lambda_n}=o(1).
\]
The Palm union bound below and Markov's inequality therefore also prove
\eqref{eq:support-dge3-palm-absolute}.

Apply Rosenthal's inequality to the centered derivatives and integrate over
\(B_R\).  Equations \eqref{eq:support-palm-derivative-moments} and
\eqref{eq:support-palm-derivative-means} yield, uniformly in \(z\),
\[
 \E\left\|
   \sum_{i=2}^nV_{n,z,i}
 \right\|_{W^{3,p}(B_R)}^p=o(1).
\]
Because \(p>d\), the embedding \(W^{3,p}(B_R)\hookrightarrow C^2(B_R)\)
and Markov's inequality imply
\[
 \sup_{z:\,\lvert |z|^2/2-b_{n,d}\rvert\le C}
 \Pbb\left\{
  \left\|\sum_{i=2}^nV_{n,z,i}\right\|_{C^2(B_R)}>\delta
 \right\}=o(1)
\]
for every \(\delta>0\).  Independence now gives the Palm union bound
\begin{align*}
 &\Pbb\left\{
   \max_{j\in\mathcal E_n(C)}
   \|\mathcal R_{n,j,A}\|_{C^2(B_R)}>\delta
 \right\}\\
 &\quad\le n\int_{\{\lvert |z|^2/2-b_{n,d}\rvert\le C\}}
  \Pbb\left\{
   \left\|\sum_{i=2}^nV_{n,z,i}\right\|_{C^2(B_R)}>\delta
  \right\}\,P(dz)=o(1),
\end{align*}
because the Gamma tail estimate gives
\(nP\{\lvert |Z|^2/2-b_{n,d}\rvert\le C\}=O_C(1)\).

Finally,
\[
 \max_i\left|
  \frac{|Z_i-\bar Z_n|^2-|Z_i|^2}{2}
 \right|=o_{\Pbb}(1).
\]
Fix \(\delta>0\).  Apart from an event of probability tending to zero,
\(|x_{n,j}|\le C\) implies \(|\widetilde x_{n,j}|\le C+\delta\), while
\(x_{n,i}\le-A\) implies
\(\widetilde x_{n,i}\le-(A-\delta)\).  Applying the raw result with
candidate bound \(C+\delta\) and lower cutoff \(A-\delta\) proves the
sample-centered assertion directly.
\end{proof}

\subsubsection{Uniform patch stability}

We need the local approximations uniformly over the random edge patches and
all candidate mixing measures. As in
Sections~\ref{sec:sharp-high-dimensional-estimates} and
\ref{sec:d-two-sharp-window}, the lower radial cutoff is removed after the
finite-sample limit.

The candidate patches and the observational cutoff play different roles.  The
candidate set is fixed by the activation threshold and remains unchanged as
the cutoff is removed.  The cutoff is used only to divide the data into a
finite exposed cloud and a lower cloud.  The following estimate controls the
lower cloud uniformly over every mixing measure on the candidate patches.
\begin{lemma}
\label{lem:support-uniform-lower-cloud}
Fix \(C,R,M,J<\infty\).  Let \(\mathcal I_n\) be any set of at most \(J\)
candidate observations with \(|x_{n,j}|\le C\), and let \(\nu_j\) be supported
in \(B_R\) for \(d\ge3\), or in \([-R,R]^2\) for \(d=2\), with
\(m=\sum_{j\in\mathcal I_n}\nu_j(\mathbb R^d)\le M\).  Form \(\nu\) by
\eqref{eq:support-patch-pushforward} and write
\(F_{n,i}=F_{n,i}(\nu)\).  Suppose \(s_n=O(1)\) for fixed \(d\ge3\), or
\(\lambda_n\to\lambda\in(0,\infty)\) for \(d=2\).

There is \(A_0=A_0(C,R,M)>C+1\) such that, for each fixed
\(A\ge A_0\), on the event \(\max_i x_{n,i}\le C\) and the usual
edge-separation event,
\begin{equation}
\label{eq:support-lower-cloud-max-square}
 \max_{i:x_{n,i}\le-A}F_{n,i}\le C_{C,R}Me^{-A},
 \qquad
 \sum_{i:x_{n,i}\le-A}F_{n,i}^2
 \le C_{C,R}M^2e^{-A}+o_{\Pbb,A}(1),
\end{equation}
apart from an event of probability tending to zero.  For every candidate
patch, one also has the derivative-envelope bound
\begin{equation}
\label{eq:support-lower-cloud-derivative-envelope}
 \max_{|\beta|\le2}\sup_u
 \frac1n\sum_{i:x_{n,i}\le-A}
 \left|\partial_u^\beta
 q_{n,\tau,a_{n,j}(u)}(i)\right|
 =O_{\Pbb,A_0,C,R}(1),
\end{equation}
uniformly for \(A\ge A_0\).
Uniformly over the same measures,
\begin{equation}
\label{eq:support-lower-cloud-linear}
 \sum_{i:x_{n,i}\le-A}F_{n,i}
 =
 \begin{cases}
   o_{\Pbb,A}(1),&d\ge3,\\[1mm]
   \displaystyle
   \sum_{j\in\mathcal I_n}\int B_{\lambda_n}(r)\,\nu_j(dr,dz)
      +O_{\Pbb,C,R,J}(Me^{-A})+o_{\Pbb,A}(1),&d=2,
 \end{cases}
\end{equation}
where
\[
 B_{\lambda_n}(r)=e^{-\lambda_n}\Phi(-r).
\]
The observations with \(-A<x_{n,i}\le C\) that do not belong to
\(\mathcal I_n\) contribute \(o_{\Pbb,A}(1)\), uniformly over the candidate
patches and measures.
\end{lemma}

\begin{proof}
Completing the square gives, for every score displacement \(a\),
\begin{equation}
\label{eq:support-lower-cloud-square-completion}
 \frac1nq_{n,\tau,a}(i)
 =\exp\left\{
   \frac{|Z_i-\gamma_n\tau|^2}{2\gamma_n}-\log n
   -\frac{\gamma_n}{2}
      \left|a-\frac{Z_i-\gamma_n\tau}{\gamma_n}\right|^2
 \right\}.
\end{equation}
For \(d\ge3\), the anchor estimate
\eqref{eq:support-high-d-anchor-error} and
\(b_{n,d}-\log n=\lambda_n-s_n\) make the first exponent at most
\(-A-s_n+o_{\Pbb}(1)\) when \(x_{n,i}\le-A\).  For \(d=2\), the radial
identity in Lemma~\ref{lem:support-d2-c2-absorption} bounds the exponent by
\(-A+O_{C,R}(1)\). Integrating
\eqref{eq:support-lower-cloud-square-completion} over a measure of mass at most
\(M\) proves the first inequality in
\eqref{eq:support-lower-cloud-max-square}.

For fixed \(d\ge3\), Lemma~\ref{lem:support-dge3-palm-c2}, followed by the
translation from \(Z_j/\gamma_n\) to
\((Z_j-\gamma_n\tau)/\gamma_n\), gives
\[
 \max_{j\in\mathcal I_n}\sup_{u}
 \frac1n\sum_{i:x_{n,i}\le-A}
 q_{n,\tau,a_{n,j}^{(\ge3)}(u)}(i)=o_{\Pbb,A}(1).
\]
The translation is legitimate in \(C^2\) because
\(|\tau|=o_{\Pbb}(1)\) and the Palm lemma controls a slightly larger compact
set.  The translated Palm field is multiplied by
\(e^{-\gamma_na\cdot\tau}\); uniformly on these patches its logarithm is
\(O_{\Pbb}(\rho_n|\tau|)=o_{\Pbb}(1)\) by
Lemma~\ref{lem:support-fixed-core-collapse}.  Positivity and Fubini now give the first line of
\eqref{eq:support-lower-cloud-linear}.
Equation \eqref{eq:support-dge3-palm-absolute} gives
\eqref{eq:support-lower-cloud-derivative-envelope} in this case.  The bound is
uniform for \(A\ge A_0\), because the absolute derivative sum is nonincreasing
as the lower cutoff \(-A\) is moved inward.

For \(d=2\), use the fixed-cutoff decomposition
Lemma~\ref{lem:d2-fixed-cutoff-decomposition}, the quantitative absorption
bound in the proof of Lemma~\ref{lem:d2-low-atom-absorption}, and the anchor
multiplier
\[
 e^{-a\cdot\Delta_n^{\rm anc}-a\cdot\bar Z_n}
 =e^{-\gamma_na\cdot\tau}=1+o_{\Pbb}(1)
\]
uniformly on the candidate patches. Together, these three estimates show that the lower field in a patch
with radial coordinate \(r\) is
\(B_{\lambda_n}(r)+O_{\Pbb,C,R}(e^{-A})+o_{\Pbb,A}(1)\), uniformly in its
direction.  Integrating over \(\nu\) proves the second line of
\eqref{eq:support-lower-cloud-linear}.
The radial identity bounds the absolute value of each derivative by a fixed
polynomial times the Gaussian angular kernel in
Lemma~\ref{lem:support-d2-c2-absorption}. For the deep part
\(\xi_i<-\ell_n^{\rm edge}\), repeat the boxwise Rosenthal--Morrey argument in
Lemma~\ref{lem:support-d2-c2-absorption} with a polynomial majorant for the
absolute derivatives; its mean is uniformly bounded by exponential tilting.
For the strip
\(-\ell_n^{\rm edge}\le\xi_i\le-A_0\), the slab-and-arc estimate in that
lemma, with the summable weight \((1+q^4)e^{-cq^2}\), is
\(O_{\Pbb,C,R}(e^{-A_0})\).  Their sum proves
\eqref{eq:support-lower-cloud-derivative-envelope} for \(d=2\).  Since the
absolute sum over \(\{x_{n,i}\le-A\}\) decreases with \(A\), the
\(O_{\Pbb,C,R}(e^{-A_0})\) bound holds simultaneously for all \(A\ge A_0\).

The linear bound is \(O_{\Pbb}(M)\) in either dimension.  Therefore
positivity and the first inequality in
\eqref{eq:support-lower-cloud-max-square} give
\[
 \sum_{i:x_{n,i}\le-A}F_{n,i}^2
 \le
 \left(\max_{i:x_{n,i}\le-A}F_{n,i}\right)
 \sum_{i:x_{n,i}\le-A}F_{n,i}
 \le C_{C,R}M^2e^{-A}+o_{\Pbb,A}(1),
\]
which is the second inequality in
\eqref{eq:support-lower-cloud-max-square}.

Finally, for fixed \(A,C\), the exposed cloud
\(\{-A<x_{n,i}\le C\}\) has tight cardinality.  Its distinct directions are
separated by \(\ell_n^{\rm sep}/\rho_n\), with
\(\ell_n^{\rm sep}\to\infty\), apart from an event of probability tending to
zero.  Completing the square and differentiating on fixed patches then show
that the response of every nonassociated exposed observation is
\(O(e^{-c(\ell_n^{\rm sep})^2})\).  Summing over the tight exposed cloud proves
the last assertion.
\end{proof}

Lemma~\ref{lem:support-uniform-lower-cloud} makes the lower-cloud contribution
uniform over all mixing measures on finitely many candidate patches. The next
proposition uses its linear, quadratic, and derivative bounds to reduce the
exact likelihood to a finite collection of local programs.
\begin{proposition}
\label{prop:support-finite-patch-expansion}
Fix \(C,R,M,J<\infty\), and let \(\mathcal I_n\) be any set of at most \(J\)
candidate observations with \(|x_{n,j}|\le C\).  Work on the event
\(\max_i x_{n,i}\le C\) and the edge-separation event.  For \(d\ge3\), let
\(\nu_j\) be supported in \(B_R\subset\mathbb R^d\); for \(d=2\), let it be
supported in \([-R,R]^2\).  Form the displacement measure using only
\(j\in\mathcal I_n\), assume \(\sum_{j\in\mathcal I_n}\nu_j(\mathbb R^d)\le M\),
and suppose the anchor satisfies \eqref{eq:support-high-d-anchor-error}.  Assume
that \(s_n=O(1)\) for fixed \(d\ge3\), or that
\(\lambda_n\to\lambda\in(0,\infty)\) for \(d=2\), and define
\[
\begin{aligned}
 \mathcal J_n^{(\ge3)}(\nu)
 &=\sum_{j\in\mathcal I_n}\log\left\{1+\int
 e^{x_{n,j}-s_n-|h|^2/2}\,\nu_j(dh)\right\}
 -\sum_{j\in\mathcal I_n}\nu_j(\mathbb R^d),\\
 B_{\lambda_n}(r)&=e^{-\lambda_n}\Phi(-r),\qquad
 C_{n,x}(r,z)=e^{-\lambda_n}e^{x-r^2/2-z^2/2},\\
 \mathcal J_n^{(2)}(\nu)
 &=\sum_{j\in\mathcal I_n}\log\left\{1+\int C_{n,x_{n,j}}\,d\nu_j\right\}
 -\sum_{j\in\mathcal I_n}\int\{1-B_{\lambda_n}(r)\}\,d\nu_j.
\end{aligned}
\]
There is \(A_0=A_0(C,R,M)<\infty\) such that, for every fixed
\(A\ge A_0\), uniformly over these measures,
\begin{equation}
\label{eq:support-finite-patch-error}
 \sup_\nu|\mathcal L_{n,\tau}(\nu)-\mathcal J_n^{(d)}(\nu)|
 =O_{\Pbb,C,R,J}\{(M+M^2)e^{-A}\}+o_{\Pbb,A}(1),
\end{equation}
where the random coefficient in \(O_{\Pbb,C,R,J}(1)\) is tight uniformly for
\(A\ge A_0\).  In particular, for every \(\delta>0\),
\begin{equation}
\label{eq:support-finite-patch-iterated}
 \lim_{A\to\infty}\limsup_{n\to\infty}
 \Pbb\left\{
  \sup_\nu|\mathcal L_{n,\tau}(\nu)-\mathcal J_n^{(d)}(\nu)|>\delta
 \right\}=0.
\end{equation}
The candidate space and \(\mathcal J_n^{(d)}\) do not depend on \(A\).  After
\(n\to\infty\) and then \(A\to\infty\), the limiting functional is the sum
of the independent local programs
\eqref{eq:support-dge3-local-program} or
\eqref{eq:support-d2-local-program} over the candidate patches.
\end{proposition}

\begin{proof}
Put \(m=\nu(\mathbb R^d)\), \(v_i=F_{n,i}(\nu)-m/n\), and split the data into
the candidate observations, the noncandidate exposed cloud
\(\{-A<x_{n,i}\le C\}\setminus\mathcal I_n\), and the lower cloud
\(\{x_{n,i}\le-A\}\).  For a candidate observation \(j\), completing the
square gives, uniformly on its associated patch,
\[
 \frac1nq_{n,\tau,a_{n,j}^{(\ge3)}(h)}(j)
 =e^{x_{n,j}-s_n-|h|^2/2+o(1)}
\]
in \(d\ge3\), and gives \(C_{n,x_{n,j}}(r,z)+o(1)\) in \(d=2\).
The separation event makes every cross-patch response \(o(1)\).  Hence the
candidate observations contribute the logarithmic terms in
\(\mathcal J_n^{(\ge3)}\) or \(\mathcal J_n^{(2)}\), uniformly over \(\nu\),
up to \(o_{\Pbb,A}(1)\).  The final assertion of
Lemma~\ref{lem:support-uniform-lower-cloud} shows that the noncandidate exposed
cloud contributes \(o_{\Pbb,A}(1)\).

Choose \(A_0\) so that \(C_{C,R}Me^{-A_0}\le1/4\), and then take
\(n\ge4M\).  On the lower cloud, the elementary inequality
\[
 |\log(1+v)-v|\le2v^2,
 \qquad v\ge-M/n,\quad |v|\le1/2,
\]
and Lemma~\ref{lem:support-uniform-lower-cloud} give
\[
 \sum_{i:x_{n,i}\le-A}\{\log(1+v_i)-v_i\}
 =O_{\Pbb,C,R,J}(M^2e^{-A})+o_{\Pbb,A}(1).
\]
Since the exposed cloud has \(O_{\Pbb,A}(1)\) observations,
\[
 \sum_{i:x_{n,i}\le-A}v_i
 =\sum_{i:x_{n,i}\le-A}F_{n,i}-m+o_{\Pbb,A}(1).
\]
Equation \eqref{eq:support-lower-cloud-linear} therefore supplies the penalty
\(-m\) in \(d\ge3\), and the penalty
\(-\sum_j\int\{1-B_{\lambda_n}(r)\}\,d\nu_j\) in \(d=2\).
Lemma~\ref{lem:support-uniform-lower-cloud} already includes the anchor error.
Combining the candidate observations, the noncandidate exposed cloud, and the
lower cloud proves \eqref{eq:support-finite-patch-error}.
\end{proof}

An elementary coercivity estimate bounds the edge mass without assuming a
fixed number of fitted atoms; it depends only on the candidate data points,
not on the mixing complexity within their patches.
\begin{proposition}
\label{prop:support-patch-coercivity}
Restrict \eqref{eq:support-anchored-likelihood} to \(J\) edge patches, and let
\(\mathcal I_n\) be their associated observations. Suppose throughout the
patches that
\[
 \frac1n\sum_{i\notin\mathcal I_n}q_{n,\tau,a}(i)\le b<1,
 \qquad
 \max_{i\in\mathcal I_n}\frac1nq_{n,\tau,a}(i)\le C_0.
\]
If \(m=nW\) is the total rescaled edge mass, then, for all large \(n\),
\[
    \mathcal L_{n,\tau}(\nu)
    \le-c_bm+J\log(1+C_0m),
    \qquad c_b=(1-b)/3.
\]
Consequently, every maximizer satisfies
\(m\le4C_0J^2/c_b^2\).
\end{proposition}

\begin{proof}
Fubini's theorem gives
\(\sum_{i\notin\mathcal I_n}F_{n,i}(\nu)\le bm\). Jensen's inequality then
implies
\[
 \sum_{i\notin\mathcal I_n}
 \log\{1-m/n+F_{n,i}(\nu)\}
 \le(n-J)\log\left\{1-\frac mn+\frac{bm}{n-J}\right\}
 \le-c_bm.
\]
At each retained observation the likelihood ratio is at most \(1+C_0m\),
which proves the displayed bound. A maximizer has nonnegative gain, and
\(\log(1+x)\le2\sqrt x\) then yields
\(m\le4C_0J^2/c_b^2\).
\end{proof}

The two-dimensional contact argument requires differentiated, rather than only
uniform, control of the inward-edge background. The next lemma upgrades the
fixed-cutoff decomposition to the \(C^2\) form needed below.
\begin{lemma}
\label{lem:support-d2-c2-absorption}
Assume \(d=2\) and \(\lambda_n\to\lambda\in(0,\infty)\), and put
\(\ell_n^{\rm edge}=4\log\rho_n\). For \(v\in S^1\), define
\[
 a_{n,v}(r,z)=(\rho_n+r)\mathsf R_{z/\rho_n}v
\]
and, for fixed \(R<\infty\), let
\begin{align*}
 \mathcal B_{n,v}(r,z)
 &=\frac1n e^{-\kappa_n|a_{n,v}(r,z)|^2}
   \sum_{i:\,\xi_i<-\ell_n^{\rm edge}}
   e^{a_{n,v}(r,z)\cdot Z_i-|a_{n,v}(r,z)|^2/2},\\
 \mathcal A_{n,v,A}(r,z)
 &=\frac1n e^{-\kappa_n|a_{n,v}(r,z)|^2}
   \sum_{i:\,-\ell_n^{\rm edge}\le\xi_i\le-A}
   e^{a_{n,v}(r,z)\cdot Z_i-|a_{n,v}(r,z)|^2/2}.
\end{align*}
Then
\begin{equation}
\label{eq:support-d2-diffuse-c2}
 \sup_{v\in S^1}
 \|\mathcal B_{n,v}-e^{-\lambda}\Phi(-r)\|_{C^2([-R,R]^2)}
 \xrightarrow{\Pbb}0,
\end{equation}
and, for every \(\delta>0\),
\begin{equation}
\label{eq:support-d2-strip-c2}
 \lim_{A\to\infty}\limsup_{n\to\infty}
 \Pbb\left\{
   \sup_{v\in S^1}\|\mathcal A_{n,v,A}\|_{C^2([-R,R]^2)}
       >\delta
 \right\}=0.
\end{equation}
\end{lemma}

\begin{proof}
Write \(a=a_{n,v}(r,z)\). The radial identity gives
\[
 \frac1n e^{-\kappa_n|a|^2}e^{a\cdot Z_i-|a|^2/2}
 =e^{\xi_i-\kappa_n|a|^2-|Z_i-a|^2/2}.
\]
Cover \(S^1\) by \(O(\rho_n)\) arcs of angular length \(O(1/\rho_n)\), and on
an arc centered at \(v_\alpha\) write
\(v=\mathsf R_{y/\rho_n}v_\alpha\), with \(y\) in a fixed interval. On the
resulting fixed boxes in \((y,r,z)\), differentiation inserts a polynomial of
bounded degree in \(Z_i-\gamma_na\), with coefficients bounded uniformly in
\(n\). Consequently, for every multi-index \(|\beta|\le3\), a summand of the
corresponding derivative of \(\mathcal B_{n,v}\) is bounded by
\(C_Re^{-\ell_n^{\rm edge}}\). Exponential tilting, as in the proof of
Proposition~\ref{prop:two-d-compact-decomposition}, also gives
\[
 \sup_{\alpha,y,r,z}\sum_i
 \E\left|\partial^\beta
   \left[\frac1n e^{-\kappa_n|a|^2}
   e^{a\cdot Z_i-|a|^2/2}
   \mathbf1_{\{\xi_i<-\ell_n^{\rm edge}\}}\right]\right|^2
 \le \frac{C_R}{\rho_n}.
\]
Here \(\partial^\beta\) denotes derivatives in the scaled local variables
\((y,r,z)\), and the supremum ranges over their fixed coordinate boxes.
The corresponding sum of \(p\)th moments is at most
\(C_Re^{-(p-2)\ell_n^{\rm edge}}/\rho_n\). For any fixed \(p>3\),
Rosenthal's inequality therefore bounds the expected \(p\)th power of the
\(W^{3,p}\) norm on one box by
\(C_R\rho_n^{-p/2}+o(\rho_n^{-p/2})\). Summing over the
\(O(\rho_n)\) boxes still gives \(o(1)\). The Sobolev--Morrey embedding
\(W^{3,p}\hookrightarrow C^2\) on each fixed box, followed by Markov's
inequality, proves uniform \(C^2\) concentration over all \(v\in S^1\).

The mean is rotation invariant and exponential tilting gives
\[
 \E\mathcal B_{n,v}(r,z)
 =e^{-\kappa_n(\rho_n+r)^2}
   \Pbb\{|G+(\rho_n+r)e_1|^2
              \le\rho_n^2-2\ell_n^{\rm edge}\},
 \qquad G\sim\mathcal N(0,I_2).
\]
The boundary expansion in the proof of
Proposition~\ref{prop:two-d-compact-decomposition} is uniform after two
derivatives in \(r\); the mean has no \(z\)-dependence. Since
\(\ell_n^{\rm edge}/\rho_n\to0\) and
\(\kappa_n\rho_n^2\to\lambda\), the last display converges in \(C^2\) to
\(e^{-\lambda}\Phi(-r)\). This proves
\eqref{eq:support-d2-diffuse-c2}.

For the strip, let \(Y_i(v)=\rho_n\operatorname{ang}(U_i,v)\). Uniformly on
the same compact set and for \(|\beta|\le2\), the radial identity and the
elementary bound of a polynomial times a Gaussian give
\[
 |\partial_{r,z}^\beta \mathcal A_{n,v,A}(r,z)|
 \le C_R\sum_{i:\,-\ell_n^{\rm edge}\le\xi_i\le-A}
 e^{\xi_i}\{1+|Y_i(v)-z|^4\}e^{-c|Y_i(v)-z|^2}.
\]
Taking the supremum over \(z\in[-R,R]\) only changes the constants and the
Gaussian exponent. The slab-and-arc occupancy proof of
Lemma~\ref{lem:d2-low-atom-absorption} applies with the summable weight
\((1+q^4)e^{-cq^2}\) in place of \(e^{-q^2/2}\). It bounds the last display
by \(C_Re^{-A}+o_{\Pbb}(1)\), simultaneously for \(|\beta|\le2\), and proves
\eqref{eq:support-d2-strip-c2}.
\end{proof}

To convert each local maximizer into an exact contact count, we need
two-derivative convergence of the fitted KKT certificate; convergence of
likelihood values alone would not exclude nearby split atoms.
\begin{proposition}
\label{prop:support-one-contact-per-patch}
Suppose an NPMLE is localized to finitely many compact patches,
its rescaled edge mass is bounded, and every associated radial excess is at
least \(\eta>0\) from the relevant threshold. After the radial cutoffs are
removed, an inactive patch contains no support point and an active patch
contains exactly one, apart from an event of probability tending to zero.
Here active means \(x_{n,j}>s_n\) for \(d\ge3\) and
\(x_{n,j}>x_{\lambda_n}\) for \(d=2\).
\end{proposition}

\begin{proof}
It is enough to argue along subsequences. The radial cutoffs and threshold
buffer allow a further subsequence on which the number of candidate patches is
constant and all marks \(x_{n,j}\) converge, say to \(x_j\). Write
\(a_{n,j}(u)=a_{n,j}^{(\ge3)}(u)\) in \(d\ge3\), and
\(a_{n,j}(u)=a_{n,j}^{(2)}(r,z)\) for \(u=(r,z)\) in \(d=2\).
The candidate index set is fixed throughout this subsequence; only the lower
observational cutoff \(A\) will later vary.  For a localized measure
\(\nu=(\nu_1,\ldots,\nu_J)\), let
\[
 M_{n,i}(\nu)=1+\frac1n\sum_j
   \int\{q_{n,\tau,a_{n,j}(u)}(i)-1\}\,\nu_j(du).
\]
Moving mass from the central atom into patch \(j\) has derivative
\begin{equation}
\label{eq:support-finite-patch-certificate}
 \Gamma_{n,j,\nu}(u)
 =\frac1n\sum_{i=1}^n
   \frac{q_{n,\tau,a_{n,j}(u)}(i)-1}{M_{n,i}(\nu)}.
\end{equation}
It is nonpositive throughout the patch and vanishes at every support point.

The limiting certificate in an active \(d\ge3\) patch is
\[
    \Gamma_j^*(h)=e^{-|h|^2/2}-1;
\]
it has a unique zero at zero and Hessian \(-I_d\). In an active
two-dimensional patch it is
\[
 \Gamma_j^*(r,z)
 =(1-B_\lambda(r))
 \left\{\frac{R_{x_j}(r,z)}{\sup_{u,y}R_{x_j}(u,y)}-1\right\},
\]
whose unique zero is \((-a_\lambda,0)\) and whose Hessian there is
\(-\{1-e^{-\lambda}\Phi(a_\lambda)\}I_2\). A threshold buffer \(\eta\)
also gives a fixed negative certificate gap in every inactive patch.

We next verify that \eqref{eq:support-finite-patch-certificate} converges to
these limits in \(C^2\) on a smaller patch, uniformly over measures in a fixed
compact mass ball.  Fix the observational cutoff \(A\).  At the candidate
observations, completing the square and differentiating three times gives
uniform \(C^3\) convergence of the associated response.  Every response from
a different exposed observation carries the factor
\(e^{-c(\ell_n^{\rm sep})^2}\), and the exposed cloud has tight cardinality.

For the lower observations in fixed \(d\ge3\),
Lemma~\ref{lem:support-dge3-palm-c2} gives, simultaneously in all candidate
patches, \(C^2\) convergence of the punctured lower field to zero.  The lemma
is centered at \(Z_j/\gamma_n\), whereas the score-displacement patch is
centered at \(Z_j/\gamma_n-\tau\). Taking the Palm estimate on a slightly
larger compact set and using \(|\tau|=o_{\Pbb}(1)\) transfers the bound to the
score-displacement patch. The accompanying multiplier
\(e^{-\gamma_na\cdot\tau}\) is \(1+o_{\Pbb}(1)\) in \(C^2\), because
\(\rho_n|\tau|=o_{\Pbb}(1)\).  Thus this is a deterministic translation of an
already uniform \(C^2\) estimate and requires no conditioning on a
sample-centered exceedance set.

In \(d=2\), Lemma~\ref{lem:support-d2-c2-absorption} gives \(C^2\)
convergence of the diffuse background to \(e^{-\lambda}\Phi(-r)\) and removes
the strip between the receding and fixed radial cutoffs in the \(C^2\)
topology.
The estimate is uniform over the patch direction, so it remains valid for the
random directions \(U_j\). The KKT field differs from the field in
Lemma~\ref{lem:support-d2-c2-absorption}
by the anchor multiplier
\[
 e^{-a\cdot\Delta_n^{\rm anc}-a\cdot\bar Z_n}
 =e^{-\gamma_na\cdot\tau}.
\]
Lemma~\ref{lem:support-fixed-core-collapse} gives
\[
 \rho_n^2|\tau|^2
 =O_{\Pbb}\left\{\frac{(\log n)^{C_d+1}}{n\kappa_n}\right\}
 =o_{\Pbb}(1),
\]
and differentiation of the multiplier in the local coordinates only
introduces further factors \(O_{\Pbb}(|\tau|)\).
It is therefore \(1+o_{\Pbb}(1)\) in \(C^2\) on every fixed patch.

It remains to pass from the unfitted fields to the fitted certificate.  Take
\(A\ge A_0(C,R,M)\) from
Lemma~\ref{lem:support-uniform-lower-cloud}.  For every lower observation,
\[
 |M_{n,i}(\nu)-1|
 \le F_{n,i}(\nu)+M/n
 \le C_{C,R}Me^{-A}+o(1)\le\frac12,
\]
uniformly over measures of mass at most \(M\).  Hence
\(|M_{n,i}^{-1}-1|\le2|M_{n,i}-1|\).  Applying
Lemma~\ref{lem:support-uniform-lower-cloud} after each of
the at most two local differentiations gives
\begin{equation}
\label{eq:support-weighted-lower-c2}
 \sup_{\nu:\,\nu(\mathbb R^d)\le M}
 \left\|
  \frac1n\sum_{i:x_{n,i}\le-A}
 \{q_{n,\tau,a_{n,j}(\cdot)}(i)-1\}
 \{M_{n,i}(\nu)^{-1}-1\}
 \right\|_{C^2}
 =O_{\Pbb,C,R,J}(Me^{-A})+o_{\Pbb,A}(1),
\end{equation}
with a tight \(O_{\Pbb}(1)\) coefficient uniform for \(A\ge A_0\).
To see this directly, note that \(M_{n,i}(\nu)\) is fixed when the prospective
contact location \(u\) is varied.  Thus all derivatives fall on the numerator,
and \eqref{eq:support-lower-cloud-derivative-envelope}, together with
\(\max_i|M_{n,i}^{-1}-1|\le CMe^{-A}+o(1)\), gives the displayed bound.  The
constant \(-1\) at derivative order zero is controlled by the same maximum.
At a candidate observation, the denominator converges to its local-program
denominator.  At a noncandidate exposed observation, every candidate response
is \(o_{\Pbb,A}(1)\), uniformly in \(C^2\); the constant \(-1\) in the
numerator contributes only \(O_{\Pbb,A}(1/n)\) after summing over the tight
exposed cloud.  Thus the complete fitted certificate
converges sequentially, first as \(n\to\infty\) and then as \(A\to\infty\),
to the displayed limiting certificates.

We now make the order of the cutoff and argmax limits explicit. Proposition
\ref{prop:support-patch-coercivity} confines the total rescaled mass to a
deterministic compact interval on events of arbitrarily high probability.
Together with the fixed compact candidate-patch union, this gives a weakly
compact space \(\mathcal M_{M,R}\) of candidate measures. Let
\(\mathcal J^\star_{\boldsymbol x}\) be the sum of the limiting local programs
for a vector \(\boldsymbol x=(x_1,\ldots,x_J)\) of candidate marks, and let
\(\nu^\star_{\boldsymbol x}\) be its optimizer. The scalar optimizations above
show that this optimizer is unique and depends continuously on
\(\boldsymbol x\) as long as the marks remain outside the threshold buffer.

Choose the patch radius \(R\) so that every such optimizer is in the interior.
For a fixed small \(\delta>0\), let
\(\mathcal U_{\boldsymbol x}\) be its radius-\(\delta\) weak neighborhood in
\(\mathcal M_{M,R}\).  The mark vectors with \(|x_j|\le C\) and distance at
least \(\eta\) from the threshold form a compact set.  Joint continuity of
the local programs and uniqueness of their optimizers therefore give the
uniform separation
\begin{equation}
\label{eq:support-argmax-gap}
 c_{\delta,\eta,C,M,R,J}
 :=\inf_{\boldsymbol x}
 \left[
  \mathcal J^\star_{\boldsymbol x}(\nu^\star_{\boldsymbol x})
  -\sup_{\nu\in\mathcal M_{M,R}\setminus\mathcal U_{\boldsymbol x}}
      \mathcal J^\star_{\boldsymbol x}(\nu)
 \right]>0.
\end{equation}
Decrease \(\delta\), if necessary, so that every active patch has mass bounded
away from zero throughout \(\mathcal U_{\boldsymbol x}\).  All the objects in
\eqref{eq:support-argmax-gap} depend only on the threshold-buffered candidate
patches, not on the observational cutoff.  Choose the deterministic cutoff
\(A\ge A_0(C,R,M)\) so large that, by
\eqref{eq:support-finite-patch-iterated}, the finite-patch approximation error
exceeds one quarter of the gap in \eqref{eq:support-argmax-gap} with
arbitrarily small limiting probability.  Then take \(n\) large.  Uniform
convergence of the local kernels shows on the complementary event that every
finite-sample maximizer belongs to \(\mathcal U_{\boldsymbol x}\): otherwise
comparison with the pushforward of \(\nu^\star_{\boldsymbol x}\) would lose at
least half that gap.  In particular, every active finite-sample patch carries
positive mass and hence contains a KKT contact.

Increase this same deterministic \(A\), if necessary, so that
\eqref{eq:support-weighted-lower-c2} exceeds one quarter of either the
inactive-patch certificate gap or the absolute value of the largest active
limiting Hessian eigenvalue with arbitrarily small limiting probability.
These constants are uniform over the threshold-buffered compact mark set,
while the candidate space does not change.  An inactive patch then has no
contact.  In an active patch,
\(C^0\) convergence confines all contacts to a small convex neighborhood of
the unique limiting zero, and \(C^2\) convergence makes the finite-sample
certificate strictly concave there.  Since the certificate is nonpositive,
strict concavity rules out two distinct zeros.  Sending the two exceptional
probabilities to zero proves the claim along every subsequence.
\end{proof}

\subsubsection{Threshold pruning and the support law}

Only observations within a fixed radial distance of the activation threshold
can support atoms. The needed deleted-field estimate is a further consequence
of the Bennett--Morrey argument in
Section~\ref{sec:sharp-high-dimensional-estimates}.
\begin{lemma}
\label{lem:support-dge3-deleted-field}
Assume \(d\ge3\) and \(s_n=O(1)\). For fixed \(a>0\), define
\[
 S_t^{(a)}=\frac1n e^{-\kappa_n|t|^2}
 \sum_{i=1}^ne^{t\cdot Z_i-|t|^2/2}
 \mathbf1_{\{|Z_i|^2/2\le\log n+\lambda_n-a\}}.
\]
If \(\zeta\in(0,1)\) and \(e^a\zeta>1\), then
\[
    \Pbb\left\{\sup_{|t|\ge\rho_n/2}S_t^{(a)}\ge\zeta\right\}\to0.
\]
\end{lemma}

\begin{proof}
Put \(R_a^2/2=\log n+\lambda_n-a\) and
\(h=R_a-|t|\). On each fixed physical annulus \(|h|\le H\),
Lemma~\ref{lem:physical-offset-fixed-t} gives
\[
 \Pbb\{S_t^{(a)}\ge\zeta\}
 \le\left(\rho_n^{-1}e^{-\lambda_n}\right)^{
     e^{a+h^2/2}\zeta-o(1)}.
\]
Lemma~\ref{lem:scaled-derivative-morrey} reduces the annulus to a net of
cardinality \(\rho_n^{d-1+o(1)}\), with \(o(1)\) oscillation. Since
\(\rho_n^{-1}e^{-\lambda_n}=\rho_n^{-(d-1)+o(1)}\) and
\(e^a\zeta>1\), the union bound tends to zero.

For \(h\ge H\), Lemma~\ref{lem:growing-offset-fixed-t} gives
\[
 \Pbb\{S_t^{(a)}\ge\zeta/2\}
 \le\exp\{-c_de^{a+h^2/4}\log\rho_n\}.
\]
A polynomial-size Euclidean net and the same derivative control make the
union bound vanish after \(H\) is large. When \(|t|\ge R_a+H\), every
truncated summand decreases along its radial ray, so the outward region reduces
to the controlled boundary. These ranges cover \(|t|\ge\rho_n/2\).
\end{proof}

The two-dimensional analogue of the deleted-field estimate has an explicit
profile gap.  It is useful to record both the global residual bound and the
stronger bound on the candidate patches, where the associated extreme
observation has itself been deleted.
\begin{lemma}
\label{lem:support-d2-residual-gap}
Assume \(d=2\) and \(\lambda_n\to\lambda\in(0,\infty)\).  Fix \(\eta>0\),
put
\[
 \mathcal I_n(\eta)=\{i:\xi_i>x_{\lambda_n}-\eta\},
 \qquad
 \mathcal R_{n,\eta}(a)
 =\frac1n\sum_{i\notin\mathcal I_n(\eta)}q_{n,\tau,a}(i),
\]
and let \(\mathcal P_n(\eta,R)\) be the union of the radius-\(R\) local
patches associated with \(\mathcal I_n(\eta)\).  Then
\begin{align}
\label{eq:support-d2-global-residual-gap}
 \sup_{|a|\ge\rho_n/2}\mathcal R_{n,\eta}(a)
 &\le1-\frac12(1-e^{-\eta})(1-e^{-\lambda})+o_{\Pbb}(1),\\
\label{eq:support-d2-candidate-residual-gap}
 \sup_{a\in\mathcal P_n(\eta,R)}\mathcal R_{n,\eta}(a)
 &\le e^{-\lambda}+o_{\Pbb,R}(1)
\end{align}
for each fixed \(R\).
\end{lemma}

\begin{proof}
Fix \(C,K,A<\infty\) and work first on
\(\max_i\xi_i\le C\) and \(|\,|a|-\rho_n|\le K\).  Write
\(a=(\rho_n+r)v\), \(|r|\le K\).  Lemma
\ref{lem:d2-fixed-cutoff-decomposition}, its quantitative absorption bound,
and the anchor estimate give, uniformly in \(r\) and \(v\),
\begin{equation}
\label{eq:support-d2-residual-decomposition}
 \mathcal R_{n,\eta}((\rho_n+r)v)
 =B_{\lambda_n}(r)
  +\sum_{\substack{-A<\xi_i\le x_{\lambda_n}-\eta}}
    C_{\lambda_n,\xi_i}(r,Y_i(v))
  +O_{\Pbb,K}(e^{-A})+o_{\Pbb,A,K}(1).
\end{equation}
For fixed \(A,C\), the exposed points have tight cardinality and, with
probability tending to one, their pairwise scaled angular distances exceed a
deterministic \(\ell_n\to\infty\).  If there are at most \(J\) exposed points,
then for any \(v\) all but at most one have \(|Y_i(v)|\ge\ell_n/2\); uniformly
for \(|r|\le K\), their total contribution is at most
\(J e^{C+K^2/2}e^{-\ell_n^2/8}\).  Tightness of \(J\) therefore makes this
remainder \(o_{\Pbb,A,C,K}(1)\), uniformly in \(v\).  Since every retained
residual mark is at most
\(x_{\lambda_n}-\eta\),
\begin{align*}
 B_{\lambda_n}(r)
 +C_{\lambda_n,x_{\lambda_n}-\eta}(r,z)
 &=e^{-\eta}
   \{B_{\lambda_n}(r)+C_{\lambda_n,x_{\lambda_n}}(r,z)\}
   +(1-e^{-\eta})B_{\lambda_n}(r)\\
 &\le e^{-\eta}+(1-e^{-\eta})e^{-\lambda_n}\\
 &=1-(1-e^{-\eta})(1-e^{-\lambda_n}).
\end{align*}
Choose \(A\) so that the absorption error is less than one quarter of the
limiting gap and then take \(n\to\infty\).  This proves the first bound on
every fixed radial window, on \(\{\max_i\xi_i\le C\}\).  As \(|r|\to\infty\),
the spike term is bounded by \(e^{C-r^2/2+o(1)}\), while
\(B_{\lambda_n}(r)\) tends to \(e^{-\lambda}\) inward and to zero outward.
The diffuse- and tail-field estimates in the proof of
Proposition~\ref{prop:d-two-radial-compactness} make these bounds uniform as
\(K\to\infty\).  Finally,
\[
 \limsup_{n\to\infty}\Pbb\{\max_i\xi_i>C\}
 =1-\exp\{-e^{-C}\}\longrightarrow0
 \qquad(C\to\infty),
\]
which removes the upper cutoff and proves
\eqref{eq:support-d2-global-residual-gap}.

If \(a\in\mathcal P_n(\eta,R)\), its associated observation belongs to
\(\mathcal I_n(\eta)\) and is absent from the residual sum.  Every exposed
residual direction is then at scaled angular distance at least
\(\ell_n-R\), so the entire exposed sum in
\eqref{eq:support-d2-residual-decomposition} is \(o_{\Pbb,A,C,R}(1)\).
The remaining profile is at most
\(\sup_rB_{\lambda_n}(r)=e^{-\lambda_n}\).  Taking first \(n\to\infty\),
then \(A,C\to\infty\), proves
\eqref{eq:support-d2-candidate-residual-gap}.
\end{proof}

Combining the edge and bulk gaps prunes subthreshold points before the
mass-coercivity bound is used. Thus the number of candidate patches is tight
independently of the lower radial cutoff.
\begin{proposition}
\label{prop:support-threshold-pruning}
Fix \(\eta>0\), and set
\[
    \vartheta_n=
    \begin{cases}
       s_n,&d\ge3,\\
       x_{\lambda_n},&d=2.
    \end{cases}
\]
In the corresponding critical window, every noncentral support point of an
NPMLE lies, with probability tending to one after the radial
cutoffs are removed, in a patch of some deterministic radius
\(R=R(\eta)<\infty\) associated with an observation whose radial excess
satisfies \(x_{n,i}>\vartheta_n-\eta\). The number of these patches is
\(O_{\Pbb}(1)\). More precisely, put
\[
 \mathcal I_n(\eta)=\{i:x_{n,i}>\vartheta_n-\eta\},
\]
and let \(\mathcal P_n(\eta)\) denote the union of their radius-\(R\) patches.
There are a deterministic \(b<1\) and
random variables \(C_{0,n}=O_{\Pbb}(1)\) such that
\[
 \sup_{a\in\mathcal P_n(\eta)}
 \frac1n\sum_{i:x_{n,i}\le\vartheta_n-\eta}
 q_{n,\tau,a}(i)\le b,
 \qquad
 \sup_{a\in\mathcal P_n(\eta)}
 \max_{i:x_{n,i}>\vartheta_n-\eta}
 \frac1nq_{n,\tau,a}(i)\le C_{0,n}.
\]
\end{proposition}

\begin{proof}
Call the sum over \(i\notin\mathcal I_n(\eta)\) the residual field. For
\(d\ge3\), every residual
observation satisfies
\[
    |Z_i-\bar Z_n|^2/2\le b_{n,d}+s_n-\eta
       =\log n+\lambda_n-\eta.
\]
Uniformly in \(i\),
\[
 \frac{|Z_i|^2}{2}
 \le\frac{|Z_i-\bar Z_n|^2}{2}
    +\max_i|Z_i-\bar Z_n|\,|\bar Z_n|
    +\frac{|\bar Z_n|^2}{2}
 =\frac{|Z_i-\bar Z_n|^2}{2}+o_{\Pbb}(1).
\]
Thus, with probability tending to one, every residual observation lies below
\(\log n+\lambda_n-\eta/2\). Put
\[
    a_\eta=\eta/2,
    \qquad
    \zeta_\eta=\frac{1+e^{-a_\eta}}2.
\]
Since \(e^{a_\eta}\zeta_\eta>1\), Lemma
\ref{lem:support-dge3-deleted-field} applied with \(a=a_\eta\) gives a gap of
\((1-\zeta_\eta)/2\) below one throughout
\(|a|\ge\rho_n/2\), with probability tending to one. To pass from the raw
field in Lemma~\ref{lem:support-dge3-deleted-field} to the residual field, note that
 \[
 \frac1n\sum_{\rm residual}q_{n,\tau,a}(i)
 \le e^{-a\cdot\Delta_n^{\rm anc}-a\cdot\bar Z_n}S_a^{(a_\eta)}.
 \]
An NPMLE may be supported in the convex hull of the rescaled observations, so
\(|a|=O_{\Pbb}(\rho_n)\); on this range
\eqref{eq:support-high-d-anchor-error} and
\(\rho_n|\bar Z_n|=o_{\Pbb}(1)\) make the prefactor \(1+o_{\Pbb}(1)\).
Thus the residual field is at most
\((1+\zeta_\eta)/2<1\) throughout the \(d\ge3\) edge region.

In \(d=2\), Lemma~\ref{lem:support-d2-residual-gap} gives the explicit global
edge gap
\[
 g_{2,\eta}=\frac12(1-e^{-\eta})(1-e^{-\lambda})>0.
\]
On the candidate patches themselves it gives the stronger residual bound
\(e^{-\lambda}+o_{\Pbb}(1)\). Thus, in the first displayed inequality of the
proposition, one may take any fixed \(b\in(e^{-\lambda},1)\) for all
sufficiently large \(n\).

The bulk estimates in Proposition~\ref{prop:bulk-mgf-exclusion} for
\(d\ge3\), and Lemma~\ref{lem:d2-bulk-exclusion} for \(d=2\), retain a gap
of order \(\kappa_n\) on the region between the fixed core and
\(\rho_n/2\). The candidate set has tight cardinality by radial Poisson
convergence, and its maximum radial excess is tight. On an event of arbitrarily
high probability, both its cardinality and its normalized peak heights are
bounded by deterministic constants. The exact square completion
\eqref{eq:support-lower-cloud-square-completion} then bounds the total
candidate contribution outside radius \(R\) of the patch centers by
\(Ce^{-cR^2}\). Choose \(R\) so this is smaller than one quarter of the
relevant residual edge gap: \((1-\zeta_\eta)/2\) for \(d\ge3\), and
\(g_{2,\eta}\) for \(d=2\). The full central field consequently has a fixed
gap outside \(\mathcal P_n(\eta)\) throughout the edge region. Equation
\eqref{eq:support-high-d-anchor-error} transfers both gaps to the central
anchor. Finally, Lemma
\ref{lem:support-no-percolation} preserves the edge gap because
\(W_n=o_{\Pbb}(1)\) and the bulk gap because
\(W_n=o_{\Pbb}(\kappa_n)\). KKT equality therefore excludes every
support point outside the candidate patches.

On the fixed candidate patches,
Lemma~\ref{lem:support-dge3-deleted-field} for \(d\ge3\) and
Lemma~\ref{lem:support-d2-residual-gap} for \(d=2\) give the first displayed
inequality with deterministic \(b<1\). The tight upper radial cutoff,
square completion, and angular separation give the second with
\(C_{0,n}=O_{\Pbb}(1)\).
\end{proof}

We can now count the patches. The following buffered sandwich is the
fixed-dimensional input to the common edge-count reduction.
\begin{proposition}
\label{prop:fixed-dimensional-edge-sandwich}
Assume \(\lambda_n\to\lambda\in(0,\infty)\) when \(d=2\), or
\(s_n^\star\to s\in\mathbb R\) when \(d\ge3\) is fixed. With
\(x_{n,i}\) defined in \eqref{eq:fixed-support-radial-excess} and
\(\vartheta_n\) as in Proposition~\ref{prop:support-threshold-pruning}, for
every fixed \(\eta>0\),
\begin{equation}
\label{eq:support-count-sandwich}
 \#\{i:x_{n,i}>\vartheta_n+\eta\}
 \le\widehat K_n-1
 \le\#\{i:x_{n,i}>\vartheta_n-\eta\},
\end{equation}
apart from an event whose limiting probability is \(O(\eta)\).
Under either critical-window assumption, the NPMLE is unique with probability
tending to one.
\end{proposition}

\begin{proof}
First assume \(\lambda_n\to\lambda\in(0,\infty)\) when \(d=2\), or
\(s_n^\star\to s\in\mathbb R\) when \(d\ge3\). Fix \(\eta>0\). Proposition
\ref{prop:support-threshold-pruning} confines every noncentral support point
to the \(O_{\Pbb}(1)\) patches associated with
\(x_{n,i}>\vartheta_n-\eta\), and supplies the two hypotheses of Proposition
\ref{prop:support-patch-coercivity}. Hence the rescaled edge mass is
\(O_{\Pbb}(1)\). After choosing an upper radial cutoff and a deterministic
bound on the number of patches, Proposition
\ref{prop:support-finite-patch-expansion} and Proposition
\ref{prop:support-one-contact-per-patch} show that every active patch contains
exactly one fitted atom and every inactive candidate patch contains none.
These are precisely the inequalities in \eqref{eq:support-count-sandwich}. Removing
the radial cutoffs costs \(o(1)\); the only remaining exceptional event is a
radial point in the threshold buffer, whose probability is \(O(\eta)\).

On the event where Propositions~\ref{prop:support-threshold-pruning} and
\ref{prop:support-one-contact-per-patch} give the stated patch localization,
the NPMLE is also unique. If \(G\) and \(H\) are two NPMLEs, then
\((G+H)/2\) is also an NPMLE.  Lemma~\ref{lem:support-fixed-core-collapse},
applied to all three measures, forces their central atoms to coincide:
otherwise the central part of \((G+H)/2\) would have positive covariance.
Propositions~\ref{prop:support-threshold-pruning} and
\ref{prop:support-one-contact-per-patch}, applied to \((G+H)/2\), likewise
force the edge contact in every active patch to coincide.  Thus all NPMLEs
are supported on the same finite contact set.

The likelihood vectors at these contacts are linearly independent.  Indeed,
on the rows indexed by the active observations, the edge-contact block has
diagonal entries of order \(n\), bounded away from zero after division by
\(n\), and off-diagonal entries \(o(n)\).  The finite-patch estimates also
give an \(O_{\Pbb}(1)\) empirical average for each of the finitely many edge
columns.  Hence there is a noncandidate row on which every edge entry is
\(o(n)\); adjoining this row and the central column produces, after scaling
the edge columns by \(n^{-1}\), an asymptotically invertible triangular
minor.  Strict concavity of the log likelihood makes the fitted likelihood
vector common to all maximizers, and linear independence then makes its
representing weights unique.  The probability of the exceptional event is
\(O(\eta)+o(1)\); letting \(\eta\downarrow0\) proves the uniqueness claim.
\end{proof}

\subsection{Growing-dimensional support recovery}

This subsection proves the support identity for every growing-dimensional
regime in Theorem~\ref{thm:main-support-law}. In the Gaussian-bump coordinates
of Subsection~\ref{sec:mv-one-atom-geometry}, the buffered reductions below
identify every noncentral fitted atom with one radial exceedance while
retaining the exact centering correction.

\subsubsection{Exact edge coordinates}

Recall the exact radial coordinate \(\xi_{n,i}^\star\), scale
\(\delta_n^\star\), and identity \eqref{eq:exact-support-radius} from
Section~\ref{sec:mv-gamma-normalization}.
For fixed \(C,J<\infty\) and \(\eta>0\), recall the buffered event
\(\mathcal B_n(C,J,\eta)\) from \eqref{eq:exact-support-buffer}.
The centered radial Poisson limit makes the complement of such buffered
events negligible after \(C,J\to\infty\) and then \(\eta\downarrow0\).
It therefore remains to prove that every positive buffered excess creates
exactly one noncentral support point.

When \(d_nL_n=o(n^2)\), the exact coordinate agrees with the
\(2\log n\)-based coordinate used in Sections~\ref{sec:radial-transition}
and~\ref{sec:proof-program}. The next
elementary comparison lets us invoke the radial estimates from those sections
without carrying two normalizations through each proof.

\begin{lemma}
\label{lem:exact-old-edge-coordinate}
Suppose \(d_nL_n=o(n^2)\) and \(s_n^\star=O(1)\).  With \(s_n\),
\(\xi_{n,i}\), and \(\delta_n=a_n/\gamma_n\) as in
\eqref{eq:s-definition} and \eqref{eq:bump-height-xi},
\[
 s_n-s_n^\star=o(1),
 \qquad
 \frac{\delta_n^\star}{\delta_n}=1+o(1).
\]
Moreover, for every fixed \(C<\infty\),
\begin{equation}
\label{eq:exact-old-edge-comparison}
 \max_{i:\,|\xi_{n,i}^\star|\le C}
 \left|\xi_{n,i}^\star-(\xi_{n,i}-s_n)\right|
 \longrightarrow_{\Pbb}0.
\end{equation}
\end{lemma}

\begin{proof}
The second assertion is immediate from
\(\delta_n^\star=\sigma_n^2a_n/\gamma_n\).  The definitions give
\[
 s_n-s_n^\star
 =\frac{\gamma_n}{a_n}
   \left(L_n-\frac{R_n^\star}{\sigma_n^2}\right).
\]
Since \(L_n-R_n^\star=O(n^{-1})\) and
\(1-\sigma_n^2=n^{-1}\), its absolute value is
\(O\{b_n/(na_n)+\gamma_n/(na_n)\}\).  This tends to zero: it is
\(O(L_n/n)\) when \(d_n=O(L_n)\), while the superlogarithmic Gamma
expansions make it \(O(\sqrt{d_nL_n}/n)=o(1)\) otherwise.

In fact, the two excess coordinates satisfy the exact identity
\begin{equation}
\label{eq:exact-old-excess-identity}
 \xi_{n,i}-s_n
 =\sigma_n^2\xi_{n,i}^\star
  -\frac{L_n-R_n^\star}{\delta_n}.
\end{equation}
Now \(L_n-R_n^\star=O(n^{-1})\).  Also
\((n\delta_n)^{-1}=o(1)\): this is immediate when \(d_n=O(L_n)\), and in
the superlogarithmic range it is
\(O\{\sqrt{d_n/L_n}/n\}=o(L_n^{-1})\).
Equation \eqref{eq:exact-old-excess-identity}, together with
\(\sigma_n^2=1-n^{-1}\), proves
\eqref{eq:exact-old-edge-comparison}.
\end{proof}


For later use, if \(c_1,\ldots,c_n>0\), define
\[
 \Gamma_c(w)=\frac1n\sum_{i=1}^n
 c_i\exp\{w\cdot q_i-|w|^2/2\}.
\]
The Gibbs variational formula gives the weighted entropy dual
\begin{equation}
\label{eq:weighted-entropy-dual}
 \log\Gamma_c(w)
 =\sup_{p\in\Delta_n}
 \left\{
  w\cdot\sum_ip_iq_i-\frac{|w|^2}{2}
  -K(p\Vert\mathbf u_n)+\sum_ip_i\log c_i
 \right\}.
\end{equation}

The entropy argument used for one-atomicity must be made stable under the
fitted denominators before it can count contacts.  The following weighted
version supplies precisely that stability.  It is stated only in the random
entropy regime, where it will be used.

\begin{lemma}
\label{lem:weighted-entropy-localization}
Suppose \eqref{eq:entropy-regime} holds, \(s_n=O(1)\), and
\(d_nL_n=o(n^2)\).  Put
\[
 q_i=\frac{Z_i-\bar Z_n}{\sqrt{\gamma_n}},
 \qquad
 R_i=\frac{|q_i|^2}{2},
 \qquad
 \delta_n=\frac{a_n}{\gamma_n},
 \qquad
 \mu_n^{\rm off}=\max_{i\ne j}|q_i\cdot q_j|,
 \qquad
 \mathfrak k_n=\frac{\|(q_i\cdot q_j)_{i,j\le n}\|_{\rm op}}{2n}.
\]
Fix \(C,J<\infty\) and \(\eta>0\), and suppose
\[
 \max_i(R_i-L_n)/\delta_n\le C,\qquad
 \#\{i:R_i>L_n-\eta\delta_n\}\le J,\qquad
 \min_i|R_i-L_n|>\eta\delta_n.
\]
Let \(A=\{i:R_i>L_n\}\).  Suppose a random vector
\(\beta=(\beta_1,\ldots,\beta_n)\) satisfies, uniformly on this event,
\begin{equation}
\label{eq:weighted-entropy-alpha-profile}
 \beta_i=-(R_i-L_n)+e_{n,i}\quad(i\in A),
 \qquad
 \max_{i\in A}|e_{n,i}|
 +\max_{j\notin A}|\beta_j|
 =:\omega_n,
\end{equation}
where
\begin{equation}
\label{eq:weighted-entropy-omega-rate}
 \omega_n
 =
 O_{\Pbb}\left\{\frac{e^{\mu_n^{\rm off}}}n+\frac{\delta_nL_n}{n}\right\}
 =n^{-1+o_{\Pbb}(1)}
 =o_{\Pbb}(\delta_n).
\end{equation}
If
\[
 F_\beta(p)
 =
 \frac12\left|\sum_i p_iq_i\right|^2
 -K(p\Vert \mathbf u_n)+\sum_i p_i\beta_i,
\]
then, with probability tending to one, every \(p\in\Delta_n\) satisfying
\(F_\beta(p)\ge0\) has one of the following two forms:
\begin{equation}
\label{eq:weighted-entropy-two-patches}
 \|p-\mathbf u_n\|_1\le n^{-1/2+o_{\Pbb}(1)},
 \qquad\text{or}\qquad
 1-p_j\le n^{-1/2+o_{\Pbb}(1)}
 \quad\text{for some }j\in A.
\end{equation}
In the first case \(|\sum_i p_iq_i|=o_{\Pbb}(1)\); in the second,
\(|\sum_i p_iq_i-q_j|=o_{\Pbb}(1)\). Moreover, if \(p=p(w)\) is the
Gibbs vector of a weighted certificate with value at least one, then \(w\)
lies within \(n^{-1/2+o_{\Pbb}(1)}\) of zero in the first case and of
\(q_j\) in the second.
\end{lemma}

\begin{proof}
Lemma~\ref{lem:random-off-diagonal} and \eqref{eq:entropy-regime} imply that
the deterministic bounds in Proposition~\ref{prop:refined-deterministic-gram}
hold with probability tending to one. Replace
\(\mu_n^{\rm off}\) and \(\upsilon_n\) by deterministic upper bounds that
hold with probability tending to one and preserve their asymptotic
separation. For readability, continue to denote these deterministic bounds by
\(\mu_n^{\rm off}\) and \(\upsilon_n\). We may then choose
deterministic \(a_n^\circ\) so that
\[
 1+\mu_n^{\rm off}\ll a_n^\circ\ll
 \min\{\log(1/\upsilon_n),\log(n\delta_n),L_n\},
\]
on an event of probability tending to one, and then set
\[
 T_n=e^{a_n^\circ},
 \qquad
 L_ny_n=\{a_n^\circ\log(n\delta_n)\}^{1/2}.
\]
These choices satisfy \eqref{eq:refined-gram-conditions}. Since
\(\max_iR_i=L_n+O(\delta_n)\), the extra diagonal term in
\eqref{eq:heavy-energy} is
\[
 O(\delta_n)\sum_i b_i^2
 =o\left(L_n\sum_i b_i^2\right)
\]
and is absorbed by the heavy-entropy term in
\eqref{eq:heavy-entropy-scalar}.

First suppose \(\max_i p_i\le1-y_n\).  Repeat the light-heavy
decomposition in Proposition~\ref{prop:refined-deterministic-gram}. Explicitly,
put
\[
 f_i=np_i,\qquad \ell_i=f_i\wedge T_n,\qquad h_i=(f_i-T_n)_+,
 \qquad b_i=\frac{h_i}{n},
 \qquad K_0=\frac1n\sum_i\phi(\ell_i),\qquad
 z=\sum_i b_i,
\]
where \(\phi(x)=x\log x-x+1\).
The active diagonal excess is canceled by the active fitted weight:
because \(b_i^2\le p_i\),
\[
 \sum_{i\in A}(R_i-L_n)b_i^2+\sum_i p_i\beta_i
 \le\omega_n.
\]
Keeping half of the slack in \eqref{eq:light-energy} and
\eqref{eq:heavy-entropy-scalar} therefore gives
\begin{equation}
\label{eq:weighted-diffuse-slack}
 F_\beta(p)
 \le
 -cK_0-ca_n^\circ z+\omega_n.
\end{equation}
Thus \(F_\beta(p)\ge0\) implies
\(K_0+a_n^\circ z=O_{\Pbb}(\omega_n)\).  Moreover,
\[
 \|p-\mathbf u_n\|_1
 \le
 C\left\{
   \left(\frac{T_n\omega_n}{a_n^\circ}\right)^{1/2}
   +\frac{\omega_n}{a_n^\circ}
 \right\}
 \le n^{-1/2+o_{\Pbb}(1)}.
\]
The last inequality uses \(T_n=n^{o(1)}\) and
\eqref{eq:weighted-entropy-omega-rate}; in particular,
\(L_n(T_n\omega_n/a_n^\circ)^{1/2}=o_{\Pbb}(1)\).
The Gram bound then gives
\(|\sum_i p_iq_i|=o_{\Pbb}(1)\), proving the first alternative.

It remains to consider \(p_i=1-y>1-y_n\).  Write
\(\sum_jp_jq_j=(1-y)q_i+y\bar q\), as in the last part of the proof of
Proposition~\ref{prop:refined-deterministic-gram}.  Put
\[
 \mathfrak H_n
 =
 L_n-\mu_n^{\rm off}-O(\delta_n)-2L_ny_n-\omega_n.
\]
The preceding choices give
\(\mathfrak H_n\sim L_n\) and
\(e^{-\mathfrak H_n}=n^{-1+o(1)}=o(\delta_n)\). Convexity,
the entropy bound \(h_{\rm bin}(y)\le y\log(e/y)\), and the fitted weights give
\begin{equation}
\label{eq:weighted-near-vertex}
 F_\beta(p)
 \le
 \alpha_i
 +y\left\{\log\frac ey-\mathfrak H_n\right\},
\end{equation}
where \(\alpha_i=R_i-L_n+\beta_i\).

If \(i\notin A\), the first term on the right of
\eqref{eq:weighted-near-vertex} is at most
\(-\eta\delta_n+\omega_n\). The second term on the right of
\eqref{eq:weighted-near-vertex} is at most
\(e^{-\mathfrak H_n}=o(\delta_n)\).
Hence \(F_\beta(p)<0\).

If \(i\in A\), then \(|\alpha_i|\le\omega_n\).  For
\(y>e^{-\mathfrak H_n/2}\),
\[
 \log(e/y)-\mathfrak H_n\le-\mathfrak H_n/3
\]
for all large \(n\).  Nonnegative free energy consequently forces
\[
 y\le e^{-\mathfrak H_n/2}+\frac{3\omega_n}{\mathfrak H_n}
 \le n^{-1/2+o_{\Pbb}(1)}.
\]
Since \(\max_j|q_j|=O_{\Pbb}(\sqrt{L_n})\), this gives
\[
 \left|\sum_jp_jq_j-q_i\right|
 =y|\bar q-q_i|=o_{\Pbb}(1),
\]
which is the second alternative in
\eqref{eq:weighted-entropy-two-patches}.

For later use, define the weighted certificate and its Gibbs vector by
\[
 \Gamma_\beta(w)
 =\frac1n\sum_{i=1}^n
   \exp\left\{w\cdot q_i-\frac{|w|^2}{2}+\beta_i\right\},
 \qquad
 p_i(w)
 =\frac{\exp\{w\cdot q_i+\beta_i\}}
        {\sum_j\exp\{w\cdot q_j+\beta_j\}}.
\]
The exact identity
\[
 F_\beta\{p(w)\}-\log\Gamma_\beta(w)
 =
 \frac12\left|w-\sum_i p_i(w)q_i\right|^2
\]
shows that the two alternatives in
\eqref{eq:weighted-entropy-two-patches} also localize \(w\) whenever
\(\Gamma_\beta(w)\ge1\). On the diffuse branch,
\eqref{eq:weighted-diffuse-slack} gives
\(F_\beta\{p(w)\}=n^{-1+o_{\Pbb}(1)}\), while the Gram estimate and
\eqref{eq:weighted-entropy-two-patches} give
\(\lvert\sum_i p_i(w)q_i\rvert=n^{-1/2+o_{\Pbb}(1)}\). On an active
branch, \eqref{eq:weighted-near-vertex} gives
\(F_\beta\{p(w)\}=n^{-1+o_{\Pbb}(1)}\), and
\(\lvert\sum_i p_i(w)q_i-q_j\rvert=n^{-1/2+o_{\Pbb}(1)}\). The displayed
identity therefore proves the asserted spatial localization.
Moreover,
\[
 \nabla^2\log\Gamma_\beta(w)
 =\operatorname{Cov}_{p(w)}(q_i)-I.
\]
On the diffuse branch the covariance norm is at most
\(2\mathfrak k_n+O\{L_n\|p-\mathbf u_n\|_1\}=o_{\Pbb}(1)\); on an active branch it is
\(O\{L_n(1-p_i)\}=o_{\Pbb}(1)\).  Thus every relevant patch is strictly
log-concave for all large \(n\).
\end{proof}

The slowly growing range needs a local version of the truncated shell
argument.  Conditional first moments are enough.  Exponential tilting
shows that the supremum of the background on a unit edge patch has mean
\(\exp\{-\eps_nL_n+O(\sqrt{d_n}+\eps_n\sqrt{L_n})\}\), which beats the
number of upper centers even when \(d_n\) grows arbitrarily slowly.

\begin{lemma}
\label{lem:slow-punctured-c3}
Suppose
\[
 d_n\longrightarrow\infty,\qquad
 d_n\le(\log L_n)^2,\qquad s_n=O(1),
\]
and put \(A_n=\sqrt{d_n}\).  There is a random set
\(\mathcal U_n\), containing every index with
\(\xi_{n,i}-s_n>-\eta\), such that, with probability tending to one,
\[
 |\mathcal U_n|\le e^{2A_n},\qquad
 \min_{\substack{i,j\in\mathcal U_n\\i\ne j}}
 |q_i-q_j|^2=4L_n\{1+o(1)\},
\]
and, for every fixed \(R<\infty\) and \(0\le k\le3\),
\begin{equation}
\label{eq:slow-punctured-c3}
 \max_{j\in\mathcal U_n}\ \sup_{|w-q_j|\le R}
 \left\|
  \nabla^k\sum_{i\notin\mathcal U_n}
  h_ie^{-|w-q_i|^2/2}
 \right\|=o_{\Pbb}(1).
\end{equation}
\end{lemma}

\begin{proof}
Work first with the independent uncentered vectors
\(Q_i=Z_i/\sqrt{\gamma_n}\), and let \(\mathcal U_n^\circ\) contain the
indices whose uncentered radial coordinate is at least \(-A_n-1\).
Write \(h_i^\circ=n^{-1}e^{|Q_i|^2/2}\) for the corresponding uncentered
bump height.
The Gamma hazard estimate and Markov's inequality give
\[
 |\mathcal U_n^\circ|\le e^{2A_n}
\]
with probability tending to one. Conditional on the radii, their
directions are independent and uniform. Since
\(\log|\mathcal U_n^\circ|=O(A_n)=o(d_n)\), spherical concentration gives
the displayed separation, first for the \(Q_i\)'s and then for the
centered \(q_i\)'s.

Conditional on \(Q_j\), \(j\in\mathcal U_n^\circ\), all remaining vectors
are independent of \(Q_j\).  Write \(w=Q_j+h\), \(|h|\le R\).  For every
multi-index \(|\alpha|\le3\), exponential tilting from
\(\mathcal N(0,\gamma_n^{-1}I)\) to
\(\mathcal N(Q_j/\gamma_n,\gamma_n^{-1}I)\) gives
\begin{align*}
 &\sum_{i\ne j}\E\left[
  \sup_{|h|\le R}
  \left|\partial^\alpha
  \{h_i^\circ e^{-|Q_j+h-Q_i|^2/2}\}\right|
  \mathbf1_{\{i\notin\mathcal U_n^\circ\}}
  \,\middle|\,Q_j\right]\\
 &\qquad\le
 C_{\alpha,R}
 \exp\left\{
  -\frac{\eps_n|Q_j|^2}{2}
  +C_{\alpha,R}\bigl(\sqrt{d_n}+\eps_n|Q_j|+1\bigr)
 \right\}.
\end{align*}
Indeed, after tilting,
\(Q_i-Q_j\) is Gaussian with mean \(-\eps_nQ_j\), covariance
\(\gamma_n^{-1}I\), and the supremum over \(h\) contributes at most a
fixed polynomial in its norm times \(e^{R|Q_i-Q_j|}\).
Uniformly over \(j\in\mathcal U_n^\circ\),
\(|Q_j|^2/2=L_n+O(A_n)\), while
\[
 \eps_nL_n=(1-o(1))\frac{d_n}{2}\log\frac{2L_n}{d_n}.
\]
This dominates
\(A_n+\sqrt{d_n}+\eps_n\sqrt{L_n}\), since
\(d_n\to\infty\) and \(d_n\le(\log L_n)^2\).  Conditional Markov's
inequality is applied to the random set through
\[
 \Pbb\{\exists j\in\mathcal U_n^\circ:B_{n,j}\}
 \le
 \sum_{j=1}^n
 \E\left[
  \mathbf1_{\{j\in\mathcal U_n^\circ\}}
  \Pbb(B_{n,j}\mid Q_j)
 \right],
\]
where \(B_{n,j}\) is the failure of the displayed derivative bound.
Since
\(\E|\mathcal U_n^\circ|=e^{(1+o(1))A_n}\), the conditional exponential
moment bound in the preceding display makes this sum \(o(1)\).  Passing from coordinate derivatives
to the tensor norm costs at most \(O(d_n^k)\) for order \(k\le3\), whose
logarithm is negligible compared with \(\eps_nL_n\).  This proves the
uncentered version of \eqref{eq:slow-punctured-c3}.

Finally,
\[
 B_n(w)=e^{-w\cdot\bar Z_n/\sqrt{\gamma_n}}B_n^\circ(w),
 \qquad
 \frac{|\bar Z_n|}{\sqrt{\gamma_n}}
 =O_{\Pbb}\{\sqrt{d_n/n}\}.
\]
On the relevant balls \(|w|=O(\sqrt{L_n})\), this multiplicative factor
is \(1+o_{\Pbb}(1)\), while each of its first four derivatives is
\(o_{\Pbb}(1)\).  The uniform
centered--uncentered radial difference is \(o_{\Pbb}(1)\), so the
one-unit cutoff buffer shows that
\(\mathcal U_n:=\mathcal U_n^\circ\) contains every centered index above
\(-\eta\).  This proves all three assertions.
\end{proof}

\subsubsection{Uniform finite-contact continuation}

The support-count argument needs a joint expansion over all active edge
patches.  The following lemma supplies it.  The hypotheses are stated in
terms of absolute derivative envelopes; this is important because the
reciprocal fitted values used in the KKT certificate need not preserve
cancellation among derivatives of different bumps.

\begin{lemma}
\label{lem:uniform-finite-edge-expansion}
Let \(L=\log n\), and let \(q_1,\ldots,q_n\in\mathbb R^{d_n}\) satisfy
\[
 \sum_{j=1}^nq_j=0,\qquad \max_j|q_j|\le C_0\sqrt L.
\]
Suppose
\[
 \frac{|q_i|^2}{2}=L+\delta v_i,\qquad 0<\delta\le2,
 \qquad \bar\delta=1\wedge\delta.
\]
Let \(\mathcal U\subseteq[n]\) contain every \(i\) with \(v_i>-\eta\),
put \(A=\{i:v_i>0\}\), and assume
\[
 |A|\le J,\qquad \eta\le v_i\le C\quad(i\in A),\qquad
 v_i\le-\eta\quad(i\in\mathcal U\setminus A).
\]
For \(0\le k\le3\), define
\begin{equation}
\label{eq:finite-edge-derivative-envelope}
\mathcal E_{i,k}(r)
 =
 \sup_{|z|\le r}\frac1n\sum_{j\ne i}
 \{1+|q_j-q_i-z|^k\}
 \exp\left\{(q_i+z)\cdot q_j-\frac{|q_i+z|^2}{2}\right\}.
\end{equation}
Fix \(0<r_n\le R_n\le1/4\), and put
\(e_{k,n}=\max_{i\in\mathcal U}\mathcal E_{i,k}(2R_n)\).  Suppose that
\begin{equation}
\label{eq:finite-edge-envelope-assumptions}
 \max_{i\in\mathcal U}\mathcal E_{i,0}(2R_n)
   =o_{\Pbb}(\bar\delta),\qquad
 \max_{i\in\mathcal U}\mathcal E_{i,1}(2R_n)
   =o_{\Pbb}(r_n),
\end{equation}
and
\begin{equation}
\label{eq:finite-edge-higher-envelope-assumptions}
 \begin{aligned}
 \max_{i\in\mathcal U}
 \{\mathcal E_{i,2}(2R_n)+\mathcal E_{i,3}(2R_n)\}
 &=o_{\Pbb}(1),\qquad
 n\bar\delta\longrightarrow\infty,
 \qquad
 r_n^2=o(\bar\delta),\\
 \frac{\bar\delta L}{n}&=o(r_n),\qquad
 r_ne_{1,n}+r_n^2e_{2,n}+r_n^3e_{3,n}
 &=o_{\Pbb}(\bar\delta).
 \end{aligned}
\end{equation}

For \(\tau\in\mathbb R^{d_n}\), \(z_i\in\mathbb R^{d_n}\), and
\(m_i\ge0\), set
\begin{equation}
\label{eq:finite-edge-candidate}
 \Pi_{\tau,z,m}
 =
 \left(1-\frac{\sum_{i\in A}m_i}{n}\right)\delta_\tau
 +\frac1n\sum_{i\in A}m_i\delta_{q_i+z_i}.
\end{equation}
Uniformly for
\[
 |\tau|\le K\bar\delta\sqrt L/n,\qquad |z_i|\le R_n,\qquad
 0\le m_i\le K\bar\delta,
\]
its relative log likelihood has the expansion
\begin{equation}
\label{eq:joint-finite-edge-expansion}
 \ell_q(\Pi_{\tau,z,m})
 =
 -\frac n2|\tau|^2
 +\sum_{i\in A}
 \left[
  \log\{1+m_i e^{\vartheta_i(\tau,z_i)}\}-m_i
 \right]
 +\mathcal R_n(\tau,z,m),
\end{equation}
where
\[
 \vartheta_i(\tau,z)
 =
 \delta v_i-\tau\cdot q_i+\frac{|\tau|^2}{2}-\frac{|z|^2}{2}.
\]
The remainder is \(o_{\Pbb}(\bar\delta^2)\), uniformly on this set.  The
same expansion may be differentiated twice in the mass variables and
three times in the spatial variables.  In particular, uniformly,
\begin{align}
\label{eq:joint-finite-edge-derivatives}
 \|\partial_m\mathcal R_n\|
   &\le O_{\Pbb}(e_{0,n}+n^{-1})=o_{\Pbb}(\bar\delta),
 &\|\partial_m^2\mathcal R_n\|&=o_{\Pbb}(1),\notag\\
 \|\nabla_z\mathcal R_n\|
   &\le O_{\Pbb}(\bar\delta e_{1,n})
     =o_{\Pbb}(\bar\delta r_n),
 &\|\partial_m\nabla_z\mathcal R_n\|
   &=O_{\Pbb}(e_{1,n})=o_{\Pbb}(r_n),\notag\\
 \|\nabla_z^2\mathcal R_n\|+
 \|\nabla_z^3\mathcal R_n\|
   &\le O_{\Pbb}\{\bar\delta(e_{2,n}+e_{3,n})\}
     =o_{\Pbb}(\bar\delta),
 &\|\nabla_\tau\mathcal R_n\|
   &=o_{\Pbb}(\bar\delta\sqrt L),
\end{align}
with all mixed blocks interpreted in the maximum block norm.

There is a critical point of \(\ell_q\) in these patches, unique near the
locations and masses displayed in
\eqref{eq:finite-edge-contact-asymptotics}--\eqref{eq:finite-edge-mass-asymptotics},
for which
\begin{equation}
\label{eq:finite-edge-contact-asymptotics}
 \tau_n
 =
 -\frac1n\sum_{i\in A}(1-e^{-\delta v_i})q_i
 +o_{\Pbb}(\bar\delta\sqrt L/n),\qquad
 z_{n,i}=o_{\Pbb}(r_n),
\end{equation}
and
\begin{equation}
\label{eq:finite-edge-mass-asymptotics}
 m_{n,i}=1-e^{-\delta v_i}+o_{\Pbb}(\bar\delta).
\end{equation}
In particular, all edge masses are positive.  If
\(\widetilde\Pi=\Pi_{\tau_n,z_n,m_n}\), then
\[
 \Gamma_{\widetilde\Pi}(\tau_n)
 =\Gamma_{\widetilde\Pi}(q_i+z_{n,i})=1,\qquad
 \nabla\Gamma_{\widetilde\Pi}(\tau_n)
 =\nabla\Gamma_{\widetilde\Pi}(q_i+z_{n,i})=0.
\]
Writing \(c_j=M_j(\widetilde\Pi)^{-1}\), one has
\begin{equation}
\label{eq:finite-edge-denominator-profile}
 \max_{j\notin A}(\log c_j)_+
 =O_{\Pbb}(\bar\delta L/n),\qquad
 \log c_i=-\delta v_i+o_{\Pbb}(\bar\delta)\quad(i\in A).
\end{equation}
Moreover, each active patch has the \(C^2\) representation
\begin{equation}
\label{eq:finite-edge-certificate-c2}
 \Gamma_{\widetilde\Pi}(q_i+y)
 =
 A_{n,i}e^{-|y|^2/2}+E_{n,i}(y),\qquad
 A_{n,i}=1+o_{\Pbb}(1),
\end{equation}
where
\[
 \|E_{n,i}\|_{C^0(B(0,R_n))}=o_{\Pbb}(\bar\delta),
 \qquad
 \|E_{n,i}\|_{C^2(B(0,R_n))}=o_{\Pbb}(1).
\]
Thus the certificate is strictly log-concave on every active patch and
has exactly one level-one point there.  On every inactive patch indexed
by \(\mathcal U\setminus A\),
\[
 \sup_{|y|\le R_n}\Gamma_{\widetilde\Pi}(q_i+y)
 \le1-c_\eta\bar\delta.
\]
\end{lemma}

\begin{proof}
Write \(k_w(j)=\exp\{w\cdot q_j-|w|^2/2\}\), \(m_+=\sum_{i\in A}m_i\),
and \(w_i=q_i+z_i\).  Factoring out the moving central component gives
the exact identity
\[
 M_j(\Pi_{\tau,z,m})=k_\tau(j)(1+x_j),\qquad
 x_j=-\frac{m_+}{n}
 +\frac1n\sum_{i\in A}m_iK_{ij},
\]
where
\[
 K_{ij}
 =\frac{k_{w_i}(j)}{k_\tau(j)}
 =\exp\left\{(w_i-\tau)\cdot(q_j-\tau)
              -\frac{|w_i-\tau|^2}{2}\right\}.
\]
At the observation belonging to the \(i\)th patch,
\begin{equation}
\label{eq:finite-edge-own-response}
 \frac{K_{ii}}n=e^{\vartheta_i(\tau,z_i)}.
\end{equation}
For \(j\notin A\), positivity and
\eqref{eq:finite-edge-envelope-assumptions} give, uniformly over the
displayed parameter set,
\[
 \max_{j\notin A}|x_j|=o_{\Pbb}(1),\qquad
 \sum_{j\notin A}x_j=-m_++o_{\Pbb}(\bar\delta^2),
 \qquad
 \sum_{j\notin A}x_j^2=o_{\Pbb}(\bar\delta^2).
\]
Indeed, after putting \(a_i=w_i-\tau\), the positive response away from
the own observation equals
\[
 \frac{m_i}{n}\sum_{j\ne i}K_{ij}
 =
 m_i e^{-a_i\cdot\tau}
 \frac1n\sum_{j\ne i}e^{a_i\cdot q_j-|a_i|^2/2}
 =o_{\Pbb}(\bar\delta^2).
\]
The negative part is \(m_+/n\) per observation.  Omitting the at most
\(J\) active observations therefore costs \(O(\bar\delta/n)\), which is
\(o(\bar\delta^2)\) by \(n\bar\delta\to\infty\).  Positivity also gives
\[
 \frac1{n^2}\sum_{j\ne i}K_{ij}^2
 \le
 \left(\frac1n\max_{j\ne i}K_{ij}\right)
 \left(\frac1n\sum_{j\ne i}K_{ij}\right)
 \le O_{\Pbb}(e_{0,n}^2),
\]
which controls the quadratic Taylor remainder. Applying the same product
bound after inserting the polynomial factors from
\eqref{eq:finite-edge-derivative-envelope} yields
\begin{align}
 |\mathcal R_n|
 &\le O_{\Pbb}\{\bar\delta e_{0,n}+e_{0,n}^2
                   +\bar\delta/n\},\notag\\
 \|\partial_m\mathcal R_n\|
 &\le O_{\Pbb}(e_{0,n}+n^{-1}),
 &\|\partial_m^2\mathcal R_n\|&=o_{\Pbb}(1),\notag\\
 \|\nabla_z\mathcal R_n\|
 &\le O_{\Pbb}(\bar\delta e_{1,n}),
 &\|\partial_m\nabla_z\mathcal R_n\|
   &\le O_{\Pbb}(e_{1,n}),\notag\\
 \|\nabla_z^2\mathcal R_n\|+
 \|\nabla_z^3\mathcal R_n\|
 &\le O_{\Pbb}\{\bar\delta(e_{2,n}+e_{3,n})\}.
 \label{eq:finite-edge-remainder-bounds}
\end{align}
Differentiation with respect to \(\tau\) produces the same relative
displacements, together with terms bounded by \(\sqrt L/n\) from the
finitely many omitted observations.  Hence
\begin{equation}
\label{eq:finite-edge-tau-remainder}
 \|\nabla_\tau\mathcal R_n\|
 \le O_{\Pbb}\{\bar\delta e_{1,n}
                +\bar\delta\sqrt L\,e_{0,n}
                +\bar\delta\sqrt L/n\}
 =o_{\Pbb}(\bar\delta\sqrt L).
\end{equation}
These estimates prove the remainder assertions in
\eqref{eq:joint-finite-edge-derivatives}.

Since \(\sum_jq_j=0\),
\(\sum_j\log k_\tau(j)=-n|\tau|^2/2\).  Taylor expansion of
\(\sum_{j\notin A}\log(1+x_j)\), together with
\eqref{eq:finite-edge-own-response}, proves
\eqref{eq:joint-finite-edge-expansion}.  Differentiating the same exact
factorization proves \eqref{eq:joint-finite-edge-derivatives}.  The
finiteness of \(|A|\) makes every bound uniform over simultaneous masses
and locations.

It remains to justify the critical point without treating the fitted
denominators as fixed.  Define the contact map
\[
 \mathcal F(\tau,z,m)
 =
 \left(
  \nabla\Gamma_\Pi(\tau),\
  \{\nabla\Gamma_\Pi(q_i+z_i)\}_{i\in A},\
  \{\Gamma_\Pi(q_i+z_i)-\Gamma_\Pi(\tau)\}_{i\in A}
 \right).
\]
Put \(m_i^0=1-e^{-\delta v_i}\),
\(\tau^0=-n^{-1}\sum_i m_i^0q_i\), and
\(\widetilde m_i^0=1-e^{-\vartheta_i(\tau^0,0)}\).  Then
\[
 |\tau^0|=O(\bar\delta\sqrt L/n),\qquad
 \widetilde m_i^0-m_i^0=o(\bar\delta),
\]
and \(m_i^0\ge c_{\eta,C}\bar\delta\).  Use the anisotropic norms
\[
 \|(f_0,\{f_i\},\{g_i\})\|_{\rm out}
 =
 \frac{n|f_0|}{\bar\delta\sqrt L}
 +\max_i\frac{|f_i|}{r_n}
 +\max_i\frac{|g_i|}{\bar\delta}
\]
and
\[
 \|(\Delta\tau,\Delta z,\Delta m)\|_{\rm in}
 =
 \frac{n|\Delta\tau|}{\bar\delta\sqrt L}
 +\max_i\frac{|\Delta z_i|}{r_n}
 +\max_i\frac{|\Delta m_i|}{\bar\delta}.
\]
The likelihood derivatives and contact equations are related by
\[
 \partial_\tau\ell_q
 =n\left(1-\frac{m_+}{n}\right)\nabla\Gamma_\Pi(\tau),
 \qquad
 \nabla_{z_i}\ell_q=m_i\nabla\Gamma_\Pi(q_i+z_i),
 \qquad
 \partial_{m_i}\ell_q
 =\Gamma_\Pi(q_i+z_i)-\Gamma_\Pi(\tau).
\]

For the quantitative estimates, put
\[
 \begin{split}
 \chi_n={}&
 \frac{e_{0,n}+n^{-1}}{\bar\delta}
 +\frac{e_{1,n}}{r_n}
 +\frac{e_{0,n}+e_{1,n}}{\sqrt L}
 +\frac Ln,\\
 \zeta_n={}&
 \frac{e_{0,n}}{\bar\delta}
 +\frac{e_{1,n}}{r_n}+e_{2,n}+e_{3,n}
 +\frac{r_ne_{1,n}+r_n^2e_{2,n}+r_n^3e_{3,n}}{\bar\delta}\\
 &+\frac1{n\bar\delta}+\frac{\bar\delta L}{nr_n}.
 \end{split}
\]
Every term is \(o_{\Pbb}(1)\) by the hypotheses.  At
\((\tau^0,0,\widetilde m^0)\), the principal part of
\eqref{eq:joint-finite-edge-expansion} is stationary in the spatial and
mass coordinates.  Its central derivative is
\(O(\bar\delta L^{3/2}/n)\).  Hence
\eqref{eq:finite-edge-remainder-bounds} and
\eqref{eq:finite-edge-tau-remainder} give the explicit residual estimate
\begin{equation}
\label{eq:finite-edge-scaled-residual}
 \|\mathcal F(\tau^0,0,\widetilde m^0)\|_{\rm out}
 =O_{\Pbb}(\chi_n)=o_{\Pbb}(1).
\end{equation}

Let \(\alpha_i=e^{-\vartheta_i(\tau^0,0)}\), which is bounded above and
below, and define
\[
 \mathcal J_n(\Delta\tau,\Delta z,\Delta m)
 =
 \left(
 -\Delta\tau-\frac1n\sum_i\alpha_iq_i\Delta m_i,\
 \{-\Delta z_i\}_{i\in A},\
 \{-\Delta m_i-\alpha_iq_i\cdot\Delta\tau\}_{i\in A}
 \right).
\]
Direct differentiation of the exact factorization gives
\begin{equation}
\label{eq:finite-edge-scaled-jacobian}
 \|D\mathcal F(\tau^0,0,\widetilde m^0)-\mathcal J_n\|_{{\rm in}\to{\rm out}}
 =O_{\Pbb}(\zeta_n)=o_{\Pbb}(1).
\end{equation}
The operator \(\mathcal J_n\) and its inverse are uniformly bounded.  Indeed,
its only nontrivial Schur correction is bounded by
\[
 \left\|\frac1n\sum_{i\in A}\alpha_i^2q_iq_i^\top\right\|_{\rm op}
 =O_J(L/n)=o(1).
\]

The own-response term requires one additional contribution in the
Jacobian-variation estimate.  On a scaled ball of radius \(\rho\), one has
\(|z_i|\le\rho r_n\), and differentiating
\(-|z_i|^2/2\) in the mass equation contributes at most
\(C\rho r_n^2/\bar\delta\) in the scaled operator norm.  The remaining
principal and remainder terms are controlled by \(\rho\) and \(\zeta_n\).
Consequently, uniformly for \(0<\rho\le1\),
\begin{equation}
\label{eq:finite-edge-scaled-jacobian-variation}
 \sup_{\|(\tau,z,m)-(\tau^0,0,\widetilde m^0)\|_{\rm in}\le\rho}
 \|D\mathcal F-D\mathcal F(\tau^0,0,\widetilde m^0)\|_{{\rm in}\to{\rm out}}
 \le
 O_{\Pbb}\left\{\zeta_n+\rho+
                  \frac{\rho r_n^2}{\bar\delta}\right\}.
\end{equation}
Since \(r_n^2=o(\bar\delta)\), the right side is
\(o_{\Pbb}(1)+O(\rho)\).

We finally verify rather than assume the existence of a vanishing Newton
radius.  Since \(\chi_n\xrightarrow{\Pbb}0\), choose integers \(N_k\uparrow\infty\)
so that
\[
 \Pbb\{\chi_n>k^{-2}\}\le k^{-1}\qquad(n\ge N_k),
\]
and set \(\varrho_n=k^{-1}\) when \(N_k\le n<N_{k+1}\).  Then
\(\varrho_n\downarrow0\) and
\(\chi_n=o_{\Pbb}(\varrho_n)\).  Equations
\eqref{eq:finite-edge-scaled-residual}--
\eqref{eq:finite-edge-scaled-jacobian-variation}, the uniformly bounded
inverse of \(\mathcal J_n\), and Newton--Kantorovich give a unique zero in
the scaled ball of radius \(\varrho_n\).  Its correction has input norm
\(o_{\Pbb}(\varrho_n)\), and therefore
\[
 |\tau_n-\tau^0|=o_{\Pbb}(\bar\delta\sqrt L/n),
 \qquad
 \max_i|z_{n,i}|=o_{\Pbb}(r_n),
 \qquad
 \max_i|m_{n,i}-\widetilde m_i^0|=o_{\Pbb}(\bar\delta).
\]
Together with the definition of the base point, these estimates give
\eqref{eq:finite-edge-contact-asymptotics} and
\eqref{eq:finite-edge-mass-asymptotics}.  Notice in particular that the
central block is \(-I+o_{\Pbb}(1)\).  The Hessian
\(\nabla^2B_n(0)\approx-\eps_nI\) is the Hessian of the unfitted bump
field; it is not the Jacobian obtained when the central atom and all
fitted denominators move together.

Stationarity in \(\tau\) and \(z_i\) gives the gradient contact equations,
while stationarity in \(m_i\) makes all displayed certificate values
equal.  Since
\(\int\Gamma_{\widetilde\Pi}\,d\widetilde\Pi=1\), their common value is
one.  For \(j\notin A\),
\[
 M_j(\widetilde\Pi)
 \ge
 \left(1-\frac{\sum_i m_{n,i}}n\right)k_{\tau_n}(j),
\]
so \(\max_{j\notin A}(\log c_j)_+=O_{\Pbb}(\bar\delta L/n)\).
The own fitted value and
\eqref{eq:finite-edge-mass-asymptotics} give
\(\log c_i=-\delta v_i+o_{\Pbb}(\bar\delta)\).

Finally, separate the \(i\)th observation in the certificate:
\[
 \Gamma_{\widetilde\Pi}(q_i+y)
 =
 c_i e^{\delta v_i}e^{-|y|^2/2}
 +\frac1n\sum_{j\ne i}c_j
   e^{(q_i+y)\cdot q_j-|q_i+y|^2/2}.
\]
The positive fitted-weight bound and the absolute envelopes imply
\eqref{eq:finite-edge-certificate-c2}.  The own Gaussian has Hessian
\(-I\) at its peak, so the certificate is strictly log-concave on a
smaller patch containing its stationary contact.  For
\(i\in\mathcal U\setminus A\), its adjusted own height is at most
\(\exp\{-\eta\delta+O_{\Pbb}(\bar\delta L/n)\}\), and the punctured
background is \(o_{\Pbb}(\bar\delta)\).  This proves the inactive-patch
claim and completes the proof.
\end{proof}

\subsubsection{Support identity when \(d_n\log n=o(n^2)\)}

The finite-contact lemma is deterministic. We now verify its hypotheses
throughout the range \(d_nL_n=o(n^2)\) and express the resulting support
identity in the exact coordinate of Theorem~\ref{thm:main-support-law}.

\begin{proposition}
\label{prop:buffered-support-below-simplex}
Suppose
\[
 d_n\longrightarrow\infty,
 \qquad d_nL_n=o(n^2),
 \qquad s_n^\star=O(1).
\]
For every fixed \(C,J<\infty\) and \(\eta>0\),
\begin{equation}
\label{eq:buffered-support-below-simplex}
 \Pbb\left\{
  \mathcal B_n(C,J,\eta),\quad
  \widehat K_n\ne1+\#\{i:\xi_{n,i}^\star>0\}
 \right\}\longrightarrow0.
\end{equation}
On \(\mathcal B_n(C,J,\eta)\), the NPMLE is also unique apart from an event
whose probability tends to zero.
\end{proposition}

\begin{proof}
Lemma~\ref{lem:exact-old-edge-coordinate} allows us to use
\(v_{n,i}=\xi_{n,i}-s_n\) in place of \(\xi_{n,i}^\star\).  On a slightly
smaller buffer the two active sets agree.  Put
\(A_n=\{i:v_{n,i}>0\}\); throughout the proof \(|A_n|\le J\).

We first describe the finite perturbation common to all the dimension
ranges.  The centered directions are not conditionally independent, so we
embed \(A_n\) in an uncentered buffered cloud.  The uniform
centered--uncentered edge-coordinate estimates in
Lemmas~\ref{lem:centered-radial-max} and
\ref{lem:centered-radial-max-high-d} show that, with probability tending
to one, every \(i\in A_n\) belongs to
\[
 \mathcal V_n=\{i:\xi_{n,i}^\circ-s_n>-\eta/2\}.
\]
This cloud has \(O_{\Pbb}(1)\) points. Conditional on its uncentered
radii, the directions are independent and uniform. Since \(d_n\to\infty\),
spherical concentration gives the following separation for the
uncentered score vectors. The common empirical-centering shift is
\(o_{\Pbb}(\sqrt{L_n})\), so it does not change their leading pairwise inner
products or distances. Consequently,
\begin{equation}
\label{eq:finite-edge-angular-separation}
 \max_{i\ne j\in A_n}
 \frac{|q_i\cdot q_j|}{L_n}=o(1),
 \qquad
 |q_i-q_j|^2=4L_n\{1+o(1)\}.
\end{equation}
The pairwise estimates in \eqref{eq:finite-edge-angular-separation} also hold
between an active point and any fixed number of other points in the upper
radial cloud.  Hence a neighborhood of \(q_i\)
sees its own bump at order one and every other retained upper bump at
\(n^{-1+o(1)}\).

Let \(\delta_n=a_n/\gamma_n\) and
\(\bar\delta_n=1\wedge\delta_n\).  The regime expansions in
Section~\ref{sec:regime-map} give \(\delta_n\le2\) for all large \(n\).
They also verify the scale assumptions in
Lemma~\ref{lem:uniform-finite-edge-expansion}. If
\(d_n=O(L_n)\), then \(\bar\delta_n\) is bounded away from zero.  If
\(L_n=o(d_n)\), then
\[
 \delta_n=\sqrt{L_n/d_n}\{1+o(1)\},
 \qquad
 n\bar\delta_n
 =\frac{L_n}{\sqrt{d_nL_n/n^2}}\{1+o(1)\}\longrightarrow\infty,
\]
because \(d_nL_n=o(n^2)\).  Thus
\(n\bar\delta_n\to\infty\) in the slowly growing, direct-truncation,
crossover, and entropy ranges alike.
After replacing \(\eta\) by a smaller fixed constant, the coordinate
comparison and the buffer imply that every nonactive index has
\(v_{n,i}\le-\eta\).  We may therefore take \(\mathcal U=A_n\) in
Lemma~\ref{lem:uniform-finite-edge-expansion}.
We verify the hypotheses of
Lemma~\ref{lem:uniform-finite-edge-expansion} uniformly over the finitely
many active indices.

Fix a small constant \(r_0<1/4\).  In the slowly growing, direct
truncation, and crossover ranges, let \(R_n=r_0\).  In the entropy range,
let \(R_n=L_n^{-1/2}\).  Define the available-radius envelopes
\[
 \widetilde e_{k,n}
 =
 \max_{i\in A_n}\mathcal E_{i,k}(2R_n),
 \qquad 0\le k\le3,
\]
and put
\[
 \Lambda_n^{\rm loc}
 =\widetilde e_{1,n}\vee\frac{\bar\delta_nL_n}{n},
 \qquad
 U_n=R_n\wedge\sqrt{\bar\delta_n}.
\]
The earlier field estimates imply, in each range, that
\begin{equation}
\label{eq:newton-radius-separation}
 \frac{\Lambda_n^{\rm loc}}{U_n}=o_{\Pbb}(1),
 \qquad
 \widetilde e_{0,n}=o_{\Pbb}(\bar\delta_n),
 \qquad
 \widetilde e_{2,n}+\widetilde e_{3,n}=o_{\Pbb}(1).
\end{equation}

In the slowly growing range, the proof of
Lemma~\ref{lem:slow-punctured-c3} bounds the sum of the absolute
derivatives of the lower cloud.  The remaining upper cloud has at most
\(e^{2\sqrt{d_n}}\) points, and angular separation bounds its first three
derivative envelopes by
\[
 e^{2\sqrt{d_n}}L_n^{3/2}e^{-2L_n+o(L_n)}=o(1).
\]
Thus \(\widetilde e_{0,n}=o_{\Pbb}(1)\) and
\(\sum_{k=1}^3\widetilde e_{k,n}=o_{\Pbb}(1)\).  Since
\(\bar\delta_n=1+o(1)\), \(U_n\) is bounded away from zero, and
\(L_n/n=o(1)\), these estimates give \eqref{eq:newton-radius-separation}.

In the direct truncation range, positivity of the Gaussian bumps and the
proof of Proposition~\ref{prop:moderate-low-cloud} give
\begin{equation}
\label{eq:direct-available-envelope}
 \widetilde e_{k,n}
 \le
 O_{\Pbb}(\delta_nL_n^{-10+k/2})+e^{-L_n+o(L_n)},
 \qquad 0\le k\le3.
\end{equation}
When \(\bar\delta_n\) vanishes, the regime condition
\(d_n(\log d_n)^2=o(L_n^3)\) and the superlogarithmic expansion imply
\(\bar\delta_n\ge L_n^{-1+o(1)}\).  Hence
\(\widetilde e_{1,n}/U_n=o_{\Pbb}(1)\) and
\((\bar\delta_nL_n/n)/U_n=o(1)\). The bounds in
\eqref{eq:direct-available-envelope} for \(k=0,2,3\) prove the other two assertions in
\eqref{eq:newton-radius-separation}.

In the crossover range, \eqref{eq:crossover-separation} gives
\begin{equation}
\label{eq:crossover-absolute-envelope}
 \widetilde e_{k,n}
 \le
 L_n^{k/2}e^{-L_n/2+O(\sqrt{L_n})}
 =o(\delta_n^m)
 \qquad(0\le k\le3)
\end{equation}
for every fixed \(m\).  Since
\(\bar\delta_n\asymp(\log L_n)/L_n\) and
\(U_n\asymp\sqrt{\bar\delta_n}\), the superpolynomial bound and
\(\sqrt{L_n\log L_n}/n=o(1)\) give
\eqref{eq:newton-radius-separation}.

Finally, in the entropy range put
\[
 \mu_n^{\rm off}=\max_{i\ne j}|q_i\cdot q_j|.
\]
The pairwise Gram bound gives
\begin{equation}
\label{eq:entropy-local-absolute-envelope}
 \widetilde e_{k,n}
 \le
 \frac{e^{\mu_n^{\rm off}+O_{\Pbb}(1)}L_n^{k/2}}n,
 \qquad 0\le k\le3.
\end{equation}
Writing \(A_n^{\rm off}=e^{\mu_n^{\rm off}}/n\), the rates in
Lemma~\ref{lem:random-off-diagonal} imply
\[
 A_n^{\rm off}=o_{\Pbb}(\delta_n),\qquad
 A_n^{\rm off}L_n=o_{\Pbb}(1),\qquad
 A_n^{\rm off}L_n^{3/2}=o_{\Pbb}(1).
\]
Since \(U_n^{-1}\le\sqrt{L_n}+\bar\delta_n^{-1/2}\),
\(A_n^{\rm off}\sqrt{L_n/\bar\delta_n}
=\{(A_n^{\rm off}/\bar\delta_n)(A_n^{\rm off}L_n)\}^{1/2}
=o_{\Pbb}(1)\).  It follows that
\(\widetilde e_{1,n}/U_n=o_{\Pbb}(1)\); also
\((\bar\delta_nL_n/n)/U_n=o(1)\), since it is at most
\(\bar\delta_nL_n^{3/2}/n+\sqrt{\bar\delta_n}L_n/n\).  The three rates
for \(A_n^{\rm off}\), applied to
\eqref{eq:entropy-local-absolute-envelope}, prove the remaining assertions
in \eqref{eq:newton-radius-separation}.

We now choose the Newton scale.  From
\eqref{eq:newton-radius-separation}, choose integers \(M_k\uparrow\infty\)
so that
\[
 \Pbb\left\{\Lambda_n^{\rm loc}/U_n>k^{-2}\right\}\le k^{-1}
 \qquad(n\ge M_k),
\]
and set \(c_n=k^{-1}\) for \(M_k\le n<M_{k+1}\).  Then \(c_n\downarrow0\)
and
\[
 \Lambda_n^{\rm loc}=o_{\Pbb}(c_nU_n).
\]
Set \(r_n=c_nU_n\).  Then
\[
 r_n\le R_n,\qquad
 \widetilde e_{1,n}=o_{\Pbb}(r_n),\qquad
 \frac{\bar\delta_nL_n}{n}=o(r_n),\qquad
 \frac{r_n^2}{\bar\delta_n}\le c_n^2\longrightarrow0.
\]
Because the envelopes decrease with the patch radius, the three inequalities
in \eqref{eq:newton-radius-separation} verify
\eqref{eq:finite-edge-envelope-assumptions}.  Moreover,
\[
 r_n\widetilde e_{1,n}=o_{\Pbb}(r_n^2)=o_{\Pbb}(\bar\delta_n),
 \qquad
 r_n^2\widetilde e_{2,n}+r_n^3\widetilde e_{3,n}
 =o_{\Pbb}(\bar\delta_n),
\]
so every condition in
\eqref{eq:finite-edge-higher-envelope-assumptions} holds.

Lemma~\ref{lem:uniform-finite-edge-expansion} therefore constructs a
measure \(\widetilde\Pi\) with one central contact \(\tau_n\) and one
contact \(w_{n,i}\) in each active patch.  Its edge weights satisfy
\begin{equation}
\label{eq:below-simplex-contact-asymptotics}
 |\tau_n|=O_{\Pbb}(\bar\delta_n\sqrt{L_n}/n),\qquad
 w_{n,i}=q_i+o_{\Pbb}(r_n),\qquad
 n\pi_{n,i}
 =1-e^{-\delta_nv_{n,i}}+o_{\Pbb}(\bar\delta_n).
\end{equation}
The buffer keeps every displayed mass positive.  In the entropy range,
retaining the explicit right side of
\eqref{eq:entropy-local-absolute-envelope} in the contact-map proof gives
\begin{equation}
\label{eq:entropy-local-contact-rate}
 n\pi_{n,i}
 =
 1-e^{-\delta_nv_{n,i}}
 +O_{\Pbb}\left\{
   \frac{e^{\mu_n^{\rm off}}}n+\frac{\delta_nL_n}{n}
 \right\}.
\end{equation}

It remains to check that the finite perturbation has not created a contact
elsewhere.  If \(c_j\) denotes the reciprocal fitted value at observation
\(j\), then
\begin{equation}
\label{eq:below-simplex-weight-perturbation}
 (\log c_j)_+
 =
 O_{\Pbb}\left\{\frac{\bar\delta_nL_n}{n}\right\}
 \quad(j\notin A_n),
 \qquad
 \log c_i=-\delta_nv_{n,i}+o_{\Pbb}(\bar\delta_n)
 \le-c_\eta\bar\delta_n
 \quad(i\in A_n).
\end{equation}
In the entropy range the contact equation gives the sharper cancellation
\begin{equation}
\label{eq:entropy-contact-cancellation}
 \max_{i\in A_n}
 \left|\delta_nv_{n,i}+\log c_i\right|
 +\max_{j\notin A_n}|\log c_j|
 =
 O_{\Pbb}\left\{
  \frac{e^{\mu_n^{\rm off}}}n+\frac{\delta_nL_n}{n}
 \right\}
 =n^{-1+o_{\Pbb}(1)}.
\end{equation}
We first control the inactive denominators without using the contact
equation.  If \(m=\sum_{i\in A_n}n\pi_{n,i}=O_{\Pbb}(\delta_n)\), then
\[
 M_j
 =
 \left(1-\frac mn\right)
 e^{\tau_n\cdot q_j-|\tau_n|^2/2}
 +\frac1n\sum_{i\in A_n}n\pi_{n,i}
 e^{w_{n,i}\cdot q_j-|w_{n,i}|^2/2}.
\]
Here
\(|\tau_n|=O_{\Pbb}(\delta_n\sqrt{L_n}/n)\), so the central term differs
from one by \(O_{\Pbb}(\delta_nL_n/n)\).  For \(j\notin A_n\), each edge
response in the last sum is at most
\(n^{-1}e^{\mu_n^{\rm off}+o_{\Pbb}(1)}\).  Consequently,
\[
 \max_{j\notin A_n}|M_j-1|
 =
 O_{\Pbb}\left\{
  \frac{\delta_nL_n}{n}
  +\frac{\delta_ne^{\mu_n^{\rm off}}}{n^2}
 \right\},
\]
which proves the inactive part of
\eqref{eq:entropy-contact-cancellation}.

We may now evaluate the exact level-one contact equation in the \(i\)th
edge patch.  If \(w_i\) is its contact and \(E_i\) is the contribution of
all observations other than \(i\) to the certificate there, then the own
term gives the exact identity
\begin{equation}
\label{eq:entropy-exact-contact-cancellation}
 \log M_i
 =
 \delta_nv_{n,i}-\frac{|w_i-q_i|^2}{2}-\log(1-E_i).
\end{equation}
The derivative envelope
\eqref{eq:entropy-local-absolute-envelope}, the \(C^2\) certificate
expansion in Lemma~\ref{lem:uniform-finite-edge-expansion}, and the local
stationarity equation give
\[
 |w_i-q_i|^2
 =O_{\Pbb}\left\{
  \frac{e^{2\mu_n^{\rm off}}L_n}{n^2}
  +\frac{\delta_n^2L_n^2}{n^2}
 \right\},
 \qquad
 E_i
 =O_{\Pbb}\left\{
  \frac{e^{\mu_n^{\rm off}}}n+\frac{\delta_nL_n}{n}
 \right\}.
\]
Indeed, a retained upper response is at most
\(e^{\mu_n^{\rm off}}/n\), and the
central displacement and lower fitted values contribute
\(O_{\Pbb}(\delta_nL_n/n)\).  Substitution in
\eqref{eq:entropy-exact-contact-cancellation} proves the active part of
\eqref{eq:entropy-contact-cancellation}.  These estimates are uniform
because \(|A_n|\le J\), and the argument is noncircular: the inactive
denominator bound was obtained before the certificate contribution
\(E_i\) was estimated.
The positive part in
\eqref{eq:below-simplex-weight-perturbation} is \(o_{\Pbb}\) of the
smallest gap used in each one-atom exclusion proof.  The following
four-region argument verifies these comparisons while proving the
weighted exclusion rather than applying the unweighted results formally.
Write
\[
 \Gamma_{\widetilde\Pi}(w)
 =
 \Gamma_I(w)+\sum_{i\in A_n}G_i(w),
 \qquad
 G_i(w)=\frac{c_i}{n}e^{w\cdot q_i-|w|^2/2},
\]
where \(I=A_n^c\) and
\[
 \Gamma_I(w)
 =\frac1n\sum_{j\in I}c_j
   e^{w\cdot q_j-|w|^2/2},
 \qquad
 D_I(w)
 =\frac1n\sum_{j\in I}
   e^{w\cdot q_j-|w|^2/2}.
\]
If
\[
 \omega_n
 =\max_{j\in I}(\log c_j)_+,
\]
then, pointwise and after every truncation used in the cited proofs,
\begin{equation}
\label{eq:weighted-inactive-domination}
 \Gamma_I(w)\le e^{\omega_n}D_I(w).
\end{equation}
Thus every Bernstein or net upper bound for the inactive field changes by
at most \(O(\omega_n)\).

There are four regions.  First consider a fixed small ball around an active
contact.  The own term satisfies the exact relation
\[
 G_i(w_{n,i})=1-E_i,
\]
where \(E_i\) is the remaining certificate contribution, rather than
being one by itself.  The \(C^2\) representation
\eqref{eq:finite-edge-certificate-c2} gives
\(\nabla^2\log\Gamma_{\widetilde\Pi}\preceq-cI\) throughout this ball.
Since the contact is stationary and has value one, it is the unique
level-one point there.

Next take a corresponding ball around an inactive upper point \(q_j\).
Its adjusted own peak is at most
\[
 \frac{c_j}{n}e^{|q_j|^2/2}
 \le
 \exp\{\omega_n+\delta_nv_{n,j}\}
 \le1-c_\eta\bar\delta_n.
\]
The punctured \(C^3\) estimates
\eqref{eq:moderate-low-cloud},
\eqref{eq:entropy-local-absolute-envelope},
\eqref{eq:slow-punctured-c3}, and
\eqref{eq:crossover-absolute-envelope},
in their respective ranges, make all remaining terms
\(o_{\Pbb}(\bar\delta_n)\).  Hence an inactive upper patch contains no
contact.

In the central region, the weighted Hessian differs from the unweighted
one by
\[
 O_{\Pbb}(\omega_nL_n)
 +O_{\Pbb}\left\{\frac{\bar\delta_nL_n^2}{n}\right\}
 =o_{\Pbb}(\eps_n).
\]
To obtain this bound, write each fitted value as \(M_j=1+x_j\).
The negative part of \(x_j\) is uniformly
\(O_{\Pbb}(\bar\delta_nL_n/n)\), while
\[
 \frac1n\sum_j(x_j)_+
 \le
 C\frac{\bar\delta_n}{n}
 +O_{\Pbb}\left\{\frac{\bar\delta_nL_n}{n}\right\};
\]
the first term follows by averaging each edge response and using
\(B_n(w_{n,i})=O_{\Pbb}(1)\). Multiplication by
\(\max_j|q_j|^2=O_{\Pbb}(L_n)\) gives the displayed covariance
perturbation.
It is therefore negative definite in the microscopic ball around the
translated central contact.  On the surrounding annuli,
\eqref{eq:weighted-inactive-domination} raises the field by at most
\(O_{\Pbb}(\omega_n)\).  We compare the active profiles with the annular
gap explicitly.  Uniformly for \(i\in A_n\),
\[
 G_i(w)
 \le
 \exp\left\{-\frac12(|q_i|-|w|)^2+o_{\Pbb}(1)\right\},
 \qquad
 |q_i|=\sqrt{2L_n}\{1+o_{\Pbb}(1)\}.
\]
Let \(g_n(r)=1-e^{-\eps_nr^2/2}\), and let \(g_{0,n}\) be the smallest
gap at the inner endpoint of the annular argument.  If
\(\eps_nr^2\le1\), then
\[
 g_n(r)\ge c\eps_nr^2\ge c g_{0,n},
 \qquad
 G_i(w)\le n^{-1+o(1)}=o(g_{0,n}).
\]
The last comparison is \(n^{-1}=o(\eps_n^3/L_n^3)\) in the slow and
direct truncation ranges and is immediate in the crossover range. If
\(\eps_nr^2\ge1\) and \(r\le|q_i|/2\), then \(g_n(r)\ge c\) and
\(G_i(w)\le n^{-1/4+o(1)}\).  Hence
\begin{equation}
\label{eq:active-central-corridor}
 \sum_{i\in A_n}G_i(w)=o_{\Pbb}\{g_n(|w|)\}
 \quad
 (r_{0,n}\le|w|\le\tfrac12\sqrt{2L_n}).
\end{equation}
Together with the annular estimates established separately in the slow,
direct-truncation, crossover, and entropy regimes, this proves that the
central contact is the only contact in the inner region.

Finally consider the complement of the central and upper patches.
Choose their fixed radius \(R=R(J)\) so large that
\(Je^{-R^2/2}<1/20\); the patches are disjoint by angular separation.
On the annulus between the implicit-function neighborhood and radius
\(R\), the punctured \(C^2\) bound and the exact level-one contact show
that the Gaussian profile generated by the nearest active observation has
its unique maximum at the constructed contact.  In the slow range this
uniform statement is precisely Lemma~\ref{lem:slow-punctured-c3}; in the
direct-truncation and crossover ranges it follows from the local \(C^2\)
estimates used to construct the active contacts.
The inactive part is again controlled by
\eqref{eq:weighted-inactive-domination}.  Conditional angular separation
implies that at most one upper profile is nearest; outside its patch that
profile loses a fixed Gaussian factor, while the sum of all other upper
profiles is \(n^{-1+o_{\Pbb}(1)}\) in the direct truncation range,
\(e^{-L_n/2+o(L_n)}\) in the crossover range, and \(o_{\Pbb}(1)\) in the
slow range by Lemma~\ref{lem:slow-punctured-c3}.  Thus the outer-region
decompositions from the unweighted proofs remain valid after the fitted
denominators are inserted: equation
\eqref{eq:weighted-inactive-domination} controls the inactive field, and
the active tails retain their original bounds.  These estimates exclude
every remaining level-one point.
More concretely, in the crossover proof
\eqref{eq:crossover-outer-net} bounds the clipped inactive field by
\(0.9+o_{\Pbb}(1)\); choose \(R\) so that the active tails total less
than \(0.05\).  In the slow proof, choose \(R\) so that those tails are
smaller than \(\rho/4\), where \(\rho\) is the fixed slack in
Lemma~\ref{lem:slow-outside-shell}; the fitted-denominator error is also
\(o_{\Pbb}(\rho)\). The direct-truncation outer estimate has
an \(o_{\Pbb}(\delta_n)\) inactive background and is stronger still.

In the entropy range,
Lemma~\ref{lem:weighted-entropy-localization}, applied with
\(\beta_j=\log c_j\), directly localizes every level-one contact to the
central patch or an active upper patch;
\eqref{eq:entropy-contact-cancellation} verifies all of the lemma's weight
hypotheses.  The strict log-concavity asserted there leaves only the
constructed contacts.

The fitted-denominator error and the tails of the active profiles are
smaller than the KKT gaps in each unweighted localization argument.  More
precisely,
\[
 \frac{L_n}{n}=o\left(\frac{\eps_n^3}{L_n^3}\right)
 \quad\text{in the slow range},\qquad
 \frac{\delta_nL_n}{n}+\delta_n\mathfrak j_n=o(g_{0,n})
 \quad\text{in the direct range},
\]
where \(g_{0,n}\) is from
Proposition~\ref{prop:moderate-inner-field}; in the crossover range,
\(\delta_nL_n/n=o(L_n^{-3})\) and
\(e^{-L_n/2+o(L_n)}=o(\delta_n^m)\) for every fixed \(m\).
Thus the weighted estimates retain the gaps and negative curvature used in
all four regions.

This verifies the global KKT inequality for the constructed measure. Strict
concavity of the log likelihood makes its fitted vector unique. The contact likelihood
vectors are linearly independent: their active-coordinate block has
diagonal entries \(n^{1+o(1)}\) and off-diagonal entries \(n^{o(1)}\),
while the central vector is bounded.  Since
 \(n^{-1}\sum_jk_{w_{n,i}}(j)=B_n(w_{n,i})=O_{\Pbb}(1)\) for each of the
finitely many edge columns, their combined column sum is \(O_{\Pbb}(n)\).
Choose a deterministic \(K_n\to\infty\) so slowly that \(K_n=o(n)\).
Then
\[
 \#\left\{j:\max_{i\in A_n}k_{w_{n,i}}(j)>K_n\right\}
 \le
 \frac1{K_n}\sum_{i\in A_n}\sum_{j=1}^n k_{w_{n,i}}(j)
 =o_{\Pbb}(n).
\]
Since \(|A_n|\le J\), with probability tending to one there is an inactive
observation at which every edge entry is at most \(K_n=o(n)\).  Adjoin that
observation as one extra row; its central entry is
\(1+o_{\Pbb}(1)\).  After scaling the edge columns by \(n^{-1}\), their
active-coordinate block is asymptotically diagonal with diagonal bounded
away from zero, while all edge entries in the added row are \(o(1)\).
The resulting square minor is therefore nonsingular.  Every
maximizing mixing measure must therefore use all and only the
\(1+|A_n|\) contacts.

Finally, pass to a subsequence on which
\(\mathfrak r_n=d_n(\log d_n)^2/L_n^3\) tends to zero, stays bounded above
and away from zero, or tends to infinity.  These cases invoke,
respectively, the slow/direct estimates (split at
\(d_n=(\log L_n)^2\)), the crossover estimates, and the entropy estimates.
They exhaust \(d_nL_n=o(n^2)\) and prove
\eqref{eq:buffered-support-below-simplex}.
\end{proof}

\subsubsection{Ordered contacts in the simplex regime}

When \(d_nL_n/n^2\) is bounded below, the centered sample is close to a
regular simplex and the isolated-bump model no longer applies.  A radial
crossing then creates a contact near an ordered Potts state.  Write
\[
 \chi_n=\frac{n-2}{n-1},
 \qquad
 p_i^{(i)}=\frac{n-1}{n},
 \qquad
 p_j^{(i)}=\frac1{n(n-1)}\quad(j\ne i).
\]
The identities \(\sum_jq_j=0\) and
\eqref{eq:potts-radial-threshold} give, exactly,
\begin{equation}
\label{eq:ordered-potts-excess}
 \sum_jp_j^{(i)}q_j=\chi_nq_i,
 \qquad
 \frac12\left|\sum_jp_j^{(i)}q_j\right|^2
 -K(p^{(i)}\Vert \mathbf u_n)
 =\chi_n^2\delta_n^\star \xi_{n,i}^\star.
\end{equation}
Thus the sign of the ordered-state free energy is the sign of the exact
radial excess, with no centering approximation.

\begin{proposition}
\label{prop:buffered-simplex-support}
Let \(n\ge3\), let \(q_1,\ldots,q_n\) be a centered triangular array, and
write
\[
 R_i=\frac{|q_i|^2}{2},
 \qquad
 u_i=\frac{R_i-R_n^\star}{\delta_n},
 \qquad
 G_n=(q_i\cdot q_j)_{i,j\le n}.
\]
Suppose \(\delta_n\downarrow0\), \(\delta_n=O(L_n/n)\), and, for fixed
\(C,J<\infty\) and \(\eta>0\),
\[
 \max_i u_i\le C,
 \qquad
 \#\{i:u_i>-\eta\}\le J,
 \qquad
 \min_i|u_i|>\eta.
\]
Assume also that
\begin{equation}
\label{eq:buffered-simplex-assumptions}
 \max_{i\ne j}|q_i\cdot q_j|=o(L_n^{-1}),
 \qquad
 \|G_n\|_{\rm op}=O(L_n),
 \qquad
 \frac{C_n}{\delta_nL_n^2}\longrightarrow0,
 \qquad
 \frac{C_n^2}{n\delta_n}\longrightarrow0,
\end{equation}
where
\[
 C_n=\max_{i\ne j}
 \left|q_i\cdot q_j+\frac{|q_i|^2}{n-1}\right|.
\]
Then, for all sufficiently large \(n\), every NPMLE has
\begin{equation}
\label{eq:buffered-simplex-count}
 |\supp(\Pi)|=1+\#\{i:u_i>0\}.
\end{equation}
There is one central contact and, for each \(i\) with \(u_i>0\), one
contact \(w_i=\chi_nq_i+o(1)\); its weight satisfies
\begin{equation}
\label{eq:buffered-simplex-mass}
 n\pi_i=\frac{n}{n-1}\delta_nu_i+o(\delta_n).
\end{equation}
All remainders are uniform when \(C,J,\eta\) in the buffer are fixed.
Moreover, the NPMLE is unique.
\end{proposition}

\begin{proof}
Put \(A=\{i:u_i>0\}\), so \(|A|\le J\), and set
\(w_i^0=\chi_nq_i\) for \(i\in A\).  The buffer and the definition of
\(R_i\) imply
\begin{equation}
\label{eq:terminal-simplex-derived-rates}
 \max_i|q_i|=O(\sqrt{L_n}),
 \qquad
 C_n=o(L_n^{-1}),
 \qquad
 \delta_nL_n=o(1),
 \qquad
 \frac{C_n^2}{n}=o(\delta_n).
\end{equation}
Indeed, the first assertion follows from
\(R_i\le R_n^\star+C\delta_n=O(L_n)\), and the second follows from the
first condition in \eqref{eq:buffered-simplex-assumptions} and
\(|q_i|^2/(n-1)=O(L_n/n)\).  The last two assertions follow directly from
the assumed bounds on \(\delta_n\) and \(C_n\).

If \(A=\varnothing\), then
\(\max_iR_i\le R_n^\star-\eta\delta_n\).  Proposition
\ref{prop:local-simplex-criterion}, with \(g_n=\eta\delta_n\), applies by
\eqref{eq:buffered-simplex-assumptions}.  Its proof is strict away from
\(p=\mathbf u_n\), so the one-atom certificate equals one only at the
origin.  The KKT conditions then force every NPMLE to be \(\delta_0\), and
the conclusion follows.  We henceforth assume that \(A\ne\varnothing\).

We begin by locating the maxima of the one-atom certificate.  If
\[
 e_{ij}=q_i\cdot q_j+\frac{|q_i|^2}{n-1},
\]
then \(|e_{ij}|\le C_n\), and direct substitution gives
\begin{align}
\label{eq:potts-response-expansion}
 k_{w_i^0}(i)
 &= (n-1)\exp\{(2\chi_n-\chi_n^2)\delta_nu_i\},\notag\\
 k_{w_i^0}(j)
 &= \frac1{n-1}
    \exp\{-(\chi_n^2+2\chi_n/(n-1))\delta_nu_i
                 +\chi_ne_{ij}\},\qquad j\ne i.
\end{align}
At \(u_i=0\) and \(e_{ij}=0\), these responses are \(n-1\) and
\((n-1)^{-1}\), respectively, and their average is one.  Moreover,
\(\sum_{j\ne i}e_{ij}=0\), by centering.  The linear row error therefore
cancels when the off-diagonal responses are summed.  Expanding
\eqref{eq:potts-response-expansion}, using \(C_n=o(1)\), gives
\begin{equation}
\label{eq:potts-local-height}
 \log B_n(w_i^0)
 =\chi_n^2\delta_nu_i
  +O_{C,J}\left(\delta_n^2+\frac{C_n^2}{n}\right)
 =\chi_n^2\delta_nu_i+o(\delta_n).
\end{equation}
The coefficient can also be read directly from the exact dual identity
\eqref{eq:ordered-potts-excess}.

The Gibbs weights at \(w_i^0\) differ from \(p^{(i)}\) by
\(O_{C,J}\{(C_n+\delta_n)/n\}\) in total variation.  Differentiating the
log-sum-exp representation therefore gives, uniformly for \(i\in A\),
\begin{align}
\label{eq:potts-local-derivatives}
 |\nabla\log B_n(w_i^0)|
 &\le K_{C,J}\frac{\sqrt{L_n}}n(C_n+\delta_n),\notag\\
 \nabla^2\log B_n(w_i^0)
 &=-I+O(L_n/n)+o(1)=-I+o(1).
\end{align}
The same estimate holds uniformly on \(B(w_i^0,L_n^{-2})\): moving by
\(L_n^{-2}\) changes every exponent by \(O(L_n^{-3/2})\), while the total
Gibbs mass outside coordinate \(i\) remains \(O(n^{-1})\).  The
Newton--Kantorovich theorem now gives a unique critical point \(\bar w_i\)
in that ball, with
\begin{equation}
\label{eq:quantitative-potts-location}
 |\bar w_i-w_i^0|
 \le K_{C,J}\frac{\sqrt{L_n}}n(C_n+\delta_n).
\end{equation}
It is a strict local maximum, and Taylor's theorem gives
\begin{equation}
\label{eq:potts-corrected-height}
 h_i:=\log B_n(\bar w_i)
 =\chi_n^2\delta_nu_i+\varepsilon_{n,i},
 \qquad
 \max_{i\in A}|\varepsilon_{n,i}|=o(\delta_n).
\end{equation}
More explicitly, the remainder is bounded by a constant times
\[
 \delta_n^2+\frac{C_n^2}{n}
 +\frac{L_n(C_n+\delta_n)^2}{n^2},
\]
which is \(o(\delta_n)\) by
\eqref{eq:terminal-simplex-derived-rates} and
\eqref{eq:buffered-simplex-assumptions}.  In particular,
\(h_i\ge c_\eta\delta_n\) for every \(i\in A\) and all sufficiently large
\(n\).

For \(i\in A\), let
\(v_i=(k_{\bar w_i}(1),\ldots,k_{\bar w_i}(n))\).  Equations
\eqref{eq:potts-response-expansion}--\eqref{eq:potts-local-derivatives}
give
\begin{equation}
\label{eq:potts-response-gram}
 A_{ij}^{(0)}
 :=\frac1{n^2}\sum_{r=1}^n\{v_i(r)-1\}\{v_j(r)-1\}
 =\beta_n^{\rm P}\mathbf1_{\{i=j\}}+e_{n,ij},
 \qquad
 \max_{i,j\in A}|e_{n,ij}|=o(1),
\end{equation}
where
\[
 \beta_n^{\rm P}=\frac{(n-2)^2}{n(n-1)}.
\]
For instance, at the unperturbed ordered state the diagonal expression is
exactly
\[
 \frac1{n^2}\left\{(n-2)^2
 +(n-1)\left(\frac1{n-1}-1\right)^2\right\}=\beta_n^{\rm P},
\]
whereas the cross expressions are \(O(n^{-1})\).  The perturbations in
\eqref{eq:potts-response-expansion} and
\eqref{eq:quantitative-potts-location} contribute \(o(1)\), uniformly over
the at most \(J\) active indices.

We next construct exact contacts and masses.  For
\(m=(m_i)_{i\in A}\) and
\(x=(\tau,(w_i)_{i\in A})\), define
\[
 \Pi_{x,m}
 =\left(1-\frac1n\sum_{i\in A}m_i\right)\delta_\tau
  +\frac1n\sum_{i\in A}m_i\delta_{w_i},
\]
and write \(\Gamma_{x,m}\) for its certificate.  Consider the stationarity
map
\[
 \mathcal S(x,m)
 =\left(
   \nabla\log\Gamma_{x,m}(\tau),
   \bigl(\nabla\log\Gamma_{x,m}(w_i)\bigr)_{i\in A}
  \right).
\]
At \(m=0\) and
\(x=x^0:=(0,(\bar w_i)_{i\in A})\), this map vanishes.  Equip the product
of location spaces with the maximum of the Euclidean block norms.  For
each fixed \(M<\infty\), the response table gives, uniformly on
\[
 |m|_\infty\le M\delta_n,
 \qquad |\tau|\le L_n^{-2},
 \qquad |w_i-\bar w_i|\le L_n^{-2},
\]
the bounds
\begin{equation}
\label{eq:terminal-location-map-bounds}
 \|D_x\mathcal S(x^0,0)^{-1}\|\le2,
 \quad
 \|D_x\mathcal S(x,m)-D_x\mathcal S(x^0,0)\|=o(1),
 \quad
 \|D_m\mathcal S(x,m)\|
 \le K_{C,J,M}\frac{\sqrt{L_n}}n.
\end{equation}
To verify the first two bounds, the diagonal location blocks are
\(-I+o(1)\) by \eqref{eq:potts-local-derivatives}, while the central block
is
\[
 \frac1n\sum_{r=1}^nq_rq_r^\top-I=-I+O(L_n/n).
\]
At zero mass, the derivative of an active equation with respect to the
central location is also negligible: if the central atom is at \(\tau\),
then
\[
 \log\Gamma_{\delta_\tau}(w)
 =\log B_n(w-\tau)-(w-\tau)\cdot\tau,
\]
so at \((\tau,w)=(0,\bar w_i)\),
\[
 D_\tau\nabla_w\log\Gamma_{\delta_\tau}(w)
 =-\nabla^2\log B_n(\bar w_i)-I=o(1).
\]
All other cross-location blocks vanish at zero mass and are
\(O_{C,J}(\delta_nL_n)=o(1)\) on the displayed neighborhood.  Finally,
one differentiation with respect to \(m_j\) produces a normalized product
of a central or ordered response and its first moment.  The response table
bounds this by \(K\sqrt{L_n}/n\), proving the last estimate in
\eqref{eq:terminal-location-map-bounds}.

The quantitative implicit-function theorem now supplies unique smooth
locations satisfying the stationarity equations.  Write
\[
 x(m)=\bigl(\tau(m),(w_i(m))_{i\in A}\bigr).
\]
They obey
\begin{equation}
\label{eq:terminal-location-map-rate}
 |\tau(m)|+\max_{i\in A}|w_i(m)-\bar w_i|
 \le K_{C,J,M}\frac{\sqrt{L_n}}n|m|_1.
\end{equation}
Define the logarithmic contact-height differences
\[
 \mathcal H_i(m)
 =\log\Gamma_{x(m),m}\{w_i(m)\}
  -\log\Gamma_{x(m),m}\{\tau(m)\}.
\]
Stationarity permits the envelope theorem.  At zero mass,
\(\mathcal H_i(0)=h_i\), and direct differentiation gives
\begin{equation}
\label{eq:terminal-mass-contact-jacobian}
 -\partial_{m_j}\mathcal H_i(0)
 =A_{ij}^{(0)}+o(1)
 =\beta_n^{\rm P}\mathbf1_{\{i=j\}}+o(1).
\end{equation}
The additional \(o(1)\), beyond
\eqref{eq:potts-response-gram}, includes the movement of the central
contact; by \eqref{eq:terminal-location-map-bounds} its contribution is
\(O(L_n/n)\).

The response averages in \eqref{eq:potts-response-gram} also control the
second derivative.  In particular,
for any \(i,j,k\in A\),
\[
 \frac1{n^3}\sum_{r=1}^n
  \{1+v_i(r)\}\{1+v_j(r)\}\{1+v_k(r)\}=O_{C,J}(1).
\]
The fitted denominators stay between \(1/2\) and \(2\) on
\(|m|_\infty\le M\delta_n\), so
\begin{equation}
\label{eq:terminal-mass-second-derivative}
 \sup_{|m|_\infty\le M\delta_n}
 \|D_m^2\mathcal H(m)\|=O_{C,J,M}(1).
\end{equation}
Consequently,
\begin{equation}
\label{eq:potts-local-program}
 \mathcal H(m)=h-A_nm+R_n(m),
 \qquad
 A_n=\beta_n^{\rm P}I_{|A|}+o(1),
 \qquad
 |R_n(m)|_\infty\le K_{C,J}|m|_\infty^2.
\end{equation}
Choose \(M\) larger than \(2C\).  Since \(u_i\in[\eta,C]\) on \(A\),
the map
\(m\mapsto A_n^{-1}\{h+R_n(m)\}\) is a contraction on a ball of radius
\(o(\delta_n)\) about \(A_n^{-1}h\).  It therefore has a unique fixed
point in \(|m|_\infty\le M\delta_n\), and
\begin{equation}
\label{eq:terminal-quantitative-mass}
 m_i
 =\frac{\chi_n^2}{\beta_n^{\rm P}}\delta_nu_i+o(\delta_n)
 =\frac{n}{n-1}\delta_nu_i+o(\delta_n).
\end{equation}
Every component is positive for all sufficiently large \(n\).  The
resulting measure \(\widetilde\Pi=\Pi_{x(m),m}\) has one stationary central
contact and one stationary contact in each active patch, all at the same
certificate level.  Since
\(\int\Gamma_{\widetilde\Pi}\,d\widetilde\Pi=1\), that level equals one.

We finally prove the global KKT inequality.  Put
\(c_r=M_r(\widetilde\Pi)^{-1}\).  The response expansion and the mass
formula imply
\begin{equation}
\label{eq:potts-fitted-weight-profile}
 |\log c_r|=O_{C,J}(\delta_nL_n/n)\quad(r\notin A),
 \qquad
 \log c_i=-\chi_n\delta_nu_i+o(\delta_n)\quad(i\in A).
\end{equation}
Indeed, \eqref{eq:terminal-location-map-rate} gives
\(|\tau(m)|=O_{C,J}(\delta_n\sqrt{L_n}/n)\) and
\(\max_i|w_i(m)-\bar w_i|=O_{C,J}(\delta_n\sqrt{L_n}/n)\).
For \(r\notin A\), the central response is
\[
 e^{\tau(m)\cdot q_r-|\tau(m)|^2/2}
 =1+O_{C,J}(\delta_nL_n/n),
\]
the total mass removed from the central atom is \(O_{C,J}(\delta_n/n)\),
and all active responses together contribute \(O_{C,J}(\delta_n/n^2)\).
Thus \(M_r=1+O_{C,J}(\delta_nL_n/n)\).  At an active observation \(i\),
\eqref{eq:potts-response-expansion} and
\eqref{eq:terminal-quantitative-mass} instead give
\[
 M_i
 =1+\frac{n-2}{n}\,
      \frac{n}{n-1}\delta_nu_i+o(\delta_n)
 =1+\chi_n\delta_nu_i+o(\delta_n).
\]
Taking negative logarithms proves
\eqref{eq:potts-fitted-weight-profile}.

For clarity, put
\[
 F_{\log c}(p)
 =\frac12\left|\sum_rp_rq_r\right|^2
  -K(p\Vert\mathbf u_n)+\sum_rp_r\log c_r.
\]
The weighted entropy dual \eqref{eq:weighted-entropy-dual} shows that
\(\Gamma_{\widetilde\Pi}(w)\ge1\) only if its Gibbs vector \(p=p(w)\)
satisfies \(F_{\log c}(p)\ge0\).  If
\(\max_rp_r\le1-L_n^{-2}\), the stability inequality
\eqref{eq:entropy-stability} may be applied after retaining the active
diagonal terms.  Indeed, writing \(r=p-\mathbf u_n\), their contribution is at
most
\[
 \sum_{i\in A}\delta_nu_i r_i^2
 \le\chi_n\sum_{i\in A}\delta_nu_i p_i+O_{C,J}(\delta_n/n^2),
\]
where the scalar inequality
\((p_i-1/n)^2\le\chi_np_i+n^{-2}\) holds for \(0\le p_i\le1\).
This is canceled by \eqref{eq:potts-fitted-weight-profile}.  The remaining
active-weight error is \(o(\delta_n)\), and the positive weight on
inactive coordinates is \(O(\delta_nL_n/n)=o(\delta_n)\).  Consequently,
\eqref{eq:entropy-stability} and
\eqref{eq:buffered-simplex-assumptions} give, for a deterministic sequence
\(\varpi_n\downarrow0\),
\[
 F_{\log c}(p)
 \le
 -\frac{c}{L_n}|p-\mathbf u_n|_1^2+\varpi_n\delta_n,
 \qquad \varpi_n\longrightarrow0.
\]
Thus nonnegative weighted free energy forces
\[
 |p-\mathbf u_n|_1=o\{\sqrt{\delta_nL_n}\}
 =o(L_n/\sqrt n).
\]
Because \(\sum_rq_r=0\), the corresponding
\(\mu(p)=\sum_rp_rq_r\) satisfies
\[
 |\mu(p)|
 \le \max_r|q_r|\,|p-\mathbf u_n|_1
 =o(L_n\sqrt{\delta_n})
 =o(L_n^{3/2}/\sqrt n)=o(1).
\]
For the Gibbs vector of an actual certificate point, the exact square
identity is
\begin{equation}
\label{eq:weighted-certificate-square-identity}
 \log\Gamma_{\widetilde\Pi}(w)
 =F_{\log c}\{p(w)\}-\frac12|w-\mu\{p(w)\}|^2.
\end{equation}
Therefore \(\Gamma_{\widetilde\Pi}(w)\ge1\) implies \(w=o(1)\) in the
diffuse case.

Suppose instead that \(p_i=1-y>1-L_n^{-2}\), and decompose \(p\) as in
the proof of Proposition~\ref{prop:local-simplex-criterion}: write
\(p_j=yr_j\) for \(j\ne i\) and
\[
 v_r=\sum_{j\ne i}\left(r_j-\frac1{n-1}\right)q_j,
 \qquad
 \mathbf u_{n-1}=\left(\frac1{n-1},\ldots,\frac1{n-1}\right).
\]
Applying the local-simplex decomposition from that proposition with
\(R_i=R_n^\star+\delta_nu_i\) gives
\begin{align}
\label{eq:weighted-potts-localization}
 &\frac12\left|\sum_rp_rq_r\right|^2-K(p\Vert \mathbf u_n)\notag\\
 &\qquad\le
 \delta_nu_i c_y^2
 -\{k_n(y)-R_n^\star c_y^2\}
 +y|c_y|C_n,
 \qquad c_y=1-\frac{ny}{n-1}.
\end{align}
Before the row error is bounded, the decomposition from the proof of
Proposition~\ref{prop:local-simplex-criterion} also retains
\begin{equation}
\label{eq:weighted-potts-conditional-entropy}
 -yK(r\Vert \mathbf u_{n-1})+\frac{y^2}{2}|v_r|^2
 \le-\{1-o(1)\}yK(r\Vert \mathbf u_{n-1}),
\end{equation}
uniformly for \(y\le L_n^{-2}\). Indeed,
\[
 |v_r|^2
 \le\frac{\|G_n\|_{\rm op}}{\log(n-1)}
       K(r\Vert \mathbf u_{n-1})
 =O\{K(r\Vert \mathbf u_{n-1})\}.
\]
We use a sharpened form of the last term. Pinsker's inequality and the
definition of \(C_n\) give
\[
 |q_i\cdot v_r|
 \le C_n\|r-\mathbf u_{n-1}\|_1
 \le C_n\{2K(r\Vert \mathbf u_{n-1})\}^{1/2}.
\]
The spectral estimate in
Proposition~\ref{prop:local-simplex-criterion} retains a fixed fraction of
the conditional entropy \(yK(r\Vert \mathbf u_{n-1})\). Young's inequality
therefore absorbs the cross term \(yc_yq_i\cdot v_r\), at the cost of
\(O(yC_n^2)\). Thus the right side of
\eqref{eq:weighted-potts-localization} may be replaced by
\begin{equation}
\label{eq:weighted-potts-localization-sharp}
 \delta_nu_i c_y^2
 -\{k_n(y)-R_n^\star c_y^2\}
 -c\,yK(r\Vert \mathbf u_{n-1})
 +O(yC_n^2).
\end{equation}
The scalar function in braces is nonnegative and, on
\(0\le y\le L_n^{-2}\), vanishes only at \(y=1/n\); its second
derivative there is \(n+O(L_n)\).  More generally, its explicit formula
gives, on this interval,
\begin{equation}
\label{eq:potts-scalar-local-stability}
 k_n(y)-R_n^\star c_y^2
 \ge
 \frac{c}{n}\{z\log z-z+1\},
 \qquad z=ny.
\end{equation}
This follows by differentiating the two sides on either side of \(z=1\);
both vanish with zero derivative at \(z=1\), while the second derivative
of the left side after the change of variables is
\(n^{-1}z^{-1}\{1+o(1)\}\).
Combining this estimate with \eqref{eq:potts-fitted-weight-profile} gives
the following dichotomy.  If \(i\notin A\), then \(u_i\le-\eta\), and the
weighted free energy is at most \(-c\eta\delta_n\) for all large \(n\).
If \(i\in A\), the active fitted weight cancels the radial term: uniformly
on the present interval,
\begin{equation}
\label{eq:potts-active-radial-cancellation}
 \delta_nu_i c_y^2+(1-y)\log c_i
 \le O_{C,J}(\delta_n/n)+o(\delta_n),
\end{equation}
because
\(\sup_{0\le y\le L_n^{-2}}\{c_y^2-\chi_n(1-y)\}
=1/(n-1)\).
Let \(\phi(z)=z\log z-z+1\).  Outside \(z=ny\in[1/2,2]\), one has
\(\phi(z)\ge c(1+z)\), so the term \(yC_n^2=zC_n^2/n\) is absorbed by
\eqref{eq:potts-scalar-local-stability}.  Inside that interval it is
\(O(C_n^2/n)=o(\delta_n)\).  Nonnegative weighted free energy therefore
implies
\[
 \phi(ny)+nyK(r\Vert\mathbf u_{n-1})=o(n\delta_n)=o(L_n).
\]
It follows that
\[
 y=o\left(\frac{L_n}{n\log L_n}\right),
 \qquad
 yK(r\Vert\mathbf u_{n-1})=o(\delta_n).
\]
Hence
\[
 y|v_r|
 \le C\{y^2K(r\Vert\mathbf u_{n-1})\}^{1/2}
 =o\{y\delta_n\}^{1/2}=o(1).
\]
Since
\(|c_y-\chi_n|=|1-ny|/(n-1)\), we obtain
\[
 \mu(p)=c_yq_i+yv_r=\chi_nq_i+o(1).
\]
Substituting \eqref{eq:potts-active-radial-cancellation} and the bounds
\(yK(r\Vert\mathbf u_{n-1})=o(\delta_n)\) and
\(y|v_r|=o(1)\) into the definition of \(F_{\log c}\) gives
\(F_{\log c}(p)=o(\delta_n)\).  Equation
\eqref{eq:weighted-certificate-square-identity} then shows that every
certificate point with value at least one lies within \(o(1)\) of
\(\chi_nq_i\).

It remains to prove uniqueness inside the localized patches.  Fix a small
constant \(r_0>0\).  In \(B(\tau,r_0)\), the Gibbs probabilities are at
most \(n^{-1+o(1)}\), and hence
\[
 \|\operatorname{Cov}_{p(w)}(q)\|_{\rm op}
 \le n^{-1+o(1)}\|G_n\|_{\rm op}=o(1).
\]
In \(B(w_i,r_0)\), the total Gibbs mass outside coordinate \(i\) is
\(n^{-1+o(1)}\), so the covariance operator is again \(o(1)\).  Therefore
\[
 \nabla^2\log\Gamma_{\widetilde\Pi}(w)
 =\operatorname{Cov}_{p(w)}(q)-I\preceq-\frac12I
\]
on every localized patch.  Each patch contains exactly one stationary
point, namely the contact already constructed there, and that point is its
unique maximum.  The diffuse and near-vertex localization arguments in this
proof now give
\(\Gamma_{\widetilde\Pi}\le1\), with equality exactly at the central and
active contacts.

The constructed measure is an NPMLE\@. Its fitted vector is unique by
strict concavity in the fitted likelihood vector.  The likelihood vectors
at its contacts are linearly independent.  Indeed, on rows indexed by
\(A\), the active-contact block has diagonal entries
\((n-1)\{1+o(1)\}\) and off-diagonal entries
\((n-1)^{-1}\{1+o(1)\}\), while the central column is \(1+o(1)\).
Choose any inactive row, which exists because \(|A|\le J<n\) for large
\(n\).  On that row the central entry is \(1+o(1)\), and every active
entry is \((n-1)^{-1}\{1+o(1)\}\).  After scaling the active columns by
\(n^{-1}\), this \((1+|A|)\)-square minor converges to an invertible
triangular matrix.  The KKT conditions for the unique fitted likelihood
vector restrict the support of every maximizing measure to these contacts;
linear independence then makes its representing weights unique.  Thus every
maximizing measure uses the same positive weights on precisely these
contacts.  Equations
\eqref{eq:buffered-simplex-count} and \eqref{eq:buffered-simplex-mass}
follow from \eqref{eq:terminal-quantitative-mass}.
\end{proof}

For the random Gaussian cloud, all assumptions of
Proposition~\ref{prop:buffered-simplex-support} hold whenever
\(d_nL_n/n^2\) is bounded below.

\begin{corollary}
\label{cor:random-buffered-simplex-support}
Suppose
\[
 d_nL_n/n^2\ge c>0,
 \qquad s_n^\star=O(1).
\]
For every fixed \(C,J<\infty\) and \(\eta>0\),
\begin{equation}
\label{eq:random-buffered-simplex-support}
 \Pbb\left\{
  \mathcal B_n(C,J,\eta),\quad
  \widehat K_n\ne1+\#\{i:\xi_{n,i}^\star>0\}
 \right\}\longrightarrow0.
\end{equation}
On \(\mathcal B_n(C,J,\eta)\), the NPMLE is also unique apart from an event
whose probability tends to zero.
\end{corollary}

\begin{proof}
Here \(L_n=o(d_n)\), and Lemma~\ref{lem:random-off-diagonal} gives
\begin{align*}
 \delta_n^\star&=\sqrt{L_n/d_n}\{1+o(1)\},&
 \mu_n^{\rm off}:=\max_{i\ne j}|q_i\cdot q_j|
   &=o_{\Pbb}(L_n^{-1}),\\
 \|G_n\|_{\rm op}&=O_{\Pbb}(L_n),&
 C_n&=O_{\Pbb}(L_n^{3/2}/\sqrt{d_n}).
\end{align*}
Consequently,
\[
 \frac{C_n}{\delta_n^\star L_n^2}=O_{\Pbb}(L_n^{-1}),
 \qquad
 \frac{C_n^2}{n\delta_n^\star}
 =O_{\Pbb}(\delta_n^\star L_n^2/n)=o_{\Pbb}(1),
\]
and \(\delta_n^\star=O(L_n/n)\to0\).  Proposition
\ref{prop:buffered-simplex-support} applies on an event whose probability
tends to one.
\end{proof}


\subsection{Proof of the main theorem}
\label{sec:edge-count-reduction}

We finish by combining the dimension-specific localization estimates with
the radial Poisson limits. This short argument proves all assertions in
Theorem~\ref{thm:main-support-law}.

For \(d_n\ge3\), let
\begin{equation}
\label{eq:exact-support-buffer}
 \mathcal B_n(C,J,\eta)
 =
 \left\{
 \max_i\xi_{n,i}^\star\le C,\quad
 \#\{i:\xi_{n,i}^\star>-\eta\}\le J,\quad
 \min_i|\xi_{n,i}^\star|>\eta
 \right\}.
\end{equation}
\begin{proof}[Proof of Theorem~\ref{thm:main-support-law}]
First let \(d=2\), and put
\(\xi_{n,i}=|Z_i|^2/2-\log n\). For every fixed \(\eta>0\),
Proposition~\ref{prop:fixed-dimensional-edge-sandwich} gives
\[
 \#\{i:\xi_{n,i}>x_\lambda+\eta\}
 \le \widehat K_n-1
 \le \#\{i:\xi_{n,i}>x_\lambda-\eta\},
\]
apart from an event whose limiting probability is \(O(\eta)\). The extremal
process of the \(\xi_{n,i}\)'s converges to a Poisson process with intensity
\(e^{-x}\,dx\). Letting \(n\to\infty\) and then \(\eta\downarrow0\) gives
\[
 \widehat K_n-1\Rightarrow\operatorname{Pois}(e^{-x_\lambda}).
\]
Taking the probability of zero gives
\(\Pbb\{\widehat K_n=1\}\to Q_2(\lambda)\).
Proposition~\ref{prop:fixed-dimensional-edge-sandwich} also shows that the
NPMLE is unique with probability tending to one.

Now suppose \(d_n\ge3\) and \(s_n^\star\to s\). We first connect support
points to radial exceedances. For every fixed \(C,J<\infty\) and \(\eta>0\),
\begin{equation}
\label{eq:universal-edge-count-sandwich}
 \#\{i:\xi_{n,i}^\star>\eta\}
 \le \widehat K_n-1
 \le \#\{i:\xi_{n,i}^\star>-\eta\}
\end{equation}
on \(\mathcal B_n(C,J,\eta)\), apart from an event whose probability tends to
zero. To verify this uniformly, argue along subsequences. If \(d_n\) is
bounded, pass to a further subsequence on which \(d_n=d\ge3\) is fixed.
Lemma~\ref{lem:fixed-d-normalization-bridge} converts the exact excess and
threshold into their fixed-dimensional versions, and
Proposition~\ref{prop:fixed-dimensional-edge-sandwich} gives the claim. If
\(d_n\to\infty\), pass to a further subsequence on which either
\(d_n\log n=o(n^2)\) or \(d_n\log n/n^2\) is bounded below. The exact support
identity follows respectively from
Proposition~\ref{prop:buffered-support-below-simplex} or
Corollary~\ref{cor:random-buffered-simplex-support}. These alternatives cover
every subsequence and prove \eqref{eq:universal-edge-count-sandwich}.  They
also show that, on \(\mathcal B_n(C,J,\eta)\), the NPMLE is unique apart from
an event whose probability tends to zero.

The centered radial exceedance count satisfies
\begin{equation}
\label{eq:exact-radial-exceedance-count}
 \#\{i:\xi_{n,i}^\star>0\}
 \ \Rightarrow\ \operatorname{Pois}(e^{-s}).
\end{equation}
For bounded \(d_n\), this follows from Lemmas
\ref{lem:radial-gamma-extremes} and
\ref{lem:fixed-d-normalization-bridge}. If \(d_n\to\infty\) and
\(d_n\log n=o(n^2)\), Proposition~\ref{prop:gamma-max} gives the uncentered
point-process limit. Lemmas~\ref{lem:centered-radial-max} and
\ref{lem:centered-radial-max-high-d} compare the centered and uncentered
radial energies uniformly over the sample. They therefore transfer counts at
thresholds separated from zero by a fixed amount; letting the separation
decrease proves \eqref{eq:exact-radial-exceedance-count}. In the complementary
range, the count assertion is part of
Proposition~\ref{prop:direct-centered-gamma-max}.

For fixed \(\eta>0\), the expected number of excesses in
\([-\eta,\eta]\) tends to
\[
 e^{-s}(e^\eta-e^{-\eta})=O_s(\eta).
\]
The radial point-process limits used in
\eqref{eq:exact-radial-exceedance-count} also make the maximum excess and
the number of excesses above \(-\eta\) tight. We may therefore choose
\(C,J\) so that
\(\mathcal B_n(C,J,\eta)\) fails, apart from an excess in
\([-\eta,\eta]\), with arbitrarily small probability. The sandwich
\eqref{eq:universal-edge-count-sandwich}, followed by \(n\to\infty\),
\(C,J\to\infty\), and \(\eta\downarrow0\), gives
\[
 \widehat K_n-1\Rightarrow\operatorname{Pois}(e^{-s}).
\]
The uniqueness conclusions in
Proposition~\ref{prop:fixed-dimensional-edge-sandwich},
Proposition~\ref{prop:buffered-support-below-simplex}, and
Corollary~\ref{cor:random-buffered-simplex-support}, together with the
preceding buffer estimate, show that the probability of nonuniqueness tends
to zero.
Again, the probability of zero gives the one-atom limit.

It remains to prove the off-critical tightness assertions in
Theorem~\ref{thm:main-support-law}. In dimension two, pass
along subsequences on which \(\lambda_n\) converges in
\((0,\infty]\). A finite limit is covered above. If
\(\lambda_n\to\infty\), monotonicity compares the one-atom event with every
fixed \(\lambda_0\); its limiting probability is
\(Q_2(\lambda_0)\uparrow1\).  For \(d_n\ge3\), any subsequence on which
\(s_n^\star\) is bounded below has a further subsequence on which
\(s_n^\star\) converges in \(\mathbb R\) or tends to \(+\infty\).
The first case is covered by the Poisson limit. In the second, monotonicity
and comparison with a fixed \(s_0\), followed by \(s_0\to\infty\), show that
\(\widehat K_n=1\) with probability tending to one.
\end{proof}

Finally, monotonicity converts the one-atom limits into the samplewise
critical-defect laws stated in the introduction.  In dimension two, the map
\[
 c_{2,n}(\eps)=\frac{\rho_n^2\eps}{2(1-\eps)}
\]
is continuous and strictly increasing. Hence, for every \(\lambda>0\),
\[
 \Pbb\{c_{2,n}(\eps_*(Z_{1:n}))\le\lambda\}
 =\Pbb\{E_n^{(1)}(\eps_n(\lambda))\}
 \longrightarrow Q_2(\lambda),
\]
where \(\eps_n(\lambda)=c_{2,n}^{-1}(\lambda)\). For \(d_n\ge3\), use
instead the critical coordinate
\[
 c_{\ge3,n}(\eps)
 =\frac{R_n^\star/[\sigma_n^2(1-\eps)]-b_n}{a_n}
\]
is also continuous and strictly increasing.  Evaluating the one-atom limit
at \(c_{\ge3,n}(\eps)=s\) gives
\(\Pbb\{S_{*,n}\le s\}\to\exp\{-e^{-s}\}\).


\subsection{The regularization path}
\label{sec:regularization-path}

This subsection proves Theorem~\ref{thm:regularization-path} and then refines
it with the mass and spatial-mark limits in
Proposition~\ref{prop:path-masses}. The fixed-parameter support identities can
be made uniform across a compact part of the critical window. To obtain a path
limit, we must also resolve what happens at each activation threshold: as the
variance defect passes the threshold associated with one radial observation,
the corresponding edge mass reaches zero transversally. We first record the
radial process that supplies these thresholds, then prove the required uniform
continuation.

For dimension two and dimensions at least three, respectively, set
\begin{equation}
\label{eq:path-radial-processes}
 \mathcal X_n^{(2)}
 =\sum_{i=1}^n\delta_{\zeta_{n,i}^{(2)}},
 \quad
 \zeta_{n,i}^{(2)}=\frac{|Z_i|^2}{2}-\log n,
 \qquad
 \mathcal X_n^{(\ge3)}
 =\sum_{i=1}^n\delta_{\zeta_{n,i}^{(\ge3)}},
 \quad
 \zeta_{n,i}^{(\ge3)}
 =\frac{|Z_i-\bar Z_n|^2/(2\sigma_n^2)-b_n}{a_n}.
\end{equation}
Thus the exact excess in \eqref{eq:exact-transition-coordinate}, evaluated
along \eqref{eq:intro-highd-path-parametrization}, is
\(\xi_{n,i}^\star(s)=\zeta_{n,i}^{(\ge3)}-s\).

The one-level exceedance limits used to prove
Theorem~\ref{thm:main-support-law} hold jointly at every finite collection of
levels. Equivalently, the whole radial cloud has the following limit.

\begin{proposition}
\label{prop:path-radial-process}
In dimension two, and for every sequence \(d_n\ge3\),
\begin{equation}
\label{eq:path-radial-ppp}
 \mathcal X_n^{(2)}\Longrightarrow\mathcal X,
 \qquad
 \mathcal X_n^{(\ge3)}\Longrightarrow\mathcal X
\end{equation}
vaguely on \(\R\), where \(\mathcal X\) has intensity \(e^{-x}\,dx\).
For each fixed \(d\ge2\), the point-process convergence remains valid when
each radial excess is marked by the direction of its observation; the
limiting directions are independent and uniform on \(S^{d-1}\).
\end{proposition}

\begin{proof}
For \(d=2\), the variables \(|Z_i|^2/2\) are standard exponential, and hence
\(n\Pbb\{\zeta_{n,i}^{(2)}>x\}=e^{-x}\) for every fixed \(x\). Counts in
disjoint bounded intervals are multinomial rare-event counts, which proves the
first convergence.

For uncentered Gamma energies, Lemma~\ref{lem:gamma-hazard} gives, uniformly
for bounded \(x\),
\[
 n\overline F_n(b_n+a_nx)\longrightarrow e^{-x}.
\]
The rare-event multinomial calculation from the \(d=2\) case proves
convergence of the uncentered point process. In the range where empirical
centering is negligible on the
\(a_n\)-scale, Lemmas~\ref{lem:centered-radial-max} and
\ref{lem:centered-radial-max-high-d} transfer the process. In the remaining
range, the proof of Proposition~\ref{prop:direct-centered-gamma-max} gives
joint upper-tail factorization uniformly when finitely many thresholds lie in
\(b_n+a_nO(1)\). Although \eqref{eq:centered-joint-tail-factorization} is
displayed there with a common threshold, its determinant-and-tilting proof
also allows bounded coordinate-dependent offsets: use the corresponding
diagonal matrix of tilt parameters in place of the scalar tilt. The
determinant comparison and tilted local limit are unchanged.
Inclusion--exclusion then gives the factorization
for disjoint intervals, so the factorial-moment criterion proves the second
convergence in \eqref{eq:path-radial-ppp}. For fixed dimension, radius and
direction are independent; the marking theorem completes the proof.
\end{proof}

We next follow the optimizer while one point crosses its activation level. In
dimension two, let \(a_\lambda>0\) be the unique solution of
\(\Phi(a_\lambda)+\phi(a_\lambda)/a_\lambda=e^\lambda\), and put
\begin{equation}
\label{eq:path-d2-mass-profile}
 b_\lambda=1-e^{-\lambda}\Phi(a_\lambda),
 \qquad
 q_{\lambda,x}
 =\frac{[1-e^{-(x-x_\lambda)}]_+}{b_\lambda},
\end{equation}
and recall that \(\Lambda\) is the inverse of
\(\lambda\mapsto x_\lambda\). For \(d_n\ge3\), define
\begin{equation}
\label{eq:path-highd-scales}
 \gamma_n(s)=\frac{\sigma_n^2(b_n+a_ns)}{R_n^\star},
 \qquad
 \delta_n(s)=\frac{\sigma_n^2a_n}{\gamma_n(s)}.
\end{equation}

The next proposition is the uniform edge reduction. Its separation
hypothesis is only a proof device; Proposition~\ref{prop:path-radial-process}
will remove it.

\begin{proposition}
\label{prop:path-uniform-activation}
Fix compact intervals
\(J\subset\operatorname{int}(J^+)\Subset(0,\infty)\) and
\(I\subset\operatorname{int}(I^+)\subset\R\), and constants
\(C,M<\infty\) and \(\eta>0\). In dimension two, work on the event
\begin{equation}
\label{eq:path-d2-buffer-event}
 \max_i\zeta_{n,i}^{(2)}\le C,
 \qquad
 \#\{i:\zeta_{n,i}^{(2)}\ge x_{\inf J^+}\}\le M,
\end{equation}
on which the values \(\Lambda(\zeta_{n,i}^{(2)})\in J^+\) are pairwise at
least \(4\eta\) apart and are at least \(4\eta\) from the endpoints of
\(J^+\). Apart from an event whose probability tends to zero, each such point
has one activation value
\(\wh\lambda_{n,i}\), and
\begin{equation}
\label{eq:path-d2-activation}
 \max_i|\wh\lambda_{n,i}-\Lambda(\zeta_{n,i}^{(2)})|=o_{\Pbb}(1).
\end{equation}
Its patch contains one fitted atom when
\(\lambda<\wh\lambda_{n,i}\) and none when
\(\lambda\ge\wh\lambda_{n,i}\), and, uniformly for \(\lambda\in J\),
\begin{equation}
\label{eq:path-d2-patch-mass}
 n\pi_{n,i}^{(2)}(\lambda)
 =q_{\lambda,\zeta_{n,i}^{(2)}}+o_{\Pbb}(1).
\end{equation}
Points above \(x_{\sup J^+}\) are active throughout \(J\), points below
\(x_{\inf J^+}\) are inactive throughout \(J\), and the mass estimate applies
to every point satisfying the second cutoff in
\eqref{eq:path-d2-buffer-event}.

For every sequence \(d_n\ge3\), work instead on the event
\begin{equation}
\label{eq:path-highd-buffer-event}
 \max_i\zeta_{n,i}^{(\ge3)}\le C,
 \qquad
 \#\{i:\zeta_{n,i}^{(\ge3)}\ge\inf I^+-1\}\le M,
\end{equation}
on which the points in \(I^+\) are pairwise at least \(4\eta\) apart and at
least \(4\eta\) from the endpoints of \(I^+\). Apart from an event whose
probability tends to zero, each point in \(I^+\) has one activation value
\(\wh s_{n,i}\), and
\begin{equation}
\label{eq:path-highd-activation}
 \max_i|\wh s_{n,i}-\zeta_{n,i}^{(\ge3)}|=o_{\Pbb}(1).
\end{equation}
Its patch contains one fitted atom when \(s<\wh s_{n,i}\) and none when
\(s\ge\wh s_{n,i}\). Fixing \(s_0\in I\) and writing
\(\bar\delta_n=1\wedge\delta_n(s_0)\), its mass satisfies, uniformly for
\(s\in I\),
\begin{equation}
\label{eq:path-highd-patch-mass}
 n\pi_{n,i}^{(\ge3)}(s)
 =\left[1-e^{-\delta_n(s)(\zeta_{n,i}^{(\ge3)}-s)}\right]_+
  +o_{\Pbb}(\bar\delta_n),
\end{equation}
with the exponential expression replaced by its linearization when
\(d_nL_n/n^2\) is bounded below, the regime in which the centered cloud is
asymptotically simplex-like. Points above \(\sup I^+\) are active throughout \(I\), points
below \(\inf I^+\) are inactive throughout \(I\), and the mass estimate applies
to every point in the second cutoff in \eqref{eq:path-highd-buffer-event}. In
both dimensions there are no other noncentral contacts, and the NPMLE is
unique throughout the original compact interval.
\end{proposition}

\begin{proof}
We first explain why the fixed-parameter estimates are uniform. In dimension
two, the local diffuse background and the response of a radial point of excess
\(x\) are
\[
 B_\lambda(r)=e^{-\lambda}\Phi(-r),
 \qquad
 C_{\lambda,x}(r,z)=e^{-\lambda}e^{x-r^2/2-z^2/2}.
\]
Their spatial derivatives through order two and their first
\(\lambda\)-derivatives are bounded uniformly when \(\lambda\) ranges over a
compact subset of \((0,\infty)\). At finite \(n\),
\[
 \partial_\lambda e^{-\kappa_n|t|^2}
 =-\frac{|t|^2}{\rho_n^2}e^{-\kappa_n|t|^2},
 \qquad \kappa_n=\lambda/\rho_n^2,
\]
so the Gaussian envelopes used for the undifferentiated field also dominate
its \(\lambda\)-derivative on \(|t|=\rho_n+O(1)\). Applying the compact decomposition in
Proposition~\ref{prop:two-d-compact-decomposition}, the lower-cloud bounds, and
the \(C^2\) estimates of Section~\ref{sec:higher-dimensional-support-counts}
on a deterministic \(\lambda\)-grid, then using these derivative bounds
between grid points, makes every fixed-parameter estimate uniform in
\(C^2(r,z)\cap C^1(\lambda)\). The radial compactness and outside-window gaps
in Proposition~\ref{prop:d-two-radial-compactness} and
Lemmas~\ref{lem:d-two-diffuse-outside}--\ref{lem:d-two-tail-outside} are
uniform for the same reason: on \(J^+\), every damping constant stays in a
fixed compact subset of \((0,\infty)\).

It remains to inspect the limiting patch program
\eqref{eq:support-d2-local-program}. Its unique spatial optimizer is
\((-a_\lambda,0)\). After the other patch variables have been optimized, the
derivative at zero rescaled mass is
\begin{equation}
\label{eq:path-d2-boundary-score}
 S_{\lambda,x}(0)=b_\lambda\{e^{x-x_\lambda}-1\}.
\end{equation}
The functions \(a_\lambda,x_\lambda,b_\lambda\) are smooth on compact
intervals, with \(x_\lambda'>0\) and \(b_\lambda>0\). At the unique zero
\(\lambda=\Lambda(x)\),
\[
 \partial_\lambda S_{\lambda,x}(0)=-b_\lambda x_\lambda',
\]
which is bounded away from zero. The spatial Hessian there is
\(-b_\lambda I_2\), and the mass curvature is strictly negative. The uniform
finite-patch expansion therefore gives one empirical zero satisfying
\eqref{eq:path-d2-activation}. Strict concavity places positive mass below
that zero and zero mass at and above it. Solving the limiting interior score
equation gives \eqref{eq:path-d2-patch-mass}. Continuity of the optimizer at
zero makes the estimate uniform through the crossing.

For \(d_n\ge3\), write
\[
 h_n(s)=\frac{\gamma_n'(s)}{\gamma_n(s)}
       =\frac{a_n}{b_n+a_ns}.
\]
If \(q_j(s)=y_j/\sqrt{\gamma_n(s)}\), differentiation of an edge kernel gives
\begin{align}
 &\left|\partial_s
 \exp\left\{(q_i(s)+z)\cdot q_j(s)-\frac{|q_i(s)+z|^2}{2}\right\}\right|
 \notag\\
 &\hspace{35mm}\le
 C h_n(s)\{\log n+|q_j(s)-q_i(s)-z|^2\}
 \exp\left\{(q_i(s)+z)\cdot q_j(s)-\frac{|q_i(s)+z|^2}{2}\right\}.
\label{eq:path-kernel-derivative-main}
\end{align}
Moreover \(h_n(s)\log n\asymp\delta_n(s)\), uniformly on \(I\). Thus the
order-zero and order-two absolute envelopes already proved in the spatial,
entropy, and support-recovery arguments also control one derivative in \(s\).
As in dimension two, a slowly refining deterministic grid makes all buffered
support identities uniform on \(I\).

Near one crossing, optimize the central contact and every other patch, leaving
only the rescaled mass \(m_i=n\pi_i\ge0\) in the crossing patch free. The
algebraic expansion in the proof of
Lemma~\ref{lem:uniform-finite-edge-expansion} is obtained before the threshold
buffer is used: its envelope assumptions do not depend on the sign of the
crossing excess. Holding this one mass free therefore gives the expansion in
Lemma~\ref{lem:uniform-finite-edge-expansion} uniformly while
\(|\zeta_{n,i}^{(\ge3)}-s|\le2\eta\), with every other patch
still buffered. In particular, the boundary score and its derivative are
\begin{align}
 \partial_{m_i}\mathcal L_{n,i}(s,0)
 &=e^{\delta_n(s)(\zeta_{n,i}^{(\ge3)}-s)}-1
   +o_{\Pbb}(\bar\delta_n),
\label{eq:path-highd-boundary-score}\\
 \partial_s\partial_{m_i}\mathcal L_{n,i}(s,0)
 &=-\delta_n(s)+o_{\Pbb}(\bar\delta_n),
\label{eq:path-highd-boundary-derivative}
\end{align}
while its second mass derivative is \(-1+o_{\Pbb}(1)\). Indeed,
\(\delta_n'=-h_n\delta_n\), so differentiating the leading exponent adds
only \(O\{h_n\delta_n|\zeta_{n,i}^{(\ge3)}-s|\}\) to
\(-\delta_n\); inequality \eqref{eq:path-kernel-derivative-main} keeps every
remainder \(o_{\Pbb}(\bar\delta_n)\) after differentiation. When
\(d_nL_n/n^2\) is bounded below, the ordered-contact argument gives
\eqref{eq:path-highd-boundary-score}--\eqref{eq:path-highd-boundary-derivative}
with the exponential linearized. More explicitly, the ordered-contact
response remainder in
\eqref{eq:potts-corrected-height} is bounded by a constant times
\[
 \delta_n^2+\frac{C_n^2}{n}
 +\frac{\log n\,(C_n+\delta_n)^2}{n^2}.
\]
Here \(\delta_n'=-h_n\delta_n\), and the Gram-error array defining \(C_n\)
is rescaled by \(\gamma_n(s)^{-1}\), so its upper Dini derivative is at most
\(h_nC_n\). The rates in \eqref{eq:terminal-simplex-derived-rates} therefore
show directly that one \(s\)-derivative leaves this remainder
\(o_{\Pbb}(\delta_n)\).

Equation \eqref{eq:path-highd-boundary-derivative} is uniformly negative.
Hence each boundary score has exactly one zero, which proves
\eqref{eq:path-highd-activation}; strict mass curvature and the spatial
certificate give one contact below the zero and none at or above it. Solving
the interior score equation gives \eqref{eq:path-highd-patch-mass}. Finally,
the uniform residual KKT gaps proved in the dimension-specific support
identities exclude contacts outside these patches. The strict concavity of
the local mass programs and the uniqueness of each spatial contact then prove
simultaneous uniqueness, including at an activation value, where the crossing
patch has zero mass.
\end{proof}

\begin{proof}[Proof of Theorem~\ref{thm:regularization-path}]
Enlarge the parameter interval slightly. Proposition
\ref{prop:path-radial-process} makes the radial maximum and the number of
points in the enlarged critical window tight. The limiting process is simple,
so after imposing fixed radial and cardinality cutoffs, its activation levels
are separated from one another and from the enlarged interval's endpoints
with probability tending to one as the separation constant decreases to zero.
Apply Proposition~\ref{prop:path-uniform-activation} first with fixed cutoffs
and separation, then let the cutoffs tend to infinity and the separation tend
to zero. A diagonal choice gives a single event of probability tending to one
on which all of its conclusions hold.

In dimension two, the support-count path is the tail-count path of the
perturbed activation measure
\(\sum_i\delta_{\wh\lambda_{n,i}}\). Equation
\eqref{eq:path-d2-activation} allows a piecewise-linear time change that moves
each \(\wh\lambda_{n,i}\) to \(\Lambda(\zeta_{n,i}^{(2)})\) and converges
uniformly to the identity. Points above the enlarged interval contribute the
same constant to both paths. After adding them, the resulting path is exactly
\(\mathcal X_n^{(2)}((x_\lambda,\infty))\). Proposition
\ref{prop:path-radial-process} and the continuous-mapping theorem for tail
counts of a simple point process prove \eqref{eq:intro-d2-path-limit} in the
\(J_1\) topology.

In dimensions at least three, equation
\eqref{eq:path-highd-activation} supplies an analogous time change. After
that time change and the addition of
the points active throughout \(I\), the path is
\(\mathcal X_n^{(\ge3)}((s,\infty))\). This proves
\eqref{eq:intro-highd-path-limit}. Simultaneous uniqueness is part of
Proposition~\ref{prop:path-uniform-activation}.
\end{proof}

Proposition~\ref{prop:path-uniform-activation} also identifies the mass and
location carried by every jump. We state
the mass limit explicitly because it is stronger than the support-count path
and shows that the fitted mixing distribution itself changes continuously even
when its support size jumps.

\begin{proposition}
\label{prop:path-masses}
Associate each noncentral fitted atom with its radial observation, and extend
its mass by zero after the corresponding patch becomes inactive. In dimension
two, for compact \(J\Subset(0,\infty)\),
\begin{equation}
\label{eq:path-d2-mass-limit}
 \max_{i:\,\zeta_{n,i}^{(2)}\ge x_{\inf J}-1}
 \sup_{\lambda\in J}
 \left|n\pi_{n,i}^{(2)}(\lambda)
       -q_{\lambda,\zeta_{n,i}^{(2)}}\right|
 \xrightarrow{\Pbb}0,
\end{equation}
after first imposing and then removing a fixed upper radial cutoff. Moreover,
\begin{equation}
\label{eq:path-d2-total-mass-limit}
 \left\{n\sum_i\pi_{n,i}^{(2)}(\lambda):\lambda\in J\right\}
 \Longrightarrow
 \left\{\sum_{x\in\mathcal X}q_{\lambda,x}:\lambda\in J\right\}
 \qquad\text{in }C(J).
\end{equation}

For \(d_n\ge3\), fix \(s_0\in I\), pass to a subsequence on which
\(\delta_n(s_0)\to\delta\in[0,\infty]\), and define
\[
 \psi_\delta(y)=
 \begin{cases}
 (1-e^{-\delta y})/(1\wedge\delta),&0<\delta<\infty,\\
 y,&\delta=0,\\
 1,&\delta=\infty,
 \end{cases}
 \qquad y>0,
 \quad \psi_\delta(y)=0\quad(y\le0).
\]
Then
\begin{equation}
\label{eq:path-highd-mass-limit}
 \max_{i:\,\zeta_{n,i}^{(\ge3)}\ge\inf I-1}
 \sup_{s\in I}
 \left|\frac{n\pi_{n,i}^{(\ge3)}(s)}{\bar\delta_n}
       -\psi_\delta((\zeta_{n,i}^{(\ge3)}-s)_+)\right|
 \xrightarrow{\Pbb}0,
\end{equation}
and the sum of these profiles converges in \(C(I)\) to
\(s\mapsto\sum_{x\in\mathcal X}\psi_\delta((x-s)_+)\). For every fixed
\(d\ge2\), the convergence is joint with the radial directions; each active
atom has the direction of its associated extreme observation asymptotically.
In dimension two its radial score location is
\(\rho_n-a_\lambda+o_{\Pbb}(1)\), uniformly away from activation.
\end{proposition}

\begin{proof}
Equations \eqref{eq:path-d2-patch-mass} and
\eqref{eq:path-highd-patch-mass} give the individual limits on events with
fixed radial and cardinality cutoffs. Lemma~\ref{lem:gamma-hazard} gives
\(a_n/b_n\to0\), and hence
\(\delta_n(s)/\delta_n(s_0)\to1\) uniformly on \(I\); this turns
\eqref{eq:path-highd-patch-mass}, after division by \(\bar\delta_n\), into the
profile \(\psi_\delta\). Proposition
\ref{prop:path-radial-process} removes these cutoffs. The profiles are jointly
continuous in the path parameter and radial mark, including at activation
where they vanish. Only finitely many limiting Poisson points contribute on a
compact parameter interval. Summation is therefore a continuous map into
\(C(J)\) or \(C(I)\), which proves the total-mass limits. The marked part of
Proposition~\ref{prop:path-radial-process}, together with the unique spatial
optimizer in each local likelihood program, gives the final assertions.
\end{proof}


\small
\bibliographystyle{alpha}
\bibliography{writing_examples/bib,references,gpt56_multivariate_variance_refs}

\end{document}
